\documentclass[a4paper]{article}


\usepackage[T1]{fontenc}
\usepackage[utf8]{inputenc}
\usepackage{ifthen}
\usepackage{pagecolor}
\usepackage[scaled]{helvet}
\usepackage{tikz}
\usepackage{titlesec}
\usepackage{pgfmath}


\usepgflibrary{fpu}
\usetikzlibrary{shapes.geometric}

\newcommand\uv[1]{\quotedblbase #1\textquotedblleft}%


\paperheight=297mm
\paperwidth=210mm
% nestandardni odsazeni (std. po 1 palci zleva a shora]
\hoffset=-1in
\voffset=-1in  \addtolength{\voffset}{10mm}
% nastaveni rozmeru textoveho tela
\textheight=272mm
\textwidth=190mm
% odsazeni spodni hrany paty od spodni hrany textoveho tela
\footskip=2mm
% pridani odsazeni zleva na sudych i lichych strankach
\oddsidemargin=10mm
\evensidemargin=10mm
% vynulovani ostatnich rozmeru.
\marginparwidth=0pt
\marginparsep=0pt
\topmargin=0pt
\headheight=0pt
\headsep=0pt


\setlength{\parindent}{0pt}


\renewcommand*{\thepage}{\color{white}--~\arabic{page}~--}
\renewcommand\familydefault{\sfdefault}
\pagecolor{black!95}
\titleformat*{\section}{\sf\bf\LARGE%
\color{yellow!70!orange}%
%\color{white}%
}


\def\FCtitle#1#2{%
  \ifthenelse{\equal{#2}{}}{%
    {\color{white}\bf Fight Club #1}%
  }{%
    {\color{white}\bf Fight Club #1 -- #2}%
  }%
}

\def\TCtitle#1#2{%
  \ifthenelse{\equal{#2}{}}{%
    {\color{white}\bf Tea Club #1}%
  }{%
    {\color{white}\bf Tea Club #1 -- #2}%
  }%
}

\def\SCtitle#1#2{%
  \ifthenelse{\equal{#2}{}}{%
    {\color{white}\bf Social Club #1}%
  }{%
    {\color{white}\bf Social Club #1 -- #2}%
  }%
}

\def\HCtitle#1#2{%
  \ifthenelse{\equal{#2}{}}{%
    {\color{white}\bf Hardware Club #1}%
  }{%
    {\color{white}\bf Hardware Club #1 -- #2}%
  }%
}


% #1 - cislo pravidla fight clubu
% #2 - text pravidla
% #3 - autor vyroku pravidla
\def\FCrule#1#2#3{Pravidlo #1: {\em #2} (vyřkl #3)}


% \def\filesize#1{%
%   \newcount\constKilo \constKilo=1024\relax%
%   \newcount\constTen  \constTen=10\relax%
%   \newcount\varX \varX=#1\relax%
%   \ifthenelse{#1 < \number\constKilo}{%
%     \number\varX\,B%
%   }{%
%     \divide\varX by \number\constKilo\relax%
%     \ifthenelse{\number\varX < \number\constKilo}{%
%       \number\varX\,KiB%
%     }{%
%       \divide\varX by \number\constKilo\relax%
%       \number\varX,\msdZ\,MiB%
%     }%
%   }%
% }
\def\filesize#1{
  \pgfkeys{/pgf/fpu=true}%
  \pgfkeys{/pgf/fpu/output format=fixed}%
  \pgfkeys{/pgf/number format/precision=2}%
  \def\pgfres{\pgfmathroundto{\pgfmathresult}\pgfmathresult}%
  %
  \newcount\varX \varX=#1\relax%
  \newcount\constKilo \constKilo=1024\relax%
  \pgfmathparse{\varX}%
  \ifthenelse{#1 < \number\constKilo}{%
    \number\varX\,B%
  }{%
    \divide\varX by \number\constKilo\relax%
    \pgfmathparse{\pgfmathresult / 1024}%
    \ifthenelse{\number\varX < \number\constKilo}{%
      \pgfres\,KiB%
    }{%
    \xdef\X{\number\varX}%
      \divide\varX by \number\constKilo\relax%
      \pgfmathparse{\pgfmathresult / 1024}%
      \ifthenelse{\number\varX < \number\constKilo}{%
        \pgfres\,MiB%
      }{%
        \pgfmathparse{\pgfmathresult / 1024}%
        \pgfres\,MiB%
      }%
    }%
  }%
}


\def\score#1{
\pgfmathsetmacro\pgfxa{#1+1}
\tikzstyle{scorestars}=[star, star points=5, star point ratio=2.25, draw,inner sep=1.3pt,anchor=outer point 3]
  \begin{tikzpicture}[baseline]
    \foreach \i in {1,...,3} {
    \pgfmathparse{(\i<=#1?"yellow!70!orange":"black!95")}
    \edef\starcolor{\pgfmathresult}
    \draw (\i*1.75ex,0) node[name=star\i,scorestars,fill=\starcolor]{};
   }
  \end{tikzpicture}
}


% #1 - filename
% #2 - oznaceni dilu fight clubu
% #3 - datum zaznamu
% #4 - seznam diskutujicich osob
% #5 - pravidlo
% #6 - popis dilu
% #7 - delka zaznamu
% #8 - velikost [B]
% #9 - uzivatelske hodnoceni
\def\FC#1#2#3#4#5#6#7#8#9{{
\hskip -4pt%
\color{black!10}
\begin{minipage}{\linewidth}
  \hrulefill\raisebox{-3pt}{\ifthenelse{\equal{#9}{}}{\score{0}}{\score{#9}}} \\[2pt]
  #2 (%
    \ifthenelse{\equal{#3}{}}{}{dne #3}%
    \ifthenelse{\equal{#7}{}}{}{; délka #7}%
    \ifthenelse{\equal{#8}{}}{}{; velikost \filesize{#8}}%
    )  \\
  Diskutují: #4 \\
  \ifthenelse{\equal{#5}{}}{}{#5 \\}
  #6
\end{minipage}
}}

% #1 - filename
% #2 - oznaceni dilu tea clubu
% #3 - datum zaznamu
% #4 - seznam diskutujicich osob
% #5 - popis dilu
% #6 - delka zaznamu
% #7 - velikost [B]
% #8 - uzivatelske hodnoceni
\def\TC#1#2#3#4#5#6#7#8{{
\hskip -4pt%
\color{black!10}
\begin{minipage}{\linewidth}
  \hrulefill\raisebox{-3pt}{\ifthenelse{\equal{#8}{}}{\score{0}}{\score{#8}}} \\[2pt]
  #2 (#3%
    \ifthenelse{\equal{#6}{}}{}{; délka #6}%
    \ifthenelse{\equal{#7}{}}{}{; velikost \filesize{#7}}%
    )  \\
  Diskutují: #4 \\
  #5
\end{minipage}
}}

% #1 - filename
% #2 - oznaceni dilu social clubu
% #3 - datum zaznamu
% #4 - seznam diskutujicich osob
% #5 - popis dilu
% #6 - delka zaznamu
% #7 - velikost [B]
% #8 - uzivatelske hodnoceni
\def\SC#1#2#3#4#5#6#7#8{{
\hskip -4pt%
\color{black!10}
\begin{minipage}{\linewidth}
  \hrulefill\raisebox{-3pt}{\ifthenelse{\equal{#8}{}}{\score{0}}{\score{#8}}} \\[2pt]
  #2 (#3%
    \ifthenelse{\equal{#6}{}}{}{; délka #6}%
    \ifthenelse{\equal{#7}{}}{}{; velikost \filesize{#7}}%
    )  \\
  Diskutují: #4 \\
  #5
\end{minipage}
}}


\def\autor#1{{\color{yellow!70!orange} #1}}

\def\Sleepless{\autor{Petr Poláček}}
\def\Luke{\autor{Pavel Dobrovský}}
\def\Janchor{\autor{Lukáš Grygar}}
\def\Blixa{\autor{Martin Bach}}
\def\Hellboy{\autor{Dan Vávra}}
\def\JOlejnik{\autor{Jan Olejník}}
\def\Farid{\autor{Farid Starman}}
\def\Tukan{\autor{Mikoláš Tuček}}
\def\Bety{\autor{Alžběta Trojanová}}
\def\Baty{\autor{Zuzana Cinková}}
\def\Cocik{\autor{Lukáš Codr}}
\def\Bludr{\autor{Honza Modrák}}
\def\JJirkovsky{\autor{Jan Jirkovský}}
\def\Cerv{\autor{Jakub Červinka}}
\def\Skeldal{\autor{Jindřích Rohlík}}
\def\Draken{\autor{Ondřej Průša}}
\def\FF{\autor{František Fuka}}
\def\KDrda{\autor{Karel Drda}}
\def\LMacura{\autor{Lukáš Macura}}
\def\MRosa{\autor{Marek Rosa}}
\def\PBulir{\autor{Petr Bulíř}}
\def\MEndrle{\autor{Miloš Endrle}}
\def\VPravda{\autor{Vašek Pravda}}
\def\JFaltus{\autor{Jarda Faltus}}
\def\Wolf{\autor{Michal Rybka}}
\def\Tlamiczka{\autor{Nikola Bornová}}
\def\PKolacek{\autor{Petr Koláček}}
\def\MPach{\autor{Martin Pach}}
\def\MJelinek{\autor{Martin Jelínek}}
\def\JHussar{\autor{Jakub Hussar}}
\def\Ulrik{\autor{Aleš Smutný}}
\def\JOrna{\autor{Jan Orna}}
\def\JSvelch{\autor{Jaroslav Švelch}}
\def\DKubec{\autor{David Kubec}}
\def\Jorssl{\autor{Jarda Šálek}}
\def\OverWatch{\autor{Adam Homola}}
\def\PKuba{\autor{Petr Kuba}}
\def\PPycha{\autor{Petr Pýcha}}
\def\OBroukal{\autor{Ondřej Broukal}}
\def\Laibach{\autor{Laibach}}
\def\Moolker{\autor{Miloš Bohonek}}
\def\FDoksansky{\autor{Filip Doksanský}}
\def\JVitik{\autor{Jan Vitík}}
\def\KMatejka{\autor{Karel Matějka}}
\def\APlechaty{\autor{Adam Plechatý}}
\def\LBasta{\autor{Lukáš Bašta}}
\def\Jorssk{\autor{Jarda Šálek}}
\def\JBeranek{\autor{Jan Beránek}}
\def\PProuza{\autor{Pavel Prouza}}
\def\AHavlova{\autor{Alžběta Havlova}}
\def\EGrillova{\autor{Eva Grillová}}
\def\VSahula{\autor{Vašek Sahula}}
\def\OBroukal{\autor{Ondřej Broukal}}
\def\ZKosticova{\autor{Zuzana Kostićová}}
\def\BHeblik{\autor{Bohdan Heblík}}
\def\Julka{\autor{Julka}}
\def\VRybar{\autor{Vašek Rybář}}
\def\JDvorsky{\autor{Jakub Dvorský}}
\def\Fiola{\autor{Filip Kraucher}}
\def\MZeman{\autor{Mirek Zeman}}
\def\DSvetlik{\autor{David Světlík}}
\def\JRejlek{\autor{Jakub Rejlek}}
\def\PDolezal{\autor{Petr Doležal}}
\def\VBrozova{\autor{Veronika Brožová}}
\def\PBarak{\autor{Pavel Barák}}
\def\MRota{\autor{Martin Rota}}
\def\MrHlad{\autor{Matěj Svoboda}}
\def\VHrincar{\autor{Vladimír Hrinčár}}
\def\JIlavsky{\autor{Jan Ilavský}}
\def\DRynes{\autor{David Ryneš}}
\def\JKral{\autor{Jirka Král}}
\def\JVrobel{\autor{Jan Vrobel}}
\def\Tafrob{\autor{Tafrob}}
\def\JValenta{\autor{Jan Valenta}}
\def\PBuday{\autor{Pavol Boday}}
\def\VGersl{\autor{Vladimír Geršl}}
\def\MMastny{\autor{Martin Mastný}}
\def\PSebor{\autor{Pavel Šebor}}
\def\fans{několik návtěvníků z řad fanouk\,\dots}
\def\Memory{\autor{Petr Dragoun}}
\def\JDiblik{\autor{Jan Diblík}}
\def\FZenisek{\autor{Filip Ženíšek}}
\def\JiP{\autor{Jiří Pavlovský}}
\def\EHarton{\autor{Eugen Harton}}
\def\VPechacek{\autor{Václav Pecháček}}
\def\MLinda{\autor{Martin Linda}}
\def\MJezek{\autor{Matouš Ježek}}
\def\JPavlasek{\autor{Jindřich Pavlásek}}
\def\MKovarik{\autor{Michal Kovařík}}
\def\TKozelek{\autor{Tomáš Kozelek}}
\def\TFaina{\autor{Tomáše Faina}}
\def\JOrszulika{\autor{Jan Orszulik}}
\def\RKajan{\autor{Ruda Kajan}}
\def\MWilczak{\autor{Martin Wilczak}}
\def\DKos{\autor{David Kos}}
\def\VKoudelka{\autor{Vláďa Koudelka}}
\def\JZeleny{\autor{Jan Zelený}}
\def\DMaslovsky{\autor{Dan Maslovský}}  
\def\PRuzicka{\autor{Petr Růžička}}
\def\JValenta{\autor{Jan Valenta}}
\def\TRepta{\autor{Tibor Repta}}
\def\MZajacik{\autor{Martin Zajačík}} 
\def\MBudik{\autor{Marek Budík}}
\def\PKnut{\autor{Peter Knut}}  
\def\MPavelka{\autor{Matěj Pavelka}}
\def\JVanecek{\autor{Jan Vaněček}}
\def\OHrabec{\autor{Ondřej Hrabec}}
\def\PSpacek{\autor{Patrik Špaček}}
\def\JHorcik{\autor{Jan Horčík}}
\def\LNizky{\autor{Lukáš Nízký}}
\def\RFriedrich{\autor{Radek Friedrich}}
\def\VVanek{\autor{Vojtěch Vaněk}}
\def\MBerlinger{\autor{Michal Berlinger}}
\def\PNagy{\autor{Peter Nagy}}
\def\GSuterova{\autor{Gabriela Šuterová}}
\def\MDuris{\autor{Marek Ďuriš}}
\def\FDusek{\autor{Filip Dušek}}
\def\PHajda{\autor{Patrik Hajda}}
\def\JHamernik{\autor{Jakub Hamerník}}
\def\Luisa{\autor{Šárka Tmějová}}
\def\MKlima{\autor{Martin Klíma}}
\def\Alladjex{\autor{Aleš Matas}}
\def\JZubec{\autor{Jindřich Zubec}}
\def\JSvak{\autor{Jiří Svák}}
\def\DVavra{\autor{David Vávra}}
\def\SKilianova{\autor{Jana Kilianová}}
\def\Shial{\autor{Gabriela Pešková}}
\def\AKlas{\autor{Antonín Klas}}
\def\TBachtik{\autor{Tomáš Bachtík}}
\def\LStojdl{\autor{Ladislav Štojdl}}
\def\HBendova{\autor{Helena Bendová}}
\def\FBayer{\autor{František Bayer}}
\def\TBarvik{\autor{Tomáš Barvík}}
\def\LHroch{\autor{Luděk Hroch}}
\def\JHajicek{\autor{Josef Hájíček}}
\def\JVasica{\autor{Jiří Vašica}}
\def\TSpousta{\autor{Tadeáš Spousta}}
\def\PMakal{\autor{Pavel Makal}}
\def\ZPovolny{\autor{Zdeněk Povolný}}
\def\JGemrot{\autor{Jakub Gemrot}}
\def\Tobi{\autor{Tobias Stolz-Zwilling}}
\def\JRucker{\autor{Jan Rücker}}
\def\RJurda{\autor{Radim Jurda}}
\def\JChlup{\autor{Jan Chlup}}
\def\PBares{\autor{Pavel Bareš}}
\def\FDufka{\autor{Filip Dufka}}
\def\OKlus{\autor{Ondřej Klus}}
\def\FDufek{\autor{Filip Dufek}}
\def\MBerlinger{\autor{Michal Berlinger}}
\def\JSlavik{\autor{Jan Slavík}}
\def\LStoidl{\autor{Ladislav Štoidl}}
\def\TomBarsweik{\autor{Tomáš Barvík}}
\def\TomBarsweik{\autor{Tomáš Barvík}}
\def\FBeyer{\autor{František Bayer}}
\def\ZyPo{\autor{Zbyněk Povolný}}
\def\JRuecker{\autor{Jan Rücker}}
\def\PSvoboda{\autor{Petr Svoboda}}
\def\MVeselovsky{\autor{Martin Veselovský}}
\def\CStrakaty{\autor{Čestmír Strakatý}}
\def\JMalcharek{\autor{Jakub Malchárek}}
\def\NAdachi{\autor{Naomi Adachi}}
\def\TTousek{\autor{Tomáš Toušek}}
\def\JSzanto{\autor{Jakub Szántó}}
\def\OBittner{\autor{Ondřej Bittner}}
\def\MHolan{\autor{Michal Holán}}
\def\MTucan{\autor{Marek Tučan}}
\def\JVana{\autor{Jiří Váňa}}
\def\VietTran{\autor{Viet Tran}}
%\def\{\autor{}}


\def\ghostofBlixa{duch \autor{Martina Bacha}}


\begin{document}



\section*{Fight Club}

{\color{white}
Fight club je podcast a vidcast webového herního serveru Games.cz.
}

\hspace{0pt}  % very 1st minipage more to the left problem hack
\FC{fight-club-001.mp3}{\FCtitle{\#1}{}}{2010-11-16}{\Sleepless, \Janchor, \Blixa, \Farid}{\FCrule{\#1}{Farid je vždycky první.}{\Janchor}}{(0:00) Klub rváčů se představuje | (5:30) Pro koho je Kinect a~bude vás šmírovat? {\em (G4TV dali na youtube video s pornoherečkou testující Kinect (project Natal))} | (15:30) S.T.A.L.K.E.R. v~televizi | (20:30) Komu vadí bazarový prodej | (26:00) Hudební hry na ústupu | (38:00) Co je férové vyčítat CoD: Black Ops? | (50:00) Zapojte se do našeho vysílání! | Soutěž o~PC verzi Arcania: Gothic~4}{0:58:14}{19240621}{1}
\FC{fight-club-002.mp3}{\FCtitle{\#2}{}}{2010-11-23}{\Sleepless, \Janchor, \Blixa, \Hellboy, \Farid}{\FCrule{\#2}{Od sedmičky dolu je to všechno podprůměr.}{\Janchor}}{Games.cz bobtná novými funkcemi | (4:00) Druhý Zaklínač bez DRM, ale s~výhružkou pro piráty? | Warez na konzolích | (18:15) Napjatá politická realita šlape na paty Homefrontu | Řím v~podání Assassin's Creed Brotherhood | Jak přistupovat k~hodnocení her | Soutěž o~Grand Theft Auto IV Complete pro Xbox 360.}{0:55:50}{17787675}{}
\FC{fight-club-003.mp3}{\FCtitle{\#3}{Posezení s klobásou}}{2010-12-01}{\Sleepless, \Janchor, \Hellboy}{\FCrule{\#3}{No comment.}{\Janchor}}{Sníh leží na pláních | Christopher Nolan chystá virtuální Inception, bude to herní experiment? | Nebudeme se bavit o~Minecraftu | Omyl Epic Mickeyho | Konzole jako brzda, nebo plyn herního průmyslu? | Milostný dopis digitální distribuci | Co nedělat na hřbitově v~Albionu | Soutěž o~Guitar Hero: Warriors of Rock | Čtenářské recenze na spadnutí}{0:55:24}{17981401}{}
\FC{fight-club-004.mp3}{\FCtitle{\#4}{Stávka}}{2010-12-08}{\Sleepless, \Janchor, \Blixa, \Hellboy}{\FCrule{\#4}{Jestli pošlete do tý soutěže něco, tak si Vás najdu!}{\Hellboy}}{Stávkokazi | Azeroth schvátilo kataklyzma | Lara zase upatlaná | Zpráva o~stavu FPS | Budoucnost her je v~prohlížeči? | Jak na videorecenze | Soutěž o~Gran Turismo 5 v~plechu}{0:46:42}{14908797}{}
\FC{fight-club-005.mp3}{\FCtitle{\#5}{}}{2010-12-16}{\Sleepless, \Janchor, \Blixa, \Hellboy}{\FCrule{\#5}{Pššš\dots}{\Janchor}}{Zimní sporty | Vylepšené komentáře | Pozvánka na automaty do IMAXu | Úspěch Minecraftu přes kopírák | Singleplayer? Ne, multiplayer je mrtvý! | Jak to bylo s~Oblivionem | Humbuk kolem oznámení na VGA Awards | Chybí hrám kvalitní střední proud? | Soutěž o~James Bond: Blood Stone pro Xbox 360}{0:41:15}{16168588}{}
\FC{fight-club-006.mp3}{\FCtitle{\#6}{}}{2010-12-23}{\Sleepless, \Janchor, \Blixa, \Hellboy}{\FCrule{\#6}{Kdo se postí, tak vidí na Štědrej večer Poláčka.}{\Janchor}}{Pan Tau versus Arnold Schwarzenegger | Jsou poměry v~herním vývoji neúnosné? A~v čem přesně? | Další kolo války Activision - EA | Jára Cimrman české herní scény | Tuzemská radiozita v~L.A. Noire | Soutěž o~červené Nintendo Wii od ConQuestu}{1:06:52}{25960084}{}
\FC{fight-club-007.mp3}{\FCtitle{\#7}{}}{2010-12-30}{\Sleepless, \Janchor, \Blixa}{\FCrule{\#7}{Jeden za všechny -- všichni za jednoho.}{\Janchor}}{Dárky, co pohladí na duši | Ohlédnutí za uplynulým rokem: překvapení i~zklamání, krize pomohla nezávislé scéně | Deformují sváteční slevové akce trh? | Soutěž o~Nintendo Wii prodloužena do příštího dílu! | Předsevzetí Games.cz}{0:48:54}{18754396}{}
\FC{fight-club-008.mp3}{\FCtitle{\#8}{}}{2011-01-06}{\Sleepless, \Janchor, \Blixa, \Hellboy}{\FCrule{\#8}{Kdo neřekne denně stokrát Jakobi, jako by nebyl.}{\Hellboy}}{Do vesmíru s~Precursors | Couvá Ubisoft z~drakonické DRM? | Zpráva o~stavu pirátství | Za hry levnější | Kam patří Mass Effect 2 | Klíč k~nelineárnímu obsahu | Soutěž o~kombo tričko + kšiltovka Playstation}{0:57:34}{23626960}{}
\FC{fight-club-009.mp3}{\FCtitle{\#9}{}}{2011-01-13}{\Sleepless, \Janchor, \Blixa, \Hellboy}{\FCrule{\#9}{Na regionální zámek je třeba paklíč.}{\Hellboy}}{LOTRO ztrojnásobilo zisky | EA věří digitálním prodejům | Regionální biče na zákazníky | Jak si stojí Gray Matter mezi adventurami a~jak si stojí adventury vůbec? | Stop grafomanii | Vysněné herní projekty | Soutěž o~Apache: Air Assault pro Xbox 360}{1:07:48}{27442420}{}
\FC{fight-club-010.mp3}{\FCtitle{\#10}{}}{2011-01-20}{\Janchor, \Blixa, \Ulrik}{\FCrule{\#10}{Když po přečtení článku nemáte tušení o~čem pojednává, přečtěte si ho znovu.}{\Janchor}}{Přivítání nového rváče | Cannon Fodder 3 made in CCCP | Hráč vydírá Runes of Magic | Proč se vrátit k~Arx Fatalis | Živý vstup z~představení Nintenda 3DS v~Amsterdamu | Tři pohledy na tři rozměry | Soutěž o~mikinu Assassin's Creed: Brotherhood}{1:02:04}{24141688}{}
\FC{fight-club-011.mp3}{\FCtitle{\#11}{}}{2011-01-27}{\Sleepless, \Janchor, \Ulrik}{\FCrule{\#11}{Brno je zlatá loď, za děvčaty z~Brna choď.}{\Janchor}}{Ikonky jako blýskání na časy | Fanfilmy, co stojí za hřích | Bude Duke Nukem Forever něco víc, než "jenom" solidní střílečka? | Outsider u~ledu | V~předvečer představení "PSP2" | Tahanice kolem práv na Fallout | Pravda o~Excaliburu | Soutěž o~Battlefield: Bad Company 2 pro PC}{1:01:32}{24658144}{}
\FC{fight-club-012.mp3}{\FCtitle{\#12}{}}{2011-02-03}{\Janchor, \Blixa, \Hellboy}{\FCrule{\#12}{I když jste herní vývojář, tak peníze pořád nejsou všechno.}{\Janchor}}{Bude novému PSP stačit pokročilá grafika? | Move na PC | Blizzard si skrze PayPal došlápnul na virtuální farmaření | Hodnocení nezávislých titulů v~kontextu herního mainstreamu | Spoiler z~finále druhé Mafie [41:40 – 43:30] | Herní novinařina jako kariéra | Dost bylo frikulínských avatarů! | Móda samoúčelného násilí | Soutěž o~českou verzi Two Worlds II od Hypermaxu}{1:12:45}{27186712}{}
\FC{fight-club-013.mp3}{\FCtitle{\#13}{}}{2011-02-10}{\Sleepless, \Janchor, \Ulrik, \Hellboy}{\FCrule{\#13}{S poctivostí a~obtížností nastavenou na easy nejdál dojdeš.}{\Janchor}}{Pozvánka na pražský sraz fanoušků Bioware | Trable s~Magickou | Povinnost volá! | Základní "práva" PC hráče | Invertovaní vyšinutci | Kdo jsou praví rytíři Jedi? | Dumáme nad smyslem recenze | Soutěž o~Epic Mickeyho na Wii}{1:10:11}{28914928}{}
\FC{fight-club-014.mp3}{\FCtitle{\#14}{}}{2011-02-24}{\Sleepless, \Janchor, \Blixa, \Farid, \Hellboy}{\FCrule{\#14}{I blbí kamarádi jsou pořád kamarádi.}{\Janchor}}{Todle \&{} todle na Games | Únik nehotové Crysis 2: jak moc průšvih a~jak moc pokrytectví | Sociální hry jsou zlo, co multiplayer? | Steam na vrcholu | Šnečí vývoj her | Soutěž o~mládeži nepřístupný dárkový balíček od Duke Nukema}{1:03:19}{25502128}{}
\FC{fight-club-015.mp3}{\FCtitle{\#15}{}}{2011-02-24}{\Sleepless, \Janchor, \Blixa, \Hellboy}{\FCrule{\#15}{Shepardova blůza dává charisma +10.}{\Janchor}}{Jak dopadl sraz bionerdů | Stal se z~Dragon Age klon Devil May Cry? | S~Jakubem Dvorským o~Machinariu pro PS3 a~dalších projektech Amanity | Potřebujeme revoluci v~číselném hodnocení? | Soutěž o~Machinarium \&{} Samorost 2}{0:52:19}{19693420}{}
\FC{fight-club-016.mp3}{\FCtitle{\#16}{}}{2011-03-03}{\Sleepless, \Janchor, \Blixa}{\FCrule{\#16}{320 na 200 by mělo stačit každému.}{\Janchor}}{Battlefield 3 vypadá k~světu | Co se děje na Game Developers Conference a~kam se celá akce vyvíjí | Zpráva o~stavu PC trhu | Gemini Rue ukazuje záda nejen adventurám | Soutěž o~výroční kolekci Super Mario Bros. pro Wii}{0:57:17}{22108948}{}
\FC{fight-club-017.mp3}{\FCtitle{\#17}{}}{2011-03-10}{\Janchor, \Blixa, \Hellboy}{\FCrule{\#17}{I zabijáci nevinných lidí mají své dny.}{\Janchor}}{Sony vs Geohot: právo, nebo prasárna? | Newellova výzva hackerům a~Steam na televizi | Zevrubně o~silných a~slabých stránkách Bulletstormu | Film podle Mafie? | Soutěž o~2100 Microsoft points pro službu Xbox Live}{0:54:30}{21643756}{}
\FC{fight-club-018.mp3}{\FCtitle{\#18}{}}{2011-03-17}{\Sleepless, \Janchor, \Blixa, \Ulrik, \Farid}{\FCrule{\#18}{Hippísácká láska nevede vždy k~optimálním výsledkům.}{\Janchor}}{Hurá do galerie! | Pět mudrců dlouze a~obšírně mudruje na téma Dragon Age II: o~kvalitách příběhu, zajímavosti postav, soubojovém systému i~celkovém zpracování – to vše v~kontextu ohlasů z~vašich řad | Soutěž o~Motorstorm Apocalypse v~neprodejné plechové krabici}{0:56:26}{23211808}{}
\FC{fight-club-019.mp3}{\FCtitle{\#19}{}}{2011-03-24}{\Sleepless, \Janchor, \Blixa, \Hellboy}{\FCrule{\#19}{200\,$\mu$Sv ještě nikoho nezabilo.}{\Janchor}}{Dojmy z~Game Expo v~Bratislavě | Co dělá Homefront líp než Call of Duty | Svítící žebříky debilizací Deus Ex? | Chlapci fandí elektronickým čtečkám | Soutěž o~knihu Assassin's Creed: Bratrstvo}{0:43:33}{17438092}{}
\FC{fight-club-020.mp3}{\FCtitle{\#20}{}}{2011-03-31}{\Sleepless, \Janchor, \Hellboy}{\FCrule{\#20}{Neber úplatky, neber úplatky nebo se z~toho zblázníš.}{\Janchor}}{Poslední škytnutí Duke Nukema | Crysis 2: úplatky, konzole, piráti, nulová asertivita\,\dots a~vůbec! | Kádrovák Metacritic.com | Danova Apokalypsa | Soutěž o~fešné tričko Killzone 3}{0:51:55}{21380128}{}
\FC{fight-club-021.mp3}{\FCtitle{\#21}{}}{2011-04-07}{\Sleepless, \Janchor, \Blixa, \Hellboy}{\FCrule{\#21}{Hardwarový klíč vynález je píč.}{\Hellboy}}{Dan Vávra je open source | Kolik se čeho prodává, kolik (a kde) se krade | Bioware se třesou nohy | Tvrdý chléb herních žurnalistů | Kolik hráčů je potřeba k~vyšroubování onlinovky | Roleplay v~singleplayeru | Hudební vstup | Soutěž o~kombo Samorost 2 + Machinarium}{0:42:51}{16884088}{}
\FC{fight-club-022.mp3}{\FCtitle{\#22}{}}{2011-04-15}{\Sleepless, \Janchor, \Blixa, \Ulrik}{\FCrule{\#22}{Tak dlouho se chodí s~dotazy do podcastu, až se Vávra utrhne.}{\Janchor}}{Jednovaječná dvojčata hrají Zaklínače 2 | Petrovy dojmy z~třetího Battlefieldu | Hlavouni si pouštějí hubu na špacír: 3D jedině s~brýlemi, Homefront je umění? | Realistický respawn nepřátel | Soutěž o~GOTY edici LittleBigPlanet pro PS3}{1:02:24}{24034408}{}
\FC{fight-club-023.mp3}{\FCtitle{\#23}{}}{2011-04-21}{\Sleepless, \Janchor, \Blixa, \Hellboy}{\FCrule{\#23}{Dělejte portály, ne válku.}{\Janchor}}{Výlet za klasickými automaty | Portal 2? Nic nečíst, nic si nepouštět - hrát! | Proč nás najednou zajímá dvojka Prey? | Omezená morálka ve hrách | Co čekat od nástupce konzole Wii | Kdo by četl "český Edge"? | Úskalí nových forem ovládání | Neberte úplatky, nebo se z~toho zblázníte | Konkurenční boj | Pryč s~růžovými brýlemi nostalgie! | Soutěž o~Operation Flashpoint: Red River pro PC}{1:19:07}{31286804}{}
\FC{fight-club-024.mp3}{\FCtitle{\#24}{}}{2011-04-28}{\Sleepless, \Janchor, \Blixa}{\FCrule{\#24}{Nemít prachy na účtu, který vám někdo nahackuje, nevadí!}{\Janchor}}{Pac-Man požírá redakci | Chystáme kufry do Los Angeles | Česká Coraabia láká na karty a~výtvarníky | Hack PlayStation Network: fakta, spekulace, důsledky | Anonymní koření diskuzí pod články | Soutěž o~luxusní červenou brašnu Super Mario}{0:52:02}{19935052}{}
\FC{fight-club-025.mp3}{\FCtitle{\#25}{}}{2011-05-06}{\Sleepless, \Janchor, \Ulrik}{\FCrule{\#25}{Nejdřív cappuccino, pak rudé víno, potom Total War.}{\Janchor}}{Čeho je moc, toho je příliš | Kauza Sony pokračuje | S~kým a~kde budeme pobíhat v~Assassin's Creed: Revelations? | Příliš mnoho PR zabíjí informace | Role originality v~hodnocení hry | Revoluce ve sklenici vody | Soutěž o~tričko \&{} epesní press kit Mortal Kombat}{0:45:26}{17234764}{}
\FC{fight-club-026.mp3}{\FCtitle{\#26}{}}{2011-05-12}{\Sleepless, \Janchor, \Blixa, \Moolker}{\FCrule{\#26}{Když mám hlad, utrhnu si malej kousek.}{\Janchor}}{Rozjíždíme uživatelské profily | Jágr? | Blýská se na nový Xbox | Za nového Hitmana serióznějšího | Co šlo vloni na dračku | Jágr! | Jak zuříme u~her | Mario Smash Football vs Mashed | Jágr!!! | Soutěž o~Arma II: Posily}{0:56:48}{22711012}{}
\FC{fight-club-027.mp3}{\FCtitle{\#27}{}}{2011-05-19}{\Sleepless, \Janchor, \Blixa, \Ulrik, \Farid}{\FCrule{\#27}{Nejlepší konzole je YouTube.}{\Janchor}}{Kam se poděl Dan | Čekání na Geralta | Quo vadis, Kinect? | Argumentace s~piráty: předem prohraná bitva | Nejasná hranice mezi výzvou a~frustrací | Je provázanost herního webu a~obchodu problém? | Soutěž o~Zaklínače 2}{0:52:05}{20142664}{}
\FC{fight-club-028.mp3}{\FCtitle{\#28}{}}{2011-05-26}{\Sleepless, \Janchor, \Blixa, \Hellboy, \Tukan}{\FCrule{\#28}{Pokaždý, když si zahrajete pirátskou verzi Zaklínače, anektuje Vás Dan Vávra.}{\Janchor}}{Představuje se nový rváč | Nový rváč představuje tuzemskou hru Mageo | Zevrubně a~necenzurovaně (ale s~minimem spoilerů) o~druhém Zaklínači: mechanismy, svět, logika, obtížnost | A~co na to Bioware? | Soutěž o~promo balíček (a potenciálně i~hru) inFAMOUS 2}{1:14:54}{29604472}{}
\FC{fight-club-029.mp3}{\FCtitle{\#29}{}}{2011-06-02}{\Sleepless, \Janchor, \Blixa, \Tukan, \Ulrik}{\FCrule{\#29}{Když nevíš, tak 7 nebo 8.}{\Janchor}}{Další týden, další nováček | Herní časopis pro celou rodinu! | Co si slibujeme od Game Fanatics | Zveme do Bio Oko na první díl akce Press Start | Balíme do Los Angeles | Carmageddon je zpátky, ale pro jakou cílovku? | Jak to mají fotbaly a~hokeje na PC | Geralt pro prezidenta | Nadsazená hodnocení zavedených značek a~naše vlastní recenzentské prohřešky | Soutěž o~Dirt 3 pro PC}{0:58:52}{23659664}{}
\FC{fight-club-030.mp3}{\FCtitle{\#30}{}}{2011-06-09}{\Blixa, \Hellboy, \Ulrik, \PBulir, \Farid}{}{E3 2011 speciál od vašich skřítků | Legenda herní scény | Proč nejsme na E3 | Wiiijůůů je zlo | Pohyb je zlo | E3 je zlo | Co na to Tintin? | Není Joker jako Joker | Soutěž o~InFamous 2}{0:54:15}{21093136}{}
\FC{fight-club-030-a.mp3}{\FCtitle{\#30 alternace}{}}{2011-06-10}{\Sleepless, \Janchor, \Tukan, \Bety}{\FCrule{\#30}{Týden v~L.A. je málo, ale tři dny výstavy jsou až až.}{\Janchor}}{Garážový přenos z~Los Angeles a~naše bezprostřední dojmy hodinu po ukončení E3: Nintendo, béčková Playstation, ale hlavně hry od X po P | Polepšil se Xcom? | Koho nadchla dvojka Prey? | A~kdo byl z~celé výstavy jelen?}{0:32:43}{13709320}{}
\FC{fight-club-031.mp3}{\FCtitle{\#31}{}}{2011-06-16}{\Sleepless, \Janchor, \Blixa, \Hellboy}{\FCrule{\#31}{Nejdůležitější je gameplay.}{\Hellboy}}{Ohlédnutí za E3: vývoj jednotlivých ročníků, srovnání s~Gamescomem, jaké lze vypozorovat trendy | Odkud se bere a~kde se zastaví vlna hackingu? | Duke Nukem Forever\,\dots hurá? | Myslí to vývojáři upřímně? | Japonská studia v~krizi | Soutěž o~art book k~druhému Zaklínači}{1:17:55}{30868864}{}
\FC{fight-club-032.mp3}{\FCtitle{\#32}{}}{2011-06-22}{\Sleepless, \Janchor, \Hellboy}{\FCrule{\#32}{Marťani jsou.}{\Janchor}}{Last call na Press Start | Na obzoru další model PS3 | Design virtuálních služeb na konzolích | Tablety jako plnohodnotná herní platforma? | Kopa dotazů: na mimozemšťany, Duka i~online hry, která nás (ne)baví | Soutěž o~Men of War - Assault Squad}{0:47:57}{19044844}{}
\FC{fight-club-033.mp3}{\FCtitle{\#33}{}}{2011-07-01}{\Sleepless, \Janchor, \Blixa, \FDoksansky}{\FCrule{\#33}{Nosferatu is dead.}{\Janchor}}{Jak bylo v~kině? | S~hlavním programátorem o~vývoji, směřování a~úskalích Carrier Command | EVE Online postihla krize | Devalvace Hvězdných válek | Lovíme z~paměti nešťastné hry | Legendární Legenda | Upírská reminiscence | Soutěž o~dvojici augmentací ruky od Sarif Industries}{1:13:47}{29347108}{}
\FC{fight-club-034.mp3}{\FCtitle{\#34}{}}{2011-07-08}{\Sleepless, \Janchor, \Blixa, \JVitik}{\FCrule{\#34}{Každý chce mít Respekt.}{\Janchor}}{Redakce na digitálních nákupech | Skyrim říká modifikacím ANO, Battlefield 3 NE | Kde se kousnul vývoj umělé inteligence ve hrách? | Soutěž o~F.E.A.R. 3 pro PS3 \&{} tričko Duke Nukem Forever}{0:52:10}{20471668}{}
\FC{fight-club-035.mp3}{\FCtitle{\#35}{}}{2011-07-14}{\Sleepless, \Janchor, \Blixa, \Hellboy, \KMatejka}{\FCrule{\#35}{Nejdůležitější je barva fontů.}{\Janchor}}{Pod pokličku vývoje Memento Mori 2 | Zpráva o~stavu adventurního žánru a~možnostech nových platforem a~trhů | Hry z~Xboxu pod Windows 8: sen nebo logický obchodní krok? | Role vydavatele a~reklamy | Kdo zařídil, aby EA koupila PopCap | Batman! | Soutěž o~poslední díl Harry Pottera pro PC}{1:14:49}{28590820}{}
\FC{fight-club-036.mp3}{\FCtitle{\#36}{}}{2011-07-21}{\Sleepless, \Janchor, \Blixa, \Ulrik}{\FCrule{\#36}{Nejlepší super kapela je Lady Gaga.}{\Janchor}}{Sweatshop baví i~apeluje na svědomí západu | Proč se těšit na Dishonored a~má taková hra šanci na masový úspěch? | Carmack brání zavedené žánry | Najít si "svého" recenzenta | Soutěž o~tričko a~nábojnici Call of Juarez}{1:03:36}{25946656}{}
\FC{fight-club-037.mp3}{\FCtitle{\#37}{}}{2011-07-28}{\Sleepless, \Janchor, \Hellboy}{\FCrule{\#37}{Největší samohana je Rihanna.}{\Janchor}}{Now in stereo (where available) | Dan představuje vizi nového studia Warhorse | Komu nestačí milion? | Méně modelek, více Ripleyových! | Rozmazlená vývojářská omladina | Soutěž o~tričko Call of Juarez: The Cartel (a hru na Xbox 360, ale pst\,\dots)}{0:42:02}{26300968}{}
\FC{fight-club-038.mp3}{\FCtitle{\#38}{}}{2011-08-04}{\Sleepless, \Janchor, \Blixa}{\FCrule{\#38}{Prasata neodejdou.}{\Janchor}}{facebook.com/fcpod | Metacritic dobyt! (a ani jsme nemuseli sníst psy) | Je kontroverze kolem třetího Diabla odůvodněná? A~co když bude hra samotná tak trochu průšvih? | Máme radost z~potvrzení Borderlands 2, ale chceme změny | Nakupujeme v~našem obchodě | Stalkujeme naše fanoušky | Soutěž o~dvojici zeldovských figurek}{0:52:21}{33615916}{}
\FC{fight-club-039.mp3}{\FCtitle{\#39}{}}{2011-08-11}{\Sleepless, \Janchor, \ghostofBlixa, \JVitik}{\FCrule{\#39}{Velký Poláček\dots\,odchází čůrat, ale hlavně Vás sleduje.}{\Janchor}}{Potkáme se u~Kolína | Kde je pravda v~otázce licence na českou fotbalovou ligu? | Vášnivé vzpomínky na Alpha Centauri a~přání plnohodnotného | Motivace k~práci v~tuzemské herní žurnalistice | Soutěž o~Virtua Tennis 4 pro PC}{0:49:18}{32748622}{}
\FC{fight-club-040.mp3}{\FCtitle{\#40}{}}{2011-08-19}{\Sleepless, \Janchor, \Farid, \Bety}{\FCrule{\#40}{EA tvrdí Muzyku.}{\Janchor}}{Hlásí se zmlsaná klubovna Kolín nad Rýnem: dvě tváře Gamescomu, basa prudí Muzyku, přehnaná očekávání | Které hry si získaly naše znavená srdce? | Jak Harvey s~Raphem vařili dort | Držíme palce odvážným MMO | Další program\dots}{0:37:22}{25583492}{}
\FC{fight-club-041.mp3}{\FCtitle{\#41}{}}{2011-08-25}{\Sleepless, \Janchor, \Blixa}{\FCrule{\#41}{Všechno, co řekneme, bude použito proti nám.}{\Janchor}}{Mono, ale živě! | Kolínské reflexe | Tanečky kolem From Dust | Odkaz prvního Deus Ex | Další roviny Human Revolution a~proč hra dostala absolutní hodnocení | Ema mele memy | Soutěž o~GamesCombo (tričko, USB klíčenky, karty do Age of Empires\dots)}{1:29:53}{37669981}{}
\FC{fight-club-042.mp3}{\FCtitle{\#42}{}}{2011-09-01}{\Sleepless, \Janchor, \Blixa, \Hellboy, \JVitik}{\FCrule{\#42}{Šest krát devět je čtyřicetdva.}{\Janchor}}{Struktura živého vysílání | Om nom nominace | Virtuální morálka | Rvačka na téma byznys Activision vs EA | Hýbou recenze s~prodejností her? | Elder Scrolls, nebo Gothic?????!!!!! | Soutěž o~Resistance 3 pro PS3 + plastové vojáčky}{0:54:31}{22171909}{}
\FC{fight-club-043.mp3}{\FCtitle{\#43}{}}{2011-09-08}{\Sleepless, \Janchor, \Blixa, \JVitik}{\FCrule{\#43}{Naše staré hodiny bijí čtyři hodiny.}{\Janchor}}{Mobilní verze Games | Průšvihy kolem vydání Dead Island i~nového Drivera | Hry s~cíleně kontroverzním námětem | Korporátní trable a~propouštění ve velkém | Bývali herní novináři v~minulosti přísnější? | Soutěž o~Bodycount pro Xbox 360}{0:43:38}{17664694}{}
\FC{fight-club-044.mp3}{\FCtitle{\#44}{}}{2011-09-15}{\Sleepless, \Janchor, \Blixa, \LMacura}{\FCrule{\#44}{I herní redaktor může být gólmanem.}{\Janchor}}{Hloubkový průzkum tuzemských herních vod | Sčítání pirátů | Syndicate jako FPS - problem? | Lukáš Macura představuje hru Gyro 13 | Jak drahý vývoj se vyplatí | Jsou "ti zlí" vždycky vydavatelé? | Fantasy liga s~Games | Herní umění, část sto padesátá | SKANDÁLNÍ historka z~TEMNÉ historie KONKURENČNÍHO redaktora!!! | Soutěž o~Sims Medieval: Pirates \&{} Nobles}{1:18:40}{31417610}{}
\FC{fight-club-045.mp3}{\FCtitle{\#45}{}}{2011-09-22}{\Sleepless, \Janchor, \Tukan}{\FCrule{\#45}{Tukana ve snu viděti - Tukana v~podcastu mýti.}{\Janchor}}{Tukan vábí na Deus Ex speciál | Mrtvo na Dead Island | Outsourcing bossů | Pravý důvod, proč THQ svěřili pokračování Homefront tvůrcům Crysis | Originální forma distribuce? Torrenty! | Hry pomáhají lékařům i~státní správě | Herní hudba inspiruje | Dojemné konce | Mobilní kútik | Soutěž o~F1 2011 pro PC}{1:05:32}{25517858}{}
\FC{fight-club-046.mp3}{\FCtitle{\#46}{}}{2011-09-30}{\Sleepless, \Janchor, \Blixa}{\FCrule{\#46}{Payne's not deed.}{\Janchor}}{Mobilní verze Games | Nakupujte od nás, podpoříte web | Bethesda u~soudů | Max Payne \&{} my | Živé reakce na trailer k~novému Syndicate | FIFA: změny k~lepšímu | Kádrujeme Half-Life | Je Postal herní punk? | Soutěž o~Driver: San Francisco a~přidružený balík výživných suvenýrů (tričko, klíčenka, samolepka\dots)}{1:28:59}{36169942}{}
\FC{fight-club-047.mp3}{\FCtitle{\#47}{Tentokrát o mobilních Games}}{2011-10-06}{\Sleepless, \Janchor, \APlechaty}{\FCrule{\#47}{Beta není demo!}{\Janchor}}{Zveme na Gameshow 2011 | Rage zlobí PC hráče | Rozčarování z~bety nového Battlefieldu | Unreal engine v~browseru a~budoucnost herních platforem | Superhrdinský kútik | Soutěž o~sběratelskou edici Cliffs of Dover}{1:06:06}{26172013}{}
\FC{fight-club-048.mp3}{\FCtitle{\#48}{O RAGE, první Mafii a~Stevu Jobsovi}}{2011-10-13}{\Sleepless, \Janchor, \Blixa, \Hellboy}{\FCrule{\#48}{Meteority nepotřebujou důvod.}{\Janchor}}{Rage, Rage, Rage! Zuříme, bavíme se, pátráme po příběhu a~zábavě v~tunelu | Multiplayer ve třetím Mass Effectu: nutné zlo, nebo dobrovolné dobro? | Trable s~progamingem | Skutečný odkaz Steve Jobse | Skutečný odkaz první Mafie | Soutěž o~F1 2011 (k dispozici všechny platformy, v~odpovědi si vyberte tu vaši) + tričko Call of Juarez}{1:39:59}{40615236}{}
\FC{fight-club-049.mp3}{\FCtitle{\#49}{O Batman: Arkham City}}{2011-10-20}{\Janchor, \Blixa, \Bety}{\FCrule{\#49}{Šusťáky neurčují charakter.}{\Janchor}}{Ticho před BlizzConem | Take On Vrtulníček | Práce a~zábava v~Batman: Arkham City | Třetí GTA slaví deset let, vzpomínáme i~na ostatní díly série | Sands of Time ftw! | Soutěž o~Red Orchestra 2}{1:21:35}{31943062}{}
\FC{fight-club-050.mp3}{\FCtitle{\#50}{O Battlefield 3 a~podezřelé ztrátě koně ve Skyrimu}}{2011-10-29}{\Sleepless, \Janchor, \Blixa, \Ulrik}{\FCrule{\#50}{Koně se také střílejí.}{\Janchor}}{Podcastová družba | Dárcovský program a~jak se bude nakládat s~vybranými penězi | Historky ze Skyrimu | Battlefield 3 pod puškohledem | Odpouštíme hrám jejich demenci? | Alternativní modely podpory Games.cz! | Soutěž o~Zaklínače 2.0 (obsahujícího i~první díl)}{1:11:37}{28926877}{}
\FC{fight-club-051.mp3}{\FCtitle{\#51}{O Uncharted 3 a~problematice pand ve WOWku}}{2011-11-03}{\Sleepless, \Janchor, \Blixa, \Farid, \Bety}{\FCrule{\#51}{Pan krab Vás všechny sleduje.}{\Janchor}}{Trable s~fps | Je třetí Uncharted "nejlepší hrou všech dob"? | Singl jako vábnička pro multiplayer | Živé ohlasy na trailer Grand Theft Auto V | Tintin lepší Nathana Drakea | Pandy vrací úder | Soutěž o~sběratelskou edici Heroes VI}{1:25:37}{34692854}{}
\FC{fight-club-052.mp3}{\FCtitle{\#52}{Spekulace s Tukanem a~o starých hrách}}{2011-11-10}{\Sleepless, \Janchor, \Blixa, \Tukan}{\FCrule{\#52}{Charismatem krysu neudoláš.}{\Janchor}}{Trable s~recenzí Modern Warfare 3 | Biometrické ovladače? Komu tím prospějete? | Kinect se chystá na PC | Co kutí Bioware | Planescape: Torment a~my | Obtížné návraty ke starým hrám | Soutěž o~Tintinova dobrodružství pro Xbox 360}{1:12:21}{30283534}{}
\FC{fight-club-053.mp3}{\FCtitle{\#53}{S Pavlem Dobrovským o TES V: Skyrim}}{2011-11-17}{\Sleepless, \Luke, \Janchor, \Blixa}{\FCrule{\#53}{Když léčivý lektvar, tak jedině Preventan.}{\Janchor}}{Pozvánka na Game Developers Session | Hack STEAMu | Nové distribuční služby v~ČR | Skyrim, Skyrim, Skyrim | Soutěž o~Red Orchestra 2 a~tričko Heroes VI}{1:42:53}{42923518}{}
\FC{fight-club-054.mp3}{\FCtitle{\#54}{O podivných nákupech a~návratu vajíčka}}{2011-11-24}{\Sleepless, \Luke, \Janchor, \Blixa}{\FCrule{\#54}{Lepší Dizzy na iPadu než pukavec k snídani.}{\Janchor}}{Koho a~na jaké platformě dnes osloví Dizzy? | Petrovy dojmy z~PS Vita | Terno v~Dark Orbit | Kladivo na spoilery | Knižní kútik | Soutěž o~mikinu a~husí brk Assassin's Creed: Revelations}{0:59:22}{24780524}{}
\FC{fight-club-055.mp3}{\FCtitle{\#55}{O cestě za děvkou orientu a~asasíny}}{2011-12-02}{\Sleepless, \Janchor, \Blixa}{\FCrule{\#55}{Není slavík jako slavík.}{\Janchor}}{Měníme hodnocení Skyrimu! | Kam se vydá nový Assassin? | Za děvkou Orientu? | Jak efektivně vyprávět herní příběh | Hlavně se nezastavovat! | Naše i~čtenářská reflexe Danova cupování Skyrimu | Neortodoxní tipy na "nejlepší strategii" | Soutěž o~kombo podložka Heroes + F1 2011 (PC) + tričko Call of Juarez | {\em 22:10 Harry Jelínek jako námět na hru}}{1:07:55}{29596914}{}
\FC{fight-club-056.mp3}{\FCtitle{\#56}{O pravidlech virtuální války}}{2011-12-08}{\Sleepless, \Janchor, \Blixa, \Farid}{\FCrule{\#56}{Rváči nosí mikrofóny a~conquéry proklatě nízko.}{\Janchor}}{Co jsme pořídili za peníze vybrané od přispěvatelů | Petr Kotvald hráčem Bastionu? | Jak se to má s~novým C\&C | Červený kříž inspiruje k~zamyšlení | Těšíme se na Legend of Grimrock | Soutěž o~kombo Tintin hra \&{} tričko + Call of Juarez: The Cartel + F1 2011}{0:48:51}{20098348}{}
\FC{fight-club-057.mp3}{\FCtitle{\#57}{O VGA a~útoku hackerů}}{2011-12-16}{\Sleepless, \Janchor, \Blixa, \LBasta}{\FCrule{\#57}{Není tank jako tank.}{\Janchor}}{Útok bratrů hackerů | Další zdroje financování Games.cz | Lze brát VGA vážně? | Handicap východoevropských her | Nový Metal Gear děsí věrné fanoušky série | Kdo zachrání druhého Stalkera, bude to Crytek? | Soutěž o~hattrick Tintin hra + tričko + CitiesXL 2012}{1:00:49}{24741880}{}
\FC{fight-club-058.mp3}{\FCtitle{\#58}{Naposledy o Skyrim}}{2011-12-23}{\Sleepless, \Janchor, \Blixa, \Hellboy}{\FCrule{\#58}{Pravda a~láska musí zvítězit nad lží a~nenávistí.}{neznámý umělec}}{Finální epická bitva na téma Skyrim: svoboda názoru, ignoranti v~Bethesdě, stromy \&{} klenby, dilemata otevřeného světa | Hry pod vánoční stromeček | Subjektivně-objektivní dekonstrukce Tintina | RIP, Galaxies | Ničitelé online světů | Jak se to má s~podplácením herních novinářů | Soutěž o~sadu LEGO Harry Potter}{1:22:52}{32890890}{}
\FC{fight-club-059.mp3}{\FCtitle{\#59}{Jak se žije a~hraje v Číně}}{2011-12-29}{\Sleepless, \Janchor, \Blixa, \Cocik}{\FCrule{\#59}{Kdo chce psa jíst, hůlky si vždy najde.}{\Cocik}}{Výlet za asijským trhem: mobilní platformy, reforma herního vývoje, co z~toho můžeme čekat v~Evropě a~kdo už se u~nás adaptoval | Free 2 Play modely, DLC\,\dots zhouba moderního hráče? | Kulturní šoky | Na co se těšíme v~příštím roce | Soutěž o~zlatou edici Sniper: Ghost Warrior a~teplákovou bundu Heroes VI}{1:00:34}{25802224}{}
\FC{fight-club-060.mp3}{\FCtitle{\#60}{Miss platí husity!}}{2012-01-06}{\Sleepless, \Luke, \Blixa, \Farid}{\FCrule{\#60}{Na Bakalu jedině s palcátem.}{\Luke}}{Jak probíhal výběr loňských Best Of | Další díl Press Start bude ve znamení IDDQD | Miss platí husity! | Dobro a~zlo v~The Old Republic | Zlo a~zlo v~SOPA | Little Big Adventure je zpátky | Soutěž o~Cities XL 2012 + tričko AC: Revelations}{1:05:56}{30510028}{}
\FC{fight-club-061.mp3}{\FCtitle{\#61}{S trojicí mušketýrů o emzácích a~nepokojích}}{2012-01-13}{\Luke, \Janchor, \Blixa}{\FCrule{\#61}{Nejlepší desková hra je Deskent.}{\Janchor}}{Ohlédnutí za anketou o~hru roku | Potenciál Xcomu od Firaxis | Vývojář čelí trestu smrti | Drama v~čínské továrně | Tuzemské RPG Nyrthos vábí na oldschool zpracování | Jak si vybíráme onlinovky (a jakou bychom si vybrali, mít neomezeně času) | Preferovaný pohled | Soutěž o~kovový přívěšek ze Skyrimu}{0:52:36}{22905080}{}
\FC{fight-club-062.mp3}{\FCtitle{\#62}{S Michalem "cokoliWolf" Rybkou}}{2012-01-19}{\Sleepless, \Luke, \Janchor, \Wolf}{\FCrule{\#62}{Myši se nikdy nesmí krmit po půnoci.}{\Luke}}{Omluvte zkraje rozháraný zvuk! | Myší problém Michala Rybky | Budoucnost tabletů jako herní platformy | První vítězství v~boji proti cenzuře internetu? | Ohlédnutí za druhým dílem akce Press Start | Jak hodnotit hry, co nejsou tak úplně hrami | Reminiscence Budokanu | Soutěž o~L mikinu Assassin's Creed: Revelations}{0:56:29}{25767004}{}
\FC{fight-club-063.mp3}{\FCtitle{\#63}{O Old Republic bez rukavic}}{2012-01-26}{\Sleepless, \Janchor, \Blixa, \Ulrik, \KDrda}{\FCrule{\#63}{Nejvýkonnější počítač je Dacia.}{\Janchor}}{Pravda o~Karlu Vojtíškovi | The Old Republic pod drobnohledem: pašeráci, otylí rytíři Jedi, důraz na příběh | Chystá se Microsoft rušit body? | Kam kráčí Resident Evil | Soutěž o~Global Ops: Commando Libya}{0:50:51}{22062058}{3}
\FC{fight-club-064.mp3}{\FCtitle{\#64}{O Half-Life 4 a~komunitních srazech}}{2012-02-02}{\Sleepless, \Luke, \Blixa, \VPravda, \JFaltus}{\FCrule{\#64}{Když gamifikovat extenze, tak jedině umělecky.}{\Luke}}{Nové koště, nové funění | Komerční sdělení: pozvánky na Gamefest a~Biocon | Zynga krade! A~co ostatní? | Síla lidu versus Valve | Zaklínač 2 2.0 | Gamifikace všude kolem nás | Sopwith a~Wings | Opět The Old Republic | Megaupload a~my | Soutěž o~Samorost, Machinarium a~F1}{1:08:43}{30948512}{}
\FC{fight-club-065.mp3}{\FCtitle{\#65}{}}{2012-02-08}{\Sleepless, \Luke, \Blixa, \Bety}{\FCrule{\#65}{Když chceš udělat hru, tak zavolej Notche.}{\Luke}}{Co nového na Games.cz | Jak hodnotíme protest hráčů Half-Life 2 | Sapkowski a~my | Alžběta a~Darkness II | Kovoví, velcí a~smrtící roboti v~Hawken | Otazníky nad zpětnou kompatibilitou PS Vita | Psychonauts 2 a~situace Double Fine | Kinect pro Windows | Uslyšíme Petra Ticháčka? | Soutěž o~Soulcalibur V~pro PS3}{0:55:49}{24701864}{}
\FC{fight-club-066.mp3}{\FCtitle{\#66}{Dejte prachy Schaferovi!}}{2012-02-16}{\Sleepless, \Luke, \Blixa, \LMacura}{\FCrule{\#66}{Není důležitý zvuk, ale informace.}{\Luke}}{Jak se díváme na PS Vita | Zevrubně o~BezOchrany.cz | Jak si vydělat jako Tim Schafer \&{} Kickstarter a~Great Satans | Jemná debata o~Mass Effect 3 | "Lukáši? Nejsi tady někde?" | Kam kráčíš Bioware? | Fenomény Rage a~pětaosmdesát procent | Soutěž o~Tekken: Hybrid na PlayStation 3}{1:09:43}{32025812}{}
\FC{fight-club-067.mp3}{\FCtitle{\#67}{Skúsime to cez vesmír!}}{2012-02-23}{\Sleepless, \Luke, \Blixa, \MRosa}{\FCrule{\#67}{Když je úplněk, pozor na upíří prasata.}{\Luke}}{Půl hodina nejen o~Miner Wars s~Markem Rosou z~Keen Software | A~zase ten Zaklínač, tentokrát milionový | Polská sci-fi do her! | Drby o~nových konzolích | Panika! Konec kamenných obchodů! | Nenápadné vývojářské celebrity | Jsou herní novináři póvl a~amatéři? | Střízlivé hodnocení diskusí na Games.cz | Soutěž o~Bodycount pro Playstation 3}{1:26:38}{42602504}{}
\FC{fight-club-068.mp3}{\FCtitle{\#68}{O mrtvých myších a~lidech}}{2012-03-01}{\Luke, \Farid, \Tlamiczka, \PKolacek, \MPach}{\FCrule{\#68}{Když chceč najít Mickey Mouse, tak zavolej Jacques Cousteauovi.}{\Luke}}{S Nikolou Bornovou (aka Tlamiczkou) o~Coral City, co dělá po nocích a~zombících | Na obranu Jennifer Helperové | Přeskakování akce v~Mass Effect 3 | Jak se stát herním scenáristou, českým i~zahraničním | Single vs multiplayer | (Ne)hraj a~div se: 39 stupňů a~Dear Esther | Tlamiczka a~borci z~Mass Effect | Bioshock Infinite 1999 a~stará hratelnost versus casual hry | Bude víc konverzí Xbox her na PC? | Soutěž o~Generator Rex pro Xbox 360}{1:23:16}{37835132}{}
\FC{fight-club-069.mp3}{\FCtitle{\#69}{Bacha na módní slova!}}{2012-03-08}{\Luke, \Blixa, \Tukan}{\FCrule{\#69}{Networkingem nic nezkazíš.}{\Luke}}{S Martinem Vaňem o~mrtvých NPCčkách, králících na traktoru a~práci na ArmA 3 | Příběhová česká RTS Age of Defenders | Staňte se podřízeným kapitána Kirka v~Artemis | Pozvánka na Indie Games CZ | O~nejprodávanějších hrách v~ČR a~SR | Copak nám dělá AHP | Chce nám Valve dát konzoli? | Ovládání her myšlenkami! | Několik slov o~Assassin's Creed 3 | Vzpomínáme a~na Petera Molyneuxe | Nejčtenější články na Games | Jak se to má s~embargy | Designéři vs uklízečky | Multiplayer v~Mass Effect 3 | Soutěž o~Kingdoms of Amalur pro Xbox 360}{1:23:53}{39470324}{}
\FC{fight-club-070.mp3}{\FCtitle{\#70}{Pivo, Fidži a~jeleni}}{2012-03-15}{\Sleepless, \Luke, \Janchor, \MJelinek}{\FCrule{\#70}{Když Fargo, tak od Cohena.}{\Janchor}}{Testujeme nový stream | S~Martinem Jelínkem zevrubně o~připravovaném tuzemském RPG Nyrthos | Hráči horují za jiný konec posledního Mass Effectu | Dalším kultovním jménem na Kickstarteru je Brian Fargo s~pokračováním Wasteland | Petr se zastává DLC | Soutěž o~výhercem vybranou starou hru z~Good Old Games}{1:18:45}{35580253}{}
\FC{fight-club-071.mp3}{\FCtitle{\#71}{Stmívání bez upírů}}{2012-03-22}{\Luke, \Janchor, \Blixa, \JHussar}{\FCrule{\#71}{Nejhorší jsou nemrtví Ewoci.}{\Janchor}}{Dobré a~špatné zprávy | Jak probíhal první ročník nové ankety Booom | Vzpomínka na LLG | Kuba Hussar představuje karetní hru Coraabia a~rozplétá její systémy i~svět | GAME se potápí, co to znamená pro tuzemský trh? | Přístup Bioware a~smysl předělávek starých her | Zombies v~adventuře? | Soutěž o~laptop brašnu Nintendo}{1:24:34}{37958873}{}
\FC{fight-club-072.mp3}{\FCtitle{\#72}{Muzejní kousek}}{2012-03-29}{\Sleepless, \Janchor, \Blixa, \JOrna}{\FCrule{\#72}{Co je na DVD-čku v Levelu, to je dobrá hra.}{\Janchor}}{Výlety za Assassinem a~novými Borderlands | Honza Orna představuje své muzeum automatů | Od sběratele k~provozovateli | Good Old Games už nejsou Old | Má cenu sbírat digitální verze her? | Osud Prey 2 visí na vlásku | Vzpomínka na LEVEL TV | Jak je to s~redakčním pokrytím online her | Soutěž o~PC verzi Alan Wake + klíčenku Arcadehry.cz}{1:42:22}{46324263}{}
\FC{fight-club-073.mp3}{\FCtitle{\#73}{Proti všem}}{2012-04-05}{\Sleepless, \Janchor, \Hellboy, \Wolf}{\FCrule{\#73}{Vyšší princip a~Justin Bieber nemůžou zůstat v opozici natrvalo.}{\Janchor}}{Chystá se herní sraz Games.cz | Velká debata kol Kickstarteru: Obsidian, Brian Fargo, Tim Schafer, Al Lowe, Justin Bieber\,\dots | Pozitivní dopad na projekty neznámých autorů? | Osa vývoje herních veteránů | Možnost "proplatit se" na vyšší level - problém, nebo prospěšný trend? | Lapkové a~katolíci | Jsou lineární hry skutečně lehčí na vývoj? | Klíč k~úspěchu na mobilních platformách | Soutěž o~Force Unleashed 2 na Xbox 360}{1:23:29}{34485581}{}
\FC{fight-club-074.mp3}{\FCtitle{\#74}{O ptácích a~jiné zvířeně}}{2012-04-12}{\Sleepless, \Luke, \Janchor, \Blixa, \JSvelch}{\FCrule{\#74}{Na ptáky jsme krátký.}{\Janchor}}{Pozvánka na dvě brněnské herní akce | Herní archeologie a~historie české herní tvorby | Vzpomínky na Svazarm i~Palachův týden | Pozice herních studií v~akademických kruzích | Jak souvisí Commodore a~morčata | Dospělé hry = dospěle se chovající hráči? | Vetřelci a~strach ve virtuálním prostoru | Bioware nakonec | Soutěž o~comics Warcraft \&{} tričko Uncharted: Golden Abyss}{1:31:04}{39271347}{}
\FC{fight-club-075.mp3}{\FCtitle{\#75}{Krok za krokem}}{2012-04-19}{\Sleepless, \Luke, \Janchor, \DKubec}{\FCrule{\#75}{Když Glass, tak Philip.}{\Janchor; pravděpodobně narážka na hudební uskupení Glass Onion}}{Co nás táhlo na starých dungeonech | Jak se s~tím popasoval Grimrock | Má Jane Jensen na to se vrátit v~plné síle? | Herní potenciál pokročilé augmentované reality | Crysis: mnoho povyku pro nic | Jak se hry potýkají s~erotikou | Proč u~nás nefrčí Nintendo | Soutěž o~Legend of Grimrock}{1:41:27}{45216209}{}
\FC{fight-club-076.mp3}{\FCtitle{\#76}{Groomování}}{2012-04-26}{\Sleepless, \Luke, \Janchor, \Farid}{\FCrule{\#76}{Kaviár předchází pád.}{\Janchor}}{Games OPEN! | Gumák slaví třicítku | Dragon's Dogma dalším signálem transformace RPG? | Jak funguje Valve | MASHED! | Soutěž o~comics World of Warcraft}{1:38:31}{44821191}{}
\FC{fight-club-077.mp3}{\FCtitle{\#77}{Cosplay a~jiné hříchy}}{2012-05-03}{\Sleepless, \Luke, \Janchor, \Blixa, \Jorssk}{\FCrule{\#77}{Šaty dělají člověka; župany rytíře Jedi.}{\Janchor}}{Poslední výzva na sobotní GAMES Open: tipy na dopravu, předpokládaný průběh | Pravda o~cosplayi u~nás | Stalker je mrtev, přežívá Survarium | Čeká sérii Call of Duty vnitřní posun? | Soutěž o~XL tričko "Mario jede na Yoshim"}{1:32:14}{40670912}{}
\FC{fight-club-078.mp3}{\FCtitle{\#78}{}}{2012-05-10}{\Sleepless, \Luke, \Janchor, \Blixa}{\FCrule{\#78}{Kdo jsi bez viny, švihni kabelem.}{\Janchor}}{Jak dopadl sraz na automatech | K~Elder Scrolls MMO zůstáváme zatím chladní | Teror v~EVE Online prospívá vývojářům | Hitman Absolution akčnější? | Největší výzva pro Beyond Good \&{} Evil | Soutěž o~špunty Panasonic}{1:25:59}{38205344}{}
\FC{fight-club-079.mp3}{\FCtitle{\#79}{Tuřín III a~další zelenina}}{2012-05-17}{\Sleepless, \Luke, \Janchor, \Blixa, \Tukan}{\FCrule{\#79}{Tuřín zásadně pod vlivem.}{\Janchor}}{Zpráva o~stavu (nejen) tabletového hraní | Na DOTA frontě zavládl mír | Maďaři fušují Rumunům do Drákuly | Secret World klepe na dveře | Nehrajeme Diablo, hrajeme tuřín! | Krvavé diamanty na DSku | Hraní pod vlivem | Soutěž o~promo verzi Sorcery pro PS Move}{1:37:10}{43301256}{}
\FC{fight-club-080.mp3}{\FCtitle{\#80}{Návštěva v pekle}}{2012-05-24}{\Luke, \Janchor, \Blixa, \KDrda}{\FCrule{\#80}{Kdo chce výzvu, hraje ruské hry.}{\Janchor}}{Diablo, Diablo, Diablo: grafika, příběh (?!), obsah, obtížnost, online DRM\,\dots | Vývojář, ten smluvně tvrdý chleba má | Koperník sem, Titan tam | Jak je na tom Guillermo del Toro | Jak se vyrovnat s~obtížnými pasážemi ve hrách | Soutěž o~MMORPG Tera + XL tričko}{1:36:38}{43975896}{}
\FC{fight-club-081.mp3}{\FCtitle{\#81}{Mlčení tučňáka}}{2012-06-01}{\Sleepless, \Luke, \Janchor}{\FCrule{\#81}{Tučňákům nenalévat.}{\Janchor}}{Žhavíme motory na E3 | Co chystá výstava, co chystáme my | Změny ve vedení THQ | Ron Gilbert se vrací v~plné síle | 38 zkrachovalých | Druhá ARMA nabírá druhý dech skrze zombies | Soutěž o~námořní dalekohled z~Risen 2}{1:21:22}{35587231}{}
\FC{fight-club-081-E3-1.mp3}{\FCtitle{E3 2012 speciál \#1}{}}{2012-06-05}{\Sleepless, \Luke, \Blixa}{}{E3, den první: tiskové konference Microsoftu, Electronic Arts, Ubisoftu, Sony | Petrův osobní problém s~Tomb Raiderem a~Splinter Cellem | Poslední z~nás | Bude Beyond vůbec hra? | Halucinace ve Far Cry 3 | Opakování matka novinářské únavy | Největší hrůzy a~co na tiskovkách nebylo}{0:33:51}{14163957}{}
\FC{fight-club-081-E3-2.mp3}{\FCtitle{E3 2012 speciál \#2}{}}{2012-06-06}{\Sleepless, \Luke, \Blixa}{}{E3, den druhý: Nintendo a~brány výstavy dokořán | Potenciál Wii U | Uhekaná Lara | Do jeskyně za Ronem Gilbertem | Piráti bez rumu?! | Jak se má Gaikai | Bude se výstava stěhovat?}{0:46:13}{18626317}{}
\FC{fight-club-081-E3-3.mp3}{\FCtitle{E3 2012 speciál \#3}{}}{2012-06-07}{\Sleepless, \Luke, \Blixa}{}{E3, den třetí: výstava v~plném proudu | Tašky a~chaos | Need for Speed: Burnout | Jak dopadl Martin v~multiplayeru Crysis 3 | Třetí Dead Space v~praxi umírněnější | Souboje v~Elder Scrolls Online | Vlastenecké okénko}{0:37:38}{15566379}{}
\FC{fight-club-081-E3-4.mp3}{\FCtitle{E3 2012 speciál \#4}{}}{2012-06-08}{\Sleepless, \Luke, \Blixa}{}{E3, den čtvrtý: zhasínáme, jdeme domů | Kdo nemá Watch Dogs, dá si tam Sleeping Dogs | Projekt "sto pé" | Napravené dojmy a~boj s~rutinou | Lepení best of samolepek | Veselé historky z~výstaviště}{0:44:03}{18791049}{}
\FC{fight-club-082.mp3}{\FCtitle{\#82}{}}{2012-06-14}{\Sleepless, \Janchor, \Hellboy, \OverWatch}{\FCrule{\#82}{Kdo jsi bez viny, hoď mrakodrapem.}{\Janchor}}{OMLUVTE OBČASNÉ PODEZŘELÉ PARANORMÁLNÍ PRASKÁNÍ | Nastavujeme zrcadlo výstavě vymknuté z~kloubů | Za kým se zdržela část losangelské výpravy? | Když na E3, tak vlakem | Kde se schovávají zajímavé hry | Watch Dogs ano, Beyond\,\dots ne? | Kterak David Cage podrazil naivní novináře | Nintendo zklamalo | Výbuchy se konečně přejedly | Noční dojezd | Soutěž o~neprodejnou encyklopedii Assassin's Creed}{1:52:46}{47922377}{}
\FC{fight-club-083.mp3}{\FCtitle{\#83}{Zvířata v zahradě}}{2012-06-21}{\Luke, \Janchor, \Blixa, \PKuba}{\FCrule{\#83}{Nejlepší skript je JavaScript.}{\Janchor}}{Jak to bylo s~Fargem | Jak se vyvíjela a~jak se má Genetická zahrada 2 | Botanicula zachraňuje tapíry | Misogynní trolling se obrátil proti trollům | Spekulativní specifikace nového Xboxu | Příliš animací, hráčova smrt? | Soutěž o~USB kartu Star Wars 1313}{1:18:33}{34429107}{}
\FC{fight-club-084.mp3}{\FCtitle{\#84}{Hudbohraní}}{2012-06-28}{\Sleepless, \Luke, \Blixa, \PPycha}{\FCrule{\#84}{Nejlepší Muzyka je Ray.}{\Luke}}{Až na dřeň o~natáčení herní hudby v~Praze s~Petrem Pýchou | Blizzard, Diablo III a~zákazníci - co ty mě, co já tobě? | Karty a~vojáčci ve Scrolls, budeme se bavit? | Notchův kult | Šílenství! Far Cry 3 posunuta! | Reapeři se nažrali a~Shepard zůstal celý | Těžko dělat soudržný a~neprůstřelný fikční vesmír | Píše nám SVětlá budoucnost! | Pes FPS | Milióny prodaných her, milióny prodaných her tam | PC konference na výstavách? | Reforma recenzí! | Soutěž o~8 GB flashku}{1:42:27}{49145920}{}
% Rossignol, Jim. This Gaming Life: Travels in Three Cities. The University of Michigan Press. ISBN 0-472-11635-5.
% Poole, Steven. Trigger Happy.
% Koster, Raph. a~Theory of Fun for Game Design.
% Ultimate History of Video Games.
\FC{fight-club-085.mp3}{\FCtitle{\#85}{GamePage Over?}}{2012-07-05}{\Sleepless, \Janchor, \Blixa, \OBroukal}{\FCrule{\#85}{Všichni jsou ve Varech.}{\Janchor}}{Léto s~automaty | Proč skončila GamePage, jak vlastně (ne)fungovala a~kam se posune dál | Vyjádření programového ředitele ČT Milana Fridricha | Postará se výnos soudního dvora EU o~revoluci v~bazarovém hraní? | Co si slibovat od rovnice Sony + Gaikai | Jak se žije finančním příspěvkům od čtenářů | Soutěž o~Steam kód hry Tiny \&{} Big}{1:34:33}{44278153}{}
\FC{fight-club-086.mp3}{\FCtitle{\#86}{Oujé!}}{2012-07-12}{\Sleepless, \Janchor, \Farid}{\FCrule{\#86}{Punk's not dead -- it's just resting.}{\Janchor}}{Propozice nové konzole zbourala Kickstarter | Jaké jsou její reálné šance? | Steamu patří zelená | Blýská se na filmy podle Assassina i~Human Revolution | Soutěž o~tričko TERA}{1:13:30}{32680275}{}
\FC{fight-club-087.mp3}{\FCtitle{\#87}{Tajný svět Karla Drdy}}{2012-07-19}{\Sleepless, \Janchor, \Blixa, \KDrda, \JOlejnik, \Laibach}{\FCrule{\#87}{Když Bach, tak Laibach.}{\Janchor}}{Připíjíme na Martina | Ouya, část druhá a~odvrácená | Rozkrýváme tajemství kvality The Secret World | Proč nehrajeme multiplayerové střílečky | Nejlepší vtip v~historii podcastu™ | Soutěž o~tričko Ghost Recon šedivé jako vlk!}{1:03:20}{29433379}{}
\FC{fight-club-088.mp3}{\FCtitle{\#88}{Nad vatrou sa blýská}}{2012-07-26}{\Sleepless, \Janchor, \Blixa, \Hellboy}{\FCrule{\#88}{Sní androidy o elektronických ovečkách?}{\Janchor}}{OMLUVTE MIZERNÝ ZVUK |  Pozvánka na arkádový turnaj | Co uhasilo Vatra Games? | Radši zadarmo, než nechat napospas pirátům | Telefonní intermezzo | Potřebují nové značky nové konzole? | Soutěž o~Ratchet \&{} Clank Trilogy pro PS3}{1:17:57}{31813809}{}
\FC{fight-club-089.mp3}{\FCtitle{\#89}{Macura ad portas!}}{2012-08-02}{\Sleepless, \Luke, \Janchor, \LMacura}{\FCrule{\#89}{Ester má ráda věci co voněj,\,\dots hlavně mýdlo s jelenem.}{\Janchor}}{Poťouchlé plány na Gamescom | Olympijský Decathlon a~jeho maraton schvalovacími procesy | Co dalšího se (ne)kutí v~Cinemaxu | Ruku do ohně za hraní | Stalker v~žoldu Bethesdy, důvod k~panice? | Přísliby remasterované verze Baldur's Gate | Kiks ubisoftího pluginu | Autoři Dear Esther se posouvají blíž k~interakci | Je u~nás prostor pro herní časopis? | Výhledy streamovacích platforem | Soutěž o~oldschool RPG Inkvizitor pro PC!}{1:30:29}{38528468}{}
\FC{fight-club-090.mp3}{\FCtitle{\#90}{Carmackgeddon}}{2012-08-10}{\Sleepless, \Luke, \Janchor, \Blixa}{\FCrule{\#90}{Za všechno může Bach.}{\Janchor}}{Žhavíme na Gamescom | Zpráva o~stavu idu | Mody jako záplaty | Kdo stojí o~virtuální realitu | Rake in Grass přichází s~karetní hrou Northmark | Soutěž o~tričko Ghost Recon a~klíčenku Driver}{1:27:45}{42066141}{}
\FC{fight-club-091-Gamescom-1.mp3}{\FCtitle{Gamescom 2012 speciál \#1}{}}{2012-08-15}{\Sleepless, \Luke, \Janchor, \Blixa}{\FCrule{\#91a}{Tady se nechodí.}{\Janchor}}{Gamescom 2012, den první | Když to nejde silou, jde to raketou | Martinova herní událost roku | Trhající se kyberpunkoví upíři | Petrovi prohráli válku inženýři | Tradiční UbiRachot | Co chystá tvůrce DayZ na příští rok? | Kyberpunkoví polští farmáři | Dabéři The Last of Us naživo}{0:26:29}{11000954}{}
\FC{fight-club-091-Gamescom-2.mp3}{\FCtitle{Gamescom 2012 speciál \#2}{}}{2012-08-16}{\Sleepless, \Luke, \Janchor, \Blixa}{\FCrule{\#91b}{Kopírování zabíjí podcast.}{\Janchor}}{Gamescom 2012, den druhý | Pět P: Petrovy Pestře Peprné PC Postřehy | Zaklínačský les levou zadní | Lukáš s~Pavlem už jsou zase ve při | Myší problém Warrena Spectora | Pomsta na Davidu Cageovi | Proč bylo Martinovi zle z~Far Cry 3 | Prázdná a~falešná výstava | Mládí vpřed? | Pachatel vykrádačky u~Larianů odhalen!}{0:32:09}{14139306}{}
\FC{fight-club-091-Gamescom-3.mp3}{\FCtitle{Gamescom 2012 speciál \#3}{}}{2012-08-17}{\Sleepless, \Luke, \Janchor}{\FCrule{\#91c}{Drda se vytře všude.}{\Janchor}}{Gamescom 2012, den třetí | Mládeži přístupno, výstava houstne | Proč tu nepotkáte menší nezávislé hry | Je Dishonored super? | Rudá orchestrální bouře na vzestupu | Rozhovory hodnotné a~rozhovory na pytel | Freemium onlinovky, kam se podíváš}{0:30:49}{13075386}{}
\FC{fight-club-092.mp3}{\FCtitle{\#92}{Ve kterém se objeví pět lidí}}{2012-08-23}{\Sleepless, \Luke, \Janchor, \Blixa, \JBeranek}{\FCrule{\#92}{Dobří holubi se vracejí.}{\Janchor}}{Novinky v~donate programu | Jaký prostor mají hry v~České televizi a~může se situace změnit? | Jak se vyvrbil prodej GAME | Zemětřesení v~Onlive | Za hry (a jejich trailery) vkusnější | Soutěž o~artbook ke druhému Zaklínači}{1:25:07}{39282325}{}
\FC{fight-club-093.mp3}{\FCtitle{\#93}{Předčasně ukončeno}}{2012-08-30}{\Sleepless, \Luke, \Janchor}{\FCrule{\#93}{Největší Mechwarrioři jsou Křemílek a~Vochomůrka.}{\Janchor}}{Games pod Lupou | Update PS Vita a~příslib nových, zajímavých titulů | Je Chris Crawford mimozemšťan? | Mechwarrior versus Hawken | WoW má utrum v~Íránu | Proč být shovívaví k~PC konverzi Dark Souls | Soutěž o~Steam kód na adventuru Deponia}{1:05:49}{31859735}{}
\FC{fight-club-094.mp3}{\FCtitle{\#94}{Snake na zelenou nečeká}}{2012-09-06}{\Luke, \Janchor, \Blixa}{\FCrule{\#94}{Snake\,\dots{} it's a~snake.}{\Janchor}}{Steamu patří zelená? | Samožerství Polygonu | Potíže s~obtížností | Těšíme se na nový Metal Gear? | Soutěž o~pivní tácky a~tričko od Larian Studios}{1:09:08}{33090571}{}
\FC{fight-club-095.mp3}{\FCtitle{\#95}{Nefiditelný Lukásh}}{2012-09-13}{\Sleepless, \Luke, \Janchor, \Bludr}{\FCrule{\#95}{Víš co hulíš?}{\Janchor}}{Rýsují se plány na stý podcast a~chystáme další speciál! | Není BIS jako BIS | Není BigPeter jako Big Picture | Chystá Microsoft předzvěst holopaluby? | Tajnosnubný comeback Chrise Robertse | Soutěž o~tričko ARMA, zápisník Carrier Command a~granát Rambo!}{1:14:27}{34466957}{}
\FC{fight-club-096.mp3}{\FCtitle{\#96}{Bude nás pět}}{2012-09-20}{\Sleepless, \Luke, \Blixa, \PProuza, \MEndrle}{\FCrule{\#96}{Dávejte si po ránu pozor na opilé velrybáře.}{\Luke}}{Jak se člověk promění ve vývojáře mobilních her | Zpráva o~stavu casual a~sociálního trhu |  Další vývoj situace zadržených českých vývojářů | Zeschuk a~Muzyka vyhořeli | WiiU nabírá konkrétní obrysy | Kdo by si koupil novou PS3? | Metro chce točit Metro! | Obsidian zbouralo Kickstarter | Soutěž o~HD kolekci Devil May Cry na PS3}{1:25:07}{37526353}{}
\FC{fight-club-097.mp3}{\FCtitle{\#97}{Karel na autopilota}}{2012-09-27}{\Sleepless, \Luke, \Janchor, \KDrda, \JOlejnik}{\FCrule{\#97}{Pro dobrý loot zásadně na zabijačku.}{\Janchor}}{Games pod Křišťálovou lupou | Lepší dramaturgii pro IGCZ | Je dvojka Guild Wars převratem v~žánru? | Virtuální fotografie v~nevirtuálních galeriích | Soutěž o~amanití kombo Samorost 2 + Machinarium}{0:59:22}{26095945}{}
\FC{fight-club-098.mp3}{\FCtitle{\#98}{Boxerka z Radiožurnálu}}{2012-10-04}{\Sleepless, \Luke, \Janchor, \Blixa, \AHavlova}{\FCrule{\#98}{Pozor na modré pipinky.}{\Janchor}}{Jak se žije hrám na vlnách rozhlasu | Matka postrachem FPS | Jubilejní stý podcast v~kině Oko | Podivná "kompilace" trilogie Mass Effect | Přesvědčí vágní projekt dvou herních legend? | Střílečka ze světa XCOMu mutuje k~horšímu | Starbreeze povolali k~nové hře režiséra komedie Policajti (Kopps) | Neznabozi of Might \&{} Magic | Soutěž o~Carrier Command pro PC nebo X360}{1:27:13}{38905812}{}
\FC{fight-club-099.mp3}{\FCtitle{\#99}{Tři už jsou skupina}}{2012-10-11}{\Sleepless, \Luke, \Janchor}{\FCrule{\#99}{Když zajíc, tak jedině s opicí.}{\Janchor}}{Další vyhořelý veterán? | Jazz Jackrabbit vs Jill of the Jungle | Genocida ve World of Warcraft | Ublížila Secret Worldu jeho jinakost? | Přešlap reklamního herce | Bad Company komediálním seriálem | Stručně k~Dishonored | Soutěž o~Resident Evil 6 pro Xbox 360}{1:03:55}{29960064}{}
\FC{fight-club-100.mp3}{\FCtitle{\#100}{O hrách s vámi}}{2012-10-22}{\Sleepless, \Luke, \Janchor, \Blixa, \Ulrik, \JOlejnik, \KDrda, \LMacura}{\FCrule{\#100}{The cake is not a~lie.}{\Janchor}}{Slavnostní přivítání | Stříháme metr do roku 2077 | Zevrubně k~Dishonored | Vývoj her jako loterie | Kultura nakupování, kultura pirátství | Soutěž o~Fallout: New Vegas a~Empire: Total War}{1:28:18}{41184326}{}
\FC{fight-club-101.mp3}{\FCtitle{\#101}{Změna je život}}{2012-10-29}{\Sleepless, \Luke, \Janchor, \PSebor}{\FCrule{\#101}{Změna je život.}{\Janchor}}{Ch-ch-ch-ch-changes! | Rozhovor s~Pavlem Šeborem o~Euro Truck Simulator 2 a~dalších aktivitích českého studia SCS  | Jedno Oko nezůstalo suché | Potřebují hry rychlejší internet? | Bioshock Infinite žije | Ohlédnutí za sérií Halo | V~Elder Scrolls Online se zablýsklo na časy | Soutěž o~Truck Driving Simulator}{1:15:36}{31938572}{}   % mising Pavel Sebor from fa CSC
\FC{fight-club-102.mp3}{\FCtitle{\#102}{Kdopak by se myši bál}}{2012-11-01}{\Luke, \Janchor, \Blixa, \OverWatch, \JOlejnik}{\FCrule{\#102}{Čtyři mikrofony nikdy nestačí.}{\Janchor}}{Disney rozpoutá nové Hvězdné války | Co je v~herní žurnalistice košer | Proč Valve nechrlí hry | Zopakuje Arakion úspěch Grimrocku? | Soutěž o~Valhalla edici King's Bounty: Warriors of the North}{1:21:09}{38442758}{}
\FC{fight-club-103.mp3}{\FCtitle{\#103}{Zápisky ze studia}}{2012-11-08}{\Sleepless, \Luke, \Janchor, \Blixa}{\FCrule{\#103}{Když se něco likeuje, tak to je vždycky psina.}{\Janchor}}{Nové funkce na Games | Zpráva o~stavu válečných stříleček | Je vývoj her předražený? | Co čekat od nového pokračování The Longest Journey | Koně (a Zaklínač) na Mars! | Soutěž o~Halo 4 pro Xbox 360 + tričko}{1:16:16}{32983562}{}
\FC{fight-club-104.mp3}{\FCtitle{\#104}{Cesty a~rohy}}{2012-11-15}{\Sleepless, \Luke, \Janchor, \Hellboy}{\FCrule{\#104}{Nepodceňujte jezevčíky.}{\Janchor}}{JRC vrací úder | Co nového se nelíbí Danovi | Jak se má projekt studia Warhorse | Čím naplnit nové Grand Theft Auto | Potenciál Crysis vs potenciál Far Cry | Bude budoucnost patřit hraní zdarma? | Routine oživuje estetiku poctivých sci-fi | Soutěž o~Need for Speed: Most Wanted a~komiks Hellblazer}{1:56:41}{50172493}{}
\FC{fight-club-105.mp3}{\FCtitle{\#105}{This is Fight Club!}}{2012-11-22}{\Sleepless, \Luke, \Janchor, \Blixa, \JJirkovsky}{\FCrule{\#105}{Občas\,\dots lagujeme.}{\Janchor}}{Jak zvýšit povědomí o~hrách? | Zákulisí Game Developers Session | Smysl podpory zadržených českých vývojářů | Není WoWko jako HoFko | WiiU otevřelo náruč nezávislému vývoji | Vesmírné hry válcují Kickstarter | Soutěž o~encyklopedii Assassin's Creed}{1:22:41}{36474350}{}
\FC{fight-club-106.mp3}{\FCtitle{\#106}{Rocky show}}{2012-11-29}{\Sleepless, \Luke, \Janchor, \EGrillova, \VSahula}{\FCrule{\#106}{Gingers have soul.}{\Janchor}}{Na čem pracují Dreadlocks bez dredů | Ken Levine vzdal multiplayer pro nový Bioshock | Ultimátní božská strategie | Bokovky studia Double Fine | Quo vadis, Doom? | Kam se (ne)posunuly fanfilmy | Soutěž o~Assassin's Creed 3 pro Xbox 360 + mikinu}{1:12:57}{32230194}{}
\FC{fight-club-107.mp3}{\FCtitle{\#107}{Baníci s latte}}{2012-12-06}{\Sleepless, \Luke, \Janchor, \Blixa, \MRosa}{\FCrule{\#107}{Baník, pičo!}{\Janchor}}{Miner Wars venku a~co bude dál? | Kamarádíme s~MojeID | První týden s~WiiU | Co chystá Bungie | Stane se z~Dishonored "obyčejná" značka? | Nová herní vlna, nebo stará neherní póza? | První vánoční tipy | Smysl projektu jakouhru.cz | Podoby interakce | Soutěž o~Far Cry 3 hru + mikinu pro prvního a~přístup do bety DOTA 2 pro druhého výherce!}{2:07:17}{56638532}{}
\FC{fight-club-108.mp3}{\FCtitle{\#108}{Psí kusy}}{2012-12-13}{\Sleepless, \Luke, \Janchor, \OverWatch}{\FCrule{\#108}{Pes, který štěká\,\dots je v klimatizaci.}{\Janchor}}{Jak to dopadlo s~přednáškami z~GDS | Nastavujeme zrcadlo VGA | Phantom Pain novým Metal Gearem | Překvapivé závěry celoevropského hráčského průzkumu | Space Hulk se vrací | Komu patří tečky ze Stalkera | Adventní doporučení \#2 | Protip na přispívání do diskuzí | Soutěž o~Steam kód Miner Wars}{1:18:23}{56446955}{2}
\FC{fight-club-109.mp3}{\FCtitle{\#109}{Návraty}}{2012-12-20}{\Luke, \Janchor, \Blixa, \Tukan}{\FCrule{\#109}{S Martinem Bachem nejsi nikdy Sama doma.}{\Janchor}}{Svátky na Games | Gamepage vstává z~popela | Jak se žije časopisu Score | Víc měst = lepší GTA? | Pět hodin v~Bioshock Infinite | Kopíruje série Dead Space "úpadek" Resident Evil? | Nová plošinovka od tvůrců Magicky | Adventní tipy pro chudé | Soutěž o~XL tričko Dishonored}{1:29:38}{64551649}{1}
\FC{fight-club-110.mp3}{\FCtitle{\#110}{Uvolněno!}}{2012-12-27}{\Sleepless, \Luke, \Janchor}{\FCrule{\#110}{Nejlepší Whiskey je Whiskey, která je tak stará, aby si mohla objednat svoji vlastní Whiskey.}{\Janchor}}{Dárky | Žáby | Péefka | Přednášky | Přehrávač | Jablka | AC3 | Speciály | Pohádky | Začátky | Kališníci | Okouni | Průměry | Soutěž!}{1:12:21}{52111078}{3}
\FC{fight-club-111.mp3}{\FCtitle{\#111}{Věštírna}}{2013-01-03}{\Sleepless, \Luke, \Janchor, \JOlejnik}{\FCrule{\#111}{Příliš mnoho zombies, Sleeplessova smrt.}{\Janchor}}{Co nás čeká v~novém roce | Třetí Press Start bude o~smrti | Těšíme se na\,\dots{}? | Trable s~The War Z | Nastartuje Steam linuxové hraní? | DRM s~lidskou tváří | Nová generace konzolí na obzoru | Bárka THQ šla ke dnu | Soutěž o~F1 2012 a~F1 Race Stars na PC}{1:12:19}{32563606}{}
\FC{fight-club-112.mp3}{\FCtitle{\#112}{Hráč a~jeho svět}}{2013-01-10}{\Sleepless, \Luke, \Janchor, \OBroukal}{\FCrule{\#112}{Volíme Karla.}{\Janchor}}{Gamepage slaví comeback | Shield nebo Piston? | Kdo si naporcuje THQ | Co zdrželo vývoj DayZ | Podrazí nízký rozpočet nohy ambiciózní Remember Me? | Soutěž o~trio triček od Ubisoftu: AC3, Ghost Recon a~Far Cry 3}{1:18:03}{35712552}{}
\FC{fight-club-113.mp3}{\FCtitle{\#113}{Ženy z polygonů}}{2013-01-17}{\Sleepless, \Luke, \Janchor, \Blixa, \ZKosticova}{\FCrule{\#113}{Nebojte se feminismu.}{\Janchor}}{Genderové studie a~hry | Šovinismus v~herním průmyslu | Jak to bude s~flashkami | Věznění vývojáři propuštěni, měly protesty smysl? | Chris Taylor sází vše na jednu kartu | Blíží se oznámení Playstation 4? | Ztracená čtvrtá epizoda Half-Life | Disney chystá hru, která míchá světy i~žánry | Soutěž o~LittleBigPlanet Karting}{1:43:21}{44365952}{}
\FC{fight-club-114.mp3}{\FCtitle{\#114}{Press Play}}{2013-01-24}{\Sleepless, \Janchor, \Blixa, \BHeblik}{\FCrule{\#114}{Nová hra od Mikrosoftu se bude jmenoval Rízident Dívl.}{\Janchor}}{Vznik, přítomnost a~směřování akce Press Start | Flashky jsou na světě | Proč pláče Chris Taylor | Šach mat pro Atari | Flashback na Flashback! | Debaty o~herním násilí | Nová krev v~nezávislé scéně | Soutěž o~2 flashky Games.cz}{1:40:11}{41683510}{}
\FC{fight-club-115.mp3}{\FCtitle{\#115}{Kdopak by se ženy bál?}}{2013-01-31}{\Luke, \Janchor, \Ulrik, \Bety}{\FCrule{\#115}{Chlupy dělaj člověka.}{\Janchor}}{Horozelecká klika na Games se rozšiřuje | Co s~načatým Spectorem? | Běloruské tanky táhnou na západ | Americký viceprezident na straně soudnosti | Matka všech online bitev | Čím nás překvapil nový Dante | Soutěž o~(nikoli véčkové) tričko Devil May Cry}{1:11:01}{30139472}{}
\FC{fight-club-116.mp3}{\FCtitle{\#116}{Zrcadlo, zrcadlo\,\dots}}{2013-02-09}{\Luke, \Janchor, \Blixa, \Moolker}{\FCrule{\#116}{Sledovat Fight club na Linuxu je známka kyberpunku.}{\Janchor}}{Nový vysílací čas! | Vítáme posilu | Stagnace RTS žánru | Nepůvodní kyberpunk | Míří Zaklínač do otevřeného světa? | Tom Hall už zase skáče přes plošinky | Mongolský kyberpunk | Soutěž o~Aliens: Colonial Marines pro PS3 nebo X360}{0:59:42}{25528872}{}
\FC{fight-club-117.mp3}{\FCtitle{\#117}{PC pod televizí}}{2013-02-13}{\Sleepless, \Luke, \Janchor, \LMacura}{\FCrule{\#117}{Pozor na netoreapery!}{\Janchor}}{Ozvěny hackathonu | Trable s~Mambou | Newell v~obýváku | Chystá se režisér Star Treku fušovat do her? | Obsidiani by rádi novou hru z~Hvězdných válek | The Longest Journey se vrací | Soutěž o~libovolnou hru z~katalogu Cinemaxu}{1:30:37}{41021898}{2}
\FC{fight-club-118.mp3}{\FCtitle{\#118}{Lamači kostí}}{2013-02-21}{\Luke, \Janchor, \Blixa, \JOlejnik}{\FCrule{\#118}{Nejlepší porno -- Lamači kostí.}{\Janchor}}{Trable s~vetřelci | Osud nového projektu od tvůrců Halo | Seriál ze světa Mortal Kombat (zmínka o filmu Lamači kostí s hodnocením 85\,\%) | Ponor do éteru | Soutěž o~System Shock 2}{1:22:32}{35035410}{3}
\FC{fight-club-119.mp3}{\FCtitle{\#119}{}}{2013-03-02}{\Sleepless, \Luke, \Janchor, \Tafrob}{\FCrule{\#119}{Není čas se ohlížet a~koukat dozadu, je třeba makat a~koukat dopředu.}{\Janchor}}{Je třeba ustříhnout Tafroba | WTF, Novinky? | Představení PS4 |  Ubisoft si to chce vyžehlit u~PC hráčů | Těšíme se na novodobé dungeony | Planescape se vrací, ale v~jaké podobě? | Soutěž o~CD Krásný ztráty + Crysis 3 pro Xbox 360 / PS3}{1:41:56}{46194421}{}
\FC{fight-club-120.mp3}{\FCtitle{\#120}{Pod palbou DRDOS}}{2013-03-06}{\Sleepless, \Luke, \Janchor, \KDrda}{\FCrule{\#120}{Když chcete sundat zpravodajský server, použijte DDoS. Když herní, použijte Drdos.}{\Janchor}}{Popel na hlavu | Pirátem mezi assassiny | Vystačí Neverwinter se zavedenou licencí? | Zmatení jazyků v~Cyberpunk 2077 | První obrázky z~nového Thiefa | Studia v~přechodu | Schizofrenní Tornquist | Soutěž o~Crysis 3 pro Xbox 360 / PS3 \&{} Tafrobovy Krásný ztráty}{1:24:27}{37408464}{}
\FC{fight-club-121.mp3}{\FCtitle{\#121}{OUYA, nebo OUNE?}}{2013-03-13}{\Luke, \Janchor, \Blixa, \Cerv, \Moolker}{\FCrule{\#121}{Bílá paní neexistuje!}{\Janchor}}{Seznamte se, OUYA | Nový videopřehrávač | Fenomenální úspěch Tormentu | Garriott se vrací | Sociální nepokoje "občanů" SimCity | Analogoví bratři od Starbreeze | Soutěž o~Metal Gear Rising: Revengeance pro X360}{1:39:22}{43240114}{2}
\FC{fight-club-122.mp3}{\FCtitle{\#122}{Rekviem za šéfa}}{2013-03-21}{\Sleepless, \Luke, \Janchor, \Blixa}{\FCrule{\#122}{Habemus CEO! (We have a~CEO!)}{\Janchor}}{Deset tisíc lajků | Zavřete oči, CEO odchází | Herní vývoj v~cloudu? | Čekání na fantomy | DLC k~Dishonored nás nechává chladnými | Autor Neverhoodu se vrací na scénu | Soutěž o~Starcraft 2: Heart of the Swarm}{1:31:26}{39710101}{1}
\FC{fight-club-123.mp3}{\FCtitle{\#123}{PES na pracovišti!}}{2013-03-28}{\Sleepless, \Luke, \Janchor, \Blixa, \Julka}{\FCrule{\#123}{The beagel has landed.}{\Janchor}}{Recenze PES | Saturace trhu velkých her | Otřesy ve Squeenix | Battlefield sází na jistotu | XCOM do kapsy | Hry jsou jako déšť | Diskutabilní BioShock Infinite a~soutěž o~verzi pro X360 / PS3}{1:32:19}{37166017}{2}
\FC{fight-club-124.mp3}{\FCtitle{\#124}{}}{2013-04-04}{\Sleepless, \Janchor, \Blixa, \Ulrik}{\FCrule{\#124}{Nová kamera zeštihluje.}{\Janchor}}{Restartujeme sebe sama | Střet herních civilizací na GDC | Phantom Pain bez obvazů a~kdo se vyzná v~chronologii Metal Gear? | Střední vývojářská třída vstává z~popela | Blizzard jde do kartiček | Co je to MOBA? | Soutěž o~Sly Cooper: Thieves in Time}{1:26:12}{38272046}{}
\FC{fight-club-125.mp3}{\FCtitle{\#125}{Tryzna za LucasArts}}{2013-04-11}{\Sleepless, \Janchor, \Blixa, \Tlamiczka}{\FCrule{\#125}{LucasArts jsou mrtví, ať žijou LucasArts\,\dots}{\Janchor}}{Nahlas! | Sbohem a~šáteček, LucasArts | Stojí osud id Software na novém Doomovi? | Vlna retro grafiky a~lo-fi estetika v~moderních hrách | Soutěž o~Cities In Motion 2}{2:03:04}{54132004}{}
\FC{fight-club-126.mp3}{\FCtitle{\#126}{Filmová vsuvka}}{2013-04-18}{\Sleepless, \Janchor, \Blixa, \VRybar}{\FCrule{\#126}{Kola se kroutí dál.}{\Janchor}}{Herní filmy tady a~teď: Need for Speed, World of Warcraft a~tuzemská Lara Croft |  Jak jsou na tom klíčové závodní série a~co nás čeká s~příchodem další generace konzolí | Z~podporovatelů regulérní investoři | Osmdesátkové okénko | Soutěž o~Army of Two 2 pro Xbox 360}{1:41:30}{44999116}{}
\FC{fight-club-127.mp3}{\FCtitle{\#127}{Studio Amanita}}{2013-04-24}{\Sleepless, \Luke, \Janchor, \Blixa, \JDvorsky}{\FCrule{\#127}{Není Molyneux jako Molydeux.}{\Janchor}}{Odkud se vzala a~co nového peče Amanita | Reportáž z~tureckého herního časopisu | Molyneux vydělává na trollingu | Gearbox si koupil Homeworld | Coraabia otevřela svou mimozemskou palubu | Soutěž o~výtisk Oyungezer + speciální edici Machinaria}{1:39:47}{39624196}{}
\FC{fight-club-128.mp3}{\FCtitle{\#128}{Rohlík s Drdolem}}{2013-05-04}{\Sleepless, \Luke, \JOlejnik, \KDrda, \Skeldal}{\FCrule{\#128}{Někdy stačí ke štěstí jen 15 kilometrů.}{\Luke}}{Výpadek zvuku | Brány Skeldalu na dotek | Vedoucí role strany | Game Dev Tycoon vytáhl na piráty\,\dots piráty! | Kinoautomat Davida Cage | XCOM střílečka žije | Soutěž o~Star Trek pro PS3}{1:40:06}{37758649}{3}
\FC{fight-club-129.mp3}{\FCtitle{\#129}{Level reloaded}}{2013-05-10}{\Sleepless, \Janchor, \Blixa, \Draken}{\FCrule{\#129}{Tenhle podcast Vám přinesl Draken.}{\Janchor}}{Draken si koupil LEVEL, co se s~ním chystá provést? | EA rozpoutá Hvězdné války | Wolfenstein s~Hendrixem | Smůla Patrice Désiletse | Soutěž o~tričko Watch Dogs}{1:32:13}{39595714}{}
\FC{fight-club-130.mp3}{\FCtitle{\#130}{Sledujme psy!}}{2013-05-13}{\Sleepless, \Luke, \Janchor, \Blixa}{\FCrule{\#130}{Není nadto si každou hru pořádně vymodit.}{\Janchor}}{Městské akce jdou do sebe, která vyhraje? | Editor pro Zaklínače a~smysluplnost modů | TIM se vrací! | Soutěž o~Leviathan: Warships}{1:19:32}{34966904}{}
\FC{fight-club-131.mp3}{\FCtitle{\#131}{Let's play Xbox One}}{2013-05-22}{\Luke, \Janchor, \Blixa, \Moolker}{\FCrule{\#131}{Amiga forever!}{\Janchor}}{Testujeme YouTube | Představení nového Xboxu: kde jsou hry, kudy se tam strkají VHSky? | Remedy se vrací skrze Quantum Break | Nintendo si duplo na výdělky z~Let's Play jejich her | Problematický vývoj Last Light | Bude Remember Me skutečně nezapomenutelná? | Hledači pokladů! | Soutěž o~kšiltovku The Last of Us + kukátko GTA V | Zpráva detektivní služby o Amiga hře Poklady zlatých sluncí}{1:38:29}{41223414}{1}
\FC{fight-club-132.mp3}{\FCtitle{\#132}{Feministka z Kickstarteru}}{2013-05-29}{\Sleepless, \Luke, \Janchor, \Blixa}{\FCrule{\#132}{Těžké být Bohem.}{\Janchor}}{Anita Sarkeesian nepřestává budit hysterii | Pokračování Neverhoodu chce vaše peníze | Molyneux udělil božství náctiletému Skotovi | S~herními horory se roztrhl pytel | Soutěž o~kompletní edici Heroes VI}{1:38:18}{45664124}{}
\FC{fight-club-133.mp3}{\FCtitle{\#133}{Žeru hry}}{2013-06-05}{\Sleepless, \Luke, \Janchor, \Moolker}{\FCrule{\#133}{Ať žije youtube.}{\Janchor}}{První konkrétní informace o~novém LEVELu | Miloš s~Pavlem balí na E3 | Prey 2 žije pod taktovkou Arkane? | První Fable se vrátí v~moderním kabátku | Zombies se chystají opět zdupat vaši zahradu |  Soutěž o~Dead Island: Riptide pro PC}{1:21:24}{36285673}{}
\FC{fight-club-134-E3-1.mp3}{\FCtitle{E3 2013 speciál \#1}{}}{2013-06-12}{\Sleepless, \Janchor, \Blixa, \Ulrik}{}{Rekapitulujeme první den E3 2013 a~vybíráme ty nejzajímavější hry a~událostí, kterým byste měli věnovat pozornost.}{1:13:34}{33889693}{}
\FC{fight-club-134-E3-2.mp3}{\FCtitle{E3 2013 speciál \#2}{}}{2013-06-13}{\Sleepless, \Janchor, \Blixa, \OverWatch}{}{Rekapitulujeme druhý den herní výstavy E3 2013 a~vybíráme ty nejzajímavější hry a~událostí, kterým byste měli věnovat pozornost.}{0:55:20}{25226315}{}
\FC{fight-club-134-E3-3.mp3}{\FCtitle{E3 2013 speciál \#3}{}}{2013-06-14}{\Janchor, \Blixa, \Ulrik, \OverWatch}{}{Rekapitulujeme třetí den E3 2013 a~vybíráme ty nejzajímavější hry a~událostí, kterým byste měli věnovat pozornost.}{0:40:12}{19084385}{}
\FC{fight-club-134-E3-4.mp3}{\FCtitle{E3 2013 speciál \#4}{FFFight club}}{2013-06-15}{\Sleepless, \Janchor, \Blixa, \FF}{\FCrule{\#134}{Děcka, za války bylo všechno lepší.}{\Janchor}}{Rekapitulace herní výstavy E3 2013, tentokrát s~Františkem Fukou, který zhodnotil nejenom výstavu, ale celý herní průmysl a~hry.}{1:24:04}{35520457}{3}
\FC{fight-club-135.mp3}{\FCtitle{\#135}{Bacha na psa!}}{2013-06-19}{\Sleepless, \Luke, \Janchor, \Moolker}{\FCrule{\#135}{Když vojáci, tak psí vojáci.}{\Janchor}}{Naposledy jsme se v~podcastu vrátili k~nedávno skončené herní výstavě E3 2013. Z~Prahy už jsme E3 probrali ze všech úhlů, tentokrát o~významné herní události uslyšíte od Miloše a~Pavla, kteří se jí přímo účastnili.}{1:38:23}{47520229}{}
\FC{fight-club-136.mp3}{\FCtitle{\#136}{Byznys je byznys}}{2013-06-26}{\Luke, \Janchor, \Blixa, \Ulrik}{\FCrule{\#136}{I live\,\dots again.}{\Janchor}}{Nový LEVEL 231 | Ordnung v~roce 1666 | Z~autora Neverhoodu se vyklubal homofob | Návrat Chaosu v~novém Enginu | Soutěž o~podepsaný signální výtisk nového Levelu}{1:18:35}{36673108}{1}
\FC{fight-club-137.mp3}{\FCtitle{\#137}{Návrat psance}}{2013-07-03}{\Sleepless, \Luke, \Janchor, \Blixa}{\FCrule{\#137}{Zrušte Levely!}{\Janchor}}{Z bláta Xboxu do louže Zyngy? | MMORPG vymknuté z~kloubů | Schafer mizerným producentem | Outcast zpátky u~svých tvůrců | Grafika znělky Fight clubu (cancoury mezi nohama) | Soutěž o~(plnou!) plechovou krabičku Watch Dogs}{1:12:57}{32580108}{3}
\FC{fight-club-138.mp3}{\FCtitle{\#138}{Nejrychlejší kolo}}{2013-07-10}{\Janchor, \Blixa, \LMacura, \VRybar}{\FCrule{\#138}{One does not simply email Valve.}{\Janchor}}{Roztříštěné Grand Theft Auto | Odvrácená strana Valve | Cesta za volant skutečného závodního auta | Soutěž o~Pravdivý příběh byzantských věrozvěstů Konstantina a~Metoděje}{1:31:39}{36208118}{2}
\FC{fight-club-139.mp3}{\FCtitle{\#139}{Volej Olej}}{2013-07-17}{\Sleepless, \Luke, \Janchor, \JOlejnik}{\FCrule{\#139}{Mobily jsou na hlavu.}{\Janchor}}{LEVEL 232 | Není Commander jako Commander | Oculus Rift branou do nových žánrů? | League of Legends ku zdraví! | Stážistou ve Valve | Soutěž o~XL tričko Destiny a~survival kit k~The Division}{1:19:17}{34923690}{3}
\FC{fight-club-140.mp3}{\FCtitle{\#140}{Všechny cesty vedou do Říma}}{2013-07-24}{\Sleepless, \Luke, \Janchor, \Ulrik}{\FCrule{\#140}{Všechny quick time eventy vedou do Říma.}{\Janchor}}{Nový LEVEL za rohem | Vyčpělost série Splinter Cell | Call of Řím | Fargo obchodníkem, Schafer troubou | Runescape žije | Soutěž o~tričko Murdered a~klíč na Legends of Dawn}{1:25:33}{39484048}{2}
\FC{fight-club-141.mp3}{\FCtitle{\#141}{Kdo uteče, prohraje}}{2013-07-31}{\Sleepless, \Luke, \Janchor}{\FCrule{\#141}{Zavřete oči, odcházíme z herního průmyslu.}{\Janchor}}{LEVEL 232 | (9:07) Phil Fish nevydýchal své kritiky | (23:48) Oblaka v~cloudu | (40:41) Duchovní nástupce Syndicate BUDE! | Soutěž o~tričko DayZ a~kód na hru Race Injection}{1:02:15}{25902578}{}
\FC{fight-club-142.mp3}{\FCtitle{\#142}{Souboj titánů}}{2013-08-07}{\Luke, \Blixa, \KDrda, \LMacura}{\FCrule{\#142}{Když čekat ve frontě, tak jedině na bosse v Everquestu.}{\Luke}}{Sleepless na dovolence | (03:32) Lukáš Macura vydává a~vítězí | (12:46) Vliv očních vad na kvalitu herního zážitku | (19:03) ArmA 3 bez kampaně! Skandál? | (27:00) John Carmack, iD Soft a~jak je vidíme | (37:10) Everquest se vrací | (39:31) Asheron's Call vs Everquest !EPIC! | (48:34) Rozjímání nad mobilními hrami | (58:06) Dotazy a~odpovědi | (1:24:08) Soutěž o~Inquisitora}{1:26:54}{37275548}{3}
\FC{fight-club-143.mp3}{\FCtitle{\#143}{Dogs Club s Danem Vávrou}}{2013-08-15}{\Sleepless, \Luke, \Blixa, \Hellboy}{\FCrule{\#143}{Když je Vlk mezi vámi, tak ho oslovujte pane Kardinále.}{\Luke}}{Balíme na Gamescom | (4:00) Výživný Humble Bundle | (11:10) Telltale chystají hru podle komiksu Fables a~co trápí adventury? | (29:12) Grand Theft Auto jde do multiplayeru | (43:38) zaklínač | (59:53) Dotazy | Soutěž o~klávesnici Battlekeyboard}{1:43:13}{42802452}{}
\FC{fight-club-143-Gamescom.mp3}{\FCtitle{Gamescom 2013 speciál}{}}{2013-08-21}{\Sleepless, \Luke, \Janchor, \Blixa, \KDrda}{}{Speciál z~GamesComu 2013 | Ale nejdřív nový LEVEL! | Poprvé na GDC | Zoo Tycoon - hra výstavy? | Shadow of the Beast - hra výstavy? | Den se Sony | Stíny a~titáni | Přeskoč, přelez, probourej, ale nepodlez}{0:30:47}{13549652}{}
\FC{fight-club-144.mp3}{\FCtitle{\#144}{Indiani z Větrova}}{2013-08-28}{\Sleepless, \Luke, \Janchor, \Fiola, \Baty}{\FCrule{\#144}{Každý je tak starý jak se cítí.}{\Janchor}}{(2:47) To nejlepší z~letošního Gamescomu | (21:34) Assassin's Creed volá jo-ho-ho! | (35:08) Online nastavování kaše a~kam se poděly akurátní datadisky? | (47:18) Návrat zlotřilého génia | Soutěž o~napěchovaný vojenský batoh Wargaming.net}{1:28:51}{36624594}{}
\FC{fight-club-145.mp3}{\FCtitle{\#145}{Otázky Petra Bulíře}}{2013-09-04}{\Sleepless, \Luke, \Janchor, \PBulir}{\FCrule{\#145}{Čim víc achievementů, tím lepší hra. Tím víc Adidas, vole.}{\Janchor{} a~dodal \Sleepless}}{(2:00) Rayman králem plošinek | (7:50) Definice nezávislosti | (11:40) Je Gone Home předražená? | (15:00) Dead Rising utíká z~interiérů | (24:50) Z~Hardware je oficiálně nový Homeworld | (31:00) Půjde Sony do vlastního "Oculusu"? | Soutěž o~tričko (L / XL) a~hru (PC) Splinter Cell: Blacklist}{1:05:27}{29478452}{}
\FC{fight-club-146.mp3}{\FCtitle{\#146}{Ptáci šli na pivo}}{2013-09-12}{\Sleepless, \Luke, \Janchor, \Blixa}{\FCrule{\#146}{Opeřencům nenalévat.}{\Janchor}}{Nezávislý server | (8:55) Ubi jde do menších her | (19:25) Staronové handheldy | (27:55) Astronomický rozpočet GTA V | (33:38) Co si slibovat od Deep Down | (41:00) Steam nabídne sdílení účtů | Soutěž o~podepsanou látkovou plachtu Tom Clancy's EndWar \&{} "řidičák" The Crew}{1:27:05}{37803336}{2}
\FC{fight-club-147.mp3}{\FCtitle{\#147}{GTA V Live!}}{2013-09-18}{\Sleepless, \Blixa, \Janchor, \VRybar}{\FCrule{\#147}{Kdo chce bejt z vás zločinec, tak ať se radši rovnou zabije.}{\Janchor}}{Ukecaný Chris Roberts | (4:20) Hrajeme živě Grand Theft Auto V | (52:00) Elite se připomíná ve stylu Hvězdných válek | (1:12:00) Spletité cesty ke kariérám v~herní novinařině | Soutěž o~Grand Theft Auto V~pro Xbox 360}{1:20:32}{36343140}{}
\FC{fight-club-148.mp3}{\FCtitle{\#148}{Pára nad hrncem}}{2013-09-25}{\Luke, \Blixa, \JOlejnik, \LMacura}{\FCrule{\#148}{Když nemáte pravidlo, řekněte: "Martine vypni to!".}{\Luke}}{LEVEL 234 | (5:10) Je virtuální mučení přes čáru? | (15:50) Valve chystá operační systém | (31:50) Apple vykročil směrem k~hráčům | (39:40) Kde končí hry a~začínají\,\dots co přesně? | Soutěž o~puzzle k~Zaklínačovi 3}{1:08:59}{29443514}{1}
\FC{fight-club-149.mp3}{\FCtitle{\#149}{Fu*k off střílečky}}{2013-10-02}{\Sleepless, \Luke, \Janchor, \Blixa, \Ulrik}{\FCrule{\#149}{Fuck off střílečky!}{\Janchor}}{Jak to bude se srazem a~výstavou For Games | (5:00) Pavel a~piráti | (15:20) Hardware od Valve | (29:45) GTA Offline | (34:30) Válečné hry s~reflexí válečných zločinů | (42:35) Audiovzkaz o slovním parazitizmu a~rétorice | Soutěž o~pirátskou vlajku Assassin's Creed IV}{1:08:24}{28711422}{2}
% Fight Club Special Diablo III (Publikováno 10. 10. 2013)
\FC{fight-club-150.mp3}{\FCtitle{\#150}{Katovna s Karlem Drdou}}{2013-10-09}{\Sleepless, \Luke, \Janchor, \KDrda}{\FCrule{\#150}{Vzhůru na Neustadt!}{\Janchor}}{Sobotní sraz | (7:00) Zpátky do Rapture | (17:00) Deus Ex se rozpíná | (25:45) Kolektivizace u~Square Enix | (35:50) Tvůrci původní Amnesie jdou do vědecké fantastiky | Soutěž o~art book Beyond}{1:17:04}{32969928}{1}
\FC{fight-club-151.mp3}{\FCtitle{\#151}{Beyond bez duše}}{2013-10-16}{\Sleepless, \Luke, \Janchor, \Blixa}{\FCrule{\#151}{David Cage je nahý.}{\Janchor}}{LEVEL 235 | (4:10) Jak se vyvedl druhý čtenářský sraz a~veletrh For Games | (10:40) Další pokus o~obrýlenou evoluci augmentované reality | (20:55) PS4 má konečně české datum i~cenu | (27:30) Watch Dogs až příští rok | (35:30) Rozpolcenost Beyond: Two Souls | Soutěž o~epesní batmanovský USB batarang}{1:22:45}{37032054}{1}
\FC{fight-club-152.mp3}{\FCtitle{\#152}{Blizzardovy nové boty}}{2013-10-23}{\Sleepless, \Luke, \Janchor, \Blixa}{\FCrule{\#152}{Dost bylo Pety, ať žije Peťa!}{\Janchor}}{Blíží se GDS | (5:25) Brazilská variace na ICO vábí i~odrazuje přeplácaností | (13:45) Blizzard vysoudil boty | (31:30) Tvůrci Total War jdou do Vetřelců | (36:30) Duch Mystu žije na Kickstarteru | Soutěž o~kód na Goodbye Deponia}{1:18:11}{33231728}{2}
\FC{fight-club-153.mp3}{\FCtitle{\#153}{Piráti bez papoušků}}{2013-10-30}{\Sleepless, \Luke, \Janchor, \Ulrik}{\FCrule{\#153}{Volte piráty!}{\Janchor}}{Piráti, kteří uspěli | (28:40) Exkluzivita vrací úder | (35:45) Co nabídne druhá herní sezóna Walking Dead? | (40:20) Na retro kolej přijíždí Confederate Express | Soutěž o~sošku Edwarda z~Assassin's Creed IV}{1:15:27}{32371530}{}
\FC{fight-club-154.mp3}{\FCtitle{\#154}{Parní Rambo}}{2013-11-06}{\Sleepless, \Luke, \Janchor, \Blixa}{\FCrule{\#154}{Hlídejte si klíče.}{\Janchor}}{Listopadové herní akce | (11:25) Trable s~rozuteklými aktivačními klíči | (18:40) Co si lze odvodit na základě dojmů z~hardwaru od Valve | (34:50) Namísto herního raději filmového Ramba | (42:00) Slibná adventura Stasis to nově zkouší přes Kickstarter | Soutěž o~Origin kód na rozšířenou verzi FIFA 14}{1:15:05}{33989300}{1}
\FC{fight-club-155.mp3}{\FCtitle{\#155}{BlizzCon s Adamem}}{2013-11-13}{\Sleepless, \Luke, \Janchor, \OverWatch}{\FCrule{\#155}{D'horreur, d'horreur\,\dots}{\Janchor}}{(7:40) Zevrubně o~BlizzConu: datadisk k~Warcraftu, film, vlastní MOBA | (28:25) FTL dostane nové scénáře | (34:10) Humble Bundle spouští obchod | (41:00) Výtvarně působivá hororovka White Night | Soutěž o~hrací karty Assassin's Creed}{1:11:35}{33468286}{2}
\FC{fight-club-156.mp3}{\FCtitle{\#156}{Představení české RPG Dex}}{2013-11-20}{\Sleepless, \Luke, \Blixa, \JJirkovsky}{\FCrule{\#156}{Nevěřte Vigo Vikernesovi.}{\Luke}}{(03:32) Dex na Kickstarteru s~Honzou Jirkovským | (30:24) Playstation 4 je tu - a? | (37:32) Star Citizen: budeme kupovat nová PC? | (45:37) Dotazy | (82:33) Soutěž o~"mystery" triko a~rohožku pod myš}{1:25:43}{38352344}{2}
\FC{fight-club-157.mp3}{\FCtitle{\#157}{}}{2013-11-27}{\Luke, \Blixa, \Moolker, \Ulrik, \MZeman, \DSvetlik}{\FCrule{\#157}{Nemíchejte papír s kachlíky.}{\Luke}}{{\em Kvůli technickým komplikacím obsahuje druhá část Fight Clubu zvukové chyby a~je tímpádem pouze pro otrlé nebo hluché. Nic s~tím nenaděláme a~omlouváme se :(} Rozhovor s~autory Databáze her | (15:52) Xbox One je tu - a? | (24:04) Voyerův ráj v~Playroom | (28:16) Šlapou nám uživatelské recenze na Steamu na paty? | (36:25) Shadowrun dostal berlínský datadisk | (41:32) Dotazy a~soutěž o~odznáčky AC4 a~tričkou BF3}{1:00:59}{27007808}{}
\FC{fight-club-158.mp3}{\FCtitle{\#158}{Teleshopping}}{2013-12-04}{\Sleepless, \Luke, \Blixa}{\FCrule{\#158}{Když podcast končí, tak se to rozpadá.}{\Luke}}{Osahání Playstation 4 | (8:05) Pečeme nový LEVEL | (13:25) Chce být Darkout novou Terrarií? | (22:05) Garry's Mod slaví třímilionové prodeje | (28:55) Double Fine staví kosmickou stanici | Soutěž o~kód na první Wasteland}{1:12:44}{30823014}{1}
\FC{fight-club-159.mp3}{\FCtitle{\#159}{PS4 pod lupou a~závody v Gran Turismo 6}}{2013-12-12}{\Sleepless, \Luke, \JOlejnik, \VRybar, \JRejlek, \PDolezal}{\FCrule{\#159}{Nejlíp vám to udělá trona.}{\Luke}}{Gran Turismo v~šestém pokračování i~"naživo" | (26:15) Hrajeme si na PlayStation 4 | Soutěž o~novinářský balíček Gran Turismo 6}{1:29:56}{40406093}{}
\FC{fight-club-160.mp3}{\FCtitle{\#160}{S Danem Vávrou o Kingdom Come}}{2013-12-19}{\Sleepless, \Luke, \Blixa, \Hellboy}{\FCrule{\#160}{Jak na Nový rok, tak po celý rok.}{\Luke}}{S Danem Vávrou o~hře Kingdom Come: Deliverance a~spoustě dalšího | Soutěž o~luxusní corebook Numenera}{1:13:47}{32123912}{1}
\FC{fight-club-161.mp3}{\FCtitle{\#161}{}}{2014-01-08}{\Sleepless, \Luke, \Ulrik, \JOlejnik}{\FCrule{\#161}{Pytel je na pytel.}{\Luke}}{Ohlášení vyhlášení her roku Best of 2013 (5:13) | Dojmy z~upatlaného ovladače Xbox One, neposlušného Kinectu a~nepřehledného menu | (19:23) Subjektivní srovnání Xbox One a~PS4 | (22:53) Hrajeme Ryse: Son of Rome | (30:28) Bouráme s~Forza Motorsport 5 | (39:46) Masakrujeme zombíky v~Dead Rising 3 | (54:36) Dotazy | (01:13) Výherce core booku Numenera z~FC \#160 | (01:13) Nová soutěže o~pirátský pytel s~motivem z~Assassin's Creed IV | (01:17) Pravidlo Fight Clubu \#161.}{1:17:47}{35847068}{2}
\FC{fight-club-162.mp3}{\FCtitle{\#162}{Vetřelec na pranýři}}{2014-01-15}{\Sleepless, \Luke, \Blixa}{\FCrule{\#162}{S vetřelcem podlahu nevytírejte.}{\Luke}}{(7:30) 2K Czech Praha zavírají krám? | (17:00) Herní Vetřelec zkusí složit reparát | (27:10) Co nás zaujalo na výstavě CES | (44:50) Hromada otázek | Soutěž o~cibulačky, tričko \&{} zápisník Assassin's Creed}{1:13:38}{33409993}{2}
\FC{fight-club-163.mp3}{\FCtitle{\#163}{Království válečného koně}}{2014-01-22}{\Sleepless, \Luke, \Janchor, \Blixa}{\FCrule{\#163}{Nenávidim metro.}{\Janchor}}{LEVEL 238 | (7:15) Přijď království Danovo! A~naše dojmy z~návštěvy u~Warhorse | (38:00) Umíráček pro Games for Windows Live | (45:15) Tvůrci Don't Starve chystají isometrickou špionážní akci | (50:30) Pustošení starého Řecka | Soutěž o~roční předplatné PlayStation+}{1:17:58}{36170941}{1}
\FC{fight-club-164.mp3}{\FCtitle{\#164}{Zpátky do kobky}}{2014-01-29}{\Sleepless, \Luke, \Janchor, \Blixa}{\FCrule{\#164}{Riftujte zásadně do kýble.}{\Janchor}}{Čtenáři vyberou hry roku | (3:50) Microsoft koupil Gears of War | (10:40) Renesance dungeonů | (27:45) SOE zařízla tři online hry, ale co peče nového? | (36:10) Ghostbusters + UFO = Ghost Control | Soutěž o~Assassin's Creed III: Liberation HD na PC }{1:23:25}{36616334}{2}
\FC{fight-club-165.mp3}{\FCtitle{\#165}{Všechno je jinak}}{2014-02-05}{\Sleepless, \Luke, \Janchor, \Blixa, KDrda}{\FCrule{\#165}{Jak to mám kurva vědět!}{\Janchor}}{(3:20) PC trh roste | (16:10) Onlinovka ReRoll chce mapovat skutečnou Zemi | (23:10) Další dungeon na obzoru, tentokrát sci-fi | (32:20) Detektivka Murdered se klube ze stínů | Soutěž o~Xbox Live kód na Orcs Must Die!}{1:08:50}{30562497}{2}
\FC{fight-club-166.mp3}{\FCtitle{\#166}{Slavnosti červů}}{2014-02-12}{\Sleepless, \Luke, \Janchor, \KDrda, \LMacura}{\FCrule{\#166}{Halapartna je základní pavoučí právo.}{\Janchor}}{(4:40) Karlovy dojmy z~prvních levelů v~Elder Scrolls Online | (37:40) Lukáš Macura má červy a~chystá se do dungeonu | (48:20) "Aféra" Flappy Bird | (61:40) Lords of the Fallen, alternativa k~Dark Souls? | (66:20) Dost bylo zombies | Soutěž o~Might \&{} Magic X}{1:37:39}{42781462}{3}
\FC{fight-club-167.mp3}{\FCtitle{\#167}{S Rybkou na suchu}}{2014-02-19}{\Sleepless, \Blixa, \Wolf}{\FCrule{\#167}{Rybka na suchu nekouše.}{\Blixa}}{(6:46) Mobilní spin-offy velkých her | (23:30) Michalův pohled na finišující kickstarter Kingdom Come | (37:50) Ohlédnutí za vývojem hardwaru | (55:20) Plný valník dotazů | Soutěž o~dvě knihy ze světa The Elder Scrolls}{1:28:10}{38148542}{3}
\FC{fight-club-168.mp3}{\FCtitle{\#168}{Návrat Farida}}{2014-02-26}{\Sleepless, \Luke, \Janchor, \Farid}{\FCrule{\#169}{Na Everest zásadně s granátem}{\Janchor}}{(5:15) Duke k~soudu? | (12:35) World of Warcraft koketuje s~pay-to-win? | (30:40) Dean Hall to u~nás balí | (40:45) Jak navrhnout moderní herní detektivku | Soutěž o~tričko (L nebo XL) Assassin's Creed IV}{1:34:05}{41251189}{}
\FC{fight-club-169.mp3}{\FCtitle{\#169}{Anatomie zloděje s Petrem Bulířem}}{2014-03-05}{\Luke, \Ulrik, \PBulir}{\FCrule{\#169}{Bez Grygara se to pravidlo vymyslet nedá!}{\Luke}}{Obšírně o~novém Thiefovi | (31:00) Stručně o~novém Batmanovi | (35:30) Kickstarter překročil vybranou miliardu | (55:30) Valník dotazů | Soutěž o~South Park: The Stick of Truth}{1:16:12}{34534525}{}
\FC{fight-club-170.mp3}{\FCtitle{\#170}{Rudá krev v zemi stínu}}{2014-03-12}{\Sleepless, \Luke, \Janchor, \Blixa}{\FCrule{\#170}{Nehrbte se.}{\Janchor}}{Blýská se na změny | (7:00) Tvůrci F.E.A.R.u to zkoušejí s~černobílou paletou | Co se to stalo s~Irrational Games? | (27:50) Třetí Zaklínač odložen na příští rok | Soutěž o~Titanfall pro PC}{1:12:43}{31506916}{}
\FC{fight-club-171.mp3}{\FCtitle{\#171}{Pád titánů z nebe}}{2014-03-19}{\Luke, \Moolker, \Ulrik, \OverWatch}{\FCrule{\#171}{Titan, který štěká, nekouše.}{\Luke}}{(5:10) Sony poodhalila svůj projekt virtuální reality | (16:30) Jak se prodává nová generace konzolí | (25:50) Titanfall, Titanfall, Titanfall | (42:25) Válka z~pohledu civilisty | Soutěž o~karty a~kapesní hodinky Assassin's Creed}{1:25:40}{37987682}{}
\FC{fight-club-172.mp3}{\FCtitle{\#172}{Jak Facebook k virtuální realitě přišel}}{2014-03-26}{\Sleepless, \Janchor, \Blixa, \JOlejnik}{\FCrule{\#172}{Aktualizujte si plužány (aka pluginy).}{\Janchor}}{LEVEL 240 | (12:45) Oculus Facebook | (31:40) Cenová válka populárních enginů? | (42:35) Assassin se vrací "domů" | (50:30) Dojmy z~GDC | Soutěž o~Steam kód Tower of Guns}{1:36:10}{42879098}{3}
\FC{fight-club-173.mp3}{\FCtitle{\#173}{Mňau!}}{2014-04-03}{\Sleepless, \Janchor, \Blixa, \Ulrik}{\FCrule{\#173}{Mňau.}{\Janchor}}{Facelift Games | (5:10) Amazon jde do her naplno | (15:00) Game Jam, který nebyl | (31:00) JRPG "pro všechny" | Soutěž o~Age of Wonders 3}{0:54:27}{23406992}{}
\FC{fight-club-174.mp3}{\FCtitle{\#174}{Hotline Miami 2 a~ti druzí}}{2014-04-09}{\Sleepless, \Luke, \Blixa}{\FCrule{\#174}{Do vesmíru jedině s hasičským přístrojem}{\Luke}}{Jak to všechno bude! | (6:15) Adam Orth se vrací oknem | (15:15) Hotline Miami se vrací ve dvou | (21:50) Scénáristka Legacy of Kain pracuje na hře podle Hvězdných válek | (34:15) Zákopčaník the Game}{0:41:05}{19093272}{}
\FC{fight-club-175.mp3}{\FCtitle{\#175}{Temné názory}}{2014-04-16}{\Luke, \Moolker, \OverWatch}{\FCrule{\#175}{S vajkingy na věčné časy!}{\Luke}}{Trable s~dárky pro novináře | (6:00) Zevrubné dojmy z~The Elder Scrolls Online | (26:30) Jak přežít v~Dark Souls 2 | (45:20) Návrat Alpha Centauri | (52:00) Nejhranější hry | (60:00) Vikinské outro}{1:04:06}{30149948}{}
\FC{fight-club-176.mp3}{\FCtitle{\#176}{LEGO hate!}}{2014-04-23}{\Sleepless, \Luke, \Moolker}{\FCrule{\#176}{Když to chrastí, je to Dobrovský.}{\Luke}}{Co přinese Dragon Age: Inquisition | Dopad zrušení české pobočky EA na hráče | Assasin's Creed nejcennější značkou Ubisoftu | Jak se dojí LEGO hry}{0:48:17}{20819270}{}
\FC{fight-club-177.mp3}{\FCtitle{\#177}{Hororový osud}}{2014-04-30}{\Luke, \JOlejnik, \Moolker}{\FCrule{\#177}{Na každého programátora alespoň deset manažerů.}{\Luke}}{Co čekat od Destiny | (18:00) "Děsivý" Daylight | (34:15) Další kolo debaty o~enginech | (45:00) E.T. vstal z~hrobu \&{} Xbox Originals}{0:56:23}{24922502}{}
\FC{fight-club-178.mp3}{\FCtitle{\#178}{VR'čení}}{2014-05-07}{\Sleepless, \Luke, \Janchor, \Blixa}{\FCrule{\#178}{Kdo maže ten jede.}{\Janchor}}{Je John Carmack zloděj? | (12:20) Musíme si promluvit o~Kevinovi | (23:10) Nikdy neříkej Never | (26:50) Matka Rus krvácí v~rytmu Streets of Rage}{0:38:45}{16418472}{}
\FC{fight-club-179.mp3}{\FCtitle{\#179}{Level up!}}{2014-05-14}{\Sleepless, \Luke, \Janchor, \Blixa}{\FCrule{\#179}{Free Pepina.}{\Janchor}}{Umíráček Kinectu? | (20:50) Šance na další Darksiders | (31:25) Nešťastná kampaň na remake Outcastu | (44:10) Jak jsem poznal svého strýčka}{0:49:55}{20836262}{2}
\FC{fight-club-180.mp3}{\FCtitle{\#180}{Bratři v triku a~košiláč}}{2014-05-21}{\Sleepless, \Luke, \Janchor, \Blixa}{\FCrule{\#180}{Nesuďte hry podle obalu.}{\Janchor}}{YouTube koupilo Twitch?! | (12:30) Robotická honitba | (21:40) Kim Čong-un v~plošinovce | (29:50) Bouře ve sklenici vody kolem Far Cry 4 | (44:35) Jazzová lovecraftovka Witchmarsh}{0:48:43}{20672912}{}
\FC{fight-club-181.mp3}{\FCtitle{\#181}{Der Käfer}}{2014-05-27}{\Sleepless, \Luke, \Blixa, \LMacura}{\FCrule{\#181}{Modrá je nová žlutá.}{\Luke}}{Vypusťte Watch Dogs! | (9:50) I~cesta může být perfektní | (13:55) Hitman jako příklad dobrého portu pro iPad | (24:30) Mají si movité firmy říkat o~peníze na Kickstarteru? | (41:55) Na dno oceánu a~ještě hlouběji}{0:52:03}{21601782}{3}
\FC{fight-club-182.mp3}{\FCtitle{\#182}{Sjednocení}}{2014-06-04}{\Sleepless, \Janchor, \Blixa}{\FCrule{\#182}{Když se podíváte do podpaží, podpaží se podívá do vás\,\dots}{\Janchor}}{E3 2014 se blíží | (8:00) Microsoft si opět zamiloval PC hráče? | (19:00) Vlna velkých odkladů | (27:20) Proč je Watch Dogs nejúspěšnější nová značka v~historii | (33:40) SOCIAL CLUB | Soutěž o~Transistor pro PC}{1:12:17}{29702422}{3}
\FC{fight-club-183.mp3}{\FCtitle{\#183}{Emergentní Proceduralita}}{2014-06-18}{\Luke, \Moolker, \Ulrik}{\FCrule{\#183}{Sprejujte jedině s kamarády.}{\Luke}}{E3 výprava se ohlíží za výstavou: dojmy, nejlepší hry, odpovědi na otázky a~soutěž o~E3 SWAG!}{1:29:19}{36939556}{2}
\FC{fight-club-184.mp3}{\FCtitle{\#184}{Mobilní disputace}}{2014-06-25}{\Sleepless, \Luke, \Blixa, \Tukan}{\FCrule{\#184}{Bez steroidů asteroid nesejmeš!}{\Luke}}{S Mikolášem o~mobilních hrách i~aplikacích | (18:00) Trable s~dotykovým ovládáním | (31:00) Atari vstává z~hrobu? | (41:30) Zodpovědnost vůči kickstarterovým podporovatelům | Soutěž o~kód na Enemy Front}{1:29:11}{38499172}{2}
\FC{fight-club-185.mp3}{\FCtitle{\#185}{Underworld}}{2014-07-02}{\Sleepless, \Luke, \Janchor}{\FCrule{\#185}{Běda mužům, kteří ženám vládnou.}{\Janchor}}{Věda žije | (10:00) Účtujeme nákupy ve výprodejích | (18:00) Ultima Underworld se vrací | (31:00) Opět rozkouskované Dreamfall Chapters | (38:00) Crytek popírá existenční problémy | (43:00) 30 versus 60 fps | Soutěž o~Sniper Elite 3 a~knížku Země bez zákona}{1:24:33}{36129308}{}
\FC{fight-club-186.mp3}{\FCtitle{\#186}{Koulovačka}}{2014-07-09}{\Luke, \Janchor, \JOlejnik, \OverWatch}{\FCrule{\#186}{Koule si musíte zasloužit.}{\Janchor}}{Gearbox předvádí vlastní MOBA | (17:30) Cliff Blezsinski chystá F2P střílečku | (26:30) Macaté rozšíření pro Hearthstone | (34:35) Progamingová liga přehodnotila svůj postoj k~hráčkám | Soutěž o~plošinovku Light}{1:13:41}{31830352}{}
\FC{fight-club-187.mp3}{\FCtitle{\#187}{Spaaaaaace!}}{2014-07-16}{\Sleepless, \Luke, \Blixa, \OverWatch}{\FCrule{\#187}{Galaxii do konzole nenarveš.}{\Luke}}{Etika youtuberů | (24:30) Free-to-play nemá budoucnost | (34:30) Pár spekulací ohledně dalšího Dishonored | (42:20) Objeví se Elite na konzolích? | Soutěž o~Divinity: Original Sin}{1:10:21}{29421880}{}
\FC{fight-club-188.mp3}{\FCtitle{\#188}{Regulátoři}}{2014-07-23}{\Luke, \Janchor, \Ulrik, \LMacura}{\FCrule{\#188}{Ať vás provází Síla.}{\Janchor}}{LEVEL bez plných her | (7:00) Nový videobsah na Games | (12:00) Databáze tuzemských vývojářů | (25:00) Duchovní nástupce Stalkera odstřihnut Kickstarterem | (34:00) Potřebuje crowdfunding regulovat | (44:45) Jak vychovat online komunitu | Soutěž o~triko Destiny a~klíčenku a~odznáček Elite: Dangerous}{1:23:40}{35392552}{2}
\FC{fight-club-189.mp3}{\FCtitle{\#189}{Lovec trofejí}}{2014-07-30}{\Luke, \Janchor, \Ulrik, \PBulir}{\FCrule{\#189}{Tak dlouho se chodí s gamepadem pro achievementy, až vás předběhne Petr Bulíř.}{\Janchor}}{Jsme paf z~Divinity: Original Sin | (19:00) John Romero věří free-to-play | (32:50) Lov achievementů a~trofejí | (45:40) Crytek potvrdil problémy | Soutěž o~speciální edici Sniper Elite 3}{1:16:02}{33540028}{2}
\FC{fight-club-190.mp3}{\FCtitle{\#190}{Letní sezona}}{2014-08-06}{\Sleepless, \Luke, \Blixa}{\FCrule{\#190}{Nemalujte Hellboye na zeď!}{\Luke}}{Proč jsme v~Levelu zrušili plnou hru | (8:30) Jak pokryjeme Gamescom a~GDC | (18:00) Utekl nový Assassin | (26:30) Activision vykázal více než půlmiliardové příjmy | (38:40) Crytek prodal Homefront | Soutěž o~Ultimu VII + datadisky}{1:17:48}{32465788}{}
\FC{fight-club-191.mp3}{\FCtitle{\#191}{GamesCom 2014 speciál}}{2014-08-14}{\Sleepless, \Luke, \Janchor, \Blixa, \Ulrik}{\FCrule{\#191}{Nejlepší hra s myšima je Mission Impossible.}{\Janchor}}{Nastavujeme zrcadlo Gamescomu | (3:00) Telefonujeme Petrovi | (11:00) Co nového bychom si přáli v~sérii Assassin's Creed | (25:20) Nová hra tvůrce Beyond Good \&{} Evil | (30:00) Pro koho bude No Man's Sky a~pro koho Elite? | (36:00) Del Toro a~Kojima spojili síly | (42:20) Ninja Theory se dušují nezávislostí, Remedy časem, sedmička HoMaM destilátem toho nejlepšího z~celé série | (53:40) Proč bude Alien: Isolation opravdu strašidelná | (62:40) Hrajte s~vlastnoručně poskládaným robotem, za strom nebo s~odvážnou myší | Soutěž o~klíč do uzavřené bety The Witcher Adventure Game}{1:26:23}{38577416}{}
\FC{fight-club-192.mp3}{\FCtitle{\#192}{Návrat exkluzivity}}{2014-08-20}{\Sleepless, \Luke, \Janchor, \Blixa}{\FCrule{\#192}{Bez chatu, mačetu.}{\Janchor}}{Jak to bude s~exkluzivitou nového Tomb Raidera | (12:20) Viktoriánský noir Bloodborne | (22:20) Plusy a~mínusy early access | (31:20) Battlefield Hardline a~militarizace americké policie | Soutěž o~early access do The Escapist | Ze z Plzně je správně\,\dots}{0:54:03}{24426844}{}
\FC{fight-club-193.mp3}{\FCtitle{\#193}{O myších a~hrách}}{2014-08-27}{\Sleepless, \Luke, \Janchor, \Blixa}{\FCrule{\#193}{Come on Amazon, light my Fire!}{\Janchor}}{Twitch byl koupen Amazonem | (16:40) Disney zavřel svou tuzemskou pobočku | (30:45) Gamescom a~jeho vliv na předobjednávky | (38:00) Jak to bylo s~hackováním PSN | Soutěž o~early access Merchants of Kaidan}{0:54:46}{23563564}{2}
\FC{fight-club-194.mp3}{\FCtitle{\#194}{Bouře ve sklenici vody}}{2014-09-03}{\Sleepless, \Luke, \Janchor, \Blixa}{\FCrule{\#194}{Sociální spravedlnost sežere vaše duše.}{\Janchor}}{Velké vlastenecké vítězství | (8:30) GamerGate je selháním komunikace | (33:30) V~Alone in the Dark už nebudete 'alone' | (41:00) Firewatch chce vyprávět příběh v~divočině | (50:40) Voláme výrobci papírových kokpitů | Soutěž o~Flockers na PC nebo Xbox One | Blixovo porušení nepsaného etického kodexu\,\dots}{1:28:21}{39362344}{3}
\FC{fight-club-195.mp3}{\FCtitle{\#195}{O skřetech a~slevách}}{2014-09-10}{\Sleepless, \Luke, \Janchor, \Blixa}{\FCrule{\#195}{Dobrou Notch.}{\Janchor}}{Trable s~bloky | (8:30) Mojang na prodej? | (21:30) Odvrácená strana slevových akcí | (38:20) Smysl remasteru Gabriela Knighta | Soutěž o~Destiny na PS4 nebo Xbox One}{1:19:16}{35914444}{2}
\FC{fight-club-196.mp3}{\FCtitle{\#196}{Prsty v boxu}}{2014-09-17}{\Sleepless, \Luke, \Janchor, \JOlejnik, \VBrozova}{\FCrule{\#196}{Matoni už není.}{\Janchor}}{Tuzemský start Xbox One | (28:00) Nový LEVEL 245 | (34:30) Notch jako George Lucas | (44:40) Pro koho jsou čtvrtí Sims | (53:15) Za obzory Forzy Horizon | Soutěž o~Forzu Motorsport 5 pro Xbox One}{1:34:37}{43709848}{2}
\FC{fight-club-197.mp3}{\FCtitle{\#197}{Shodan z Googlu}}{2014-09-24}{\Sleepless, \Luke, \Janchor, \Blixa}{\FCrule{\#197}{Lepší Fargo v hrsti nežli Schafer na střeše.}{\Janchor}}{Wasteland 2 sklízí úspěch | (24:00) Co udělal Tim Schafer špatně | (47:15) Sbohem, Titáne | (56:00) Evoluce Steamu | Soutěž o~NHL 15 na Xbox One}{1:33:09}{43849096}{}
\FC{fight-club-198.mp3}{\FCtitle{\#198}{Ptáci, lišky, jeleni a~psíci}}{2014-10-01}{\Sleepless, \Luke, \Janchor, \OverWatch}{\FCrule{\#198}{Kdo nekliká s námi, kliká proti nám.}{\Janchor}}{Plány na výroční vysílání | (8:10) Janchor varuje před Clicker Heroes -- nezapínat! | (11:00) Destiny osekaná pár měsíců před vydáním? | (23:30) Němé tváře ve hrách | (39:30) Proč frčí survival | Soutěž o~joystick Logitech Extreme 3D Pro}{1:31:02}{42329428}{2}
\FC{fight-club-199.mp3}{\FCtitle{\#199}{V izolaci}}{2014-10-08}{\Sleepless, \Luke, \Janchor}{\FCrule{\#199}{Ve vesmíru váš debilní křik nikdo neuslyší.}{\Janchor}}{Dvoustovka za rohem | (7:45) Síly i~slabiny virtuálního vetřelce | (30:30) Fallout na kolečkách | (37:00) Hozená černá rukavice | (48:10) Pravidla pro kurátory | Soutěž o~Diablo 3: Reaper of Souls pro Xbox One}{1:30:29}{39251572}{1}
\FC{fight-club-200.mp3}{\FCtitle{\#200}{88 grilovaných kachen}}{2014-10-15}{\Sleepless, \Luke, \Janchor, \Blixa, \Hellboy, \JOlejnik, \Farid, \OverWatch, \PBulir, \KDrda, \LMacura, \Moolker}{\FCrule{\#200}{Dvoustovka je nová osmdesátosmička.}{\Blixa}}{Dort a~šampus | (15:00) Miloš Bohoněk a~Lukáš Macura | (25:00) Ty vole, Petr Bulíř | (52:00) Farid Starman a~Karel Drda | (79:00) Honza Olejník | (109:00) Adam Homola a~Dan Vávra | Soutěž o~tričko Kingdom Come}{2:59:57}{76924060}{2}
\FC{fight-club-201.mp3}{\FCtitle{\#201}{Zakážeme vás!}}{2014-10-22}{\Luke, \Janchor, \Blixa}{\FCrule{\#201}{All hail Gaben!}{\Janchor}}{Sobotní sraz | Nový LEVEL | (11:40) Jade Raymond odchází z~Ubisoftu | (18:30) Je odstranění hry adekvátní reakcí na výhrůžky? | (32:50) Tvůrce Little Big Adventure se vrací se survival projektem | (38:25) Smysl remaků amigáckých klasik | Soutěž o~kód na remaster Gabriela Knighta}{1:37:19}{37269412}{1}
\FC{fight-club-202-b.mp3}{\FCtitle{Bonus round}{Vox popoli, vox dei}}{2014-10-25}{\Sleepless, \Luke, \Janchor, \Blixa, \JOlejnik, \OverWatch, \KDrda, \fans}{\FCrule{\#200 bonus round}{Amy na to má lidi!}{\Janchor}}{Co třeba Kingdom Come (komposty a~včelíny)? | (15:00) Další očekávané tuzemské hry | (22:00) Zaklínačovy slovanské kořeny | (30:00) Trable her z~východu | (43:00) Jakou obtížnost lze tolerovat | (54:00) Volnost do adventur! | (61:00) Co kutí Blizzard}{1:12:11}{32530444}{3}
\FC{fight-club-202.mp3}{\FCtitle{\#202}{Wing-Wing}}{2014-10-29}{\Sleepless, \Luke, \Janchor}{\FCrule{\#202}{Nekrmte impuriální troly.}{\Janchor}}{Jak dopadla veselka | (8:10) Lucasarts konečně na GOG | (18:30) Dva pohledy na remasterovaného Gabriela Knighta | (25:30) Lords of the Fallen Souls? | (34:50) Valkyrie Chronicles se objeví na PC | Soutěž o~FIFA 15 na PS4}{1:11:48}{28534516}{2}
\FC{fight-club-203.mp3}{\FCtitle{\#203}{Sekanica}}{2014-11-05}{\Sleepless, \Luke, \Janchor, \OverWatch}{\FCrule{\#203}{Není hříva jako oháňka.}{\Janchor}}{GTA z~pohledu vlastních očí | (14:30) EA ruší svou MOBA hru, co oznámí Blizzard? | (25:15) Hádáme se s~umělou inteligencí | Soutěž o~lístek na GDS v~Praze}{1:01:22}{27416104}{}
\FC{fight-club-204.mp3}{\FCtitle{\#204}{Na just!}}{2014-11-12}{\Sleepless, \Luke, \Blixa, \PBarak}{\FCrule{\#204}{Na každou Alfu Romeo patří zemědělský pluh.}{\Luke}}{Jak bylo a~jak bude na GDS | (23:55) Nová mobilní verze Games a~nový LEVEL | (32:00) Trable s~embargem na AC: Unity | Snaha o odvedení pozornosti of Gamergate -- narážka na Hovory z Lán | (44:20) Abandonware na cestě k~legalitě | (53:40) Just Cause se vrací | Soutěž o~Forza Horizon 2 pro Xbox One / 360}{1:31:54}{39896152}{}
\FC{fight-club-205.mp3}{\FCtitle{\#205}{Deník horala}}{2014-11-19}{\Sleepless, \Luke, \Janchor}{\FCrule{\#205}{Nedráždi slona bosou nohou.}{\Janchor}}{Povedený ročník GDS | (11:00) Far Cry ve stopách Assassina | (32:10) Elite striktně online | (45:05) Kvalitní scénář nadevše | Soutěž o~shmup The Blue Flamingo}{1:26:00}{36524644}{}
\FC{fight-club-206.mp3}{\FCtitle{\#206}{Naprosto retardovaný}}{2014-11-26}{\Luke, \Blixa, \Ulrik, \MRota}{\FCrule{\#206}{Nejlepší cestování je ve vlaku.}{\Luke}}{Youtuberova cesta ke slávě | (11:00) Zkušenosti z~kooperace v~Unity | (27:00) Steam zpřesňuje pravidla | (36:50) Česká hra Dark Train zaujala originalitou | (43:00) Záplava dotazů | Soutěž o~Age of Wonders III}{1:45:05}{47183740}{2}
\FC{fight-club-207.mp3}{\FCtitle{\#207}{Řád a~neřád}}{2014-12-10}{\Sleepless, \Luke, \Blixa}{\FCrule{\#207}{Pravidlo je dobrý.}{\Luke}}{Nový LEVEL | (11:00) Jak bylo v~Londýně na prezentaci Sony | (36:40) Zelda předvádí svět, Zaklínač výlet do toho svého odkládá | Soutěž o~Far Cry 4 pro PC}{1:18:58}{32684452}{2}
\FC{fight-club-208.mp3}{\FCtitle{\#208}{Rapture na Moravě}}{2014-12-17}{\Sleepless, \Luke, \Janchor}{\FCrule{\#208}{Ondrášů dělé!}{\Janchor}}{Ruský chlapec se ptá svého otce. "Tati, dáš mi pět rublů na mléko a~chleba?" Otec odpoví. "Dvacet rublů? Na co sakra potřebuješ padesát rublů!" | Sváteční provoz | (7:30) Amiibo pod stromeček | (18:50) K~Oculusu přibydou i~virtuální ruce | (40:50) Jaké hry k~Ježíšku | Soutěž o~herní zápisník od Sony}{1:16:19}{31252804}{2}
\FC{fight-club-209.mp3}{\FCtitle{\#209}{3+1}}{2015-01-07}{\Sleepless, \Luke, \Janchor, \Blixa}{\FCrule{\#209}{Piráti zásadně po schodech.}{\Janchor}}{Ankety, kam se podíváš | (10:30) Talos Principle nachytala piráty | (19:25) Vzhůru ke hvězdám | (28:00) iOS zaplaven hrami | (36:20) Omezení prodejů na vlastní triko}{1:10:33}{28737880}{}
\FC{fight-club-210.mp3}{\FCtitle{\#210}{Středověcí inženýři}}{2015-01-14}{\Sleepless, \Luke, \LMacura}{\FCrule{\#210}{Inženýři jsou nejlepší.}{\Luke}}{Kolik obálek bude mít jubilejní LEVEL? | (11:20) Jak to bude s~prodejem her prostřednictvím vlastních stránek | (24:20) Nejlepší herní ovladače | (35:15) Speedruny pro charitu | (44:05) Středověcí inženýři}{1:18:45}{32085124}{2}
\FC{fight-club-211.mp3}{\FCtitle{\#211}{Remasterováno}}{2015-01-21}{\Sleepless, \Luke, \Janchor, \Blixa, \JOlejnik, \Draken}{\FCrule{\#211}{Pravdu má Olejník.}{\Janchor}}{20 let Levelu | (19:20) Je kritika českých her oprávněná? | (29:25) Sid Meier se vytasil s~flotilou kosmických lodí | (35:50) Které remastery mají smysl | Soutěž o~10 kódů do bety Heroes of the Storm}{1:22:58}{36456604}{}
\FC{fight-club-212.mp3}{\FCtitle{\#212}{Elitáři}}{2015-01-28}{\Sleepless, \Janchor, \Blixa, \Ulrik}{\FCrule{\#212}{Přes den zelinářem, v noci pirátem.}{\Janchor}}{Třikrát Elite: Dangerous | (35:10) Proč Ubisoft mazal koupené hry | (45:40) Windows 10 a~hry | HALDA DOTAZŮ}{1:45:01}{42574264}{1}
\FC{fight-club-213.mp3}{\FCtitle{\#213}{Pod zemí}}{2015-02-04}{\Sleepless, \Luke, \Janchor, \OverWatch}{\FCrule{\#213}{Podcast macht frei.}{\Janchor}}{Ultima Underworld se vrací ve velkém stylu | (19:10) Renoir nebude jen další černobílou plošinovkou | (27:15) Oprašte VHS kazety! | (33:40) Týden s~Blizzardem}{1:14:46}{31879384}{}
\FC{fight-club-214.mp3}{\FCtitle{\#214}{Filmostroj}}{2015-02-11}{\Sleepless, \Luke, \Blixa, \MrHlad}{\FCrule{\#214}{Když točit film, tak v Kladně.}{\Luke}}{Jak dopadne druhý pokus o~zfilmování Hitmana | (8:00) A~co první pokus o~Minecraft, Ratcheta \&{} Clanka nebo Splinter Cell? | (28:45) Čekání na Warcraft | (47:20) Jsou hry příliš složité na naučení? | (54:20) Trable s~hodnocením}{1:31:03}{37725028}{}
\FC{fight-club-215.mp3}{\FCtitle{\#215}{16 odstínů her}}{2015-02-18}{\Sleepless, \Luke, \Cerv}{\FCrule{\#215}{S Amigou na věčné časy.}{\Luke}}{Vzpomínáme s~Červem | (11:10) Jak funguje Startovač | (17:50) Temné duše, temnější kobky | (29:25) Hrajeme s~Červem | (36:20) Otázky, samé otázky}{1:11:25}{32655832}{1}
\FC{fight-club-216.mp3}{\FCtitle{\#216}{Blizzard free edition}}{2015-02-25}{\Luke, \Janchor, \Blixa, \OverWatch}{\FCrule{\#216}{Nevěřte pohádkám.}{\Janchor}}{Námořní LEVEL 250 | (8:25) Když si veteráni základají nezávislá studia | (20:45) Co čekat od GDC v~San Franciscu | (36:00) Pro koho je nová verze Nintenda 3DS | (41:35) Lze soucítit s~Peterem Molyneuxem?}{1:24:50}{37508848}{1}
\FC{fight-club-217.mp3}{\FCtitle{\#217}{Sám před vraty}}{2015-03-04}{\Luke, \Janchor, \Ulrik}{\FCrule{\#216}{Žijte dlouho a~blaze.}{\Janchor}}{GDC hardwarová | (14:30) Nová konkurence pro Oculus Rift, parní stroje a~hydry | (35:45) Vzpomínky na Spocka}{1:27:16}{39402880}{}
\FC{fight-club-217-GDC.mp3}{\FCtitle{GDC 2015 speciál}{}}{2015-03-04}{\Sleepless, \Blixa}{}{Martin a~Petr se hlásí z~Game Developers Conference v~San Franciscu se svými postřehy a~názory.}{0:31:39}{12849604}{}
\FC{fight-club-218.mp3}{\FCtitle{\#218}{Kouzelná tlačítka}}{2015-03-11}{\Sleepless, \Luke, \Blixa, \LMacura}{\FCrule{\#218}{Kdo šetří, má za Notche.}{\Luke}}{Shrnutí GDC od navrátilců | (8:00) Bude se v~Turecku banovat Minecraft? | (22:50) Descent se vrací | (32:00) Enginy zdarma z~pohledu vývojáře | (55:00) Amazon zlanařil autory zrušeného Prey 2 | (66:45) Valník dotazů}{1:46:46}{45303748}{2}
\FC{fight-club-219.mp3}{\FCtitle{\#219}{Výpary pixelů}}{2015-03-18}{\Sleepless, \Luke, \Blixa, \PBulir}{\FCrule{\#219}{Level TV byla vždycky napřed.}{\Luke}}{Co hraje Petr Bulíř | (18:00) Nintendo zbrojí do útoku na mobily | (34:00) Do kin se valí Pixely}{1:06:36}{29181868}{1}
\FC{fight-club-220.mp3}{\FCtitle{\#220}{Vyvádění s Honzou Vrobelem}}{2015-03-25}{\Sleepless, \Luke, \Janchor, \JVrobel}{\FCrule{\#220}{Nejoblíbenější jídlo Jana Vrobela -- svíčková.}{\Janchor}}{Nový LEVEL 251 | (6:30) Zenový Vrobel vyznává lásku Bloodborne | (34:00) Konami vs Kojima | (51:10) Perspektivy virtuálních světů | Soutěž o~vymazlenou edici Bloodborne}{1:17:36}{34648324}{1}
\FC{fight-club-221.mp3}{\FCtitle{\#221}{Guilty pleasure}}{2015-04-01}{\Sleepless, \Luke, \Janchor, \OverWatch}{\FCrule{\#221}{Když guilty pleasure, tak jenom trošku.}{\Janchor}}{Uklízečem v~Do Not Cross | (15:00) Crowfall crowluje crowdfundingu | (26:40) Co jsme se dozvěděli z~interního auditu Valve}{1:15:09}{33003952}{}
\FC{fight-club-222.mp3}{\FCtitle{\#222}{Muž, chlapec a~myš}}{2015-04-08}{\Sleepless, \Luke, \Blixa}{\FCrule{\#222}{Lepší než žádné pravidlo je jen divácké pravidlo.}{\Luke}}{Herní bulvár od Tiscali (booom.cz) | (4:00) Oznámení nového Deus Ex | (16:30) Brány Skeldalu se otevřou potřetí | (23:00) Finále kampaně Dungeons of Aledorn | (30:00) Klady a~zápory Nvidia SHIELD | (41:30) Co vyčmuchal Steam Spy | Soutěž o~sběratelskou edici Blackhole}{1:19:28}{33428788}{}
\FC{fight-club-223.mp3}{\FCtitle{\#223}{Ta naše kapela}}{2015-04-15}{\Sleepless, \Luke, \Blixa, \JOlejnik}{\FCrule{\#223}{Fight club je makový a~ne takový.}{\Luke}}{Úvod s~GTA ve 4K | (6:30) Nahý zadek v~Halo 2 | (15:00) Návrat Kultury kniplů? | (33:00) Chválíme třetí Samorost | (38:00) Ron Gilbert říká, že adventurám škodí moc příběhu | (47:30) Dotazy z~mailu i~chatu}{1:09:34}{29293628}{2}
\FC{fight-club-224.mp3}{\FCtitle{\#224}{Staré lásky nerezaví}}{2015-04-22}{\Sleepless, \Luke, \Janchor, \Bety}{\FCrule{\#224}{Když zdi, tak zaflusaný.}{\Janchor}}{Nový LEVEL 252 | (6:10) Quo vadis, Mass Effect? | (18:25) Jak bylo na Star Wars Celebration | (53:15) Jak mohl vypadat nový Alan Wake | (44:00) Shadow of the Beast vypadá všelijak | Soutěž o~Shadowrun Returns a~Dragonfall - Director's Cut}{1:30:45}{38572792}{2}
\FC{fight-club-225.mp3}{\FCtitle{\#225}{Janda ve Vlašimi}}{2015-04-29}{\Sleepless, \Luke, \Janchor, \Blixa}{\FCrule{\#225}{Radši vařit z vody, než-li platit mody.}{\Janchor}}{Just Cause 3 se předvádí v~nejlepším světle, jak dopadne Mad Max? | (12:50) Silent Hills zrušeno | (25:55) Fiasko s~placený mody na Steamu | (57:40) Vliv všudypřítomných nápověd na herní design | Soutěž o~progamerské brýle}{1:28:18}{38090716}{}
\FC{fight-club-226.mp3}{\FCtitle{\#226}{Ouné!}}{2015-05-06}{\Luke, \Janchor, \Blixa}{\FCrule{\#226}{Nejlepší prognostik je Martin Bach.}{\Janchor}}{Pozvánka na víkend | (7:40) Ouya nad hrobem | (18:20) Velký potenciál Hololens | (39:45) Je nová služba od GOG konkurencí Steamu? | Soutěž o~klíče do Heroes of the Storm}{1:15:37}{32994268}{}
\FC{fight-club-227.mp3}{\FCtitle{\#227}{Assassin's Skate}}{2015-05-13}{\Sleepless, \Luke, \Janchor, \JOlejnik}{\FCrule{\#227}{Není Stig jako Stiga.}{\Janchor}}{Assassin's Creed dělá brajgl v~Londýně | (20:30) Kam se vydá nový Tony Hawk | (30:40) Co nového může přinést dvojka Audiosurfu | Soutěž o~Total War: Shogun 2 - Fall of the Samurai}{1:15:56}{32717464}{}
\FC{fight-club-228.mp3}{\FCtitle{\#228}{Zaklínačování}}{2015-05-20}{\Sleepless, \Luke, \Blixa, \OverWatch}{\FCrule{\#228}{Každý hráč trhá kalendář.}{\Luke}}{Nový LEVEL se Zaklínačem | (11:50) Jak vypadá a~jak se hraje nový Zaklínač | (28:25) Space Engineers nabízejí zdrojový kód | (38:45) Dobří bardové se vracejí | (47:40) Divinity: Original Sin dostane rozšířenou edici}{1:14:10}{30978016}{}
\FC{fight-club-229.mp3}{\FCtitle{\#229}{Keepmaster}}{2015-05-27}{\Sleepless, \Luke, \Ulrik, \LMacura}{\FCrule{\#229}{Necheatujte nebo zemřete.}{\Luke}}{Co se kutí v~Cinemaxu | (16:45) Co kutí 2K | (22:00) Je potřeba rychlosti opravdu potřeba? | (34:10) Kontroverze v~Guild Wars aneb jak se vypořádat s~hackery | Soutěž o~retro tričko Necromania}{1:16:30}{33995068}{2}
% 2015-06-07
\FC{fight-club-230.mp3}{\FCtitle{\#230}{Návrat emzáků}}{2015-06-03}{\Sleepless, \Luke, \Janchor, \OverWatch}{\FCrule{\#230}{Aleš -- nejlepší imaginární přítel Petra Poláčka.}{\Janchor}}{Co si slibovat od nového Falloutu | (12:00) Co si slibovat od nového XCOM | (24:40) Problematické vracení her na Steamu}{1:25:51}{39027904}{1}
\FC{fight-club-231.mp3}{\FCtitle{\#231}{Klid před bouří}}{2015-06-10}{\Sleepless, \Luke, \OverWatch}{\FCrule{\#231}{Kdo nastoupí k Bohmece, z toho bude zombie.}{\Luke}}{Skvělé prodeje nového Zaklínače | (20:50) Žhavíme na E3 | (40:50) Duchovní otec Castlevanie drtí Kickstarter}{1:37:55}{43596124}{}
\FC{fight-club-231-E3-1.mp3}{\FCtitle{E3 2015 speciál \#1}{}}{2015-06-17}{\Janchor, \Ulrik, \OverWatch, \LMacura}{}{Hodnotíme ty nejzajímavější hry oznámené na tiskových konferencích Microsoftu, Sony, EA a~Ubisoftu.}{0:59:04}{26466388}{}
\FC{fight-club-231-E3-2.mp3}{\FCtitle{E3 2015 speciál \#2}{}}{2015-06-18}{\Janchor, \VHrincar, \JIlavsky, \MRota}{}{O dosavadním průběhu E3 2015 a~nejzajímavějších hrách výstavy s~našimi hosty.}{1:40:47}{44341612}{2}
\FC{fight-club-231-E3-3.mp3}{\FCtitle{E3 2015 speciál \#3}{}}{2015-06-19}{\Janchor, \JOlejnik, \KDrda}{}{Třetí díl povídání o~E3 2015 s~našimi hosty.}{1:03:57}{27286540}{3}
\FC{fight-club-231-E3-4.mp3}{\FCtitle{E3 2015 speciál \#4}{}}{2015-06-20}{\Janchor, \OverWatch, \FF, \Baty}{}{Poslední edice speciálního E3 Fight Clubu.}{1:19:13}{31603336}{3}
\FC{fight-club-232.mp3}{\FCtitle{\#232}{Nebijte pejsky!}}{2015-06-24}{\Sleepless, \Luke, \Moolker}{\FCrule{\#232}{Nebuďte Last nebo budete Lost.}{\Luke}}{Nový LEVEL s~Battlefrontem | (8:40) Rozporuplné dojmy z~Video Games Live | (22:30) Ohlédnutí za E3: hry, hry, hry! | (53:00) Neortodoxní herní filmy}{1:28:28}{41498080}{}
\FC{fight-club-233.mp3}{\FCtitle{\#233}{Plachty napnuty!}}{2015-07-01}{\Sleepless, \Luke, \Ulrik, \OverWatch}{\FCrule{\#233}{Francouzštinu nechte na Homolovi.}{\Luke}}{Tale of Tales zavřeli krám | (24:10) PlayStation opanovala evropský trh | (29:55) Renoir je na Kickstarteru | (40:00) Uvolněná letní atmosféra}{1:34:00}{42247960}{}
\FC{fight-club-234.mp3}{\FCtitle{\#234}{S pánem i bez něj}}{2015-07-08}{\Luke, \Janchor, \Blixa, \LMacura}{\FCrule{\#234}{Mluv o Pánu Ježíši s dalšími lidmi.}{\Janchor}}{Vývojáři lámou hůl nad mobilními hrami | (24:15) Síla jejího příběhu | (40:15) Je PewDiePie nemorálně bohatý?}{1:21:05}{36883780}{}
\FC{fight-club-235.mp3}{\FCtitle{\#235}{Alfa vývojář}}{2015-07-15}{\Sleepless, \Blixa, \Hellboy, \Ulrik}{\FCrule{\#235}{Díky Iwata-san.}{\Blixa}}{S Danem, o~všem! | Je na tom Nintendo špatně? | (21:40) JRC změnilo majitele | (34:55) Co nového v~Kingdom Come}{2:01:45}{53712340}{3}
\FC{fight-club-236.mp3}{\FCtitle{\#236}{Gundam sem, Gundam tam}}{2015-07-22}{\Sleepless, \Luke, \Janchor, \OverWatch}{\FCrule{\#236}{Polidštěte si svoje Schaffery.}{\Janchor}}{Oprašování zašlých svitků | (17:45) Bluesové kulky | (22:45) Online svět, mířící k~neodvratné apokalypse | (29:50) Dokument osvětluje peripetie herního vývoje | Valník otázek}{1:34:01}{43558792}{}
\FC{fight-club-237.mp3}{\FCtitle{\#237}{IT Mafie}}{2015-07-29}{\Sleepless, \Luke, \Ulrik, \OverWatch}{\FCrule{\#237}{Nebojte se mafie.}{\Luke}}{Co bychom si přáli od třetí Mafie | (23:40) Xzone fúzuje s~Gameexpress | (39:00) Cesta na Blizzcon vede přes Prahu}{1:28:55}{40340032}{}
\FC{fight-club-238.mp3}{\FCtitle{\#238}{Na GamesCom a~zase zpátky}}{2015-08-12}{\Sleepless, \Luke, \Blixa, \DRynes}{\FCrule{\#238}{Chceš-li hardcode hráčem být, musíš tetování mít.}{\Luke}}{Vášně kolem třetí Mafie | (22:20) Návrat Dark Souls | (29:50) Quantum Break pochyb (skoro) zbavený | (40:00) Petr mluví o~Mount \&{} Blade 2 | (47:00) Dračí radovánky ve Scale Bound | (55:00) Dvojnásobná dávka dotazů}{1:41:44}{44310344}{2}
\FC{fight-club-239.mp3}{\FCtitle{\#239}{16 odstínů podcastu}}{2015-08-19}{\Luke, \Ulrik, \OverWatch}{\FCrule{\#239}{Stavění zdí ve vesmíru, to se pěkně prodraží.}{\Luke}}{Odchod Clinta Hockinga z~Amazonu, odhalení crowdfundingové platformy Fig a~další šokující odhalení}{1:23:58}{35846024}{}
\FC{fight-club-240.mp3}{\FCtitle{\#240}{Přiznání Poláčka}}{2015-08-26}{\Sleepless, \Luke, \Ulrik}{\FCrule{\#240}{Zloděj her jezdí zásadně na koloběžce.}{\Luke}}{Těšíme se na Divinity 2 | (14:30) Herní ekologie v~ECO | (23:00) Hudební magie v~Branách Skeldalu | (31:20) Útok titánů | (40:30) Mraky dotazů}{1:04:51}{28596380}{1}
\FC{fight-club-241.mp3}{\FCtitle{\#241}{\#MBPPPD}}{2015-09-02}{\Sleepless, \Luke, \Blixa}{\FCrule{\#241}{Kuželky hrejte zásadně s trpaslíky.}{\Luke}}{Chceme nové funkce Kickstarteru | (6:00) Trpasličí kuželky | (12:00) Zlatý valoun jménem Witcher 3 | (27:30) SecuROM a~Safedisc udeřily ze záhrobí | (35:20) Dotazová sekce | (47:10) Zásek u~Metal Gearu}{1:16:21}{34749068}{}
\FC{fight-club-242.mp3}{\FCtitle{\#242}{Metař metal meta}}{2015-09-09}{\Luke, \Blixa, \JOlejnik, \Ulrik}{\FCrule{\#242}{Sprchujte se častěji než Snake.}{\Luke}}{Kde bude FC \#250? | (05:57) Mikrotransakce v~MGS V~a velkých hrách | (19:32) Theme park Ubi Land | (28:08) Bulvarizace herních médií | (43:53) Simulátor mravence již brzy na vašich monitorech | (51:09) Dotazy}{1:34:10}{40810136}{3}
\FC{fight-club-243.mp3}{\FCtitle{\#243}{Král a~babety ve studiu}}{2015-09-16}{\Luke, \Bety, \Baty, \JKral}{\FCrule{\#243}{Vivaldiho do každého domu.}{\Luke}}{03:57 Baty a~DayZ | 09:05 Jirka a~Until Dawn | 16:49 Babety a~Youtube | 24:41 Jirka na Vidconu | 33:09 Bety na PAXu | 56: 44 Dotazy}{1:34:58}{42391148}{}
\FC{fight-club-244.mp3}{\FCtitle{\#244}{Za oponou}}{2015-09-23}{\Sleepless, \Blixa, \JOlejnik, \JVrobel}{\FCrule{\#244}{Nejlepší hobby má Honza Vroby.}{\Blixa}}{Speciální moderátor a~nový LEVEL | (9:10) Průnik her a~reality v~Železné oponě | (28:00) Situace kolem Konami | (42:00) Jak dobrou AI ve hrách chceme? | (57:30) Bez oslího můstku k~Aurionu | (1:06:00) Dotazy z~mailu i~chatu}{1:42:13}{42815814}{2}
\FC{fight-club-245.mp3}{\FCtitle{\#245}{Tour de Studio}}{2015-09-30}{\Sleepless, \Luke, \Ulrik}{\FCrule{\#245}{Na čarodějnice kladivo a~na Poláčka krabici.}{\Luke}}{Armikrog je zde | (06:57) Unbox je o~krabici | (11:35) Strategie Battletech | (22:46) Naporcovaný nový Hitman | (36:36) Dotazy}{1:27:24}{41017748}{1}
\FC{fight-club-246.mp3}{\FCtitle{\#246}{Pravěk}}{2015-10-07}{\Sleepless, \Luke, \Blixa, \OverWatch}{\FCrule{\#246}{Poslouchejte pana Humla nebo vás zašlápne mamut, vy holoto.}{\Luke}}{Pac-Man s~raptory | (06:17) Road to Blizzcon a~dojmy z~Overwatch | (20:11) Naše reakce na Far Cry Primal | (27:23) Předplatíte si hry na Humble Monthly? | (36:15) Česká Novus Inceptio spuštěna | (40:22) Dotazy}{1:05:19}{29149823}{}
\FC{fight-club-247.mp3}{\FCtitle{\#247}{Force is strong with us!}}{2015-10-14}{\Sleepless, \Luke, \OverWatch, \MRota}{\FCrule{\#247}{Bez polštáře nelez mezi rváče.}{\Luke}}{Nový Level! | (03:36) Náš názor na Heroes VII | (10:09) Jak hodnotíme výstavu For Games | (17:09) Zážitky ze Star Wars: Battlefront | (35:00) Underworld Ascendant? Fight! | (44:06) Očekávání hrůz v~The Park | (50:32) Dotazy}{1:31:54}{43183827}{1}
\FC{fight-club-248.mp3}{\FCtitle{\#248}{Návraty do budoucnosti}}{2015-10-21}{\Sleepless, \Luke, \Blixa, \DRynes}{\FCrule{\#248}{Své plutonium dobře ukrývejte. a~Vy víte proč\,\dots}{\Luke}}{Překvapení pro Petra | (03:30) Návrat do budoucnosti | (06:36) HALO 5 | (15:49) One Life, The Flock a~Upsilon Circuit | (31:28) Nintendo NX | (42:30) Dotazy - SW: Battlefront, virtuální realita, Zaklínač, Metal Gear Solid V~a další}{1:19:07}{34350178}{2}
\FC{fight-club-249.mp3}{\FCtitle{\#249}{Dárky strýčka Sonyho}}{2015-10-28}{\Sleepless, \Luke, \JOlejnik, \PBulir}{\FCrule{\#249}{XBox patří do lednice.}{\Luke}}{Destiny podle Petra Bulíře | (09:27) JRPG podle Petra Bulíře | (14:41) Nová hra Davida Cage: Detroit | (20:44) Wild od Michela Ancela | (25:15) Robinson: The Journey a~virtuální realita | (31:46) Gran Turismo Sport | (35:57) Horizon: Zero Dawn | (41:50) Stručně o~Assassin's Creed: Syndicate | (48:01) Dotazy}{1:23:21}{37766979}{}
\FC{fight-club-250.mp3}{\FCtitle{\#250}{Brněnský speciál}}{2015-11-07}{\Sleepless, \Luke, \Blixa, \JOlejnik, \Tafrob, \VGersl, \MMastny}{\FCrule{\#250}{Dost bylo pravidel!}{\Luke}}{06:27 Kam kráčíte české hry? | (18:59) Vláďa Geršl (Future Factory) \&{} Martin Mastný (Novus Inceptio) o~svých zkušenostech | (49:08) Hardcore vs Casual vs Mainstream a~co my s~tím | (72:00) Problematika délky her | (87:00) Na co se těšíme \&{} Hra roku \&{} Tafrob}{1:36:23}{42534248}{2}
\FC{fight-club-251.mp3}{\FCtitle{\#251}{Mouchou chtěl bych být\,\dots}}{2015-11-11}{\Sleepless, \Luke, \OverWatch, \JValenta}{}{Fall of the Fly: Hra Honzy Valenty | (23:29) Film Warcraft | (34:39) Overwatch o~placeném Overwatch | (43:41) Dotazy}{1:16:54}{36075319}{}
\FC{fight-club-252.mp3}{\FCtitle{\#252}{Project Červ}}{2015-11-18}{\Sleepless, \Luke, \Blixa, \Cerv}{\FCrule{\#252}{Pravidla budou.}{\Blixa}}{Nekonečný labůžo \&{} Nový Level | (10:40) Nový Sensible Soccer | (20:54) Fallout 4 dle Červa | (34:10) Nadšení z~InfiniFactory | (40:24) Návrat Everquestu v~Project 1999 | (50:23) Dotazy}{1:24:00}{38912135}{2}
\FC{fight-club-253.mp3}{\FCtitle{\#253}{Delegace z Číny}}{2015-11-25}{\Sleepless, \Luke, \Blixa, \Cocik}{\FCrule{\#253}{Když nemůžete do Číny, tak si prostě počkejte až Čína dojede za váma.}{\Luke}}{Čočík o~Číně | (58:30) Dotazy}{1:31:22}{46588889}{3}
\FC{fight-club-254.mp3}{\FCtitle{\#254}{Krásné české hry}}{2015-12-02}{\Sleepless, \Luke, \PBuday}{\FCrule{\#254}{Pozor na trans a~trauma.}{\Luke}}{Shrnutí GDS 2015 | (12:54) Western Blood Will be Spilled | (15:26) Mašinky Logistics | (21:18) Gerontofilní Planetoid Pioneers | (24:04) Rubačka Company of Vikings | (24:45) Rytmik Ultimate od Cinemaxu | (27:54) Experiment Dimension Brothers | (29:59) Ghost Theory pro VR | (37:31) Granátovačka Grenade Madness | (41:38) Chameleon Run pro twitch gamery | (46:14) Remaster System Shock | (58:22) Dotazy}{1:29:47}{43177279}{}
\FC{fight-club-255.mp3}{\FCtitle{\#255}{Pláč a~smích}}{2015-12-09}{\Luke, \Blixa, \Fiola}{\FCrule{\#255}{Nerušte odběry!}{\Luke}}{Blackhole pro konzole | (12:02) Hodnocení Dreamfall Chapters | (19:16) Psychonauts 2 | (23:58) Jak funguje FIG.co? | (38:31) Crowdfunding u~nás | (53:50) Remaster Day of the Tentacle | (58:08) Dotazy}{1:35:17}{43741499}{}
\FC{fight-club-256.mp3}{\FCtitle{\#256}{Ztracen v alfě}}{2015-12-16}{}{\FCrule{\#256}{X je nové Z.}{\Luke}}{Představení Levelu 259 | (09:23) Dojmy ze Star Citizen alphy 2.0 | (23:45) Lezecký simulátor The Climb pro virtuální realitu | (39:25) Dojmy ze strategie Valhalla Hills | (43:34) Dojmy z~karetní strategie Cards \&{} Castles | (52:18) Dojmy z~Xenoblades Chronicles X | (58:51) Dotazy}{1:29:57}{38010858}{2}
\FC{fight-club-257.mp3}{\FCtitle{\#257}{Urnový háj}}{2015-12-23}{\Sleepless, \Luke, \Blixa, \OverWatch}{\FCrule{\#257}{Jůlie, pojď sem, pojď sem, pojď sem, pojď sem,\,\dots}{\Luke}}{Co budeme hrát přes svátky | (09:47) Nová hra s~barbarem Conanem? | (20:12) Homeworld: Deserts of Kharak | (25:05) Filmové okénko: Star Wars: Síla se probouzí | (63:56) Dotazy}{1:35:47}{39934423}{}
\FC{fight-club-258.mp3}{\FCtitle{\#258}{Bob \&{} Bobek Interactive}}{2016-01-06}{\Sleepless, \Luke, \Ulrik, \DRynes}{\FCrule{\#258}{Když začnete hýbat kamerou, tak pozor na pancéřované mloky.}{\Luke}}{Co jsme hráli přes svátky | (05:21) Naši favoriti za rok 2015 | (14:05) David žije v~Rainbow 6 Siege | (22:42) Film podle Dragon's Lair | (27:30) Soulož ve Čtyřlístku | (32:01) Aleš má radost z~UFC 2 | (38:41) Dotazy | {\em (85:55) Tatarku na sebe lejt nebudeme$\cdot$}\/)}{1:35:46}{37574672}{2}
\FC{fight-club-259.mp3}{\FCtitle{\#259}{Moravský les}}{2016-01-13}{\Sleepless, \Luke, \Ulrik, \OverWatch}{\FCrule{\#259}{Držte se a~važte si turistického značení!}{\Luke}}{Oculus Rift je tu! | (24:56) Honza Kavan o~novince Someday You'll Return | (44:05) Football Manager 2016 | (50:25) The Fall of the Dungeon Guardians | (59:52) Dotazy}{1:21:08}{32834912}{}
\FC{fight-club-260.mp3}{\FCtitle{\#260}{Slavíme!}}{2016-01-20}{\Sleepless, \Luke, \Blixa, \Ulrik}{\FCrule{\#260}{Poslouchejte pořád pořad Fight club.}{\Luke}}{10 let podcastování | (04:20) Level 260 | (15:37) Iron Maiden dělají hru! | (22:28) Japonská brutalitka Danganronpa: Trigger Happy Havoc | (33:02) Cenzura ve Westport Independent | (40:58) Sony se pokusila patentovat "let's play" | (46:12) Dotazy}{1:12:49}{31194176}{1}
\FC{fight-club-261.mp3}{\FCtitle{\#261}{Pohoda nebo smrt}}{2016-01-27}{\Luke, \JOlejnik, \OverWatch, \PBulir}{\FCrule{\#261}{Drbati tygra je třeba.}{\Luke}}{Oculus Rift a~HTC Vive na vlastní kůži | (13:21) Petr Bulíř o~The Witness | (30:55) Česká hra Ghost Theory | (38:27) Česká hra Planet Nomads | (44:23) Český dokument New Game o~českých vývojářích | (50:57) Netflix a~EA Access | (59:23) Dotazy}{1:27:59}{39625232}{1}
\FC{fight-club-262.mp3}{\FCtitle{\#262}{Rytmizování}}{2016-02-03}{\Luke, \Blixa, \Ulrik, \OverWatch, \LMacura}{\FCrule{\#262}{Kdo zlobí, ten dostane banánem.}{\Luke}}{Lukáš Macura představuje Rytmik | (16:35) Aleš Smutný komentuje změny v~Hearthstone | (33:18) Adam Homola rozebírá předběžný přístup na her u~GOG  | (48:44) Martin Bach vysvětluje E3 s/bez Electronic Arts | {\em (58:00) Ignorovaný Gamestar byl na Moravě populární}\/ | (57:18) Dotazy}{1:28:02}{39659792}{2}
\FC{fight-club-263.mp3}{\FCtitle{\#263}{Fiolův klub}}{2016-02-10}{\Sleepless, \Luke, \Baty, \Fiola}{\FCrule{\#263}{Dobrý den je, když je Fiola spokojený.}{\Janchor}}{Fight o~XCOM 2 | (18:45) Dream job: Baty testuje DayZ (a~o~testování her obecně) | (39:48) Vývojářská odysea Blackhole | (53:11) Amazon rozdává (cry)engine | (61:01) Dotazy}{1:47:58}{50190728}{2}
\FC{fight-club-264.mp3}{\FCtitle{\#264}{Červí díra}}{2016-02-17}{\Luke, \Blixa, \Ulrik, \Cerv}{\FCrule{\#264}{Dávejte si pozor na\,\dots bad sektory!}{\Luke}}{Hlášení o Levelu | (05:13) Nové herní projekty na Startovači | (09:26) Textové RPG Sanctuary | (13:38) U Fuky doma \&{} SpaceCom | (17:04) Dvakrát Descent: Underworld a~Overload | (23:34) Assassin's Creed letos vynechá a~nás to nemrzí | (30:33) (Ne)Exkluzivita a~naše očekávání Quantum Break | (44:25) Červ a~Aleš hrají Warhammer: Vermintide | (51:13) Dotazy}{1:28:35}{39987068}{1}
\FC{fight-club-265.mp3}{\FCtitle{\#265}{Warren Bros.}}{2016-02-24}{\Sleepless, \Luke, \Blixa, \Ulrik}{\FCrule{\#265}{Pozor na kavárenské fotbalisty!}{\Luke}}{Level 261 je tady! | (10:00) Dino Dini vrací fotbal hráčům + fotbaly na mobily | (23:39) Warren Spector se vrací k System Shock | (35:43) HTC Vive: debata o ceně a~virtuální realitě | (43:33) Dotazy}{1:15:54}{33414260}{}
\FC{fight-club-266.mp3}{\FCtitle{\#266}{Back to the Microsoft}}{2016-03-02}{\Sleepless, \Luke, \OverWatch}{\FCrule{\#266}{Hry nejsou jenom pro hráče.}{\Luke}}{HoloLens jde mezi vývojáře | (10:42) Budoucnost konzolí a~her podle Microsoftu | (28:15) Rozruch kolem That Dragon, Cancer | (41:39) Dotazy | {\em (63:00) Nové brány Skelalu}\/}{1:08:41}{31065008}{}
\FC{fight-club-267.mp3}{\FCtitle{\#267}{Kumáni v zádech}}{2016-03-09}{\Blixa, \OverWatch, \MRota}{\FCrule{\#267}{Když diverzita, tak jedině ve VR.}{\Blixa}}{Jdeme na Supermana | (05:30) Hrajeme Factorio a~betu Kingdom Come | (14:50) Konec Lionheadu | (24:15) Intel Extreme Masters a~budoucnost progamingu | (37:30) Game Developers Conference 2016 - na co se (ne)těšit | (47:00) Dotazy}{1:19:51}{34542197}{}
\FC{fight-club-268.mp3}{\FCtitle{\#268}{Bez dozoru}}{2016-03-16}{\Luke, \JOlejnik, \Ulrik, \OverWatch}{\FCrule{\#268}{Bez dožíněk se nedá žít.}{\Luke}}{O čem bude Oddworld: Soulstorm? (03:52 ) | CryEngine zdarma (07:26) | Filmový Portal od JJ Abramse? (27:28 ) | Zrušení Everquest Next (34:44) | Velké zmatení nad remakem System Shock (39:51) | Dotazy (44:26)}{1:17:14}{34031804}{}
\FC{fight-club-269.mp3}{\FCtitle{\#269}{Gifové hnízdo}}{2016-03-23}{\Sleepless, \Luke, \Blixa}{\FCrule{\#269}{Kdo negifuje s námi, gifuje proti nám.}{\Luke}}{Co jsme viděli na GDC 2016 v San Franciscu | (25:33) Lawbreakers už není free2play | (38:28) Nová česká hra Hobo simuluje život bezdomovce | (46:06) Dotazy}{1:08:03}{30731756}{1}
\FC{fight-club-270.mp3}{\FCtitle{\#270}{Šlehačková výzva}}{2016-03-30}{\Sleepless, \Blixa, \Ulrik, \OverWatch}{\FCrule{\#270}{Šlehačkou se nepoléváme.}{\Blixa}}{Čína kupuje LEVEL? | (12:05) Co bude PlayStation 4.5? | (23:40) Potíže vývojářů s letsplay videi | (38:30) Velká debata o blokování reklamy | (60:20) Spousta dotazů}{1:34:59}{38455142}{1}
\FC{fight-club-271.mp3}{\FCtitle{\#271}{Repete!}}{2016-04-06}{\Sleepless, \Luke, \Ulrik}{\FCrule{\#271}{Nejlepší sidekick Kratose je\,\dots Mario.}{\Luke}}{Mass Effect: Andromeda | (13:12) Onimushi + Ninja Gaiden = Nioh | (24:26) Mindfuck jménem Miegakure | (31:37) God of War 4 | (43:53) Dotazy | {\em (60:35) Dotaz na datum konání srazu platinového předplatného}\/}{1:09:50}{31289216}{2}
\FC{fight-club-272.mp3}{\FCtitle{\#272}{Skotský král}}{2016-04-13}{\Sleepless, \Luke, \Blixa, \OverWatch}{\FCrule{\#272}{Chcete horor -- jeďte do Vlašimi!}{\Luke}}{Leslie Benzies vs Rockstar | (18:59) Codies jedou! | (28:20) Bude Phantasmal horor? | (34:00) Dlouhé dotazy}{1:21:33}{35291012}{}
\FC{fight-club-273.mp3}{\FCtitle{\#273}{Maf III e}}{2016-04-20}{\Sleepless, \Luke, \Blixa, \Ulrik}{\FCrule{\#273}{Prosim vás, ty hry nepalte ani nevařte.}{\Luke}}{Nový Level 263 | (11:09) Mafie III a~my | (22:49) Pyre od Supergiant Games | (31:02) Nová hra Deana Halla Out Of Ammo | (41:10) Dotazy}{1:20:32}{34862108}{1}
\FC{fight-club-274.mp3}{\FCtitle{\#274}{Grafici útočí}}{2016-04-27}{\Blixa, \Ulrik, \Memory, \JDiblik}{\FCrule{\#274}{Co se stalo v dark zoně, zůstane v dark zoně.}{\Blixa}}{Aktuálně ke změnám v Heartstone | (15:50) Naděje a~obavy kolem Warcraft filmu | (25:45) The Division s odstupem času | (40:00) Kdo (ne)bude na E3? | (47:30) Dotazy}{1:23:22}{31949626}{2}
\FC{fight-club-275.mp3}{\FCtitle{\#275}{Lvohlavouni}}{2016-05-04}{\Sleepless, \Luke, \Ulrik, \PBulir}{\FCrule{\#275}{May the pixel be with you.}{\Luke}}{Nový Ratchet \&{} Clank | (14:19) Salt \&{} Sanctuary | (20:26) Dawn of War III | (31:56) Vzpomínáme na Lionhead | (44:34) Dotazy | {\em (72:03) Superholub a~Star Wars Day}\/}{1:15:07}{32356400}{1}
\FC{fight-club-276.mp3}{\FCtitle{\#276}{Překladatelský oříšek}}{2016-05-11}{\Luke, \Blixa, \Ulrik, \FZenisek}{\FCrule{\#276}{Pod bestií je největší stín.}{\Luke}}{Filip Ženíšek o překládání Zaklínače a~her obecně | (69:00) Oznámení Battlefield 1 | (79:00) Konec Disney Infinity | (97:00) Dotazy}{2:02:51}{54556160}{2}
\FC{fight-club-277.mp3}{\FCtitle{\#277}{Tři avokáda pro Alžbětu/Domácí zvířata útočí}}{2016-05-18}{\Sleepless, \Luke, \OverWatch, \Bety}{\FCrule{\#277}{A sedmého dne Bůh odpočíval a~hrál izometrické hry.}{\Luke}}{Brigador | (05:04) Zaklínač 3: O vínu a~krvi | (24:20) Uncharted 4 a~Alžběta | (36:36) Dishonored 2 | (49:11) Shawden | (50:30) Dotazy}{1:54:13}{55485752}{2}
\FC{fight-club-278.mp3}{\FCtitle{\#278}{Dům plný Doomu}}{2016-05-25}{\Luke, \Blixa, \Ulrik, \JOlejnik}{\FCrule{\#278}{Git normal.}{\Luke}}{Nový Level 264 a~historka o Romerovi a~Level TV | (12:03) Nový Doom a~aféra "git gud" | (24:26) Starý a~nový Doom | (38:06) The Way kombinuje Another World, Flashback a~Heart of Darkness | (44:26) Dotazy}{1:20:44}{36019724}{1}
\FC{fight-club-279.mp3}{\FCtitle{\#279}{Jirka Pavlovský Unleashed}}{2016-06-01}{\Sleepless, \Luke, \Ulrik, \JiP}{\FCrule{\#279}{Neplést Čáryfuka s Mrakoplašem.}{\Luke}}{Pavlovský z řetězu urvaný (filmy, hry a~komixy) | 65:16 Dotazy}{1:42:52}{49175528}{}
\FC{fight-club-280.mp3}{\FCtitle{\#280}{Klubové umění}}{2016-06-08}{\Luke, \Ulrik, \OverWatch, \DRynes}{\FCrule{\#280}{Schnoucí barvu pozorujte pouze ve více lidech.}{\Luke}}{Film Warcraft a~my | (36:43) Overwatch a~David Ryneš | (49:00) Fable Fortune a~další karetní hry | (56:44) Dotazy}{1:18:43}{36233780}{2}
\FC{fight-club-281-E3-1.mp3}{\FCtitle{E3 2016 speciál \#1}{}}{2016-06-14}{\Sleepless, \Janchor, \PBulir, \EHarton}{}{Souhrn tiskových konferencí a~prvního dne E3 2016.}{1:29:17}{39073696}{2}
\FC{fight-club-281-E3-2.mp3}{\FCtitle{E3 2016 speciál \#2}{}}{2016-06-15}{\Luke, \Blixa, \Moolker, \VHrincar}{}{(05:49) Steep | (12:23) Quake Champions | (23:58) Inside | (27:06) God of War | (29:32) Dual Universe | (34:38) Nová Kojimova hra | (37:54) Ghost Recon: Wildlands | (43:26) Resident Evil 7 | (43:33) Star Trek: Bridge Crew | (52:18) Titanfall 2 \&{} Battlefield 1 | (62:56) Mafie 3 + Dotazy}{1:30:27}{40824824}{1}
\FC{fight-club-281-E3-3.mp3}{\FCtitle{E3 2016 speciál \#3}{}}{2016-06-18}{\Janchor, \Ulrik, \JOlejnik, \MRota}{}{Watch Dogs 2 bereme s rezervou | (23:30) Čekáme ještě na The Last Guardian? | (28:00) Forza a~Horizon | (40:00) Souboje v Bannerlord, Final Fantasy a~Scalebound | (60:30) Dishonored 2 | (76:00) Dotazy ze všech stran}{1:40:30}{42749980}{3}
\FC{fight-club-282.mp3}{\FCtitle{\#282}{Hm s Lukášem Macurou!}}{2016-06-23}{\Luke, \Blixa, \OverWatch, \LMacura}{\FCrule{\#282}{Nakonec\,\dots stejně všechno koupí Tencent.}{\Luke}}{Hrneme E3 2016 na hromadu | (05:38) Nová Zelda | (12:24) Lesk a~bída VR na E3 | (25:07) Aféra G2A a~dalších přeprodejců kódů | (46:43) Co je to Tencent? Co ještě koupí? | (55:28) Dotazy}{1:19:14}{35064968}{1}
\FC{fight-club-283.mp3}{\FCtitle{\#283}{Návraty}}{2016-06-29}{\Sleepless, \Luke, \Ulrik}{\FCrule{\#283}{Pozor, ale opravdu pozor, na bloudění v mlze.}{\Luke}}{Nový Level 265 | (07:52) Tahovka Phoenix Point od autora UFO: Enemy Unknown a~Chaos Reborn | (16:59) Kniha Brány Skeldalu a~gamebooky | (23:26) Remake System Shock | (36:37) Bojová umění v Absolver | (44:08) Dotazy}{1:21:34}{37968764}{1}
\FC{fight-club-284.mp3}{\FCtitle{\#284}{Nesponzorováno}}{2016-07-13}{\Sleepless, \Luke, \Ulrik}{\FCrule{\#284}{Pokémoni neexistují.}{\Luke}}{Harvey Smith zasahuje | (01:27) Pokémon GO a~my | (16:51) Právní mračna nad youtubery | (29:04) Představení Dark Fear | (31:31) Představení Beat Cop | (35:38) Dotazy}{1:12:50}{30913520}{2}
\FC{fight-club-285.mp3}{\FCtitle{\#285}{Potvory do kapsy}}{2016-07-20}{\Luke, \Blixa, \OverWatch, \DRynes}{\FCrule{\#285}{Bez Anči a~Pikachu ani ránu.}{\Luke}}{Hrajeme Pokémon GO | (27:54) Mini NES a~další nové retro konzole | (39:47) Viktor Antonov opouští Arkane | (48:39) Dotazy}{1:17:52}{31479728}{2}
\FC{fight-club-286.mp3}{\FCtitle{\#286}{Drinking Game}}{2016-07-27}{\Luke, \Blixa, \OverWatch}{\FCrule{\#286}{Dávejte si pozor na svoje data.}{\Luke}}{Myší dírou v Ghost of a~Tale | (06:53) Jak to je s Ninteno a~Pokémon Go | (13:58) Budoucnost Nintenda (NX) | (28:47) Quake Champions a~chaos | (43:00) Dotazy (dlouhý blok o Level TV included)}{1:23:32}{34694960}{2}
\FC{fight-club-287.mp3}{\FCtitle{\#287}{Platón v sudu}}{2016-08-03}{\Luke, \JOlejnik, \VPechacek}{\FCrule{\#287}{Co je lesklého,\,\dots to může být i shiny.}{\Luke}}{Vašek Pecháček se představuje | (03:28) Galaxy in Turmoil končí | (14:46) GamesCom zpřísňuje bezpečnost | (22:57) Rio v Overwatch | (31:55) Dotazy}{1:06:54}{26992724}{}
\FC{fight-club-288.mp3}{\FCtitle{\#288}{Když chybí Petr Bulíř}}{2016-08-10}{\Luke, \Blixa, \Ulrik, \VPechacek}{\FCrule{\#288}{Jde to i s ochranou.}{\Luke}}{Mladá krev na GamesComu | (13:30) Pavel je spokojený s Abzû | (18:47) Kontroverze kolem No Man's Sky | (30:00) Denuvo prolomeno? | (47:02) Dotazy | {(73:20) Na jakém nejpodivnějším místě jste hráli hry?}}{1:24:00}{34748564}{2}
\FC{fight-club-289.mp3}{\FCtitle{\#289}{S Mr. Hladem: Mizerný režisér David Cage}}{2016-08-17}{\Sleepless, \Luke, \Ulrik, \MrHlad}{\FCrule{\#289}{Tenhle Big Boss není pro starý.}{\Luke}}{Představení nového Levelu | (09:03) Mr. Hlad a~film Warcraft | (19:24) Šance filmového Assassin's Creed | (25:43) Filmové Division, Resident Evil, Zaklínač, Tomb Raider a~další | (44:25) Reakce na Metal Gear Survive | (51:43) Battlefield a~Aleš | (55:20) Dotazy}{1:22:59}{34947140}{2}
\FC{fight-club-290.mp3}{\FCtitle{\#290}{Ozvěny Gamescomu}}{2016-08-24}{\Janchor, \VPechacek, \Bety, \EHarton}{\FCrule{\#290}{Květiny jsou víc než pušky.}{\Janchor}}{Gamescom | Deus Ex: Mankind Divided | Dotazy | Pravidlo Klubu rváčů}{1:28:52}{38241520}{1}
\FC{fight-club-291.mp3}{\FCtitle{\#291}{Za oponou Prahy v Deus Ex}}{2016-08-31}{\Janchor, \OverWatch, \VPechacek, \FZenisek}{\FCrule{\#291}{Dokupte papírové!}{\Janchor}}{Dojmy z PlayStation VR | Čeština v Deus Ex: Mankind Divided | Dotazy ({\em Jindřich Rohlík}\/)| Pravidlo Klubu rváčů}{1:49:38}{48941208}{3}
\FC{fight-club-292.mp3}{\FCtitle{\#292}{Věštění budoucnosti}}{2016-09-07}{\Sleepless, \Janchor, \Ulrik, \OverWatch}{\FCrule{\#292}{I s malou trofejí jde dělat velkej nohejbal.}{\Janchor}}{{\em (0:20) Polévka s koriandrem | (1:10) Press vs. Devs 2016}\/ | (9:10) System Shock 3 | (46:40) Free to play verze EVE Online | Dotazy | Pravidlo Klubu rváčů}{1:41:21}{40438888}{2}
\FC{fight-club-293.mp3}{\FCtitle{\#293}{Mafiáni}}{2016-09-14}{\Janchor, \Blixa, \Ulrik, \VPechacek}{\FCrule{\#293}{Čůráci, čůráci everywhere.}{\Janchor}}{Aktuální dění okolo PlayStation 4 | Zásadní změna systému uživatelských recenzí na Steamu | Dojmy z hraní Mafia III | Dotazy | Pravidlo Klubu rváčů}{1:17:48}{30394260}{2}
\FC{fight-club-294.mp3}{\FCtitle{\#294}{Sportovní snění}}{2016-09-21}{\Sleepless, \Janchor, \OverWatch, \VPechacek}{\FCrule{\#294}{Vraťte nám hokej, ať už je od nás pokej.}{\Janchor}}{Pozvánky na GDS 2016 a~Gamer Pie | (12:00) Nová FIFA a~NHL | (36:15) Nové kauzy kolem Steamu | (50:20) Nové dotazy}{1:38:08}{37392192}{1}
\FC{fight-club-295.mp3}{\FCtitle{\#295}{Hráč z laboratoře}}{2016-10-05}{\Luke, \Ulrik, \VPechacek, \OHrabec}{\FCrule{\#295}{Mějte styl.}{\Luke}}{Ondřej Hrabec o výzkumech hráčských stylů a~závislosti | (58:16) Úpravy v Hearthstone | (73:00) VR od Google | 80:00 Dotazy | {\em Minuta ticha}}{2:02:20}{49372772}{1}
\FC{fight-club-296.mp3}{\FCtitle{\#296}{Krabí maso}}{2016-10-12}{\Luke, \Blixa, \OverWatch, \VPechacek}{\FCrule{\#296}{Chlebíčky jsou základ Fight Clubu.}{\Luke}}{Aktuální Level a~newsletter Games.cz | (09:06) Mafie 3 | (27:46) Kde jsou hry pro rodiny a~děti?! | (40:51) VR Lone Echo | (46:02) Dotazy}{1:06:35}{29378588}{2}
\FC{fight-club-297.mp3}{\FCtitle{\#297}{Jádroví lidé}}{2016-10-19}{\Sleepless, \Luke, \Blixa, \PSpacek}{\FCrule{\#297}{Obarvujte realitu.}{\Luke}}{Nový Level (268) je tu! | (08:24) Go west! Aneb, jak se vyvíjí hry za velkou louží | (18:51) Patrik Špaček o remasterování Indiana Jones and the Fate of Atlantis | (52:54) Red Dead Redemption 2 | (63:00) Ujíždíme z komunistického Československa! | (70:21) Dotazy}{1:27:09}{42365372}{1}
\FC{fight-club-297-GP.mp3}{\FCtitle{Gamer Pie festival speciál}{}}{2016-10-22}{\Sleepless, \Luke, \Blixa, \VPechacek, \JHorcik}{\FCrule{Gamer Pie}{Když jste v Brně, tak si dávejte pozor na Chuchla.}{\Luke}}{Česká herní scéna | Herní média vs. youtubeři vs. videoherní průmysl | Co přinese virtuální realita}{1:16:09}{31097520}{}
\FC{fight-club-298.mp3}{\FCtitle{\#298}{Recenzní politka}}{2016-10-26}{\Sleepless, \Luke, \Ulrik, \OverWatch, \PBarak}{\FCrule{\#298}{Měcháček to zařídí.}{\Luke}}{Pozvánka na GDS 2016 | (29:38) Nintendo Switch | (53:10) Bethesda vs. novináři | (66:39) Dotazy}{1:37:23}{39371072}{3}
\FC{fight-club-299.mp3}{\FCtitle{\#299}{ČeRV ve VR}}{2016-11-02}{\Luke, \OverWatch, \VPechacek, \Cerv}{\FCrule{\#299}{Co Červ nedohrál, no to prostě ani nehrál.}{\Luke}}{ČeVR: Červ a~virtuální realita | (33:39) Nové hry od Bohemia Interactive: Project Argo a~Ylands | (43:02) Oťukáváme Gwent a~jsme nadšení | (47:59) Bud Spencer a~Terrence Hill ve hře | (49:29) Dotazy}{1:13:35}{33190088}{2}
\FC{fight-club-300.mp3}{\FCtitle{\#300}{Tři a~půl kola}}{2016-11-12}{\Sleepless, \Luke, \Janchor, \Blixa, \Ulrik, \OverWatch, \Moolker, \Farid, \Wolf, \JOlejnik, \JVrobel, \RFriedrich}{\FCrule{\#300}{This is Fight Club!}{\Luke}}{{\bf Sešli jsme se v kině 64 U Hradeb} | 00:22 Původní Klub rváčů promlouvá o životě (Lukáš Grygar, Martin Bach, Petr Poláček a~Farid Starman) | (21:40) Další generace nastupuje (Pavel Dobrovský, Aleš Smutný, Adam Homola a~Honza Olejník) | 37:58 Tajná Games / Level TV | (42:45) Střet titánů (Radek Friedrich, Miloš Bohoněk, Michal Rybka a~Honza Vrobel) | (74:00) Tak na další stovku!}{1:19:45}{33979424}{3}
\FC{fight-club-301.mp3}{\FCtitle{\#301}{Cesta vlakem}}{2016-11-16}{\Sleepless, \Janchor, \VPechacek, \VVanek}{\FCrule{\#301}{Nejlepší hry jsou hry, který začínají na písmeno D.}{\Janchor}}{Představení Levelu 269 | (04:00) Půlhodinka o Dark Train | (28:00) Dishonored 2 a~jeho příbuzní | (59:00) Otázky a~odpovědi}{1:29:11}{36565812}{2}
\FC{fight-club-302.mp3}{\FCtitle{\#302}{Srdíčka pro Adama}}{2016-11-23}{\Sleepless, \Luke, \Ulrik, \OverWatch}{\FCrule{\#302}{Kdo si počká, tomu Vlasta zavolá.}{\Luke}}{Český survival mezi dinosaury Claw Hunter | (09:31) Salvaged VR: Trochu jiný Space Hulk | (17:01) Star Citizen upřímný | (24:02) Koho podporuje EA a~koho koťátka | (34:46) Dotazy | (50:10) Větší než malá dávka Tyranny}{1:10:05}{31001612}{}
\FC{fight-club-303.mp3}{\FCtitle{\#303}{No Sky for Man}}{2016-11-30}{\Luke, \Blixa, \OverWatch, \VPechacek}{\FCrule{\#303}{Ať žije ICQ!}{\Luke}}{Facebook šlape do her | (15:51) Hovory o No Man's Sky | (33:10) Odkud vysílá Signal from Tölva? | (43:25) Dotazy}{1:12:35}{28623704}{}
\FC{fight-club-304.mp3}{\FCtitle{\#304}{Lízací klub}}{2016-12-07}{\Sleepless, \Luke, \Ulrik, \Bety}{\FCrule{\#304}{Bez perzistentní, emergentní imerze to prostě nejde!}{\Luke}}{Co si přejeme v Last of Us 2 | (40:46) Kojima trollí v Death Stranding | (34:08) Hodnocení Dead Rising 4 | (54:37) DOTAZY}{1:30:59}{40081532}{}
\FC{fight-club-305.mp3}{\FCtitle{\#305}{Indiáni jedou!}}{2016-12-14}{\Luke, \OverWatch, \Fiola, \Baty}{\FCrule{\#305}{Bez Baty tu komunitu ani nezkoušejte.}{\Luke}}{Jaké mobilní hry hrajeme | (06:25) Insiderský pohled na Playstation Experience | (11:23) Dinosauři a~kovbojové v Dino Frontier | (25:00) The Steep: Sníh je tu! | (38:34) The Last Guardian | (45:00) Dotazy}{1:10:54}{30369596}{}
\FC{fight-club-306.mp3}{\FCtitle{\#306}{Vaaaaaaanoceeeeeeee!}}{2016-12-21}{\Blixa, \Janchor, \Ulrik, \OverWatch, \DRynes}{\FCrule{\#306}{Ta nejkrásnější emoce je erekce.}{\Janchor}}{Nový LEVEL a~jeho předvánoční předplatné | (9:55) Co přinesl rok 2016? | (37:33) Filmové okénko o Rogue One | (48:10) Líbí se nám Super Mario Run? | (51:20) Přichází Aleš | (62:45) Dotazy z mailu a~chatu}{1:37:39}{38620204}{3}
\FC{fight-club-307.mp3}{\FCtitle{\#307}{Sépiové studio}}{2017-01-04}{\Sleepless, \Luke, \OverWatch, \VPechacek}{\FCrule{\#307}{Každou středu musíte poslouchat Fight Club\dots\,prosím.}{\Luke}}{Mass Effect: Andromeda | (08:52) Divinity: Original Sin 2 | (12:02) Mount \&{} Blade: Bannerlord | (19:10) Church in the Darkness | (24:06) Echo | (27:40) What Remains of Edith Finch | (31:03) Nioh | (34:25) Absolver | (35:31) Return of the Obra Dinn | (40:00) Tides of Numenera | (44:21) Persona 5 | (46:33) Tokyo 42 | (51:08) DOTAZY}{1:12:46}{27905000}{}
\FC{fight-club-308.mp3}{\FCtitle{\#308}{Tak to končí\,\dots}}{2017-01-11}{\Luke, \Ulrik, \VPechacek}{\FCrule{\#308}{Buďte hodní i na cizí pejsky.}{\Luke}}{Zrušení Scalebound | (09:22) Zrušení Everquest Landmark | (18:08) Psí bezdomovec v Home Free | (25:30) Co je to GeFroce NOW | (35:38) Dotazy}{1:18:56}{35248856}{}
\FC{fight-club-309.mp3}{\FCtitle{\#309}{Switchujeme}}{2017-01-18}{\Janchor, \Blixa, \OverWatch, \VPechacek}{\FCrule{\#309}{Nejlepší Switch je Moskvič.}{\Janchor}}{Nový mikrofon a~hledání programátorů pro Games.cz | (5:00) Jak realistická bude hra realpolitics? | (18:30) Zajímavá čísla ze Steamu | (27:35) Rýpání v Nintendu Switch | (51:00) Dotazy ze všech stran}{1:19:11}{29263608}{2}
\FC{fight-club-310.mp3}{\FCtitle{\#310}{Lovecký lístek}}{2017-01-25}{\Sleepless, \Luke, \Ulrik, \OverWatch}{\FCrule{\#310}{Always look at the bright side of life.}{\Luke}}{Level \#271 brzy i u vás! | (13:09) Masakr zvířátek v The Hunter: Call of the Wild \&{} Feral | (25:35) Banán pro Banner Saga 3? | (34:53) Call of the Cthulhu pana Lovecrafta | (45:30) Dotazy}{1:21:45}{34357820}{1}
\FC{fight-club-311.mp3}{\FCtitle{\#311}{Jedeme do Afriky!}}{2017-02-01}{\Luke, \Ulrik, \OverWatch, \VPechacek}{\FCrule{\#311}{Konec všem ptákovinám!}{\Luke}}{Dojmy Adama a~Aleše z For Honor | (16:26) Beautiful Desolation a~herní scéna v Africe | (30:10) Penisy v Conan Exiles | (42:24) Dotazy}{1:02:19}{26230892}{2}
\FC{fight-club-312.mp3}{\FCtitle{\#312}{Ať žije Molyneux!}}{2017-02-08}{\Sleepless, \Luke, \Blixa, \Ulrik}{\FCrule{\#312}{Každý má rád koupelnu.}{\Luke}}{Facebook vs. Bethesda | (22:07) Názory na hru World of Castles | (32:08) Jak to je s Psychonauts 2 | (48:35) Dotazy}{1:13:04}{30535448}{}
\FC{fight-club-313.mp3}{\FCtitle{\#313}{Vizuální zklamání}}{2017-02-15}{\Luke, \Blixa, \Ulrik, \OverWatch}{\FCrule{\#313}{Emergentně launchujte jen ve springu.}{\Luke}}{Ze zákulisí České hry roku 2016 | (20:07) Vzpomínky na začátky internetu | (39:13) Virtually VR, Legion 1917 a~další tituly z konference White Nights | (49:37) Důsledky zrušení Greenlightu | (62:35) Dotazy | {\em vzpomínky na ilegal}}{1:31:58}{38318360}{2}
\FC{fight-club-314.mp3}{\FCtitle{\#314}{Lovec platinovek s Petrem Bulířem}}{2017-02-22}{\Luke, \OverWatch, \PBulir}{\FCrule{\#314}{Lokalizujte s citem.}{\Luke}}{Nový Level 272 | (05:49) Nioh | (20:00) Horizon: Zero Dawn | (31:41) Yakuza 0 | (43:36) Dotazy}{1:20:38}{33847160}{3}
\FC{fight-club-315.mp3}{\FCtitle{\#315}{Neovlivnitelní}}{2017-03-01}{\Luke, \Ulrik, \OverWatch, \VPechacek}{\FCrule{\#315}{Čtení není mrtvé!}{\Luke}}{(03:16) Dojmy z Mass Effect: Andromeda | (15:14) Dojmy z Torment: Tides of Numenéra| (34:00) Dotazy}{0:59:56}{23952488}{}
\FC{fight-club-316.mp3}{\FCtitle{\#316}{Degustace zakázána}}{2017-03-08}{\Blixa, \OverWatch, \Moolker, \LMacura}{\FCrule{\#316}{Cartridge nežereme!}{\Blixa}}{Úvod ze Středozemě | (04:00) Nintendo Switch a~nová Zelda | (35:35) Vychází The Keep na PC | (42:50) Jaké bylo GDC 2017? | (61:00) Čeští vývojáři jednají se státem | (74:50) Dotazy}{1:34:40}{36231280}{}
\FC{fight-club-317.mp3}{\FCtitle{\#317}{Se sukní i bez}}{2017-03-15}{\Janchor, \DRynes, \OverWatch, \GSuterova}{\FCrule{\#317}{I sukně může být trofej.}{\Janchor}}{Něco málo o Games TV | (04:30) Hry na NI(er) | (22:00) Nový Styx pro Adama | (26:30) Gauneři v herní podobě | (32:40) Kde je pravda kolem Ghost Recon Wildlands? | (48:00) Dotazová smršť}{1:23:44}{32416518}{3}
\FC{fight-club-318.mp3}{\FCtitle{\#318}{Žlutej pes}}{2017-03-22}{\Sleepless, \Luke, \OverWatch, \VPechacek}{\FCrule{\#318}{Posunujte stropy!}{\Luke}}{Petr v Americe | (08:54) Unavowed od Wad Jet Eye Games | (22:31) Vypadá to dobře! Slovenská Sacred Fire | (27:54) MMORPG Camelot Unchained a~Crowfall | (41:49) 10 let PlayStation 3 v Evropě | (56:21) Dotazy}{1:27:14}{37020632}{}
\FC{fight-club-319.mp3}{\FCtitle{\#319}{Oftahistoričnost}}{2017-03-29}{\Sleepless, \Luke, \Blixa, \Ulrik}{\FCrule{\#319}{Nekousejte si nehty, nebo dopadnete jako Vašek.}{\Luke}}{Level 273 s virtuální realitou a~rozhovorem s Ronem Gilbertem | (10:39) Remaster Planescape: Torment, chceme ho? | (16:09) Pomořský horor Narcosis a~studentské projekty | (29:41) Fandom: Herní novinařina nebo komunitní služba? | (38:32) Co by mohlo být v nové Call of Duty | (49:54) Dotazy}{1:24:18}{35213936}{}
\FC{fight-club-320.mp3}{\FCtitle{\#320}{Politika}}{2017-04-05}{\Sleepless, \Luke, \OverWatch, \PBarak}{\FCrule{\#320}{Nepodporujte předprodeje!}{\Luke}}{Prorůstání her a~politiky s Pavlem Barákem | (33:52) Petr a~Adam juchají nad Zeldou | (48:48) Dean Hall představil Stationeers - vyjde mu to? | (59:20) {\bf Remaster Starcraftu} + další hry od Blizzardu | (69:51) Dotazy}{1:50:07}{47438132}{}
\FC{fight-club-321.mp3}{\FCtitle{\#321}{Virtuální realita v praxi}}{2017-04-12}{\Luke, \Blixa, \VPechacek, \MDuris}{\FCrule{\#321}{Piš jako slyš.}{\Luke}}{(01:10) Jak to funguje v herně VR s Markem Ďurišem | (33:13) Kladivo na youtubery | (48:05) Star Wars Battlefront II | (53:17) Frostpunk od 11 bit | (57:34) Dotazy}{1:12:57}{30443360}{2}
\FC{fight-club-322.mp3}{\FCtitle{\#322}{Hekající skála}}{2017-04-19}{\Janchor, \Blixa, \Ulrik, \OverWatch}{\FCrule{\#322}{Hekající skálu nenapodobíš.}{\Janchor}}{Nový LEVEL 274 je blízko | (08:00) Na co se těšíme v Battlefront 2 |  (18:30) Nová konkurence pro Steam z Číny? | (31:00) Pro koho je Julka-Laylee | (41:00) Dotazy z mailu a~chatu}{1:04:12}{26269236}{2}
\FC{fight-club-323.mp3}{\FCtitle{\#323}{Opičí destilát}}{2017-04-26}{\Sleepless, \Luke, \Ulrik, \VPechacek}{\FCrule{\#323}{Pirátů není nikdy dost.}{\Luke}}{Ancestors od Patrice Désiletse (Assassin's Creed) | (13:37) Esport na Asian Games | (20:59) Figment a~herní muzikály | (30:59) Dotazy}{1:00:23}{25062224}{}
\FC{fight-club-324.mp3}{\FCtitle{\#324}{Ve vyhnanství}}{2017-05-03}{\Sleepless, \Blixa, \OverWatch, \VPechacek}{\FCrule{\#324}{Kdo neskáče, není klokan.}{\Sleepless}}{Nový Quake v provizorním studiu | (16:30) Nadšení z nového Nintenda 2DS | (27:00) Co čekat od Darksiders 3? | (36:30) Dvě novinky z Marsu | (44:30) Dotazy leakují horem dolem}{1:11:42}{27217816}{2}
\FC{fight-club-325.mp3}{\FCtitle{\#325}{Prej je to tak!}}{2017-05-10}{\Luke, \Janchor, \OverWatch, \VPechacek}{\FCrule{\#325}{Ať žije prase.}{\Luke}}{Bavíme se o Prey | (47:15) Dotazy}{1:06:14}{27918212}{2}
\FC{fight-club-326.mp3}{\FCtitle{\#326}{Návrat Geralta}}{2017-05-17}{\Sleepless, \Luke, \Ulrik, \OverWatch}{\FCrule{\#326}{Nepodceňujte nejlepší novinkáře.}{\Luke}}{{\em (00:25) Nové studio | (01:35) LEVEL párty}\/ | Poslední výstřel Hitmana | (12:49) Assassin's Creed v Egyptě, Far Cry 5 a~my | (23:28) Zaklínač se vrací v seriálu od Netflix | (31:30) Řecký soud o českých vývojářích | (37:12) Co hrajeme | (45:10) Dotazy}{1:03:23}{26753684}{}
\FC{fight-club-327.mp3}{\FCtitle{\#327}{Slow TV}}{2017-05-24}{\Sleepless, \Luke, \Blixa, \Ulrik}{\FCrule{\#327}{Bacha na mrkajícího Bacha.}{\Luke}}{Level 275 is here! | (08:39) 2D Dark Souls jménem Blasphemous | (14:46) Kam jdeš, Unity? | (23:18) Ovládání her myšlenkami | (32:17) Ochmatali jsme New Nintendo 2DS XL | (39:51) Xbox Game Pass je tu! | (43:48) Dotazy}{1:12:34}{30531920}{2}
\FC{fight-club-328.mp3}{\FCtitle{\#328}{Uctívači ještěrek}}{2017-05-31}{\Sleepless, \Ulrik, \OverWatch, \VPechacek}{\FCrule{\#328}{Kdo chce uspat Adama, cačne mluvit o Warhammeru.}{\Sleepless}}{Rozebíráme Far Cry 5 | (18:15) Dojmy z hraní Total War: Warhammer 2 | (35:00) Overwatch slaví první výročí | (50:31) | (1:03:20) Odpovídáme  na dotazy {\em | (1:25:36) Neznají Legendy Polskie}}{1:31:15}{36681748}{1}
\FC{fight-club-329.mp3}{\FCtitle{\#329}{Need for oves}}{2017-06-07}{\Sleepless, \Luke, \Blixa, \OverWatch}{\FCrule{\#329}{Počkejte si až to vyjde.}{\Luke}}{Co čekáme od E3 2017 | (15:04) Valve opouští scénáristé | (20:00) Fízlácká simulace Police Stories | (24:40) Kopeme do Need for Speed: Payback | (38:10) Dotazy}{1:07:08}{28265828}{}
\FC{fight-club-329-E3-1.mp3}{\FCtitle{E3 2017 speciál \#1}{}}{2017-06-14}{\Janchor, \OverWatch, \JOlejnik, \KDrda}{}{-}{1:20:10}{40817565}{2}
\FC{fight-club-329-E3-2.mp3}{\FCtitle{E3 2017 speciál \#2}{}}{2017-06-15}{\Janchor, \JVrobel, \JIlavsky, \VHrincar}{}{-}{1:27:27}{41626737}{}
\FC{fight-club-330.mp3}{\FCtitle{\#330}{UJD}}{2017-06-21}{\Sleepless, \Blixa, \Ulrik}{\FCrule{\#330}{Když flashka, tak dřevěná.}{\Sleepless}}{Martin a~Aleš hodnotí své dojmy z výstavy E3 2017, řeší největší překvapení i největší zklamání. {\em Medvěd je nový luk\,\dots :)}}{1:42:55}{43463768}{2}
\FC{fight-club-331.mp3}{\FCtitle{\#331}{SNES dnes}}{2017-06-28}{\Sleepless, \Luke, \Blixa, \OverWatch}{\FCrule{\#331}{Nestaňte se obětí své hry.}{\Luke}}{Nový Level 276 | (09:50) Odchod autora Prey za scény | (20:31) Zákaz a~povolení OpenIV pro GTA IV | (34:50) Mini SNES má super hry | (46:50) Dotazy}{1:16:31}{32448776}{}
\FC{fight-club-331-E3-3.mp3}{\FCtitle{E3 2017 speciál \#3}{}}{2017-07-04}{\Sleepless, \Janchor, \OverWatch, \VPechacek}{}{Rekapitulace a~hodnocení letošní E3 člelny redakce, kteří zůstali v Praze a~sledovali všechno zpovzdálí.}{1:12:48}{30852700}{}
\FC{fight-club-332.mp3}{\FCtitle{\#332}{Vesmírné okno}}{2017-07-12}{\Sleepless, \Luke, \OverWatch, \VPechacek}{\FCrule{\#332}{Když se koukneš do dvora, bouchne Ti tam mrtvola.}{\Luke}}{Vesmírný výlet pana DeGrasse Tysona | (17:19) Deník zombokalypsy v Zafehouse Diaries 2 | (27:01) Válka druhá světová v Steel Division Normandy a~Hell Let Loose | (38:34) Nekromantská radost v Diablo III | (50:50) DOTAZY}{1:12:05}{30141572}{}
\FC{fight-club-333.mp3}{\FCtitle{\#333}{Méně hymen, více RPG}}{2017-07-19}{\Sleepless, \Janchor, \OverWatch, \VPechacek}{\FCrule{\#333}{Slávie je pro psa.}{\Janchor}}{Sparta Praha jako jediný CZ tým ve FIFA 18 | (18:00) Ubisoft koketuje s Battle Royale | (30:10) Co nabídne Ataribox? | (36:30) Otec Mass Effectu se vrací do BioWare | (45:45) DOTAZY}{1:13:23}{29095324}{1}
\FC{fight-club-334.mp3}{\FCtitle{\#334}{Hrdinova cesta}}{2017-07-26}{\Luke, \Ulrik, \OverWatch, \BHeblik}{\FCrule{\#334}{Nedělejte si to složitější.}{\Luke}}{Press Start na Letní filmové škole | (15:11) Hry a~příběhy | (25:49) Náročnost ve hrách | (36:50) Fortnite je konečně tady! | (46:22) Zkusili jsme My Time At Portia | (50:57) Dotazy}{1:16:44}{29095452}{2}
\FC{fight-club-335.mp3}{\FCtitle{\#335}{My tady děti nežereme!}}{2017-08-02}{\Sleepless, \Blixa, \OverWatch, \VPechacek}{\FCrule{\#335}{My tady děti nežerem.}{\VPechacek}}{Adam, hraje Pyre | (14:45) Tencent utrácí v Evropě | (29:20) Pillars of Eternity II je nový Mass Effect | (37:30) Ready Player One | (47:00) Už zase ten Star Citizen | (58:00) Dotazy}{1:30:21}{33554838}{1}
\FC{fight-club-336.mp3}{\FCtitle{\#336}{Zlé lootboxy}}{2017-08-09}{\Luke, \Ulrik, \OverWatch, \VPechacek}{\FCrule{\#336}{Hippocampus atrofuje.}{\Luke}}{Epizodické hraní není pro každého! | (22:23) Přijdou herní "netflixy"? | (35:02) Placené předměty v (ne)placených hrách | (51:47) Dotazy}{1:14:38}{31597972}{1}
\FC{fight-club-337.mp3}{\FCtitle{\#337}{Je suis mrkev!}}{2017-08-16}{\Sleepless, \Luke, \OverWatch, \VPechacek}{\FCrule{\#337}{Angličtinu na vás, holoto!}{\Luke}}{Lawbreakers nic nebrejkuje | (15:16) VR Arktika.1 stojí za vidění | (30:31) EVE Valkyrie jde do světa | (36:15) Gamifikované fake news | (49:42) Dotazy}{1:16:42}{31876684}{}
\FC{fight-club-338.mp3}{\FCtitle{\#338}{GamesCom 2017 speciál}}{2017-08-24}{\Sleepless, \Blixa, \OverWatch, \PHajda}{\FCrule{\#338}{Angela mutuje.}{\Sleepless}}{Vyšel nový LEVEL 277 | (09:00) GamesCom 2017 a~naše tipy: Haimrik | (17:00) Petr versus Biomutant | (24:45) Nové Age of Empires a~Anno | (30:15) Cuphead a~Sociable Soccer | (41:15) Ancestors | (45:45) Česká stopa: Hobo, Novus Inceptio \&{} Beat Saber | (55:50) Patrik popisuje Kingdom Come: Deliverance | (62:30) Frostpunk od tvůrců This War of Mine | (68:15) Přednášky Dana Vávry a~Tima Sweeneyho | (72:30) Dotazy}{1:38:25}{34355231}{}
\FC{fight-club-339.mp3}{\FCtitle{\#339}{Návrat MMORPG}}{2017-08-30}{\Luke, \Ulrik, \OverWatch, \PHajda}{\FCrule{\#339}{Das ist Aleš.}{\Luke}}{Warhammer 40k: Inquisitor | (08:52) Richard Garriott a~návrat MMORPG | (18:12) Pokračování Rune: Ragnarok se hlásí | (28:22) Adam a~Starcraft Remastered | (37:20) Dotazy}{1:17:05}{33991288}{}
\FC{fight-club-340.mp3}{\FCtitle{\#340}{Polské karty}}{2017-09-06}{\Luke, \Ulrik, \OverWatch, \Bety}{\FCrule{\#340}{Nemrkejte zbytečně.}{\Luke}}{Warhammer 40k: Inquisitor | (08:52) Richard Garriott a~návrat MMORPG | (18:12) Pokračování Rune: Ragnarok se hlásí | (28:22) Adam a~Starcraft Remastered | (37:20) Dotazy}{1:17:05}{37312144}{}
\FC{fight-club-341.mp3}{\FCtitle{\#341}{Opevněný klub}}{2017-09-13}{\Luke, \Blixa, \PHajda, \JHamernik}{\FCrule{\#341}{Mrzký peníz nesmrdí.}{\Luke}}{Představení World of Castles s Jakubem Hamerníkem | (33:12) Žravá střílečka od Ubisoftu Atomega | (38:32) Augmentovaná realita od Applu | (50:50) Bach + Rabbids | (61:00) Dotazy}{1:18:12}{34548496}{}
\FC{fight-club-342.mp3}{\FCtitle{\#342}{Integrace s božským Karlem}}{2017-09-20}{\Sleepless, \Ulrik, \KDrda, \VPechacek}{\FCrule{\#342}{Kooperace řeší integraci.}{\Sleepless}}{Entrée Karla Drdy a~příchod Wargamingu do Prahy | (15:15) Divinity: original Sin 2 se prodává jak máslo | (34:55) Hráči bombardují recenze na Steamu | (47:55) Bruce Straley opouští Naughty Dog | (51:00) Dotazy}{1:27:09}{34916632}{2}
\FC{fight-club-343.mp3}{\FCtitle{\#343}{Změna je život}}{2017-09-27}{\Sleepless, \Blixa, \Ulrik}{\FCrule{\#343}{Ať žije šéfredaktor, ať už je to kdokoliv.}{\Blixa}}{Střídání stráží | Legends of Azulgar vážně nejsou fantasy (05:00) | Co je to ten Ataribox? (34:55) | Left Alive a~skvadra elitních Japonců (45:20) | Dotazy (51:30)}{1:19:07}{32114484}{}
\FC{fight-club-344.mp3}{\FCtitle{\#344}{Aligátoři a~písky}}{2017-10-04}{\Sleepless, \Janchor, \Ulrik, \OverWatch}{\FCrule{\#344}{Není krokodýl jako Aligator.}{\Janchor}}{Red Dead Redemption 2 | Assassin's Creed Origins vás bude učit (17:30) | Co hrajeme nebo bychom chtěli (32:05) |Spooousta dotazů (35:24)}{1:23:16}{35076148}{}
\FC{fight-club-345.mp3}{\FCtitle{\#345}{Támhle jsou dveře!}}{2017-10-11}{\Sleepless, \OverWatch, \VPechacek, \LMacura}{\FCrule{\#345}{Když něvíš, kde jseš, tak jseš v Brně.}{\VPechacek}}{Rytmik | Wolfenstein a~náckové (18:40) | Cuphead je těžká\,\dots{} ale je to zábavné? (29:32) | Brno, Gamerpie a~my (46:30) | Dotazy (56:30)}{1:23:26}{35196820}{1}
\FC{fight-club-346.mp3}{\FCtitle{\#346}{Slabá síla}}{2017-10-18}{\Luke, \OverWatch, \VPechacek, \PHajda}{\FCrule{\#346}{Každý týden jiný satan.}{\Luke}}{EA ruší očekávaný Star Wars titul (09:23) | Jak se líbí ELEX  (24:34) | Adam hodnotitel: Ghost Recon Wildlands: Ghost War (36:25) | Dotazy (47:01)}{1:14:24}{31873084}{1}
\FC{fight-club-347.mp3}{\FCtitle{\#347}{Kompatibilita vítězná}}{2017-10-25}{\Sleepless, \Luke, \Ulrik, \OverWatch}{\FCrule{\#347}{Xylofon do každé rodiny.}{\Luke}}{(01:35) Unboxing Levelu \#279 | (04:50) Hudební adventura Beat the Game | (17:20) Zpětná kompatibilita Xbox | (26:19) Single player hry vs hry jako služby | (46:12) Co budeme v pátek hrát? | (51:02) Dotazy}{1:04:02}{26591056}{1}
\FC{fight-club-348.mp3}{\FCtitle{\#348}{Regnum Regina}}{2017-11-01}{\Ulrik, \OverWatch, \VPechacek}{\FCrule{\#348}{Networkinguj dokud nepadneš.}{\VPechacek}}{Regina chce dobrovolníky | (12:10) Adam chce velký Blizzcon | (16:00) Sony překvapila | (46:50) Násilné hry! | (1:06:02) Dotazy}{1:20:09}{35899153}{}
\FC{fight-club-349.mp3}{\FCtitle{\#349}{X Box One X ASAP!}}{2017-11-08}{\Sleepless, \Luke, \OverWatch, PBarak}{\FCrule{\#349}{Sežeňte si kontext.}{\Luke}}{(02:08) Pavel Barák o GDS2017: Větší a~lepší! | (29:05) Splněný Adamův sen: Návštěva Blizzcon 2017 | (45:31) X box One X: Ano nebo ne? | (56:44) HOMAM + Civilizace = Battle of Polytopia | (60:45) Dotazy}{1:21:59}{37037356}{}
\FC{fight-club-350.mp3}{\FCtitle{\#350}{Bacha na Teraflopy}}{2017-11-18}{\Janchor, \Ulrik, \OverWatch, \Luisa}{\FCrule{\#350}{Skřety v krabicích nechceme.}{\Janchor}}{(02:10) Co s těmi lootboxy? | (33:44) Respawn už patří Electronic Arts | (43:40) Dotazy}{1:17:45}{33820914}{}
\FC{fight-club-351.mp3}{\FCtitle{\#351}{Gambleboxy}}{2017-11-22}{\Sleepless, \Luke, \OverWatch, \VPechacek}{\FCrule{\#351}{Lootboxy jsou temná strana síly.}{\Luke}}{(01:22) Level \#280 | (02:36) Sklizeň z GDS 2017 | (34:10) Need for Speed, Battlefront a~lootboxy | (64:45) Dotazy}{1:30:24}{39680692}{}
\FC{fight-club-352.mp3}{\FCtitle{\#352}{Bětka Force}}{2017-12-03}{\Sleepless, \Luke, \OverWatch, \Bety}{\FCrule{\#352}{Pozor na nahé vlkodlaky.}{\Luke}}{(04:53) Málo rychlosti a~zběsilosti: Sonic Forces | (17:45) Hitman nevyhynul! | (22:20) Hellblade ukazuje cestu | (43:00) Bad cop! L. A. Noire na Switch | (49:12) Óda na Ode | (52:35) Dotazy}{1:19:43}{35207188}{2}
\FC{fight-club-353.mp3}{\FCtitle{\#353}{Stavitelé mostů}}{2017-12-09}{\Luke, \OverWatch, \PHajda, \Luisa}{\FCrule{\#353}{Kdo chce být Mario, musí do Tokya.}{\Luke}}{(01:57) Mario Kart IRL | (09:31) Postavte si most s GLaDOS | (14:58) Jak zabít boha v Assassin's Creed: Origins |(23:30) Patrik má radost: Division se daří | (31:37) Dotazy}{1:03:13}{26415844}{}
\FC{fight-club-354.mp3}{\FCtitle{\#354}{V hlavě pana Kojimy}}{2017-12-14}{\Luke, \Ulrik, \VPechacek, \Luisa}{\FCrule{\#354}{Nekupujte hry v Day One.}{\Luke}}{(03:47) WTF Death Stranding | (15:16) Falling Sky podle vzoru Twin Peaks | (22:23) Dojmy z bitvy v Ancestors | (37:19) Návraty ke starým hrám | (53:39) Dotazy}{1:27:48}{37640284}{2}
\FC{fight-club-355.mp3}{\FCtitle{\#355}{Tady je Conanovo}}{2017-12-21}{\Sleepless, \Luke, \Ulrik, \OverWatch}{\FCrule{\#355}{Abyste si ty svátky fakt užili, jsou jenom jednou za rok.}{\Luke}}{(01:41) Nový Level 281 | (07:10) Co bude s Conanem a~odkazem R. E. Howarda ve hrách? Funcom ví! | (31:56) Komiksová Gorogoa nám dělá radost! | (34:50) Moleman Longplay: Dokument o maďarských hrách | (41:58) Volná debata o roce 2018 a~2017 | (50:39) Dotazy}{1:13:21}{30943292}{}
\FC{fight-club-356.mp3}{\FCtitle{\#356}{Hon na predátora}}{2018-01-03}{\Luke, \Ulrik, \OverWatch, \PHajda}{\FCrule{\#356}{Nežerte kočky, nosí to smůlu.}{\Luke}}{(01:18) Best of\,\dots{} články jsou tu! | (03:44) WHO definuje závislost na hrách | (18:44) Predátor v Ghost Recon: Wildlands | (28:48) Ogre Tactics: Let Us Cling Together | (30:43) Adam opět dohrává Zeldu (opět) | (32:38) Patrik hraje Bridge Constructor Portal  | (35:51) Pavel ztracen v Egyptě | (43:20) Patrik a~deskovka This War of Mine | (51:49) Hry, na které se těšíme v roce 2018 | (59:31) Dotazy}{1:22:18}{34234304}{2}
\FC{fight-club-357.mp3}{\FCtitle{\#357}{Čučel}}{2018-01-11}{\Sleepless, \Janchor, \OverWatch, \VPechacek}{\FCrule{\#357}{Nejlepší pražská čtvrť je Sačdol.}{\Janchor}}{Servisní okénko | (06:40) Sneaky Adam | (13:10) Češi ve světě, Atentát a~ti další | (44:20) Dotazy všeho druhu}{1:10:53}{47810689}{2}
\FC{fight-club-358.mp3}{\FCtitle{\#358}{Remcání na remastery}}{2018-01-17}{\Luke, \Ulrik, \OverWatch, \PHajda}{\FCrule{\#358}{Káčo, Káčo, Mamma Mia!}{\Luke}}{(01:38) Level \#282 | (08:29) Robot Cache: Kryptoměna a~distribuční systém od Briana Farga | (20:17) Two Point Hospital: Nástupce Theme Hospital | (28:05) Strategická bajka Tooth and Tail | (37:12) Remaster Dark Souls | (44:35) Dotazy}{1:11:51}{31515872}{}
\FC{fight-club-359.mp3}{\FCtitle{\#359}{Core Loop}}{2018-01-24}{\Luke, \Ulrik, \OverWatch, \VPechacek}{\FCrule{\#359}{Obalte si to papírem.}{\Luke}}{(02:35) Zážitky z PUBG a~Fortnite | (24:47) Filmy podle her (Tomb Raider a~Duke Nukem aj) | (38:20) Papír a~hry: Nintendo LABO | (45:10) Největší bitva v EVE Online (která nevyšla) | (55:12) Dotazy}{1:31:30}{38611420}{}
\FC{fight-club-360.mp3}{\FCtitle{\#360}{Příští pětiletka}}{2018-01-31}{\Luke, \Ulrik, \OverWatch, \VPechacek}{\FCrule{\#360}{Vzpomínky na budoucnost nejsou nebezpečné.}{\Luke}}{(02:42) Microsoft versus kamenné obchody | (28:38) Ubisoftí SAM a~inteligentní asistenti pro hráče | (37:43) Umělá inteligence dělá hry | (48:13) Facebook jde do stremování esportů | (58:06) Dotazy}{1:28:38}{38011552}{}
\FC{fight-club-361.mp3}{\FCtitle{\#361}{Ponory}}{2018-02-07}{\Luke, \OverWatch, \VPechacek, \PHajda}{\FCrule{\#361}{Nedívejte se dlouho do hlubin, nebo se hlubiny podívají do vás.}{\Luke}}{(06:32) Global Game Jam 2018: Přes 8 000 her! | (16:05) CRPG Book: 400 RPGček na 500 stránkách | (18:51) Dojmy z Age of Empires Definite edition | (29:43) Ve světě monster: Monster Hunter World | (43:00) Ponory v Subnautice | (60:43) Dotazy}{1:24:38}{36707324}{}
\FC{fight-club-362.mp3}{\FCtitle{\#362}{O Kingdom Come s Martinem Klímou}}{2018-02-14}{\Sleepless, \Luke, \PHajda, \MKlima}{\FCrule{\#362}{Tvé království již přišlo.}{\Luke}}{Povídání o vývoji a~vydání Kingdom Come: Deliverance -- od začátku do konce}{1:35:25}{41690464}{}
\FC{fight-club-363.mp3}{\FCtitle{\#363}{Original Syn}}{2018-02-21}{\Sleepless, \Luke, \Blixa, \Ulrik}{\FCrule{\#363}{Neprovokujte zbytečně.}{\Luke}}{(02:57) Level \#283 | (06:25) Virtuální turistika v Assassin's Creed | (12:44) Kingdom Come v neherních médiích | (20:58) Remaster Burnout Paradise | (25:11) THQ Nordic kupuje Koch Media - a~co? | (40:27) Turecký dystopický kyberpunk ManMade | (46:23) Rozpaky z Shroud of the Avatar | (52:51) Dotazy}{1:26:04}{37094468}{}
\FC{fight-club-364.mp3}{\FCtitle{\#364}{Pandy, Lary a~Vilgefortz}}{2018-02-28}{\Janchor, \Ulrik, \OverWatch, \PHajda}{\FCrule{\#364}{Namísto Vikanders si denjte radši Nikander.}{\Janchor}}{(3:50) Filmy, filmy a~seriál | (25:00) Slitherine ve slevě | (36:00) ESRB vs mikrotransakce | (45:50) Dotazy}{1:19:30}{31175388}{2}
\FC{fight-club-365.mp3}{\FCtitle{\#365}{Optimalizované deskovky}}{2018-03-07}{\Luke, \Ulrik, \PHajda, \Alladjex}{\FCrule{\#365}{AmeriTrash vládne - Euro sucks.}{\PHajda}}{Povídání o deskových hrách, tipy na nejlepší kusy a~demytizace - 72:13 Dotazy}{1:32:09}{39021028}{1}
\FC{fight-club-366.mp3}{\FCtitle{\#366}{Schrödingerův šnek}}{2018-03-14}{\Luke, \OverWatch, \VPechacek, \Luisa}{\FCrule{\#366}{Periferie vymlacujte s citem.}{\Luke}}{(01:54) Nadšení z Into the Breach | (10:41) Perzonalizace hrdinů | (18:22) Co připravuje Nintendo | (26:44) MiniDoom II: 2D krvárna | (28:40) Odkaz Heretica se jmenuje Amid Evil | (32:13) Indický surrealismus Under a~Porcelain Sun | (38:32) Dotazy}{1:07:09}{28414580}{}
\FC{fight-club-367.mp3}{\FCtitle{\#367}{Čistka v redakci}}{2018-03-21}{\Luke, \Ulrik, \OverWatch, \PHajda}{\FCrule{\#367}{Japonské plodící hry forever.}{\Luke}}{(03:31) Level \#284: Čučl rulez! | (08:31) GDS 2018: Autoři her podvádí hráče! | (21:12) GDS 2018: Potřebujeme herní odbory? | (34:13) Ve stínu Tomb Raidera | (38:04) V-Rally 4 je tu! | (40:43) Scavengers: Koopetitivní potenciál | (43:03) Nezávislé veselice na Nintendo Switch | (47:42) Stárnutí ve virtuálních světech | (56:27) Dotazy}{1:16:21}{32605700}{}
\FC{fight-club-368.mp3}{\FCtitle{\#368}{Vrátili se z Ameriky}}{2018-03-28}{\Sleepless, \Blixa, \OverWatch, \VPechacek}{\FCrule{\#368}{Ať žije čínská kavárna.}{\VPechacek}}{Kluci se vrátili z GDC a~radostně sdělují své dojmy!}{1:27:59}{55598236}{}
\FC{fight-club-369.mp3}{\FCtitle{\#369}{Duben, s Kratem budem}}{2018-04-04}{\Ulrik, \OverWatch, \VPechacek, \PHajda}{\FCrule{\#369}{Kratos je Nietzsche bez kníru.}{\Ulrik}}{Dubnový výhled a~spousta novinek, plus pravidelné dotazy}{1:04:16}{41669800}{}
\FC{fight-club-370.mp3}{\FCtitle{\#370}{Defaultní podcast}}{2018-04-11}{\Luke, \Ulrik, \OverWatch, \PHajda}{\FCrule{\#370}{Patro nahoře neexistuje.}{\Luke}}{(01:55) Valve a~soukromí | (10:46) Toxické komunity | (27:11) Radical Heights | (36:01) Průlet novinkami | (54:45) Dotazy}{1:18:43}{33802340}{}
\FC{fight-club-371.mp3}{\FCtitle{\#371}{Jarní přemítání}}{2018-04-18}{\Janchor, \Ulrik, \OverWatch, \Luisa}{\FCrule{\#371}{Hmm, prostě\,\dots}{\Janchor}}{(01:50) Nový formát Board Game Club | (03:00) Longplay XCOM 2 s herní redakcí | (07:00) Problémové změny Gwentu | (13:08) Další karetní hry: Artifact od Valve, Elder Scrolls: Legends | (14:10) Fenomén Battle Royale -- Přibude do COD: Black Ops? | (21:00) Nový stupeň hodnocení do Levelu -- Hmm, prostě\,\dots Meh. | (21:20) EA se poučili | (28:00) Retro novinky -- konečné řešení pro pavouky v Thiefovi | (29:30) Stížnost na Sama Fishera ve Wildlands | (35:50) Rampage a~hodnocení filmů na předlohy videoher na Rotten Tomatoes | (43:45) Dotazy}{1:06:55}{28716876}{}
\FC{fight-club-372.mp3}{\FCtitle{\#372}{Festivalová sezóna s Michalem Berlingerem}}{2018-04-25}{\Sleepless, \Luke, \Ulrik, \MBerlinger}{\FCrule{\#372}{Prachový králíček na hrad!}{\Luke}}{(01:04) Nový Level 285 | (05:49) Ocenění Someday You'll Return na konferenci Reboot Develop | (15:00) Ancestors od autora Assassin's Creed | (19:26) Co přinese Game Day na Anifilmu v Třeboni | (30:27) Under Leaves a~Hravouka od Michala Berlingera | (41:07) Seriál Zaklínač od Netflixu a~film podle Chuchla | (51:27) Dotazy}{1:21:02}{34355984}{2}
\FC{fight-club-373.mp3}{\FCtitle{\#373}{Jen houšť a~větší hry!}}{2018-05-02}{\Luke, \OverWatch, \VPechacek, \PHajda}{\FCrule{\#373}{Ion nebo Iron -- hlavně, že je Maiden.}{\Janchor}}{(3:00) God of War je boží! Anebo...? | (13:33) Otazníky nad Total War | (28:50) Invalidní SteamSpy se vrací | (33:15) Těšíme se na Phoenix Point! | (39:51) Nintendo LABO | (45:57) Dotazy | (55:28) Nový trailer na Red Dead Redemption 2 | (60:17) Pokračování dotazů}{1:06:13}{28359284}{}
\FC{fight-club-374.mp3}{\FCtitle{\#374}{Frakční klub}}{2018-05-09}{\Luke, \OverWatch, \VPechacek, \Luisa}{\FCrule{\#374}{Predátor není vetřelec.}{\Luke}}{(01:36) Au! Nintendo Switch Online | (11:26) Hardcore české RPG Lost Hero | (21:22) Česká inženýřina Volcanoids | (32:18) Frakce v RPGčkách a~my | (53:46) Dotazy}{1:23:57}{35378044}{2}
\FC{fight-club-375.mp3}{\FCtitle{\#375}{Ijáčkův klub}}{2018-05-16}{\Luke, \Ulrik, \PHajda, \Luisa}{\FCrule{\#375}{Bez S.T.A.L.K.E.R.a není sranda.}{\Luke}}{(04:07) Rage nad Rage 2 | (18:11) Stalker 2 se velmi pomalu blíží! | (28:56) Smích a~pláč: Cliff Blezsinski zavírá Boss Key | (32:40) Pády legendárních vývojářů | (49:49) Dotazy}{1:22:45}{34991800}{}
\FC{fight-club-376.mp3}{\FCtitle{\#376}{Replikant v klubu}}{2018-05-23}{\Sleepless, \Luke, \OverWatch, \VPechacek}{\FCrule{\#376}{Digitální hřiště do každé rodiny.}{\Luke}}{(01:21) Nový Level \#286 | (05:59) Česká hra roku 2018 již dnes! | (12:52) Erotika na Steamu a~parenting | (34:46) Soumrak PlayStation 4 | (47:25) Možná budoucnost: Cloudové hraní na Switch | (52:35) Dotazy}{1:20:33}{35469268}{}
\FC{fight-club-377.mp3}{\FCtitle{\#377}{V menšině}}{2018-05-30}{\Blixa, \OverWatch, \PHajda, \Luisa}{\FCrule{\#377}{Znejte svůj web.}{\OverWatch}}{(00:00) Jaká byla Česká hra roku | (06:00) Další české hry! | (16:00) Battlefield V a~jeho trable | (32:20) Detroit: Become Human | (37:00) Dotazy}{1:12:46}{30893272}{}
\FC{fight-club-378.mp3}{\FCtitle{\#378}{Závodní klub}}{2018-06-06}{\Luke, \Ulrik, \OverWatch, \PHajda}{\FCrule{\#378}{Čokoládové nalýže pro všechny.}{\Luke}}{(04:07) Rage nad Rage 2 | (18:11) Stalker 2 se velmi pomalu blíží! | (28:56) Smích a~pláč: Cliff Blezsinski zavírá Boss Key | (32:40) Pády legendárních vývojářů | (49:49) Dotazy}{1:13:01}{31401232}{}
\FC{fight-club-378-E3-1.mp3}{\FCtitle{E3 2017 speciál \#1}{}}{2018-06-11}{\Sleepless, \Janchor, \JIlavsky, \Luisa}{}{První speciál k E3 s hostem Jánem "Splitem" Ilavským}{1:17:35}{32565486}{}
\FC{fight-club-378-E3-2.mp3}{\FCtitle{E3 2017 speciál \#2}{}}{2018-06-12}{\Sleepless, \Janchor, \PBulir, \LMacura}{}{Druhý speciál věnující se E3 s hosty Petrem Bulířem a~Lukášem Macurou}{1:07:03}{28362306}{}
\FC{fight-club-378-E3-3.mp3}{\FCtitle{E3 2017 speciál \#3}{}}{2018-06-13}{\Blixa, \Ulrik, \OverWatch, \Luisa}{}{Třetí speciál věnující se E3 s hostem Martinem Bachem}{0:53:42}{21709506}{}
\FC{fight-club-378-E3-4.mp3}{\FCtitle{E3 2017 speciál \#4}{}}{2018-06-14}{\Janchor, \Moolker, \VPechacek, \DVavra}{}{Čtvrtý speciál věnovaný E3 s hosty Milošem Bohoňkem a~Davidem Vávrou {\em | (34:00) Malby Axela Sauerwalda}}{1:12:54}{30084438}{}
\FC{fight-club-379.mp3}{\FCtitle{\#379}{PostEtrojkový klub s Alžbětou}}{2018-06-20}{\Luke, \Bety, \VPechacek, \PHajda}{\FCrule{\#379}{Čím špinavější tím lepší.}{\Luke}}{(03:22) Nejhorší prezentace E3: Last of Us 2 | (07:37) Cyberpunk 2077 | (19:09) E3 sběr: Planet Alpha, Control, Beyond Good \&{} Evil 2, Sekiro, Plague Tale, Sinking City, Dying Light, Skull \&{} Bones, Anthem a~další | (59:02) Dotazy}{1:20:03}{34328552}{}
\FC{fight-club-380.mp3}{\FCtitle{\#380}{Zvídavý klub s Martinem Klímou}}{2018-06-27}{\Sleepless, \Luke, \OverWatch, \MKlima}{\FCrule{\#380}{Každý musí do Ratají.}{\Luke}}{Život ve Warhorse, úspěch Kingdom Come, nová DLC, E3 a~spousta informací a~dotazů}{1:18:19}{33401608}{}
\FC{fight-club-381.mp3}{\FCtitle{\#381}{Zobací klub}}{2018-07-04}{\Luke, \Ulrik, \OverWatch, \VPechacek}{\FCrule{\#381}{Zabraňte kolonizaci!}{\Luke}}{(01:53) Dokáží nezávislí vývojáři těžit z velkých producentů? | (20:31) Problém kolonizace ve videohrách | (49:45) Co právě hrajeme (Octopath Traveler, Ultimate General a~další) | (60:04) Dotazy}{1:19:47}{34499464}{}
\FC{fight-club-382.mp3}{\FCtitle{\#382}{Bilanční klub}}{2018-07-11}{\Luke, \Ulrik, \OverWatch, \PHajda}{\FCrule{\#382}{Pravý Spiderman žije v Liberci.}{\Luke}}{(01:30) Rozlučka s Petrem | (04:44) K jakým letošním hrám se chceme vrátit | (14:16) Prozatímní faily roku 2018 | (20:27) Na co se těšíme | (43:28) Dotazy}{1:14:46}{31268104}{}
\FC{fight-club-383.mp3}{\FCtitle{\#383}{Petrova rozlučka se svobodou}}{2018-07-18}{\Sleepless, \Janchor, \OverWatch, \Farid, \VPechacek}{}{(03:45) Pozdravy starých známých | (06:10) Hry roku 2002 | (36:50) Fotbalová nostalgie bez patosu | (58:37) Co bylo a~bude | (87:45) Není slávy bez dortů}{1:29:29}{57161226}{3}
\FC{fight-club-384.mp3}{\FCtitle{\#384}{Fandovský klub}}{2018-07-25}{\Luke, \Ulrik, \VPechacek, \Luisa}{\FCrule{\#384}{Vašek tančí o půlnoci.}{\Luke}}{(02:27) Streamování Scarlett Cloud pro Xbox od Microsoftu | (12:35) Hejtujeme Vampyr nebo ne? | (15:49) Return of the Tentacle a~šance fandovských her | (38:39) Dobrovský ve Fortnite | (50:41) Gnog je tu! | (51:56) Dotazy}{1:21:21}{33592804}{}
\FC{fight-club-385.mp3}{\FCtitle{\#385}{Žíznivý klub}}{2018-08-01}{\Luke, \Ulrik, \OverWatch, \VPechacek}{\FCrule{\#385}{Co na kazetě -- to na osmibitu.}{\Luke}}{(06:04) No Man's Sky Next, For Honor a~další staronová překvapení | (30:20) Láska k mytologii World of Warcraft | (45:02) Ztraceni v Grimmwood | (50:18) Dotazy}{1:20:58}{34777672}{}
\FC{fight-club-386.mp3}{\FCtitle{\#386}{Klub na páru}}{2018-08-08}{\Luke, \Ulrik, \JOlejnik, \OverWatch}{\FCrule{\#386}{Jezděte metrem.}{\Luke}}{(06:22) Fallout 76 bez Steamu! | (17:57) Soumrak Steamu! | (28:39) Návrat Aleše do Hearthstone a~World of Warcraft | (39:19) Diskuze o budoucnosti MMORPG | (50:31) Dotazy}{1:20:22}{34153360}{2}
\FC{fight-club-387.mp3}{\FCtitle{\#387}{Petr Bulíř v klubu}}{2018-08-15}{\Luke, \OverWatch, \VPechacek, \PBulir}{\FCrule{\#387}{Mutantí hlavy dávejte na zeď jako trofeje.}{\Luke}}{(02:11) Trofeje Petra Bulíře | (17:58) Na co se Petr těší | (20:51) Fallout 76 po QuakeConu | (27:11) Doom Eternal po QuakeConu | (32:23) Rage 2 po QuakeConu | (40:52) Plagiátorství v herní novinařině | (58:10) Dotazy}{1:18:27}{33773828}{}
\FC{fight-club-388.mp3}{\FCtitle{\#388}{Koření musí proudit!}}{2018-08-22}{\Janchor, \Ulrik, \VPechacek, \JSvak}{\FCrule{\#388}{Spáč se musí probudit.}{\Janchor}}{(2:00) Gamescom | (28:21) Možná přijde i Kwisatz Haderach | (35:11) The Banner Saga je u konce | (1:02:37) Dotazy}{1:27:22}{36431454}{2}
\FC{fight-club-389.mp3}{\FCtitle{\#389}{Kolínský klub}}{2018-08-29}{\Luke, \Ulrik, \OverWatch, \Luisa}{\FCrule{\#389}{Na Fiolu se nesahá.}{\Luke}}{(04:43) Gamescom a~Devcom | (06:19) Nový Level \#288 | (07:44) To nejlepší z GamesComu: Cyberpunk 2077, Assassin's Creed: Odyssey, ANNO 1800, Dying Light 2 + Bad Blood, 11-11, Twin Mirror a~další | (52:40) Dotazy}{1:18:51}{33723880}{}
\FC{fight-club-390.mp3}{\FCtitle{\#390}{Zaklínačský casting}}{2018-09-05}{\OverWatch, \VPechacek, \PHajda, \Luisa}{\FCrule{\#390}{Evičko, pojď hrát Yennefer.}{\VPechacek}}{(02:00) Spider-man | (18:00) Henry Cavill a~seriál Zaklínač | (32:00) Two Point Hospital | (37:00) Dotazy 1 | (42:00) Team Triss nebo team Yenn? | (44:30) Nové LotR MMORPG | (58:30) Dotazy 2}{1:28:08}{37804566}{}
\FC{fight-club-391.mp3}{\FCtitle{\#391}{Toulky s Larou}}{2018-09-12}{\Ulrik, \OverWatch, \PHajda, \JSvak}{\FCrule{\#391}{Pokud nám na donate pošlete půl mega, tak si zaplatíme Chrise Avelona.}{\OverWatch}}{(01:00) Shadow of the Tomb Raider | (26:00) Adorace Chrise Avellona | (42:30) Dotazy}{1:31:22}{38828722}{}
\FC{fight-club-392.mp3}{\FCtitle{\#392}{Vzpomínky na oslíky}}{2018-09-19}{\Ulrik, \OverWatch, \VPechacek, \Luisa}{\FCrule{\#392}{Nezabíjejte ženy, hrozně křičí.}{\VPechacek}}{(01:24) PlayStation Classic | (17:10) Vyvíjíme si vlastní hru! | (21:30) Assassin's Creed Odyssey | (39:40) Umíte plavat? a~jiné otázky | (55:00) Kdo tady dělá sekundární PR? | (01:11:00) Čtení ve hrách | (01:24:30) Naše oblíbené cutscény | (01:35:30) Naše oblíbené soundtracky}{1:46:49}{45173874}{}
\FC{fight-club-393.mp3}{\FCtitle{\#393}{Svět bez vlků mezi námi}}{2018-09-26}{\Janchor, \Luisa, \PHajda, \VPechacek}{\FCrule{\#393}{Když se koni svraští varlata, vemte si čepici.}{\VPechacek}}{(02:00) Nový Level | (04:30) Red Dead Redemption 2 | (31:30) Konec Telltale Games | (01:00:00) Dotazy}{1:39:03}{41956186}{}
\FC{fight-club-394.mp3}{\FCtitle{\#394}{Potterova odysea}}{2018-10-03}{\Ulrik, \Luisa, \OverWatch, \VPechacek}{\FCrule{\#394}{Sparta je ta víc homo část Řecka a taky Prahy.}{\VPechacek}}{(01:00) Fight Club 400? | Co hrajeme: (04:00) Life is Strange 2 (09:40) Spider-Man (21:00) Assassin's Creed Odyssey (38:00) Pathfinder: Kingmaker | (46:00) RPG ze světa Harryho Pottera | (54:30) Streamování her od Google | (01:02:30) GOG slaví 10. narozeniny | (01:14:00) Dotazy}{1:45:50}{45671754}{}
\FC{fight-club-395.mp3}{\FCtitle{\#395}{Hledání ztracené cesty}}{2018-10-10}{\OverWatch, \VPechacek, \PHajda, \JSvak}{\FCrule{\#395}{My pliveme na housenky.}{\VPechacek}}{(01:00) Fight Club 400! | (05:00) PlayStation 5 | (22:00) Assassin's Creed Odyssey a~mrkvička | (45:00) Pathfinder: Kingmaker | (01:07:30) Dotazy}{1:30:22}{38883378}{}
\FC{fight-club-396.mp3}{\FCtitle{\#396}{Disco na Discordu}}{2018-10-17}{\Ulrik, \VPechacek, \Luisa, \JSvak}{\FCrule{\#396}{Všichni chtějí byt Regis.}{\VPechacek}}{(02:00) Co teď hrajeme? | (08:00) Star Citizen | (17:00) Witcher casting | (32:40) Redakce točí vlastní zaklínačský seriál | (37:00) Discord jede | (54:00) Creaks od Amanity | (01:01:00) Dotazy}{1:39:40}{42274938}{}
\FC{fight-club-397.mp3}{\FCtitle{\#397}{Rozbíjíme trůny}}{2018-10-24}{\OverWatch, \VPechacek, \PHajda, \Luisa}{\FCrule{\#397}{Nešahejte Adamovi na banán.}{\VPechacek}}{(01:00) Co se děje na Games | (06:30) Starlink: Battle for Atlas | (21:00) Thronebreaker: The Witcher Tales | (39:30) Reigns: Game of Thrones | (49:00) Dotazy }{1:28:07}{38097786}{}
\FC{fight-club-398.mp3}{\FCtitle{\#398}{Podzim s kovbojem}}{2018-10-31}{\Ulrik, \OverWatch, \JOlejnik, \Luisa}{\FCrule{\#398}{V osmdesátemčtvrtém takovéhle koňské koule ještě nebyly.}{\OverWatch}}{(01:00) Henry Cavill jako Geralt | (02:30) Red Dead Redemption 2 | (56:00) Westerny | (01:01:00) Dotazy}{1:26:06}{39062046}{2}
\FC{fight-club-399.mp3}{\FCtitle{\#399}{Rok bez ďábla}}{2018-11-07}{\Ulrik, \OverWatch, \VPechacek, \PHajda}{\FCrule{\#399}{Fight Club příště jenom na mobilech.}{\VPechacek}}{(00:30) Smutný Fight Club | (01:30) Co hrajeme? | (16:00) Blizzcon | (55:00) Dotazy}{1:20:34}{36624126}{}
\FC{fight-club-400.mp3}{\FCtitle{\#400}{Sešli jsme se v továrně}}{2018-11-10}{\Sleepless, \Blixa, \JOlejnik, \OverWatch, \KDrda, \VPechacek, \PHajda, \Luisa, \JSvak}{\FCrule{\#400}{Když se všichni společně sejdou u fondue na náměstí Svatopluka, děkujem.}{\VPechacek{} aj.}}{Jubilejní Fight Club \#400 s mnoha hosty i tématy}{2:26:30}{59489964}{2}
\FC{fight-club-401.mp3}{\FCtitle{\#401}{Retrobraní}}{2018-11-14}{\Janchor, \VPechacek, \PHajda, \JZubec}{\FCrule{\#401}{Teď si nás Microsoft koupí všechny.}{\VPechacek}}{(01:00) Jaký byl Fight Club 400? | (04:00) Co hrajeme? | (23:00) Akvizice herních studií | (35:00) PC Classic | (41:00) Beta Falloutu 76 | (44:40) Česká hra Volcanoids | (48:30) DayZ | (53:00) Víme něco o Final Fantasy? | (55:00) Microsoft na nákupech | (01:16:00) Dotazy}{1:32:57}{41697786}{}
\FC{fight-club-402.mp3}{\FCtitle{\#402}{Severský bůh se světelným mečem}}{2018-11-21}{\Ulrik, \OverWatch, \VHrincar, \PHajda}{\FCrule{\#402}{Když hra na mobil, tak s chameleonem.}{\OverWatch}}{(01:00) Beat Saber a jak se vydává hra u Sony | (39:25) Sony končí s E3 | (58:00) Dotazy}{1:37:52}{46187670}{}
\FC{fight-club-403.mp3}{\FCtitle{\#403}{Konce konců}}{2018-11-28}{\Ulrik, \OverWatch, \Luisa}{\FCrule{\#403}{Když herní ceny, tak jedine na Games a letos už v prosinci.}{\OverWatch}}{(01:00) Servisní okénko | (06:00) Co nás čeká a nemine? | (32:00) The Game Awards nominace a spekulace | (01:03:00) Dotazy}{1:31:13}{41595726}{}
\FC{fight-club-404.mp3}{\FCtitle{\#404}{Nenalezen}}{2018-12-05}{\OverWatch, \PBarak, \VPechacek, \Luisa}{\FCrule{\#404}{Když je v baráku Barák, tak je vždycky legrace.}{\VPechacek}}{(01:00) Pozvánka na GDS 2018 | (13:30) Co teď hrajeme a co hraje náš host? | (25:30) Změny cenové politiky Steamu | (48:20) Dotazy}{1:16:47}{35013738}{}
\FC{fight-club-405.mp3}{\FCtitle{\#405}{}}{2018-12-12}{\OverWatch, \VPechacek, \PHajda, \Luisa}{\FCrule{\#405}{Nad Himalájí se blízká, hromy divo bijů.}{\VPechacek}}{}{1:28:08}{40615086}{}
\FC{fight-club-406.mp3}{\FCtitle{\#406}{Děti nezávislosti}}{2018-12-19}{\Ulrik, \OverWatch, \PHajda, \Luisa}{\FCrule{\#406}{Když děti tak jedině indie.}{\OverWatch}}{Kromě provozních věcí kolem her roku byly tématem jen vaše dotazy a naše plány na Vánoce, přejeme příjemné prožití vánočních svátků a vše nejlepší do nového roku:)}{1:27:09}{40236510}{}
\FC{fight-club-407.mp3}{\FCtitle{\#407}{Novoroční předsevzetí}}{2019-01-02}{\Ulrik, \OverWatch, \VPechacek, \PHajda, \Luisa, \JSvak}{\FCrule{\#407}{Bez Jidáše Poslední večeři neuděláte.}{\VPechacek}}{(01:00) Jak jsme trávili svátky? | (27:00) Na co se těšíme v roce 2019? | (38:30) Přebarvený Chuchel | (50:00) Projekt Meditations | (55:30) Dotazy}{1:28:08}{40158966}{}
\FC{fight-club-408.mp3}{\FCtitle{\#408}{Fičení PlayStationů}}{2019-01-09}{\Ulrik, \OverWatch, \PHajda, \Luisa}{\FCrule{\#408}{Až Cage udělá pokračování Deatroitu, tak to bude Become Patrik.}{\OverWatch}}{(01:00) Co vás na Games čeká a nemine? | (12:00) PlayStation 4 se dobře daří | (21:00) Co na to Microsoft? | (25:00) A co je sakra Mad Box? | (31:00) Zpětná kompatibilita | (38:00) My a módy | (40:00) Porody a genitálie | (46:30) Russia 2055 a Death Stranding | (49:00) Dotazy}{1:09:51}{32387142}{}
\FC{fight-club-409.mp3}{\FCtitle{\#409}{Sbírání brouků}}{2019-01-16}{\OverWatch, \VPechacek, \Luisa, \SKilianova}{\FCrule{\#409}{Nejlepší kouzelnické triky najdete na PornHubu.}{\VPechacek}}{(01:00) Co teď hrajeme? | (07:00) Kdo je to Yuffie? | (22:00) Jediný pravý pixelart | (30:00) Open world Star Wars nebude | (40:00) Pinkertoni proti Red Dead Redemption 2 | (44:00) Randy Pitchford byl zlobivý kluk | (52:00) Asasíni musí mít děti | (01:00:00) Novinky ve filmech podle her | (01:05:00) Dotazy}{1:29:21}{41790306}{}
\FC{fight-club-410.mp3}{\FCtitle{\#410}{Rezidentní expert}}{2019-01-23}{\Ulrik, \OverWatch, \DRynes, \Luisa}{\FCrule{\#410}{Kdo v tomhle podcastu nemluví o hrách, je hovado!}{\VPechacek}}{(03:00) Resident Evil 2 | (19:00) Ace Combat 7 | (26:00) At the Gates | (32:00) Destiny 2, The Division, Hunt: Showdown | (37:30) Hry a závažná témata | (01:03:30) Dotazy}{1:20:48}{36970806}{2}
\FC{fight-club-411.mp3}{\FCtitle{\#411}{Exodus Metra}}{2019-01-30}{\Ulrik, \VPechacek, \PHajda, \Luisa}{\FCrule{\#411}{Co řeknete, když vaše vysněná hrá nevyjde na Steamu? Boofish.}{\VPechacek}}{(01:00) Co hrajeme? | (16:30) Konec exkluzivit Quantic Dream | (21:00) Metro Exodus exkluzivně pro Epic Store | (51:00) Koncert Games vol. 2 | (01:00:30) Dotazy}{1:20:31}{37677054}{}
\FC{fight-club-412.mp3}{\FCtitle{\#412}{Filmové legendy}}{2019-02-06}{\OverWatch, \MrHlad, \VPechacek, \Luisa}{\FCrule{\#412}{Nejlepší filmový kritici jsou zároveň nejlepší kuchaři.}{\VPechacek}}{(01:00) Co hrajeme? | (03:00) Mr. Hlad taky hraje! | (16:30) Apex Legends je pecka | (26:30) The Division 2 by taky mohla být pecka | (32:00) Hrozné a méně hrozné herní filmy | (50:30) Bude hrozný i Zaklínač od Netflixu? | (01:00:00) Primárně filmový dotazník}{1:42:28}{48113958}{}
\FC{fight-club-413.mp3}{\FCtitle{\#413}{S Alžbětou Trojanovou}}{2019-02-13}{\Ulrik, \OverWatch, \Bety, \PHajda}{\FCrule{\#413}{Pokud chcete někoho zapalovat tak jednom pavouky.}{\OverWatch}}{(01:00) Máme nový vizuál | (03:00) New Game+? | (12:30) THQ Nordic koupilo Warhorse | (21:30) Metro Exodus | (41:30) Demo Anthem | (01:00:00) The Division 2 | (01:05:00) Dotazy I. | (01:19:00) Zaklínačské zastavení | (01:33:00) Dotazy II.}{1:42:28}{46545870}{}
\FC{fight-club-414.mp3}{\FCtitle{\#414}{Metro Exodus}}{2019-02-20}{\Ulrik, \Luisa, \JSvak, \PHajda}{\FCrule{\#414}{Jak by řekl Artiom...}{\Ulrik}}{(01:00) Všichni hrajeme Metro Exodus | (58:30) Dotazy}{1:09:10}{30813402}{}
\FC{fight-club-415.mp3}{\FCtitle{\#415}{S Gabrielou "Shial" Peškovou}}{2019-02-27}{\OverWatch, \Ulrik, \PHajda, \Shial}{\FCrule{\#415}{Pokud chcete ušetřit na cosplay, tak Half-Life 3.}{\OverWatch}}{(01:00) Cosplay, kam se podíváš | (15:00) Shial ale hraje i hry! | (28:00) Last Remnant | (30:00) Crackdown 3 | (35:30) Trials a Yoshi | (39:00) Apex Legends | (41:00) Anthem a budoucnost Bioware | (52:30) Duna | (58:30) Dotazy}{1:20:09}{41275830}{}
\FC{fight-club-416.mp3}{\FCtitle{\#416}{S Martinem "Markusem" Mastným o The Wild Age}}{2019-03-06}{\VPechacek, \PHajda, \Luisa, \MMastny}{\FCrule{\#416}{Hratelnost... pfff, příběh... brrr, hlavní jsou lesy.}{\VPechacek}}{(01:00) The Wild Age | (18:00) Novus Inceptio | (24:30) Když Martin nevyvíjí, hraje! | (34:30) Rallye! | (41:30) Stíhačky, horory, Civilizace | (48:00) Průšvih jménem Anthem | (56:30) Xbox One bez mechaniky? | (01:02:00) Dotazy}{1:20:09}{36308406}{}
\FC{fight-club-417.mp3}{\FCtitle{\#417}{S Petrem Poláčkem}}{2019-03-13}{\VPechacek, \Luisa, \PHajda, \Sleepless}{\FCrule{\#417}{Jedna Bohemka je málo.}{\VPechacek}}{(01:00) Co teď Petr dělá? | (06:00) GDC 2019 a očekávání | (14:30) DiRT Rally 2.0 | (19:30) A co teď Petr hraje? | (29:30) Smrdí nový Level? | (37:00) Harry Potter: Wizards United | (45:00) Kdy budou exkluzivity pro PlayStation? | (49:30) Fotbalové okénko | (51:00) Karetní hry | (59:00) Dotazy}{1:31:11}{40965006}{}
\FC{fight-club-418.mp3}{\FCtitle{\#418}{0 Google Stadia}}{2019-03-20}{\VPechacek, \OverWatch, \PHajda, \JSvak}{\FCrule{\#418}{Chcete review bombing? Recenze na Games.cz jsou bombo.}{\VPechacek}}{(01:00) Co hrajeme? | (12:00) Google Stadia | (44:00) Review bombing | (51:30) Sony nekupuje Take-Two | (55:00) Streameři jako forma propagace | (01:03:00) Halo na PC | (01:07:30) Dwarf Fortress | (01:11:00) Dotazy}{1:25:32}{38418582}{}
\FC{fight-club-419.mp3}{\FCtitle{\#419}{O GDC 2019}}{2019-03-27}{\VPechacek, \Blixa, \Sleepless, \Ulrik}{\FCrule{\#419}{Nobelova cena míru patří Rockstaru.}{\VPechacek}}{(01:00) GDC 2019: (03:00) Google Stadia a streamování | (17:00) Růst Epic Store | (37:00) Historky z natáčení | (58:00) Vampire: The Masquerade - Bloodlines 2 | (01:01:30) Dotazy}{1:27:03}{37956774}{}
\FC{fight-club-420.mp3}{\FCtitle{\#420}{S Davidem Vávrou}}{2019-04-03}{\VPechacek, \Luisa, \OverWatch, \DVavra}{\FCrule{\#420}{Známka dobré hry... WSAD.}{\VPechacek}}{(01:00) Co obnáší QA? | (17:00) Co hrajeme? | (26:30) Proč je Anthem tak špatný? | (56:30) Borderlands 3 | (59:30) Artifact a karetní hry | (01:14:00) Gamer Pie | (01:18:00) Dotazy}{1:32:30}{41642058}{}
\FC{fight-club-421.mp3}{\FCtitle{\#421}{O obtížných hrách}}{2019-04-10}{\VPechacek, \Luisa, \Ulrik, \JSvak}{\FCrule{\#421}{Take it easy and play it hard.}{\VPechacek}}{(01:00) Co hrajeme? | (06:00) Bojíme se o Dragon Age 4 | (26:00) Je Sekiro moc těžké? | (01:01:00) Dotazy}{1:27:56}{40977426}{}
\FC{fight-club-422.mp3}{\FCtitle{\#422}{O Star Wars Jedi: Fallen Order}}{2019-04-17}{\OverWatch, \Ulrik, \PHajda, \JSvak}{\FCrule{\#422}{V 8K 60 fps s 3D audiem a lagem 0,8 vteřiny a ray-tracingem budeme dělat i Fight Club... za deset let.}{\OverWatch}}{(01:00) Co hrajeme a chystáme? | (04:00) PlayStation 5 | (54:30) Star Wars Jedi: Fallen Order | (01:10:30) The Elder Scrolls: Blades | (01:19:00) Dotazy}{1:39:33}{45194718}{}
\FC{fight-club-423.mp3}{\FCtitle{\#423}{O Notre-Dame a Imperator: Rome}}{2019-04-24}{\VPechacek, \OverWatch, \Luisa, \Ulrik}{\FCrule{\#423}{Uživatelský recenze, ne ne, uživatelský pindy.}{\VPechacek}}{(01:00) Co jsme hráli o Velikonocích? | (17:00) Assassin’s Creed Unity a review bombing | (30:00) Imperator: Rome | (40:30) Apex ztrácí popularitu, Fortnite ničí vývojáře | (58:00) Dotazy}{1:30:04}{41011050}{}
\FC{fight-club-424.mp3}{\FCtitle{\#424}{O Borderlands 3 a bitvě o Winterfell}}{2019-05-02}{\VPechacek, \Ulrik, \PHajda, \Luisa}{\FCrule{\#424}{Mikrotransakce nejsou mikrotransakce.}{\VPechacek}}{(01:00) Co hrajeme? | (17:30) Borderlands 3 | (31:00) Veteráni z BioWare | (38:00) Filmový Sonic a Saints Row | (44:00) Dotazy | (53:00) Game of Thrones okénko}{1:22:19}{37366338}{}
\FC{fight-club-425.mp3}{\FCtitle{\#425}{O vysněných hrách a Assassin's Creed s vikingy}}{2019-05-09}{\VPechacek, \Luisa, \OverWatch, \Ulrik}{\FCrule{\#425}{Podporujte svoje herní novináře na Kickstarteru.}{\VPechacek}}{(03:30) Další Assassin's Creed? | (15:00) Vaškův pád Konstantinopole | (38:30) Šárčino jednorožčí pólo | (46:30) Adamův válečný tycoon | (58:00) Alešova andělská bojovka | (01:13:30) Dotazy}{1:29:55}{41902230}{}
\FC{fight-club-426.mp3}{\FCtitle{\#426}{O Ghost Recon: Breakpoint}}{2019-05-15}{\VPechacek, \JSvak, \Ulrik, \OverWatch}{\FCrule{\#426}{Kdo bydlí s Alešem zemře}{\VPechacek}}{(9:30) Ghost Recon: Breakpoint a survivaly | (31:30) Red Dead Online a GTA Online | (42:30) Hardware na E3 a downgrady | (1:02:30) Dotazy}{1:21:43}{36555942}{}
\FC{fight-club-427.mp3}{\FCtitle{\#427}{O GRIDu a přátelích Sony a Microsoftu}}{2019-05-22}{\Ulrik, \PHajda, \OverWatch, \AKlas}{\FCrule{\#427}{Když telefonovat, tak s gamingovým headsetem}{\Ulrik}}{Obsah bude :)}{1:18:17}{36152224}{}
\FC{fight-club-428.mp3}{\FCtitle{\#428}{S tvůrcem Golem VR}}{2019-05-29}{\VPechacek, \PHajda, \Luisa, \TBachtik}{\FCrule{\#428}{Když chcete dědit, běžte na Golema.}{\VPechacek}}{(08:00) Golem VR, Arachnoid a spol. | (39:00) Nilfgaardská zbroj v Zaklínači od Netflixu | (45:00) SW: KOTOR film? | (53:00) Kingdom Come, Beat Saber, Arma 3 | (01:07:00) Dotazy}{1:21:45}{37117554}{}
\FC{fight-club-429.mp3}{\FCtitle{\#429}{O významech Death Stranding a Conanovi}}{2019-06-05}{\VPechacek, \Ulrik, \JSvak, \AKlas}{\FCrule{\#429}{Nejlepší je crushovat, drivovat a hearovat lamentation.}{\VPechacek}}{(06:00) Death Stranding | (21:00) Barbar Conan Unconquered | (01:01:30) Baldur's Gate III | (01:19:00) Dotazy}{1:40:32}{46105914}{}
\FC{fight-club-430.mp3}{\FCtitle{\#430}{O E3 a jen o E3}}{2019-06-12}{\VPechacek, \Ulrik, \PHajda, \OverWatch}{\FCrule{\#430}{Některé babičky vaří květákové mozačky a jiné rozstřelují ty lidské. (Vypůjčené z Watchdogs.)}{\VPechacek}}{(02:00) Microsoft | (11:00) Ubisoft | (27:00) EA a Square Enix | (32:00) Marvel's Avengers a jiná zklamání | (47:00) Nejlepší hry výstavy | (52:00) Nenaplněné predikce | (01:10:00) Dotazy}{1:28:46}{39644382}{}
\FC{fight-club-431.mp3}{\FCtitle{\#431}{O návratu z E3}}{2019-06-19}{\VPechacek, \Luisa, \PHajda, \OverWatch}{\FCrule{\#431}{Tři tildy jsou haluz.}{\VPechacek}}{(06:30) Psychonauts 2, Watch Dogs Legion | (14:00) Cyberpunk 2077 | (23:30) Vampire: The Masquerade - Bloodlines 2 | (28:00) Avengers | (31:30) Ghost Recon Breakpoint | (45:00) Dotazy}{1:07:04}{30691938}{}
\FC{fight-club-432.mp3}{\FCtitle{\#432}{Harry Potter: Wizards Unite}}{2019-06-26}{\VPechacek, \Luisa, \PHajda, \JSvak}{\FCrule{\#432}{Nejlepší novou základní postavou je nechápající Tomáš.}{\VPechacek}}{V redakci jsme propadli nejnovější mánii z žánru augmentovaných realit - Harry Potter: Wizards Unite. Stane se druhým Pokémon GO?}{1:11:50}{30575370}{}
\FC{fight-club-433.mp3}{\FCtitle{\#433}{O dalších projektech CD Projektu}}{2019-07-03}{\VPechacek, \OverWatch, \Ulrik, \PHajda}{\FCrule{\#433}{Fight Club je lepší než Movie Zone.}{\VPechacek}}{(01:00) Omluva za zvukové potíže | (05:00) Co teď hrajeme? | (17:00) Seriálové okénko | (40:00) Komiksové okénko | (48:00) CD Projekt si protiřečí | (01:02:00) Hry na cesty | (01:15:00) Dotazy}{1:30:09}{40580706}{}
\FC{fight-club-434.mp3}{\FCtitle{\#434}{O Aledornu}}{2019-07-10}{\OverWatch, \Ulrik, \PHajda, \LStoidl}{\FCrule{\#434}{Něco končí, něco začíná a něco pokračuje dál.}{\OverWatch}}{(01:00) Co hrajeme? | (13:00) Switch Lite | (21:00) Aledorn | (51:00) Dotazy}{1:24:12}{39112950}{}
\FC{fight-club-435.mp3}{\FCtitle{\#435}{Nejen o příliš dlouhých hrách}}{2019-07-17}{\VPechacek, \OverWatch, \Ulrik, zJSvak}{\FCrule{\#435}{Je naprosto geniální a prostě nejlepsí pravidlo, který kdy zaznělo a byla by naprostá strašná škoda, kdyby vypadnul zvuk, takže jdeme na to...}{\VPechacek}}{(01:00) Co hrajeme? | (13:00) Změny v Apex Legends, tahový Pathfinder | (27:00) NeoRetro? | (38:00) Příliš dlouhé hry | (01:09:00) Dotazy}{1:35:24}{43646250}{}
\FC{fight-club-436.mp3}{\FCtitle{\#436}{S launchem Kickstarteru Extinction Protocol}}{2019-07-24}{\Ulrik, \PHajda, \Luisa, \TBachtik}{\FCrule{\#436}{Pozor na špunty.}{\Ulrik}}{(01:00) Extinction Protocol míří na Kickstarter | (41:30) Co hrajeme? | (01:02:00) Dotazy}{1:14:51}{34579758}{}
\FC{fight-club-437.mp3}{\FCtitle{\#437}{S Martinem Bachem}}{2019-07-31}{\VPechacek, \Ulrik, \OverWatch, \Blixa}{\FCrule{\#437}{Dobré matchování vede k dobrému mečování.}{\VPechacek}}{(01:00) Co hrajeme? | (29:30) Ultimate Team vydělává EA víc než FIFA | (34:00) Předplatné |  (49:00) Kasíno v GTA | (01:00:00) Zaklínačský a LotR seriál | (01:22:00) Dotazy}{1:49:47}{50349774}{}
\FC{fight-club-438.mp3}{\FCtitle{\#438}{O politické kritice her}}{2019-08-07}{\VPechacek, \OverWatch, \Luisa, \JZubec}{\FCrule{\#438}{Nejlepší anekdoty pojednávají o mrtvých dětech.}{\VPechacek}}{(01:00) Popkulturní okénko | (13:00) Greedfall a Dark Envoy | (19:00) Způsobují hry násilí? | (37:00) EVE Online a Starcraft | (54:30) Ooblets a kyberšikana | (01:12:00) Dotazy}{1:32:25}{44900814}{}
\FC{fight-club-439.mp3}{\FCtitle{\#439}{O hrách, na které se těšíme}}{2019-08-14}{\VPechacek, \PHajda, \Luisa, \OverWatch}{\FCrule{\#439}{Naše omáž ti blowne mind.}{\VPechacek}}{(01:00) S čím teď bojujeme a co hrajeme? | (18:30) Na co se těšíme v srpnu a září? | (41:30) Call of Duty: Modern Warfare | (48:00) Need for Speed Heat | (56:00) Dotazy}{1:21:08}{38545410}{}
\FC{fight-club-440.mp3}{\FCtitle{\#440}{O Gamescomu s Davidem Rynešem}}{2019-08-21}{\OverWatch, \Luisa, \JSvak, \DRynes}{\FCrule{\#440}{Já mám tady takovou krátkou literární ukázku od velice uznávaného autora a Nobelovy ceny za mír a tak prostě... Františka Kotlety z jeho série Bratrstvo krve, tentokrát z nového díu Stalingrad. Tady je předmluva: "Mámě. Té, co každému z nás řekne jeho první příběh". Kapitola se jmenuje Šváby. První čtyři řádky: "Střelím tu píču do kundy. Kurva, buzno, kundu jí nech. Tak ji vošukáš do prdele. Nejsem buzna, abych šukal do prdele."}{\OverWatch}}{(01:00) Co nás nejvíc zaujalo na Gamescomu? | (31:00) A co teď budeme hrát? | (37:00) Technologické novinky - streamování, ray tracing | (01:06:00) Dotazy}{1:34:28}{45884370}{3}
\FC{fight-club-441.mp3}{\FCtitle{\#441}{O Gamescomu, Control a Ancestors: The Humankind Odyssey}}{2019-08-28}{\Ulrik, \PHajda, \JSvak, \Luisa}{\FCrule{\#441}{Fuck Darksiders.}{\Ulrik}}{(01:00) Zážitky z Gamescomu | (38:00) Ancestors: The Humankind Odyssey | (49:30) Control | (01:08:00) Dotazy}{1:31:11}{43254066}{}
\FC{fight-club-442.mp3}{\FCtitle{\#442}{S Karlem Drdou}}{2019-09-04}{\Ulrik, \PHajda, \VPechacek, \KDrda}{\FCrule{\#442}{Letošní Gamescom byl Vietcong.}{\Ulrik}}{(00:30) Nestandardnosti | (03:00) Co nového ve Wargamingu? | (26:00) Co se Karlovi nelíbí a co hraje? | (40:00) Cyberpunk 2077 a jeho cutscény a multiplayer | (49:00) Zbytky z Gamescomu | (54:00) Dotazy}{1:27:35}{41916220}{2}
\FC{fight-club-443.mp3}{\FCtitle{\#443}{S Helenou Bendovou}}{2019-09-11}{\VPechacek, \Luisa, \Ulrik, \HBendova}{\FCrule{\#443}{Medové koláčky jedině ve tvaru myšky.}{\VPechacek}}{(01:00) Co hrajeme a proč je to WoW Classic? | (18:30) Co je nového v redakci Games? | (20:00) A co je nového v game studies? | (01:01:00) A v počítačových hrách? | (01:04:00) Dotazy}{1:24:15}{42268782}{1}
\FC{fight-club-444.mp3}{\FCtitle{\#444}{O nejlepších českých hrách}}{2019-09-18}{\VPechacek, \Luisa, \OverWatch}{\FCrule{\#444}{Když na vás bafne démon, hoďte po něm hovno.}{\VPechacek}}{(04:00) Co hrajeme? | (14:00) Na co se těšíme? | (19:00) Novinky z domova a ze světa | (30:00) Nejlepší české hry všech dob | (01:06:00) Dotazy}{1:15:47}{37966530}{1}
\FC{fight-club-445.mp3}{\FCtitle{\#445}{O předplatném a The Last of Us: Part II}}{2019-09-25}{\VPechacek, \OverWatch, \Luisa, \PHajda}{\FCrule{\#445}{Pro Boha jsou Japonci ptáci.}{\VPechacek}}{(02:00) Co hrajeme? | (30:00) The Last of Us 2 a State of Play | (44:00) Call of Duty: Modern Warfare a exkluzivity | (01:01:00) Dotazy}{1:20:55}{39642078}{}
\FC{fight-club-446.mp3}{\FCtitle{\#446}{O Ghost Recon: Breakpoint}}{2019-10-02}{\OverWatch, \Ulrik, \Luisa, \PHajda}{\FCrule{\#446}{Epitomem Fight Clubu je příliš mnoho pravidel.}{\OverWatch}}{(01:00) Za co Šárka utrácí | (09:00) Ghost Recon: Breakpoint | (25:00) Herní zklamání | (40:00) Co hrajeme? | (49:30) Dotazy}{1:10:56}{34732650}{}
\FC{fight-club-447.mp3}{\FCtitle{\#447}{S Lokim z Beat Games}}{2019-10-09}{\OverWatch, \Ulrik, \Luisa, \VHrincar}{\FCrule{\#447}{Páteř sice není vidět, ale je důležitá.}{\OverWatch}}{(01:00) Co nového v Beat Games? | (16:00) Setkání s Johnem Carmackem | (38:00) PlayStation 5 | (01:02:30) Blizzard nemá páteř | (01:28:40) Dotazy}{1:46:17}{52518198}{}
\FC{fight-club-448.mp3}{\FCtitle{\#448}{O Fortnite a League of Legends}}{2019-10-16}{\OverWatch, \PHajda, \Ulrik, \Luisa}{\FCrule{\#448}{Riot má sice projekty A až F, ale my máme projekt FC.}{\OverWatch}}{(05:00) Fortnite zamrzlo | (14:00) Co hrajeme: Disco Elysium, Switcher, GRID, Luigi’s Mansion 3 atd. | (54:00) League of Legends a novinky od Riotu | (01:14:00) Dotazy}{1:29:27}{42808386}{}
\FC{fight-club-449.mp3}{\FCtitle{\#449}{PDXCon 2019}}{2019-10-23}{\VPechacek, \OverWatch, \Luisa, \Ulrik}{\FCrule{\#449}{Nikdy neříkejte, že se nemůže nic stát.}{\VPechacek}}{(01:00) PDXCon: Crusader Kings 3 | (21:30) Stellaris | (29:50) Battletech | (36:00) BlizzCon 2019 | (57:00) Dotazy}{1:23:59}{39117270}{}
\FC{fight-club-450.mp3}{\FCtitle{\#450}{Rozhodují kostky}}{2019-10-30}{\VPechacek, \PHajda, \Luisa, \OverWatch}{\FCrule{\#450}{Nikdy nevěřte dementním kostkám.}{\VPechacek}}{450. Fight Club vychází ve speciální, randomizované edici, kde probíráme třeba velkou porci herních odkladů nebo návrat EA na Steam. V blíže neurčeném pořadí.}{1:23:28}{41690370}{}
\FC{fight-club-451.mp3}{\FCtitle{\#451}{O Death Stranding, Red Dead Redemption 2 na PC a BlizzConu}}{2019-11-06}{\OverWatch, \JSvak, \Luisa, \PHajda}{\FCrule{\#451}{Death Stranding je sice walking simulátor, ale my jsme talking simulátor.}{\OverWatch}}{(01:00) Red Dead Redemption 2 na PC | (18:30) Death Stranding | (48:00) BlizzCon: Diablo 4, Overwatch 2 atd. | (01:06:00) Dotazy}{1:24:39}{42275658}{}
\FC{fight-club-452.mp3}{\FCtitle{\#452}{S tvůrci Velvetist inspirované první Mafií}}{2019-11-13}{\OverWatch, \JSvak, \FBeyer, \TomBarsweik}{\FCrule{\#452}{Indie vývojáři se mají sametově.}{\VPechacek}}{(02:00) Nová česká hra Velvetist | (22:00) Mafia 4 | (43:00) Jaké remaky a remastery bychom si přáli? | (01:06:00) Dotazy}{1:18:41}{42077694}{}
\FC{fight-club-453.mp3}{\FCtitle{\#453}{S Pavlem Barákem}}{2019-11-20}{\OverWatch, \JSvak, \Ulrik, \PBarak}{\FCrule{\#453}{Můžeš věřit v Krista Dia, nepojede ti Stadia.}{\VPechacek}}{(01:00) Game Developers Session 2019 | (14:00) The Game Awards 2019 | (25:00) Launch Google Stadia | (49:00) Dotazy}{1:08:20}{35396994}{}
\FC{fight-club-454.mp3}{\FCtitle{\#454}{S tvůrci strategie Nebuchadnezzar}}{2019-11-28}{\OverWatch, \Ulrik, \LHroch, \JHajicek}{\FCrule{\#454}{Úrodný půlměsíc je na Facebooku.}{\VPechacek}}{(0:20) Snad už skočí ban na YT | (01:30) Česká strategie Nebuchadnezzar | (41:40) Stíhají vývojáři i hrát? | (01:00:00) Dotazy}{1:12:31}{36903810}{}
\FC{fight-club-455.mp3}{\FCtitle{\#455}{Máme rádi Sílu}}{2019-12-04}{\VPechacek, \JSvak, \Luisa}{\FCrule{\#455}{Když jsi lepší zvíře, jsi lepší člověk.}{\VPechacek}}{(00:10) Jsme na Mixéru | (01:00) Co nového v redakci? | (04:30) Co hrajeme a na co se (ne)díváme? | (13:00) Star Wars Jedi: Fallen Order | (36:00) GDS 2019 | (55:30) Dotazy}{1:05:01}{34277142}{}
\FC{fight-club-456.mp3}{\FCtitle{\#456}{BioShock, Darksiders, Sony a Nintendo}}{2019-12-11}{\OverWatch, \Ulrik, \Luisa, \PHajda}{\FCrule{\#456}{Když Game Awards, tak ne teď aleaž pěkně mezi svátkama.}{\OverWatch}}{(01:00) Prosinec na Games.cz | (05:30) Co hrajeme? (Darksiders Genesis) | (20:00) Nový BioShock | (35:00) State of Play – Resident Evil 3, Babylon’s Fall, Ghost of Tsushima… | (52:00) Dotazy}{1:03:11}{32424726}{}
\FC{fight-club-457.mp3}{\FCtitle{\#457}{Předvánoční freestyle}}{2019-12-18}{\VPechacek, \JSvak, \Ulrik, \Luisa, \PHajda, \OverWatch}{\FCrule{\#457}{Děleje to, co dělá švec se svým kopytem.}{\VPechacek}}{Poslední Fight Club letošního roku... zcela bez pravidel. A bez časových značek.}{1:17:02}{39746946}{}
\FC{fight-club-458.mp3}{\FCtitle{\#458}{Bětka a Geralt}}{2020-01-08}{\VPechacek, \Ulrik, \OverWatch, \Bety}{\FCrule{\#458}{Trápí vás hlad, smutek, bída? Řešením je genocida.}{\VPechacek}}{(01:00) Na co se těšit v roce 2020? | (38:45) Co jsme hráli přes Vánoce? | (42:00) Bětka nesnáší Zaklínače od Netflixu | (59:00) Duna | (01:03:00) Dotazy}{1:26:23}{42668670}{}
\FC{fight-club-459.mp3}{\FCtitle{\#459}{Pán prstenů - Návrat Gluma}}{2020-01-15}{\OverWatch, \Ulrik, \OverWatch, \Luisa}{\FCrule{\#459}{Ať žije re-žije.}{\VPechacek}}{(02:00) Lord of the Rings: Gollum a naši oblíbení herní Páni prstenů | (35:20) Vynikající hry, které nás nezaujaly | (56:40) Nový Xbox nebude mít exkluzivity? | (01:06:30) Dotazy}{1:26:48}{43442598}{}
\FC{fight-club-460.mp3}{\FCtitle{\#460}{O Hobo, Cyberpunku a odkladech}}{2020-01-22}{\VPechacek, \Luisa, \OverWatch, \JVasica}{\FCrule{\#460}{Oblíbené jídlo Hobo: Tough Life — dort z pejska a kočičky.}{\VPechacek}}{(02:30) Návštěva od Hobo: Tough Life | (28:30) Jarní odklady očekávaných her | (49:40) Bilancujeme s hrami roku 2019 | (01:06:00) Dotazy}{1:17:15}{37869078}{}
\FC{fight-club-461.mp3}{\FCtitle{\#461}{Coronavirus party}}{2020-01-29}{\OverWatch, \Luisa, \Ulrik}{\FCrule{\#461}{Když to tak půjde ještě dál, tak místo "Netflix and chill" budeme dávat "Netflix and pill".}{\OverWatch}}{(03:00) Co hrajeme? | (11:00) Remastery a remaky, které se nepovedly | (21:00) Netflix a hry | (48:00) Koronavirus ve hrách i mimo ně | (58:00) Dotazy}{1:20:33}{36758068}{}
\FC{fight-club-462.mp3}{\FCtitle{\#462}{Warcraft 3: Reforged se nepovedl}}{2020-02-06}{\VPechacek, \Ulrik, \JSvak, \Luisa}{\FCrule{\#462}{První titul pro Vaška, druhý pro Slávii.}{\VPechacek}}{(01:00) Vašek je magistrem | (09:30) Warcraft 3: Reforged není dobrý remaster | (31:30) Z Rockstaru odchází Dan Houser | (46:00) Mass Effect 2 slaví 10. narozeniny | (01:02:00) Dotazy}{1:23:13}{42247780}{}
\FC{fight-club-463.mp3}{\FCtitle{\#463}{O Anthem, GeForce Now a The Sims}}{2020-02-12}{\VPechacek, \JSvak, \Luisa, \OverWatch}{\FCrule{\#463}{Místo umíráčku hymna znovuzrození.}{\VPechacek}}{(06:30) Změny v Anthem | (13:00) Looter shooter Outriders | (22:30) GeForce Now opouští Activision | (31:30) Vyplácí se předplatné a streamování? | (47:00) The Sims a herní styl | (01:01:00) Dotazy}{1:12:00}{31823368}{}
\FC{fight-club-464.mp3}{\FCtitle{\#464}{Mythic Quest, hry ve hře a budoucnost E3}}{2020-02-19}{\VPechacek, \PHajda, \OverWatch, \Bety}{\FCrule{\#464}{Bětka Trojanová, novinářka, moderátorka, cosplayerka, největší noční můra vévody z Wellingtonu}{\VPechacek}}{(00:30)(04:00) Budoucnost E3 a herních výstav | (21:00) Dreams a další hry ve hrách | (43:20) Seriál Mythic Quest | (58:00) Dotazy}{1:21:55}{35667916}{}
\FC{fight-club-465.mp3}{\FCtitle{\#465}{O dobru, zlu a herních komediích}}{2020-02-26}{\VPechacek, \Ulrik, \Luisa, \OverWatch}{\FCrule{\#465}{Brno, země korony české.}{\VPechacek}}{(03:30) Jak koronavirus ovlivní herní průmysl? | (18:30) Morálka ve hrách | (47:00) Zábava, vtipy a humor: kvalitní herní komedie | (01:00:00) Dotazy}{1:12:05}{32439488}{2}
\FC{fight-club-466.mp3}{\FCtitle{\#466}{O 20. narozeninách PlayStation 2 a remaku FF7}}{2020-03-04}{\Bety, \Luisa, \Ulrik, \OverWatch}{\FCrule{\#466}{Dejte si bacha na ftáky noci.}{\Bety}}{(01:00) Bětka obohacuje naši redakci | (12:30) Remake Final Fantasy VII | (30:30) 20. výročí PlayStation 2 | (42:30) Dotazy}{1:25:45}{42810908}{}
\FC{fight-club-467.mp3}{\FCtitle{\#467}{O Horizon: Zero Dawn na PC a zrušené E3}}{2020-03-11}{\VPechacek, \PHajda, \Luisa, \OverWatch}{\FCrule{\#467}{Dávejte na sebe pozor a hodně zdraví.}{\VPechacek}}{(01:00) Jdeme do karantény? | (05:00) E3 2020 se ruší | (17:30) Horizon: Zero Dawn míří na PC | (32:00) Call of Duty: Warzone a naše vzpomínky na CoD | (50:00) Dotazy}{1:13:25}{36153284}{}
\FC{fight-club-468.mp3}{\FCtitle{\#468}{Improvizace za časů koronaviru}}{2020-03-18}{\Bety, \Ulrik, \VPechacek, \Luisa}{\FCrule{\#468}{Gulag je zvláštní komunitní zázitek.}{\Bety}}{(01:00) Jak se nám daří v karanténě? | (24:00) Call of Duty: Warzone | (33:30) Co bychom chtěli od nového Zaklínače? | (42:00) Baldur’s Gate 3 | (58:00) Představení PS5}{1:38:46}{54110192}{}
\FC{fight-club-469.mp3}{\FCtitle{\#469}{O Half-Life: Alyx a Animal Crossing}}{2020-03-25}{\Bety, \Ulrik, \VPechacek, \OverWatch}{\FCrule{\#469}{Báječná komunita na Spyware.}{\Bety}}{(04:00) Half-Life: Alyx | (13:00) Hry pro herní začátečníky | (36:00) Animal Crossing: New Horizons | (48:00) Máme Discord a dotazy}{1:25:10}{41509076}{}
\FC{fight-club-470.mp3}{\FCtitle{\#470}{Mount \& Blade II: Bannerlord}}{2020-04-01}{\Bety, \Ulrik, \VPechacek, \JSvak}{\FCrule{\#470}{Děti do sedel nepatří. Bohužel.}{\VPechacek}}{(09:30) Jirkovy dojmy z Half-Life: Alyx | (36:30) Mount \&{} Blade II: Bannerlord | (01:04:30) Jaké hry by měly změnit svůj ustálený žánr? | (01:20:30) Dotazy}{1:35:00}{45108284}{}
\FC{fight-club-471.mp3}{\FCtitle{\#471}{O ovladači k PS5 a remacích RE3 a FF7}}{2020-04-08}{\Bety, \VPechacek, \JSvak, \Ulrik, \Luisa}{\FCrule{\#471}{Kyž chceš jako mořský koník býti, musís svoje děti pod jazykem mýti.}{\Bety}}{(02:00) Jak pokračuje karanténa? | (30:00) Ovladač DualSense | (58:30) Remaky a remastery aktuálních dní: RE3, FF7, MW2 | (01:13:30) Soundtracky | (01:28:30) Dotazy}{1:43:46}{44457872}{}
\FC{fight-club-472.mp3}{\FCtitle{\#472}{O Gears Tactics, XCOM: Chimera Squad a Duně}}{2020-04-15}{\VPechacek, \PHajda, \Ulrik, \OverWatch}{\FCrule{\#472}{Délka Adamova vousu je přímo úměrná počtu snědených luskounů.}{\VPechacek}}{(01:00) Rouškoviny | (09:00) Co jsme hráli v karanténě? | (22:00) Dobročinnost herních studií | (30:30) Gears Tactics, XCOM: Chimera Squad, Crusader Kings 3 | (01:02:30) Duna}{1:28:44}{40426916}{}
\FC{fight-club-473.mp3}{\FCtitle{\#473}{O roztodivných simulátorech}}{2020-04-22}{\OverWatch, \Luisa, \PHajda, \JSvak}{\FCrule{\#473}{Když chcete jít na půnoční mši, nechoďte tam s Patrikem.}{\OverWatch}}{(01:00) Redakční aktuality | (06:00) Divné, divnější a vlastně i docela normální simulátory |  (01:04:30) Dotazy}{1:21:30}{37794272}{}
\FC{fight-club-474.mp3}{\FCtitle{\#474}{O RPG deskovce Euthia: Torment of Resurrection}}{2020-04-29}{\VPechacek, \PHajda, \Ulrik, \TSpousta}{\FCrule{\#474}{Nechceš-li vypadat jako Rus, musíš mít na roušky dobrý vzkus.}{\VPechacek}}{(04:00) Euthia: Torment of Resurrection | (27:20) Co hrajeme? Nová mechanika v MtG: Arena | (43:00) XCOM: Chimera Squad a Gears Tactics, Gwent, SnowRunner | (54:00) Dotazy}{1:12:00}{27573044}{}
\FC{fight-club-475.mp3}{\FCtitle{\#475}{O Assassin's Creed Valhalla a TLoU2}}{2020-05-06}{\VPechacek, \Luisa, \Ulrik, \OverWatch}{\FCrule{\#475}{Nekonečná zábava na zeleném pažitu? Ne raději kostičky.}{\VPechacek}}{(01:00) Nový trailer na The Last of Us: Part 2 | (16:30) Okénko pro Animal Crossing | (21:30) Assassin’s Creed Valhalla | (50:20) Dotazy}{1:21:22}{32226764}{}
\FC{fight-club-476.mp3}{\FCtitle{\#476}{O češtinách ve hrách s Filipem Ženíškem}}{2020-05-13}{\VPechacek, \Luisa, \FZenisek}{\FCrule{\#476}{Chcete poznat ve hře Filipa, pusťte si dnešní podcast na repeat ke spánku. To bude creepy.}{\VPechacek}}{Do studia jsme si pozvali Filipa Ženíška, který překládal poslední dva díly Assassin's Creed. Společně probereme, proč Valhalla tentokrát nebude mít české titulky.}{1:38:51}{42236554}{}
\FC{fight-club-477.mp3}{\FCtitle{\#477}{O Mafia: Trilogy}}{2020-05-20}{\VPechacek, \PHajda, \Luisa, \OverWatch}{\FCrule{\#477}{Musíme se navzájem krejt, TOME.}{\VPechacek}}{Remake Tony Hawk. Oznámením Mafia Trilogy se nám témata do dnešního Fight Clubu okamžitě vyřešila. Zradilo nás internetové připojení, pravidlo Fight Clubu \#477 zní: Musíme se navzájem krejt, TOME.}{1:07:56}{30421562}{}
\FC{fight-club-478.mp3}{\FCtitle{\#478}{O nové podobě Games.cz}}{2020-05-27}{\VPechacek, \PHajda, \OverWatch, \Ulrik}{\FCrule{\#478}{No tak dej si — nový Gamesi.}{\VPechacek}}{(02:00) Games.cz mají nový design | (37:00) Novinky ve (video)obsahu | (47:00) Adventurní zamyšlení | (01:15:00) Dotazy}{1:26:16}{39963445}{}
\FC{fight-club-479.mp3}{\FCtitle{\#479}{Tencent (ne)kupuje Bohemku}}{2020-06-03}{\VPechacek, \PHajda, \Ulrik}{\FCrule{\#479}{Kdo chce psa bít, zastaví o vrtící ocas.}{\VPechacek}}{(01:00) Tencent nekupuje Bohemia Interactive | (19:00) Hrajeme si se zvířátky | (47:00) Dotazy}{1:19:20}{36148273}{}
\FC{fight-club-480.mp3}{\FCtitle{\#480}{S Petrem Poláčkem z nového studia}}{2020-06-10}{\VPechacek, \Sleepless, \Ulrik, \Luisa}{\FCrule{\#480}{Když nevíš, nekomentuj.}{\VPechacek}}{(01:00) Máme nové studio a starého hosta | (05:00) Jak Tencent nekoupil Bohemku | (14:00) Jak probíhá běžný den v Bohemce a co se tam chystá? | (52:00) Na co se Petr těší, co by chtěl hrát \&{} jiné dotazy}{1:22:16}{34178821}{}
\FC{fight-club-481.mp3}{\FCtitle{\#481}{O PlayStation 5 a jeho exkluzivitách}}{2020-06-17}{\Ulrik, \JSvak, \Luisa}{\FCrule{\#481}{Televize jedině do vázy.}{\Ulrik}}{(01:00) Technické okénko | (09:00) Koupíme si PlayStation 5? | (35:00) A co si na něm zahrajeme? | (01:01:00) Dotazy}{1:29:49}{51030823}{}
\FC{fight-club-482.mp3}{\FCtitle{\#482}{O Star Wars: Squadrons}}{2020-06-24}{\VPechacek, \OverWatch, \Luisa, \Bety}{\FCrule{\#482}{Brad, pot a pixel.}{\VPechacek}}{(03:00) Oznámení z EA Play: SW Squadrons | (15:00) Skejty ve hrách | (22:00) It Takes Two | (37:00) Bětka a The Last of Us: Part 2 | (54:00) Dotazy}{1:18:46}{35411497}{}
\FC{fight-club-483.mp3}{\FCtitle{\#483}{O výročí Diabla 2}}{2020-07-01}{\Ulrik, \OverWatch, \Luisa}{\FCrule{\#483}{Stay awhile and subscribe.}{\OverWatch}}{(01:00) Co nového na Games.cz? | (10:00) 20. výročí Diabla 2 | (01:05:00) Dotazy}{1:20:20}{44049596}{}
\FC{fight-club-484.mp3}{\FCtitle{\#484}{O monetizačních modelech a RPG světě World of Darkness}}{2020-07-08}{\VPechacek, \PHajda, \Ulrik, \Luisa}{\FCrule{\#484}{Aleš — mistr jemného ochlupení.}{\VPechacek}}{(01:00) Popkulturní a novinkové okénko |  (24:00) Trackmania a zvláštní způsoby monetizace her | (53:00) World of Darkness | (01:07:00) Dotazy}{1:19:06}{40944337}{}
\FC{fight-club-485.mp3}{\FCtitle{\#485}{S Pavlem Makalem o UbiForward}}{2020-07-15}{\VPechacek, \PHajda, \OverWatch, \PMakal}{\FCrule{\#485}{Šunková či salámová, nejlepší je bagera.}{\VPechacek}}{Hlavním tahákem uplynulých dnů nebylo nic jiného než salva oznámení a novinek z akce Ubisoft Forward. Ten nakonec nebyl tak macatý jako klasická E3 tiskovka, ale přesto nabídl hned několik zajímavých her v čele s krásnými videi k Far Cry 6. A právě Ubisoft Forward a všechno kolem něj bude jedním z hlavních témat dnešního podcastu.}{1:12:25}{34772106}{}
\FC{fight-club-486.mp3}{\FCtitle{\#486}{O prvních záběrech z Mafie}}{2020-07-22}{\VPechacek, \JSvak, \OverWatch, \PMakal}{\FCrule{\#486}{Zvěstujte lakadémonským, že Kassandra je lepší.}{\VPechacek}}{(00:00) úvod | (01:00) První záběry z remaku Mafie | (14:00) Ubisoft nemá rád ženy | (47:00) Co chystá Microsoft pro Xbox Series X? | (59:00) Death Stranding na PC | (01:06:00) Dotazy}{1:11:28}{34424953}{}
\FC{fight-club-487.mp3}{\FCtitle{\#487}{O Xbox Games Showcase se Zypem}}{2020-07-29}{\Bety, \PMakal, \VPechacek, \ZyPo}{\FCrule{\#487}{Když se Pavel zakuklí, tak se máme na co těšit.}{\VPechacek}}{(00:00) úvod | (03:00) Shrnutí Xbox Games Showcase | (39:00) Co hrajeme a čteme? | (48:00) Jaké to je být herním novinářem? | (01:06:00) Dotazy}{1:17:58}{32938729}{}
\FC{fight-club-488.mp3}{\FCtitle{\#488}{Google sabotuje Attentat 1942}}{2020-08-05}{\Ulrik, \PMakal, \JGemrot}{\FCrule{\#488}{38-45-89}{\Ulrik}}{O tom, jak obtížné je dostat na Google Play hru o atentátu na Heydricha si popovídáme s hlavním programátorem Charles Games, Jakubem Gemrotem.}{1:16:51}{34420381}{}
\FC{fight-club-489.mp3}{\FCtitle{\#489}{O herním PR s Tobim}}{2020-08-12}{\Bety, \Ulrik, \Tobi}{\FCrule{\#489}{Když herní PR, nebojte se být punkový, ale hlavně pozitivně.}{\Bety}}{S naším hostem jsme na příkladu Kingdom Come zevrubně probrali, jak se dělá herní PR.}{1:35:05}{54713558}{2}
\FC{fight-club-490.mp3}{\FCtitle{\#490}{O bitvě Epicu s Applem a Googlem}}{2020-08-19}{\VPechacek, \Luisa, \PMakal, \OverWatch}{\FCrule{\#490}{Tvoje máma má ráda totální válku.}{\VPechacek}}{(00:00) Úvod, co hrajeme? | (02:40) Civilizace 6 na mobilu | (07:30) Total War Saga: Troy | (20:20) Mortal Shell a Tony Hawk | (32:00) Peaky Blinders a AC: Odyssey | (38:50) Epic vs. Apple a Google | (01:07:03) Dotazy}{1:18:23}{40945309}{}
\FC{fight-club-491.mp3}{\FCtitle{\#491}{O Creaks a Amanitě s vývojáři z Amanity}}{2020-08-26}{\VPechacek, \PMakal, \RJurda, \JChlup}{\FCrule{\#491}{Nejlepší je, když to skřípe.}{\VPechacek}}{Hry z produkce českého studia Amanita Design jsou světovým unikátem a říct to můžeme i o jeho nejnovějším kousku Creaks. A právě o tomto studiu a jeho tvorbě bude dnes řeč.}{1:01:40}{32747538}{}
\FC{fight-club-492.mp3}{\FCtitle{\#492}{S Janem Rückerem o testování her}}{2020-09-02}{\VPechacek, \Bety, \JRuecker}{\FCrule{\#492}{Kdyby ho testoval Honza, tak už tady Covid nemáme.}{\VPechacek}}{Jak se před vydáním i po něm testovalo české Kingdom Come? Ve studiu přivítáme šéfa testerů ve Warhorse a popovídáme si o tom, jak se hledají bugy.}{1:17:26}{41289577}{}
\FC{fight-club-493.mp3}{\FCtitle{\#493}{O Xboxu Series X a S}}{2020-09-09}{\Bety, \VPechacek, \PMakal, \JSvak}{\FCrule{\#493}{Průduchy a incest y jsou dobrá věc.}{\VPechacek}}{(00:00) Úvod: Co hrajeme – Wasteland 3 | (05:00) Microsoft Flight Simulator 2020 | (13:30) Iron Harvest | (18:10) Tony Hawk Pro Skater 1+2 | (29:30) Xbox Series X a S | (1:00:10) Dotazy}{1:27:00}{44686789}{}
\FC{fight-club-494.mp3}{\FCtitle{\#494}{S Jirkou Pavlovským o Marvelu v komiksu i ve hrách}}{2020-09-16}{\VPechacek, \PHajda, \PMakal, \JiP}{\FCrule{\#494}{Myšpulín je tatínek, Fifinka je maminka, Bobík zuří, Pinďa hoří.}{\VPechacek}}{(00:00) Úvod: Jirka Pavlovský a vydavatelství CREW | (04:55) Hry podle komiksů nejen od Marvelu | (36:20) Komiksy podle her | (45:00) Naše první komiksy | (1:00:00) Dotazy}{1:15:17}{38348953}{}
\FC{fight-club-495.mp3}{\FCtitle{\#495}{O nákupu Zenimaxu a Bethesdy Microsoftem a nedostatku PS5}}{2020-09-23}{\VPechacek, \Ulrik, \Luisa, \PMakal}{\FCrule{\#495}{Když máš Bethesdu, můžeš si dovolit být mrňavej a měkej.}{\VPechacek}}{Microsoft se tenhle týden postaral o pořádný poprask, když oznámil akvizici firmy ZeniMax za 7,5 miliardy dolarů. Tahle firma má pod palcem Bethesdu, tedy vydavatele stojícího za hrami jako Doom, Dishonored a především The Elder Scrolls nebo Fallout.}{1:04:39}{32539957}{}
\FC{fight-club-496.mp3}{\FCtitle{\#496}{S naším hodnocením Mafia: Definitive Edition}}{2020-09-30}{\VPechacek, \PMakal, \Luisa}{\FCrule{\#496}{Fight Club není cuddle club.}{\VPechacek}}{Uplynulý týden ovládla v herním světě Mafie a my na ni tedy máme názor nad rámec recenze, videorecenze a dalších videí.}{1:06:01}{34540972}{}
\FC{fight-club-497.mp3}{\FCtitle{\#497}{S Pavlem Barešem, českým autorem sci-fi}}{2020-10-07}{\VPechacek, \PMakal, \PBares}{\FCrule{\#497}{Hnědý kabát je bližší než hnědá košile.}{\PMakal}}{Podcast serveru Games.cz tentokrát přivítá hosta ze světa literatury, Pavla Bareše, autora série Kronos či superhrdinského thrilleru Meta.}{}{}{}
\FC{fight-club-498.mp3}{\FCtitle{\#498}{O filmech podle Monster Huntera, Kingdom Come: Deliverance a anime Dragon's Dogma}}{2020-10-14}{\Bety, \VPechacek, \JSvak, \OverWatch, \PMakal}{\FCrule{\#498}{Zpropadená cenzura Fight Club dnes vypnula.}{\Bety}}{Vracíme se do preventivní karantény, ale o Fight Club vás neochudíme ani tentokrát. Probírat budeme čerstvé filmové a seriálové adaptace herních značek.}{}{}{}
\FC{fight-club-499.mp3}{\FCtitle{\#499}{O problémech Vampire: The Masquerade - Bloodlines 2 a herním kalendáři}}{2020-10-21}{\PMakal, \Ulrik, \VPechacek, \Luisa}{\FCrule{\#499}{Lockdown nebo rajs hlavní je mít plnej špajz.}{\Luisa}}{Mimo aktuálního dění z herního světa probereme i to, co chystáme na 500. jubileum Fight Clubu během zuřící pandemie.}{}{}{}
\FC{fight-club-501.mp3}{\FCtitle{\#501}{O dalším odkladu Cyberpunku 2077 a o Watch Dogs: Legion}}{2020-10-28}{\Bety, \Ulrik, \PMakal, \VPechacek}{\FCrule{\#501}{Buď jak Fassbender a nepoužívej hole při golfu.}{\Bety}}{Výroční Fight Club \#500 z pandemických důvodů přeskakujeme, ale nebojte, při první příležitosti si ho vynahradíme. | (03:00) Watch Dogs: Legion | (23:00) Jak dopadne Assassin’s Creed Valhalla? | (59:00) Kolikátý už odklad Cyberpunku 2077? | (01:11:00) Control na Switchi, náhodné novinky a dotazy}{}{}{}
\FC{fight-club-502.mp3}{\FCtitle{\#502}{Před příchodem PlayStation 5 a Xboxu Series X/S}}{2020-11-04}{\VPechacek, \Luisa, \Ulrik, \PHajda}{\FCrule{\#502}{Správná babička vaří rajskou a paří Banana Grams.}{\VPechacek}}{Nesnášíme loučení, jak praví jedna česká píseň. I proto se PS4 a XONE ještě tak úplně nezbavujeme. Nebo ano? | (00:00) Úvod | (02:00) Co hrajeme? | (16:00) Loučíme se s PS4 a Xboxem One | (24:00) Xbox Series X/S, nebo PlayStation 5? | (33:30) Baterky?! | (35:40) Ekonomické hraní | (44:40) Konzole jako multimediální centrum | (54:00) Podivný fenomén Genshin Impact | (01:10:00) Dočkáme se do konce roku ještě nějakých hitů?}{}{}{}
\FC{fight-club-503.mp3}{\FCtitle{\#503}{O Assassin's Creed Valhalla a nových Xboxech}}{2020-11-11}{\VPechacek, \PMakal, \JSvak, \JSlavik}{\FCrule{\#503}{Kdo užívá nový X-Box, ten si to užívá.}{\VPechacek}}{Zavítal k nám nový díl Assassin's Creed a taky start nové konzolové generace v podobě Xbox Series X.}{}{}{}
\FC{fight-club-504.mp3}{\FCtitle{\#504}{O výuce herního vývoje na středních školách}}{2020-11-18}{\Bety, \VPechacek, \OKlus, \FDufek}{\FCrule{\#504}{Když něco chceš jdi direkt za direktorem.}{\VPechacek}}{V podcastu Fight Club přivítáme hned 2 hosty: Ondřeje Kluse a Filipa Dufku, zakladatele oboru Game Art z brněnské Střední školy umění a designu. Jak se tvorba her učí a co budou absolventi oboru po maturitě umět?}{}{}{}
\FC{fight-club-505.mp3}{\FCtitle{\#505}{Čerstvé zkušenosti s PlayStation 5}}{2020-11-25}{\Bety, \Ulrik, \PMakal, \VPechacek}{\FCrule{\#505}{Pralinka, rozinka, čokoška, nejlépe chutná předkožka.}{\VPechacek}}{PS5 si mohlo osahat 5 členů naší redakce a všichni z něj máme dojmy. Co se Sony povedlo a co zkazilo?}{}{}{}
\FC{fight-club-506.mp3}{\FCtitle{\#506}{O předvánočním dění v redakci}}{2020-12-02}{\Bety, \PMakal, \Luisa, \Ulrik, \VPechacek}{\FCrule{\#506}{Ať jsi paní nebo pan, nejsi víc než Angolan.}{\VPechacek}}{Začal nám advent, což v herním světě znamená jediné - do konce roku nás mimo Cyberpunku už nečekají žádné velké pecky. Jak to snášíme?}{}{}{}
\FC{fight-club-507.mp3}{\FCtitle{\#507}{Vše co chcete vědět o Cyberpunku 2077}}{2020-12-09}{\Bety, \Luisa, \Ulrik, \PMakal}{\FCrule{\#507}{Když chceš sbalit lesbu, dej si šluka helia.}{\Bety}}{Fight Club nám začíná jen pár hodin před oficiálním vydáním Cyberpunku 2077 - pokecáme tak o všem, co byste před hraním měli vědět.}{}{}{}
\FC{fight-club-508.mp3}{\FCtitle{\#508}{O nejlepších hrách roku 2020}}{2020-12-16}{\Bety, \Luisa, \VPechacek, \PMakal}{\FCrule{\#508}{Komu se nepostaví z Liary, ten si nezaslouží Grey Poupon}{\VPechacek}}{Cyberpunk 2077 je konečně venku, ale naše rozhodování o nejlepších hrách za uplynulý rok teprve začíná.}{}{}{}
\FC{fight-club-509.mp3}{\FCtitle{\#509}{Vánoční speciál s cukrovím}}{2020-12-23}{\VPechacek, \Ulrik, \Luisa, \PMakal}{\FCrule{\#509}{Šťastné krásné Vánoce, nezabíjej pásovce.}{\VPechacek}}{Nesem vám podcasty, poslouchejte! Vánoce jsou za dveřmi a s nimi i speciální epizoda Fight Clubu z našich svátečních domovů.}{}{}{}
\FC{fight-club-510.mp3}{\FCtitle{\#510}{Co videohry čeká v roce 2021?}}{2021-01-06}{\Bety, \PMakal, \VPechacek, \Ulrik, \OverWatch}{\FCrule{\#510}{Pavla umlčíš jedině prachama.}{\VPechacek}}{Nový rok je tu, což kromě nového Fight Clubu znamená i nové hry. Jaké herní pecky nás letos čekají?}{}{}{}
\FC{fight-club-511.mp3}{\FCtitle{\#511}{Nejočekávanější hry roku 2021 pohledem redakce}}{2021-01-13}{\Bety, \PMakal, \Ulrik, \Luisa}{\FCrule{\#511}{Pavel je elitní eskort, úsměv hodí až od pětistovky výš.}{\Bety}}{Přiznání (alkoholik, kafoholik, tři lajny koksu). Které tituly naplánované na letošní rok nejvíc vyhlížíme? 511. díl podcastu věnujeme očekávaným hrám roku 2021.}{}{}{}
\FC{fight-club-512.mp3}{\FCtitle{\#512}{Jak dopadlo čtenářské hlasování o hře roku 2020}}{2021-01-20}{\Bety, \PMakal, \Ulrik, \Luisa, \VPechacek}{\FCrule{\#512}{Ve Fight Clubu si zazpíváme až na to budeme mít... vejce.}{\VPechacek}}{My už jsme svůj názor na nejlepší hry roku 2020 vyjádřili, ale teď je řada na vás, našich čtenářích}{}{}{}
\FC{fight-club-513.mp3}{\FCtitle{\#513}{Bojíme se s Resident Evil Village a The Medium}}{2021-01-27}{\Bety, \PMakal, \Luisa, \Ulrik}{\FCrule{\#513}{Když na vás přijde duch, nebojte se. Stačí ho jenom počůrat.}{\Bety}}{Hororový žánr má u nás v redakci své příznivce i odpůrce. Bojíme se u her rádi? A pokud ne, tak proč ne? Probereme to v dnešním Fight Clubu.}{}{}{}
\FC{fight-club-514.mp3}{\FCtitle{\#514}{Nepovedený vstup Googlu a Amazonu do herního vývoje}}{2021-02-03}{\Bety, \PMakal, \OverWatch, \VPechacek}{\FCrule{\#514}{CHceš vypadat dobře, čti Vogue, ne Mein Kampf.}{\VPechacek}}{Herní studia Amazonu jsou prodělečná i přes milardové investice, Google to své rovnou zavírá. Proč k vývoji her nestačí jen peníze?}{}{}{}
\FC{fight-club-515.mp3}{\FCtitle{\#515}{Když se míchají hry s patenty a autorským právem}}{2021-02-10}{\Bety, \PMakal, \Luisa, \PHajda, \VPechacek}{\FCrule{\#515}{Demokracie je v krizi. Bětka osvícenou císařovnou.}{\VPechacek}}{Měl by na herní mechanismy mít nárok ten, kdo s nimi přijde jako první? A jde vůbec vlastnit něco tak abstraktního? Nejen patenty řešíme v 515. díle podcastu Fight Club, kde dojde řeč i na dema z končícího Steam Game Festivalu.}{}{}{}
\FC{fight-club-516.mp3}{\FCtitle{\#516}{Vikinský Valheim ovládl Steam}}{2021-02-17}{\Bety, \VPechacek, \PMakal, \Ulrik}{\FCrule{\#516}{Pavel zabij jelena, svět závistí zelená.}{\VPechacek}}{Vikingové jsou vděčné herní téma, jak dokazuje God of War, AC Valhalla a nejčerstvěji také Valheim. Proč nás to táhne na sever?}{}{}{}
\FC{fight-club-517.mp3}{\FCtitle{\#517}{Diablo II Resurrected, The Burning Crusade Classic a další novinky z BlizzConline}}{2021-02-24}{\PMakal, \Luisa, \JSvak, \VPechacek}{\FCrule{\#517}{Nevyměňuj staré ďábly za nové.}{\VPechacek}}{O víkendu proběhl BlizzConline 2021, kde Blizzard odhalil svůj chystaný herní line-up. Co čeká v brzké době Diablo, Hearthstone či WoWko?}{}{}{}
\FC{fight-club-518.mp3}{\FCtitle{\#518}{Nový obor herní design na FAMU s Michalem Berlingerem}}{2021-03-03}{\Bety, \VPechacek, \MBerlinger}{\FCrule{\#518}{Možná přijde i Dan Vávra.}{\VPechacek}}{Pražská FAMU po letech příprav otevírá obor herní design a my máme na jednoho z jeho pedagogů Michala Berlingera spoustu zvídavých otázek.}{}{}{}
\FC{fight-club-519.mp3}{\FCtitle{\#519}{Jak bude hraní vypadat v roce 2031?}}{2021-03-10}{\Bety, \Ulrik, \VPechacek, \PMakal}{\FCrule{\#519}{Pavel se seká, ale okýnka nerozprcává.}{\VPechacek}}{V roce 2011 jsme si asi nepředstavovali, že rok 2021 strávíme v karanténě. Jak bude herní svět vypadat za dalších 10 let?}{}{}{}
\FC{fight-club-520.mp3}{\FCtitle{\#520}{Co jsme hráli jako děti?}}{2021-03-17}{\Bety, \JSvak, \VPechacek, \PMakal}{\FCrule{\#520}{Sláva kresleným.}{\VPechacek}}{Co nás prapůvodně přivedlo k hraní her a časem k naší vysněné práci? Co, kde a jak jsme hráli poprvé?}{}{}{}
\FC{fight-club-521.mp3}{\FCtitle{\#521}{Jsou dnešní hry až příliš dlouhé?}}{2021-03-24}{\Bety, \PMakal, \Ulrik, \OverWatch}{\FCrule{\#521}{I s krátkým kašpárkem se dá udělat divadlo.}{\Bety}}{Navazujeme na tematický komentář od Pavla, který řeší, jsou hry dnes příliš dlouhé, nebo ne. Článek na webu vyvolal pořádnou diskuzi ostatně jak samotný, text tak i názory čtenářů si můžete přečíst zde: https://games.tiscali.cz/komentar/hry-jsou-nesmyslne-dlouhe-stoji-o-to-jen-hrstka-oddanych-507423 Co myslíte vy? Přidejte se k naší diskuzi o délce her!}{}{}{}
\FC{fight-club-522.mp3}{\FCtitle{\#522}{Patch pro Cyberpunk 2077, vylepšený Zaklínač a další novinky}}{2021-03-31}{\Bety, \JSvak, \OverWatch, \PMakal}{\FCrule{\#522}{Když se vám něco povede jednou neznamená, že se vám to povede i podruhé.}{\Bety}}{Téma o hrách, co nemáme rádi, jsme museli dnes z organizačních důvodů přesunout na příští týden (nebojte, že byste o to přišli). Místo toho tak probereme, co všechno si pro nás připravuje CD Projekt. Ať už nově vydaný patch pro Cyberpunk, tak novinky ohledně multiplayeru nebo  vylepšeného Zaklínače pro novou generaci konzolí.}{}{}{}
\FC{fight-club-523.mp3}{\FCtitle{\#523}{Proč nesnášíme milované hry?}}{2021-04-07}{\Bety, \VPechacek, \OverWatch, \PMakal}{\FCrule{\#523}{Žluté barety nepřátelům Dishonored.}{\VPechacek}}{Některé hry se staly ikonami a všichni je zbožňují. Občas ale moc nechápeme proč a co na nich vlastně všichni vidí. U kterých absolutně netušíme, proč je ostatní staví na piedestal?}{}{}{}
\FC{fight-club-524.mp3}{\FCtitle{\#524}{Nejlepší a nejhorší herní konce}}{2021-04-14}{\Bety, \VPechacek, \PMakal, \Luisa}{\FCrule{\#524}{Dobrej elf, mrtvej elf.}{\VPechacek}}{Všechno krásné jednou končí, ale ne vždycky úplně tak, jak bychom si představovali. Předem varujeme, že se v tomhle dílu z principu nebudeme vyhýbat všemožným spoilerům!}{}{}{}
\FC{fight-club-525.mp3}{\FCtitle{\#525}{Vysněné remaky a remastery}}{2021-04-21}{\VPechacek, \PMakal, \Ulrik, \Luisa}{\FCrule{\#525}{Nejsi hejtr, dokud o tom nenapíšeš knihu.}{\VPechacek}}{Pořád ještě existují hry, které se nedočkaly moderní předělávky a jsou skvělé, ale na dnešní poměry už ošklivé a zastaralé. Které klasiky by podle nás zasloužily předělat?}{}{}{}
\FC{fight-club-526.mp3}{\FCtitle{\#526}{Jak nám herní designéři kazí hraní?}}{2021-04-28}{\Bety, \OverWatch, \PMakal, \Luisa}{\FCrule{\#526}{Jestli nechceš, aby kolem tebe kroužila Šárka, musíš sakra přidat do kroku.}{\Bety}}{Mizerné minimapy, zdlouhavé úvody, quick time eventy... Čím nám herní designéři ztrpčují požitek z her a které mechanismy by měly zmizet v propadlišti dějin a už se nikdy nevrátit?}{}{}{}
\FC{fight-club-527.mp3}{\FCtitle{\#527}{Pohřešovaná pokračování oblíbených herních sérií}}{2021-05-05}{\VPechacek, \JSvak, \Ulrik, \OverWatch}{\FCrule{\#527}{Kenedy ty běhno, seberu ti stehno.}{\VPechacek}}{Stýská se nám po některých hrách - ve své době byly úspěšné, ale kdovíproč už nedostaly žádná další pokračování. Po kterých návratech herních sérií neúspěšně voláme?}{}{}{}
\FC{fight-club-528.mp3}{\FCtitle{\#528}{Naše největší herní úspěchy}}{2021-05-12}{\PMakal, \Ulrik, \JSvak, \Luisa}{\FCrule{\#528}{Pro PlayStation 5 si k nám do redakce nechoďte.}{\PMakal}}{Achievementy, trofeje, úspěchy. Sbíráme, nebo neřešíme? A na který dosažený milník jsme nejvíc pyšní? O tom všem si povíme v dalším díle podcastu Fight Club redakce Games.cz!}{}{}{}
\FC{fight-club-529.mp3}{\FCtitle{\#529}{Hráli jsme Mass Effect: Legendary Edition}}{2021-05-19}{\VPechacek, \PMakal, \Ulrik, \Luisa}{\FCrule{\#529}{I když se blbě ovládá Pavlo Mako milujeme.}{\VPechacek}}{Shepardi se vrátili a my se vrátili k Shepardům. Jaké jsou naše dojmy z remasterované trilogie Mass Effect? Je to pro nás pořád nejlepší herní sci-fi?}{}{}{}
\FC{fight-club-530.mp3}{\FCtitle{\#530}{Bulánci ovládli Startovač}}{2021-05-26}{\VPechacek, \PMakal, \Ulrik, \Luisa}{\FCrule{\#530}{Soul Calibr ma příběh kalibru Souls.}{\VPechacek}}{Čeští Bulánci si na přelomu tisíciletí podmanili školní učebny IVT a teď převálcovali všechny ostatní kampaně na Startovači. Jak na buclaté polštářky vzpomínáme my? A co nového je na poli velkolepých strategií od Paradoxu?}{}{}{}
\FC{fight-club-531.mp3}{\FCtitle{\#531}{Nejhorší a nejlepší začátky her}}{2021-06-02}{\VPechacek, \Bety, \Luisa, \PMakal}{\FCrule{\#531}{Začátek dobrý, všechno dobré. Dokud nepřijde indián.}{\VPechacek}}{Některé hry mají začátek a konec. Koncům už jsme se věnovali, ale které herní úvody jsou podle nás mistrně zvládnuté, ať už vypravěčsky, scenáristicky anebo nějak úplně jinak?}{}{}{}
\FC{fight-club-532.mp3}{\FCtitle{\#532}{Co čekáme od E3 2021?}}{2021-06-09}{\VPechacek, \Ulrik, \PMakal, \Luisa}{\FCrule{\#532}{Kolik bodů z deseti dostane letošní největší konference... Eee 3.}{\VPechacek}}{Každoroční svátek hráčů je skoro tady! Co si od letošní digitální E3 slibujeme, co čekáme, na co se těšíme a jak ji plánujeme pokrývat?}{}{}{}
\FC{fight-club-533.mp3}{\FCtitle{\#533}{Dojmy z E3 2021}}{2021-06-16}{\VPechacek, \PHajda, \Bety, \PMakal}{\FCrule{\#533}{Patrikovo těšení — vývojářova smrt.}{\VPechacek}}{}{}{}{}
\FC{fight-club-534.mp3}{\FCtitle{\#534}{Naše vzpomínky a zážitky z minulých E3}}{2021-06-23}{\PMakal, \Ulrik, \PHajda, \Luisa}{\FCrule{\#534}{Figurky? My chceme jedině peníze.}{\PMakal}}{Letos výstava E3 proběhla poprvé v historii jen digitálně - my ale pamatujeme i ročníky, kdy se herní průmysl sešel na výstavišti i fyzicky. Jak na výlety do USA v redakci vzpomínáme?}{}{}{}
\FC{fight-club-535.mp3}{\FCtitle{\#535}{Špatné i skvělé tutoriály}}{2021-06-30}{\PMakal, \Luisa, \Ulrik, \PHajda}{\FCrule{\#535}{Moby Dobby jsou mor dnešní doby.}{\PMakal}}{Herní tutoriály jsou nedílnou součástí her. Ve většině případů jsou ale neskutečně otravný prvek, který přeskočíme a pak se divíme, že netušíme, co dál dělat. Jaké hry mají ale skvělé tutoriály, a kde jsme naopak návod, jak na to, od srdce nesnášeli?}{}{}{}
\FC{fight-club-536.mp3}{\FCtitle{\#536}{O herních klientech, ale i o novém Switchi}}{2021-07-07}{\VPechacek, \Luisa, \Ulrik}{\FCrule{\#536}{D'Artagnan do pusy nepatří.}{\VPechacek}}{Dnes se vracíme ke klasickému Fight Clubu, kde nebudeme řešit jedno téma, ale popovídáme si hned o několika různých událostech. Ať už jde o představení nových Windows, OLED Switch a nebo třeba, co říkáme na různé herní klienty jako Steam nebo GOG.}{}{}{}
\FC{fight-club-537.mp3}{\FCtitle{\#537}{O novinkách kolem Zaklínače}}{2021-07-14}{\Bety, \JSvak, \VPechacek, \Ulrik}{\FCrule{\#537}{Jirka — chytrá liška s tmavou barvou bříška.}{\VPechacek}}{O víkendu se konal WitcherCon - událost, kde se spojil Netflix a CD Projekt a společně přinesly novinky jak o seriálovém Zaklínači, tak i tom herním. Probereme tak to nejzajímavější z akce.}{}{}{}
\FC{fight-club-538.mp3}{\FCtitle{\#538}{O tom, jak se žije herním novinářům}}{2021-07-21}{\VPechacek, \Ulrik, \JSvak, \PMakal}{\FCrule{\#538}{Nové popravčí nástroje objednávejte jedině od Valve.}{\VPechacek}}{Dnešní téma bude ještě dozvuk Steam Decku, popovídáme si o chystaném XDefiant, nové týmové střílečce podle Toma Clancyho, a taky rozebreme, jak se vlastně žije takovému běžnému českému hernímu novináři.}{}{}{}
\FC{fight-club-539.mp3}{\FCtitle{\#539}{O žalobě proti Blizzardu a remake Dead Space}}{2021-07-28}{\Bety, \PMakal, \Luisa}{\FCrule{\#539}{Buďte na lidi hodní a nebuďte lulani.}{\Bety}}{Dnešní téma}{}{}{}
\FC{fight-club-540.mp3}{\FCtitle{\#540}{O herní hudbě}}{2021-08-04}{\PMakal, \Bety, \JSvak, \PHajda}{\FCrule{\#540}{Na koncert Low Roar nechoďte ani s Karlem Drdou.}{\PMakal}}{Naše oblíbené herní soundtracky, melodie či písničky.}{}{}{}
\FC{fight-club-541.mp3}{\FCtitle{\#541}{O morálních volbách}}{2021-08-11}{\PMakal, \Ulrik, \Luisa, \JSvak}{\FCrule{\#541}{Pozor na to jak přehýbáte svůj pytlík.}{\PMakal}}{Tentokrát probereme hry, u kterých jsme se pořádně zapotili, když jsme se museli rozhodnout. Které hry to zvládly dobře a které ne?}{}{}{}
\FC{fight-club-542.mp3}{\FCtitle{\#542}{O Diablu II a dalších tématech}}{2021-08-18}{\VPechacek, \Luisa, \PMakal, \JSvak, \Bety}{\FCrule{\#542}{Když vidíš holuba, zabij ho.}{\VPechacek}}{Jako každý týden se i tento ve středu sejdeme ve studiu a probereme aktuální palčivá témata z herního světa. Tentokrát se zaměříme na naše dojmy z bety Diablo II: Resurrected, ale mrkneme i na patch pro Cyberpunk a probereme další novinky týdne.}{}{}{}
\FC{fight-club-543.mp3}{\FCtitle{\#543}{O Gamescomu 2021 a dalších tématech}}{2021-08-25}{\VPechacek, \Luisa, \Ulrik, \PMakal}{\FCrule{\#543}{Šárko a Aleši, jediná výtka, na naší kameře nemáte lýtka.}{\VPechacek}}{Tenhle týden se zaměříme na Gamescom 2021, ale i na to, co jsme v poslední době hráli a další volná témata.}{}{}{}
\FC{fight-club-544.mp3}{\FCtitle{\#544}{V čem jsou nová CRPG lepší než stará dobrá klasika (a naopak)}}{2021-09-01}{\VPechacek, \Ulrik, \JSvak}{\FCrule{\#544}{Za kolik dní bude recenze Pathfindera? Hoďte si 8k10.}{\VPechacek}}{Naši redakční experti na RPG s D\&{}D systémy Vašek, Jirka a Aleš si ve Fight Clubu popovídají o tom, kam se žánr v průběhu let posunul, jaké nové vychytávky je baví a co jim v dnešních CRPG pro změnu chybí.}{}{}{}
\FC{fight-club-545.mp3}{\FCtitle{\#545}{O hrách, co nechceme hrát, ale jejich světy nás fascinují}}{2021-09-08}{\PMakal, \JSvak, \Bety, \Luisa}{\FCrule{\#545}{K některým příběhům se musíte dopracovat až na záchodě.}{\PMakal}}{Probereme téma našeho společného článku o herních světech, co se nám líbí, ale hry bychom nehráli. Krom toho budeme také diskutovat o jiných tématech}{}{}{}
\FC{fight-club-546.mp3}{\FCtitle{\#546}{O PlayStation Showcase, Deathloop a herním podzimu}}{2021-09-15}{\Ulrik, \Bety, \Luisa, \PMakal}{\FCrule{\#546}{Moops do her patří.}{\Bety}}{Probereme, co všechno se oznámilo na streamu Sony z minulého týdne a co si o tom myslíme, pokecáme o Deathloop a taky, jaké další hry nás čekají tenhle podzim.}{}{}{}
\FC{fight-club-547.mp3}{\FCtitle{\#547}{O herním podzimu}}{2021-09-22}{\Ulrik, \PHajda, \Luisa, \PMakal}{\FCrule{\#547}{Pokud chcete Fight Club 500, vyčistěte si dutiny.}{\PMakal}}{V minulém díle se na to nedostalo, ale tentokrát se skutečně podíváme na to, co nás čeká v nadcházejících týdnech a měsících.}{}{}{}
\FC{fight-club-548.mp3}{\FCtitle{\#548}{Opět o Diablu II, ale i dalších tématech}}{2021-09-29}{\Bety, \VPechacek, \Luisa, \JSvak, \PMakal}{\FCrule{\#548}{Než-li každodenní sebehana, lepši každodenní sebeláska.}{\Bety}}{Minulý týden vyšlo Diablo II: Resurrected, tenhle týden má pro nás Pavel recenzi, ale dojmy, které může probrat s Vaškem a Jirkou. Diablo ovšem nebude jediné téma našeho Fight Clubu :) | Zde slíbený Vaškův podcast, proč je Mass Effect 3 špatný: https://www.youtube.com/watch?v=fHVCshwwTe8}{}{}{}
\FC{fight-club-549.mp3}{\FCtitle{\#549}{O Far Cry 6 a celé sérii Far Cry}}{2021-10-06}{\Ulrik, \VPechacek, \JSvak, \PMakal}{\FCrule{\#549}{Kakal v redakci kakal, než ho pan Makal spapal.}{\VPechacek}}{Tenhle týden vychází Far Cry 6, a tak si nejenom popovídáme o novém dílu, ale i zavzpomínáme na starší díly. | Zde slíbený Vaškův podcast, proč je Mass Effect 3 špatný: https://www.youtube.com/watch?v=fHVCshwwTe8}{}{}{}
\FC{fight-club-550.mp3}{\FCtitle{\#550}{Vzpomínání na sérii GTA}}{2021-10-13}{\VPechacek, \Luisa, \PMakal}{\FCrule{\#550}{Tři protagonisté GTA 6: Klaus, Zeman a Havel.}{\VPechacek}}{Rockstar konečně potvrdil, že nás čeká vydání remasterované verze GTA III, GTA Vice City a  GTA San Andreas. Je tak na čase si na tyhle hry zavzpomínat.}{}{}{}
\FC{fight-club-551.mp3}{\FCtitle{\#551}{O Duně ve všech jejích podobách}}{2021-10-20}{\Ulrik, \VPechacek, \JSvak, \PMakal}{\FCrule{\#551}{Dřív Paul se Chani nabazí, než Jirka lidských podpaží.}{\VPechacek}}{Ve čtvrtek má premiéru filmová Duna, a tak si připomeneme jak filmy, tak samozřejmě i hry, ale i původní knížky.}{}{}{}
\FC{fight-club-552.mp3}{\FCtitle{\#552}{O herních adaptacích komiksů}}{2021-10-27}{\Ulrik, \Luisa, \PMakal}{\FCrule{\#552}{Mrtvoly jedině v podobě sídla pro skřítky.}{\Ulrik}}{Hry a komiksy jdou spolu ruku v ruce už nějaký pátek. Zavzpomínejte s námi na ty nejlepší a nebo i ty nejhorší.}{}{}{}
\FC{fight-club-553.mp3}{\FCtitle{\#553}{O herních hororech}}{2021-11-03}{\Ulrik, \Luisa, \PMakal}{\FCrule{\#553}{Opravdový horor je stěhování.}{\Ulrik}}{Halloween je za námi, ale to neznamená, že se pořád nemůžeme trochu postrašit!}{}{}{}
\FC{fight-club-554.mp3}{\FCtitle{\#554}{O Elden Ringu a soulsovkách}}{2021-11-10}{\Ulrik, \VPechacek, \JSvak, \PMakal}{\FCrule{\#554}{Pavel Makal kooperuje v zásadě sám.}{\VPechacek}}{Pavel si vyzkoušel Elden Ring na vlastní kůži, a tak si popovídáme, jak se mu to líbilo, probereme ale i celkově hry ze série Souls.}{}{}{}
\FC{fight-club-555.mp3}{\FCtitle{\#555}{O nepovedených startech slavných her}}{2021-11-24}{\Ulrik, \VPechacek, \JSvak, \PMakal}{\FCrule{\#555}{Tucking... meh, docking... eee, ale furt lepší než GTA Remaster Trilogy.}{\VPechacek}}{Ne vždycky se hrám start povede. Pojďme zavzpomínat na ty nejhorší launche v historii, ale popovídat si o nedávném fiasku s GTA Trilogy Definitive Edition a nebo Battlefield 2042.}{}{}{}
\FC{fight-club-556.mp3}{\FCtitle{\#556}{O TV seriálech pro fandy fantasy a sci-fi}}{2021-12-01}{\Ulrik, \Bety, \VPechacek, \PMakal}{\FCrule{\#556}{Jsi matka? Zpracuje Tě Betky mamky taťka.}{\VPechacek}}{Na podzim se urojilo hned několik zajímavých televizních seriálů. Ať už skvěle zpracovaný Arcane nebo třeba Kolo času či Nadace. Tentokrát tak ve Fight Clubu nebudeme probírat hry, jako spíš na co se koukat za dlouhých dní.}{}{}{}
\FC{fight-club-557.mp3}{\FCtitle{\#557}{NFT ve hrách a výhledy do roku 2022}}{2021-12-08}{\VPechacek, \Luisa, \JSvak}{\FCrule{\#557}{K čemu server, když máte autobus.}{\VPechacek}}{Ubisoft představil svou novou platformu Quartz, a tak se pobavíme o tom, jaká je podle nás budoucnost NFT ve hrách, ale i o tom, co nás v herním světě čeká v prvních měsících 2022.}{}{}{}
\FC{fight-club-558.mp3}{\FCtitle{\#558}{The Game Awards}}{2021-12-15}{\Ulrik, \Bety, \VPechacek, \PMakal}{\FCrule{\#558}{Nesouložte s koňmi ani sestrami. Radši topte nemluvňata v hnoji.}{\VPechacek}}{Jak probíhalo udílení The Game Awards? S čím jsme souhlasili, co nás překvapilo? Jaké hry nás čekají?}{}{}{}
\FC{fight-club-559.mp3}{\FCtitle{\#559}{Vánoční: Šťastné a veselé feat. Vaječňák}}{2021-12-22}{\Ulrik, \VPechacek, \Luisa, \JSvak, \PMakal}{\FCrule{\#559}{Krásné sobotní noc svátky bez Jirky a Vaška a Honzi a kdo? Proč? Co? Zač? Díky.}{\VPechacek{} + \Luisa{} + \JSvak{} + \Ulrik{} + \PMakal}}{V redakci zavládla sváteční atmosféra a tak poslední Fight Club tohoto roku pojmeme velmi oddechově. Takže probereme hlavně to, co budeme hrát o Vánocích a na co se těšíme v příštím roce.}{}{}{}
\FC{fight-club-560.mp3}{\FCtitle{\#560}{Novoroční: Na co se těšíme v letošním roce}}{2022-01-05}{\Ulrik, \VPechacek, \PHajda, \JSvak}{\FCrule{\#560}{Nekřičte smrt fašistickým psům, nebo se Patrik uchlastá.}{\VPechacek}}{Máme nový rok plný dvojek, ale také spousty her, co nás čekají. Na co se těšíme nejvíc?}{}{}{}
\FC{fight-club-561.mp3}{\FCtitle{\#561}{Hry roku 2021}}{2022-01-12}{\Ulrik, \VPechacek, \Luisa, \PMakal}{\FCrule{\#561}{Obrýlený Pavel, rivalova smrt.}{\VPechacek}}{Na jednu stranu se teď těšíme, co nového přinese rok 2022, na tu druhou je na místě zrekapitulovat, co nejlepšího se událo v tom předcházejícím. A proto, jako každý rok, vyhlašujeme Best of 2021! Ve Fight Clubu se tak pobavíme o tom, co je pro nás hrou roku 2021 a proč.}{}{}{}
\FC{fight-club-562.mp3}{\FCtitle{\#562}{Microsoft koupil Activision Blizzard}}{2022-01-19}{\Ulrik, \Luisa, \JSvak, \PMakal}{\FCrule{\#562}{Mmmm...}{\JSvak}}{V úterý Microsoft nečekaně oznámil, že koupil vydavatele Activision Blizzard za závratných 70 miliard dolarů. Tahle nečekaná akvizice rozhodně rozhýbala vody herního světa. Co z toho plyne? O tom všem si budeme povídat v novém Fight Clubu}{}{}{}
\FC{fight-club-563.mp3}{\FCtitle{\#563}{Jsou dnešní hry až moc dlouhé?}}{2022-01-26}{\Bety, \VPechacek, \JSvak, \PMakal}{\FCrule{\#563}{Co se říká, když potkáte Patrika? Hi, Daddy!}{\VPechacek}}{Jsou dnešní hry opravdu až tak dlouhé, nebo se nám to jenom zdá? A kdo z toho má radost...}{}{}{}
\FC{fight-club-564.mp3}{\FCtitle{\#564}{Naše dojmy z Dying Light 2?}}{2022-02-02}{\Bety, \VPechacek, \PMakal}{\FCrule{\#564}{Chcete být jako Pavel, přidejte se k fašistické paramilitaristické organizaci.}{\VPechacek}}{Jak moc se těšíme na Dying Light 2? Jaká hra bude?}{}{}{}
\FC{fight-club-565.mp3}{\FCtitle{\#565}{Vzpomínáme na Horizon Zero Dawn}}{2022-02-09}{\VPechacek, \Luisa,\JSvak, \PMakal}{\FCrule{\#565}{Trouba, vhodné místo pro kapustičky i Jiří Svák.}{\VPechacek}}{Co je třeba si všechno připomenout před chystaným vydáním Horizon Forbidden West? Vzpomínáme na první díl.}{}{}{}
\FC{fight-club-566.mp3}{\FCtitle{\#566}{Jak se hraje Horizon Forbidden West?}}{2022-02-16}{\OverWatch, \Luisa, \PMakal}{\FCrule{\#566}{Vynášení ortelů GTA 6 nechceme.}{\Luisa}}{Horizon Forbidden West za chvíli vychází, a tak nám Adam přišel popovídat o všem, co se nevešlo do recenze.}{}{}{}
\FC{fight-club-567.mp3}{\FCtitle{\#567}{S Pavlem o Elden Ring do detailu}}{2022-02-23}{\Ulrik, \PHajda, \PMakal}{\FCrule{\#567}{Pokud želvičku, tak jedině s Tiarou.}{\Ulrik}}{Elden Ring za chvíli vychází, a tak nám o hře Pavel do detailu povypráví. Recenzi si můžete přečíst zde: https://games.tiscali.cz/recenze/elden-ring-recenze-nejambicioznejsi-souls-hry-529377}{}{}{}
\FC{fight-club-568.mp3}{\FCtitle{\#568}{S Patrikem o Shadow Warrior 3 a Gran Turismo 7}}{2022-03-02}{\VPechacek, \PHajda, \JSvak, \PMakal}{\FCrule{\#568}{Tapety mají rady prvočísla.}{\VPechacek}}{Tenhle týden vychází hned dvě hry, Shadow Warrior 3 a Gran Turismo 7, tak se na obě pořádně podíváme.}{}{}{}
\FC{fight-club-569.mp3}{\FCtitle{\#569}{Co je nového v popkultuře}}{2022-03-09}{\Ulrik, \Bety, \Luisa, \PMakal}{\FCrule{\#569}{Nepijte v kině a nebo alespoň nesmrďte.}{\Bety}}{Místo herního tématu si zase dáme rekapitulaci toho, co jsme v poslední době viděli a nebo četli. Třeba se dozvíte pár dobrých tipů, co si pustit a čeho se naopak vyvarovat.}{}{}{}
\FC{fight-club-570.mp3}{\FCtitle{\#570}{O protiválečných hrách a dalších tématech}}{2022-03-16}{\Ulrik, \VPechacek, \JSvak, \PMakal}{\FCrule{\#570}{Vyfič Brumbič. Čumim na Fampič.}{\VPechacek}}{Tentokrát si budeme povídat o hrách, které kritizují válku. Krom toho taky probereme, jak je to s obtížností v Elden Ring, ale i co čekáme od chystaného streamu Hogwarts Legacy}{}{}{}
\FC{fight-club-571.mp3}{\FCtitle{\#571}{Vše, co víme o novém Zaklínači: engine, možný příběh a další...}}{2022-03-23}{\Bety, \VPechacek, \PMakal}{\FCrule{\#571}{Moje dcera se vyspala s mojí matkou a teď zabíjím monstra.}{\VPechacek}}{O Ghostwire: Tokyo si nakonec nepopovídáme, místo toho tak bude řeč o potvrzeném pokračování ze světa Zaklínače. Co si pro nás CD Projekt nachystal.}{}{}{}
\FC{fight-club-572.mp3}{\FCtitle{\#572}{PlayStation Plus v novém}}{2022-03-30}{\Ulrik, \VPechacek, \JSvak, \PMakal}{\FCrule{\#572}{Welcome to Jomo Dojo.}{\VPechacek}}{Sony představilo novou podobu svého předplatného PlayStation Plus. Podařilo se mu vytvořit kompetentní konkurenci pro Xbox Game Pass? Povíme si v dnešním Fight Clubu.}{}{}{}
\FC{fight-club-573.mp3}{\FCtitle{\#573}{O české hře Matcho s Filipem „Fiolou“ Kraucherem}}{2022-04-06}{\Ulrik, \VPechacek, \PMakal, \Fiola}{\FCrule{\#573}{Alešvi prstíky: zdroj mrtvolného ticha.}{\VPechacek}}{České studio FiolaSoft konečně představilo svou spoj3 střílečku Matcho, a tak jsme si do studia pozvali Fiolu, aby nám o ní něco víc povyprávěl.}{}{}{}
\FC{fight-club-574.mp3}{\FCtitle{\#574}{Nejdivnější herní periferie}}{2022-04-13}{\Ulrik, \PHajda, \Luisa, \PMakal}{\FCrule{\#574}{Pokud nový počítač, pak jedině s kličkou.}{\Ulrik}}{Nepraktické chlupaté gamepady, balanční podložky, bubínky nebo posilovací kruhy. Netradičních způsobů, jak ovládat hry, vzniklo nepřeberné množství, a tak si tentokrát ve Fight Clubu probereme ty nejpodivnější herní periferie, s nimiž jsme se za své hráčské kariéry setkali}{}{}{}
\FC{fight-club-575.mp3}{\FCtitle{\#575}{O našich láskách ve videohrách}}{2022-04-20}{\Ulrik, \Bety, \VPechacek, \Luisa}{\FCrule{\#575}{Rozdíl mezi břízou a Morigan? Tu druhou beztrestně nepočůráš.}{\VPechacek}}{Nejsou skuteční a stejně nám občas nedají spát. Kdo je naše největší virtuální láska?}{}{}{}
\FC{fight-club-576.mp3}{\FCtitle{\#576}{O Poldovi 7 a pirátství s Petrem Svobodou}}{2022-04-27}{\Bety, \VPechacek, \PSvoboda}{\FCrule{\#576}{Jack Sparrow patří na šibenici.}{\VPechacek}}{Médii proběhla informace o tom, jak se daří Poldovi 7 na pirátské scéně. O nepříjemné situaci si popovídáme s producentem hry Petrem Svobodou.}{}{}{}
\FC{fight-club-577.mp3}{\FCtitle{\#577}{O Hvězdných válkách na všechny způsoby}}{2022-05-04}{\Ulrik, \VPechacek, \JSvak, \PMakal}{\FCrule{\#577}{Kolik vlastníš Fallen Orderů, tolikrát jsi Svákem.}{\VPechacek}}{Je 4. května a s tím svátek tzv. Star Wars Day. Proto v našem Fight Clubu zavzpomínáme na staré hry ze světa Hvědných válek, ale také si popovídáme o nových projektech, co se právě teď chystají. Řeč bude samozřejmě i o filmech a seriálech. Zkrátka Star Wars Day jak se patří.}{}{}{}
\FC{fight-club-578.mp3}{\FCtitle{\#578}{EA končí s FIFOU}}{2022-05-11}{\Ulrik, \VPechacek, \PHajda, \Luisa}{\FCrule{\#578}{Nenaolejuješ-li Patrika, naolejuje Patrika Aleš.}{\VPechacek}}{Dneska si budeme povídat o budoucnosti her FIFA, ale i o dalších značkách, co pomalu umírají...}{}{}{}
\FC{fight-club-579.mp3}{\FCtitle{\#579}{Jaké herní série bychom rádi viděli ve středověku}}{2022-05-18}{\Ulrik, \Bety, \VPechacek, \Luisa}{\FCrule{\#579}{Královna kamene nezná slitování.}{\VPechacek}}{Změna plánu, hry podle filmů si necháme na příště, teď si dáme téma inspirované drby, že Bloober připravoval hru Aliens zasazenou do středověku. Jaké jiné hry bychom rádi viděli ve středověku?}{}{}{}
\FC{fight-club-580.mp3}{\FCtitle{\#580}{Vzpomínky na herní žně}}{2022-05-25}{\Ulrik, \Bety, \PHajda, \PMakal}{\FCrule{\#580}{Fight Club za 88 let? Jedině s opičími neštovicemi.}{\Bety}}{Wolfenstein: The New Order slaví tenhle týden výročí a spolu s tím jsme se rozhodli zavzpomínat na krásný rok 2014, který přinesl množství skvělých her.}{}{}{}
\FC{fight-club-581.mp3}{\FCtitle{\#581}{O hraní na mobilech}}{2022-06-01}{\Bety, \VPechacek, \JSvak, \PMakal}{\FCrule{\#581}{Skin do Fortnite? Staňte se raději rytířem.}{\VPechacek}}{Diablo Immortal je za rohem, tak je na čase si pokecat o tom, jak jsme na tom s mobilním hraním. Kdo je z redakce největší pařan a proč Pavel utratil za Candy Crush nehorázné množství peněz...}{}{}{}
\FC{fight-club-582.mp3}{\FCtitle{\#582}{Co čekáme od letošní "E3"?}}{2022-06-08}{\Ulrik, \Bety, \PHajda, \PMakal}{\FCrule{\#582}{Trpíte na modré koule? Můžou za to herní studia.}{\Bety}}{E3 je zrušená, ale to neznamená, že se nechystá spousty herních streamů, které nám snad představí nové hry, a nebo odhalí dosud neviděné záběry z titulů chystaných. Co bychom tedy na tomhle festivalu her rádi viděli?}{}{}{}
\FC{fight-club-583.mp3}{\FCtitle{\#583}{Naše dojmy z "E3 2022"}}{2022-06-15}{\Ulrik, \Bety, \Luisa, \PMakal}{\FCrule{\#583}{Starfield zjevně nemá tep na prsu doby.}{\Bety}}{Hlavní prezentace z letošního festivalu her máme za sebou. Tak jaké jsou naše dojmy?}{}{}{}
\FC{fight-club-584.mp3}{\FCtitle{\#584}{O demech z akce Steam Next Fest}}{2022-06-22}{\Bety, \VPechacek, \PHajda, \PMakal}{\FCrule{\#584}{Bereme koruny, eura, orgány, sperma. Hlavně darujte.}{\VPechacek}}{Vyzkoušeli jsme si dema k indie hrám v rámci akce Steam Next Fest. Co jsme si zahráli a na co se těšíme? To všechno si povíme v dnešním Fight Clubu.}{}{}{}
\FC{fight-club-585.mp3}{\FCtitle{\#585}{O smutných příbězích herního světa}}{2022-06-29}{\Bety, \PHajda, \Luisa, \PMakal}{\FCrule{\#585}{Pokud se chcete zabít, zavolejte linku bezpečí 11 61 11.}{\Bety}}{Spec Ops The Line slaví 10 let od svého vydání. Vydání se ale nesetkalo s úspěchem, hra komerčně propadla a studio skončilo. Jaké další smutné osudy herní svět zažil?}{}{}{}
\FC{fight-club-586.mp3}{\FCtitle{\#586}{O tom, kdo je nejhorší herní postava}}{2022-07-13}{}{\FCrule{\#586}{Když ženy, tak jedině v plátové zbroji.}{\Bety}}{Popovídáme si o tom, koho bychom nejraději ve hře vůbec neměli. A nebo měli, ale aspoň s páskou přes pusu.}{}{}{}
\FC{fight-club-587.mp3}{\FCtitle{\#587}{O hrách, co nás uklidňují}}{2022-07-20}{\Bety, \Ulrik, \JSvak, \PHajda}{\FCrule{\#587}{Kanci a hroši... s těma si nezahrávejte.}{\Bety}}{Jaké hry nás dostávají do zenu?}{}{}{}
\FC{fight-club-588.mp3}{\FCtitle{\#588}{O světech, kam bychom rádi na dovolenou}}{2022-07-27}{\Ulrik, \PHajda, \Luisa, \PMakal}{\FCrule{\#588}{Když na párty, tak transportérem. Když bradavky, tak s jazykem.}{\Ulrik}}{Jaký herní svět nejvíc láká k tomu, abychom tam strávili dovolenou?}{}{}{}
\FC{fight-club-589.mp3}{\FCtitle{\#589}{CD Projekt slaví 20 let}}{2022-08-03}{\Ulrik, \VPechacek, \Luisa, \PMakal}{\FCrule{\#589}{Nejlepší zákusek, Pavlův banán v Alešově čokoládě.}{\VPechacek}}{Tvůrci série Zaklínač, karetního Gwentu nebo RPG Cyberpunk 2077 slaví 20 let od založení studia. Jak na jejich tvorbu vzpomínáme my?}{}{}{}
\FC{fight-club-590.mp3}{\FCtitle{\#590}{O špatných misích v jinak dobrých hrách}}{2022-08-10}{\Ulrik, \Bety, \PMakal}{\FCrule{\#590}{Vývojáři, neberte nám naše hračky.}{\Bety}}{Zavzpomínáme na všechny ty mizerné mise, co nám zkazily jinak dobrý zážitek ze hry.}{}{}{}
\FC{fight-club-591.mp3}{\FCtitle{\#591}{O historkách z Gamescomu}}{2022-08-17}{\VPechacek, \Luisa, \PMakal}{\FCrule{\#591}{Za zvratky na Rýně v německém Klíně Pavel je na vině a Šárka na víně.}{\VPechacek}}{Naše vzpomínky z Gamescomu, dobré i špatné.}{}{}{}
\FC{fight-club-592.mp3}{\FCtitle{\#592}{Co jsme viděli a vyzkoušeli na Gamescomu 2022}}{2022-08-26}{\Bety, \Luisa, \PMakal}{\FCrule{\#592}{Nešlapejte na německý včely.}{\Bety}}{Gamescom 2022: Pavel a Šárka přijeli plni dojmů. Jaké hry si zahráli, které prezentace viděli? Co jim prozradili vývojáři... To všechno v našem Fight Clubu}{}{}{}
\FC{fight-club-593.mp3}{\FCtitle{\#593}{Jsme staří}}{2022-08-31}{\Luisa, \JSvak, \PMakal}{\FCrule{\#593}{Je nutno dát na vědomí, ze Morrowind je jedna z her všech dob.}{\JSvak}}{Letos slaví řada slavných sérií velká výročí. Pojďme se podívat, kdo všechno si letos sfouknul pomyslnou svíčku na oslaveneckém dortu. | (00:00) loading | (01:08) úvod | (03:03) co právě hrajeme | (08:57) co všechno letos slaví výročí | (10:49) Diablo na PC slaví 25 let | (15:50) Final Fantasy VII slaví 25 let | (17:04) Contra slaví 30 let | (18:47) Neir Automata slaví 5 let | (19:29) Hotline Miami slaví 10 let | (20:57) Fatal Frame slaví 20 let | (22:48) Turok slaví 20 let | (24:39) Street Fight II, Tekken a Mortal Kombat | (25:46) Morrowind slaví 20 let | (32:52) Star Wars Jedi Knight slaví 20 let | (36:05) Wolfenstein 3D slaví 30 let | (41:27) Warcraft III: Reign of Chaos slaví 20 let | (44:52) Metal Gear série slaví 35 let | (49:54) Bioshock slaví 15 let | (58:24) Borderlands 2 slaví 10 let | (1:03:17) Hitman 2 slaví 20 let | (1:05:57) Fallout slaví 25 let | (1:07:31) Portal slaví 15 let | (1:12:11) Pong slaví 50 let | (1:12:26) Kane \&{} Lynch: Dead Men slaví 15 let | (1:14:02) GTA slaví 25 let | (1:16:47) Dune II slaví 30 let | (1:19:49) závěr}{}{}{}
\FC{fight-club-594.mp3}{\FCtitle{\#594}{Hry \&{} Velká francouzská revoluce}}{2022-09-07}{\Ulrik, \VPechacek, \Luisa, \PMakal}{\FCrule{\#594}{Pistáciovou do prdele si strč.}{\VPechacek}}{Tento týden vychází Steelrising, kde Ludvík XVI. potlačil revoluci díky svým robotickým loutkám. Jaké další hry se vydávají do tohohle zajímavého období francouzské historie? Povíme si v tomto díle Fight Clubu.}{}{}{}
\FC{fight-club-595.mp3}{\FCtitle{\#595}{Budoucnost Assassin's Creed}}{2022-09-14}{\Ulrik, \PHajda, \Luisa, \PMakal}{\FCrule{\#595}{Když do postele, tak jedině s Ichibanem.}{\Ulrik}}{Ubisoft v sobotu představil ne jednu, ale rovnou čtyři hry ze série Assassin's Creed. Co říkáme na Mirage a na Codename: Jade, Red a Hexe?}{}{}{}
\FC{fight-club-596.mp3}{\FCtitle{\#596}{Herní únik století? Rockstar neuhlídal GTA VI}}{2022-09-21}{\Ulrik, \JSvak, \PHajda, \PMakal}{\FCrule{\#596}{Sváku vypni to.}{\Ulrik}}{Internet v neděli zachvátily obrázky a videa z nadcházejícího GTA VI. Proč je předčasné z uniklé verze hry cokoliv soudit a proč bude výsledek vypadat úplně jinak? A může tenhle únik Rockstar nějak poškodit dlouhodobě? Probíráme ve Fight Clubu \#596.}{}{}{}
\FC{fight-club-597.mp3}{\FCtitle{\#597}{Vracíme se k Cyberpunku 2077}}{2022-10-05}{\Bety, \VPechacek, \Luisa, \JSvak}{\FCrule{\#597}{Pokud nemáte peníze na koks, pořiďte si ušní infekci.}{\VPechacek}}{Cyberpunk 2077 se dva a půl roku po svém vydání vyšvihl na přední špičky v prodejích, prošel řadou vylepšení, jak si tedy vede teď a stojí za to se k němu vrátit?}{}{}{}
\FC{fight-club-598.mp3}{\FCtitle{\#598}{Série Fallout slaví 25 let}}{2022-10-12}{\Ulrik, \Bety, \Luisa, \PMakal}{\FCrule{\#598}{Když potřebujete povzbudit, dejte si kyselý želé od Hariba.}{\Bety}}{První Fallout vyšel před 25 lety. Tak si pojďme zavzpomínat na celou sérii!}{}{}{}
\FC{fight-club-599.mp3}{\FCtitle{\#599}{Co chceme od nového Silent Hillu}}{2022-10-19}{\Bety, \Luisa, \PMakal}{\FCrule{\#599}{Nový Silent Hill jedině s fousatým chlapíkem, co chodí s teenagerem.}{\Bety}}{Večer nás čeká odhalení nového Silent Hillu (který bude nejspíše remake). Co bychom rádi slyšeli? A chceme se vlastně do zkaženého světa tohohle hororu vracet?}{}{}{}
\FC{fight-club-600.mp3}{\FCtitle{\#500 + \#600}{Speciál z nového studia}}{2022-11-02}{\Ulrik, \OverWatch, \Bety, \VPechacek, \Luisa, \PMakal}{\FCrule{\#1100}{Mokvavý a plesnivý Jiří strašně chybí.}{\Bety}}{Nechceme nic zakřiknout, ale snad se nám konečně poštěstí mít s vámi Fight Club speciál, ve kterém oslavíme hned 500 a 600 díl z našeho nového studia. | (00:00) znělka | (00:26) - úvodní slovo, odhalení studia, odpovědi na otázky | (23:44) sestřih evoluce Fight Clubu | (25:22) zpátky ve studiu | (47:09) přeceňované, nebo podceňované? | (1:15:41) sestřih top momentů z Fight Clubu | (1:31:09) zpátky ve studiu | (1:36:03) soutěž Co na to Gamesáci | (2:02:23) Odpovědi na vtipné otázky, co vás donutí se zamyslet | (2:27:09) vzkazy od externistů | (2:29:08) znělka Fight Club ala seriál z 80. let | (2:29:47) vánoční přání za rok 2020 | (2:33:22) zpátky ve studiu | (2:34:00) Herní kvíz | (2:44:24) poděkování, rozloučení, pravidlo}{}{}{}
\FC{fight-club-601.mp3}{\FCtitle{\#601}{O novinkách, co právě hrajeme}}{2022-11-09}{\Ulrik, \Bety, \Luisa, \PMakal}{\FCrule{\#601}{Chces Fight Club metch, dostaneš prdy v lahvičce!}{\Bety}}{God of War Ragnarök ale také Call of Duty: Modern Warfare II, to všechno a víc bude tématem našeho Fight Clubu.}{}{}{}
\FC{fight-club-602.mp3}{\FCtitle{\#602}{Vzpomínáme na sérii Splinter Cell}}{2022-11-16}{\VPechacek, \Luisa, \PMakal}{\FCrule{\#602}{FOrtnite je sračka. Nejlepší battle royal vypukne v neděli.}{\VPechacek}}{Splinter slaví 17. listopadu výročí 20 let od vydání prvního dílu. Ideální čas na to zavzpomínat, co si o téhle sérii myslíme, zda nám chybí a co bychom rádi příapdně od nového dílu.}{}{}{}
\FC{fight-club-603.mp3}{\FCtitle{\#603}{Co říkáme na nominace TGA 2022}}{2022-11-23}{\Bety, \PHajda, \Luisa, \PMakal}{\FCrule{\#603}{Alana je láska s rozkošným pršáčkem a hlavně dobrým srdcem.}{\PMakal}}{Máme tu vyhlášení cen Golden Joystick 2022 a vedle toho The Game Awards, ve zkratce TGA, oznámily minulý týden nominace v řadě kategorií. Společně projdeme, co nás překvapilo, a kde se naopak ničemu nedivíme. | (00:00) znělka | (00:40) přivítání | (02:30) Co poslední dobu hrajeme? | (04:35) Bětce se ztratil prsten z Hobita | (06:38) God of War: Ragnarok | (26:07) South Park The Stick of Truth | (27:24) Co máme vyhlídnuto? | (30:34) Evil West | (32:53) Warhammer 40k: Darktide | (34:53) Pentiment | (36:44) Vampire Survivors | (38:27) Golden Joystick Awards Best Storytelling | (40:37) Still Playing Award | (41:08) Best Visual Design | (43:14) Dotaz fanouška: Nominace na hru roku? | (45:17) Studio of the Year | (46:11) Best Game Expansion | (46:51) Best Early Access Launch | (47:57) Best Indie Game | (49:06) Best Multiplayer Game | (50:16) Best Audio | (52:55) Best Game Trailer | (54:32) Best Game Community | (55:13) Best Gaming Hardware | (56:08) Breakthrough Awar, Critics' Choice Award, Best Performer | (58:07) Nintendo, PC, PlayStation, Xbox Game and Most Wanted Game of the Year | (1:00:30) Ultimate Game of the Year | (1:00:58) The Game Awards | (1:17:22) Rahul Kohli a Bětka | (1:18:48) Dotaz od fanouška: Co říkáme na Micka Gordona a Bethesdu ? | (1:20:10) Loučení a Pravidlo \#603}{}{}{}
\FC{fight-club-604.mp3}{\FCtitle{\#604}{Nejlepší a nejhorší herní puzzly}}{2022-11-30}{\Bety, \VPechacek, \PHajda, \PMakal}{\FCrule{\#604}{Zastavit zlou Lilith? Ne, raději knihu.}{\VPechacek}}{Jaké hádánky ve hrách máme rádi, co naopak nesnášíme? A samozřejmě přijde i řeč na to, jaké hry právě hrajeme. | (00:00) Znělka | (00:17) Přivítání | (01:03) Záhada ve studiu!? | (02:11) Marvel's Midnight Suns | (17:45) Dotaz od fanouška: Trailer Super Mario Bros | (19:34) Dotaz od fanouška: Bude doporučení deskovek? | (20:36) Pavel hraje Evil West | (28:09) Dotaz od fanouška: Bude recenze The Callisto Protocol ? | (28:37) Bětka hrála God of War a Pentiment | (30:58) Vašek hraje Victoria 3 | (34:48) Somerville a Vampire Survivors | (37:07) Dnešní téma puzzly | (37:53) Názor Pavla/Bětky na Puzzly | (41:34) Názor Patrika na puzzly | (42:19) Puzzly Evil west/ Marvel's Midnight Suns/ A Plague Tale: Requiem/ Half-Life 2 | (45:11) Puzzly, který nenávidí Pavel/Bětka | (49:05) Puzzly v Uncharted/Tomb Raider/ God of War | (55:05) Dotaz od fanouška: Jaký gear měl Pavel v God of War? | (56:50) Jaké puzzly jsou dobré pro Pavla? | (57:40) Silent hill Puzzly | (58:30) Dotaz od fanouška: Puzzly ve VR? | (59:43) Patrik doporučuje sérii We Were | (1:01:14) Vašek a puzzly | (1:04:47) Escape room ve Falloutu 76 | (1:06:38) The Incredible Machine/Laser puzzly | (1:09:11) Sherlock Holmes a puzzly | (1:10:14) Loučení a Pravidlo \#604}{}{}{}
\FC{fight-club-605.mp3}{\FCtitle{\#605}{Kontroverzní osobnosti herního průmyslu}}{2022-12-07}{\Bety, \VPechacek, \PMakal}{\FCrule{\#605}{Bětky se vždycky dotýkejte více prsty.}{\VPechacek}}{V herním průmyslu se najde nejeden člověk s pochroumanou pověstí. Popovídejte si s námi o těch nejhorších případech. | (00:00) Úvod | (01:44) Co poslední dobu hrajeme? | (02:29) The Callisto Protocol | (05:18) Marvel's Midnight Suns | (08:14) Need for Speed Unbound | (10:46) Victoria 3 | (12:55) Call of Duty: Modern Warfare II | (21:29) Dotaz od fanouška: Dívali jste se na tu "podivnou" Dead Island 2 prezentaci? | (22:23) Dotaz fanouška: Názor na hodnocení 10/10 Indianu na hru The Callisto Protocol? | (24:00) Hlavní téma Kontroverzní osobnosti herního průmyslu | (26:15) Tommy Tallarico | (32:06) Vaškův oblíbený kontroverzní vývojář | (33:40) Peter Molyneux | (36:41) Dotaz od fanouška: Nemyslíte si, že kontroverzní osobnost u herního vývojáře je záruka úspěchu hry? Např. Vávra, Kojima a Molyneux | (41:44) David Cage | (43:21) Randy Pitchford | (45:02) Marty Stratton | (52:39) Blizzard | (55:00) Chris Avellone | (58:17) Rozloučení a pravidlo \#605}{}{}{}
\FC{fight-club-606.mp3}{\FCtitle{\#606}{To dobré a špatné z The Game Awards 2022}}{2022-12-14}{\Bety, \PHajda, \Luisa, \PMakal}{\FCrule{\#606}{Microsofte, zakra doval exkluzivity!}{\Bety}}{Pokud chcete shrnutí toho nejdůležitějšího The Game Awards, či třeba pokec o tom, jaké nově představené hry se nám líbily nejvíce, nezapomeňte se podívat na dnešní Fight Club. Jako vždycky se také budeme věnovat hrám, které právě teď hrajeme. | (00:00) úvod | (01:57) Need For Speed Unbound | (10:00) Co hraje Bětka? | (14:25) Co hraje Šárka? | (15:35) Duck Simulator | (16:31) Co hraje Pavel? | (18:45) High on Life | (24:26) Téma The Game Awards 2022 | (24:50) Dotaz od fanouška: Crysis remaster | (25:26) Co nás šokovalo na The Game Awards 2022? | (31:55) Co se nám líbilo na The Game Awards 2022? | (32:44) Hellboy | (34:56) Death Stranding 2 | (36:00) Banishers: Ghosts of New Eden | (38:42) Judas | (41:10) Armored Core VI | (43:06) Immortals of Aveum | (46:04) Hades 2 | (48:14) Star Wars Jedi: Survivor | (49:30) Cyberpunk 2077: Phantom Liberty | (52:58) Diablo IV | (54:17) Dotaz od fanouška: Co říkáme na Crime Boss: Rockay City? | (56:12) Bayonetta | (57:49) Co potěšilo Pavla a na co se těší? | (59:30) Transformers: Reactivate | (1:00:33) Behemoth | (1:02:17) Space Marine 2 | (1:02:59) Co nás zamrzelo? | (1:09:51) Dotaz od fanouška: Jestli jsme zkoušeli Forspoken? | (1:13:00) Dotaz od fanouška: Jestli jsme zkoušeli The Witcher 3: Wild Hunt Next-Gen? | (1:13:28) Dotaz od fanouška: Kdy jdeme na Avatara 2? | (1:17:53) Rozloučení a pravidlo \#606}{}{}{}
\FC{fight-club-607.mp3}{\FCtitle{\#607}{Šťastné, herní a veselé!}}{2022-12-21}{\Ulrik, \VPechacek, \Luisa, \JSvak, \PMakal}{\FCrule{\#607}{Jiří míří pryč, tak fič.}{\VPechacek}}{Máme tu poslední Fight Club roku 2022, což znamená, že se budeme bavit o zásadních událostech uplynulého roku, očekávaných hrách toho příštího, ale taky o tom, jak hodláme trávit vánoční svátky a co očekáváme pod stromečkem. | (00:00) Úvod | (01:58) Kdo opouští Gamescz? | (02:25) Kdo přináší dárky Ježíšek, Santa nebo Satan? | (08:43) Jak slavíme Vánoce? | (11:56) Dotaz od fanouška: Jak vypadá naše kancelář? | (19:49) Co budeme dělat o Vánocích? | (28:01) Avatar 2 | (34:52) Na jaký filmy se těšíme? | (39:32) Co budeme dělat na Boží hod? | (43:57) Dotaz od fanouška: Jakou hru jsme si pořídili ve slevách a jaký hry budeme hrát? | (55:09) Hogwarts Legacy | (57:20) Co budeme dělat na Silvestra? | (1:03:59) Fun Fact od Vaška | (1:06:37) Jakou kolej si vybereme v Hogwarts Legacy? | (1:12:26) Hráli jsme Witcher 3 remaster? | (1:15:57) Dotaz od fanouška: Doporučení na nějakou party/rodinou hru na PS5 | (1:18:15) Dotaz od fanouška: Kingdom Come: Deliverance | (1:24:49) AI Chatbot | (1:31:25) Budeme zvát zajímavého hosty? | (1:36:17) Konec a pravidlo 607}{}{}{}
\FC{fight-club-608.mp3}{\FCtitle{\#608}{Naše výhledy na rok 2023}}{2023-01-04}{\Ulrik, \Bety, \Luisa, \PMakal}{\FCrule{\#608}{Ne každa J.K.Rawlingová je Rus, ale každý Rus je J.K.Rawlingová.}{\VPechacek}}{Jaké byly naše Vánoce a na co se těšíme v letošním roce? | (00:00) Úvod | (03:57) Jak jsme slavili Silvestr? | (13:26) Co jsme viděli? | (18:20) Jaké hry vyjdou v lednu? | (22:44) Jaké hry vyjdou v únoru? | (33:44) Jaké hry vyjdou v březnu? | (38:44) Jaké hry vyjdou v dubnu? | (39:24) Jaké hry vyjdou v květnu? | (43:46) Jaké hry vyjdou v červnu? | (44:53) Jaké hry vyjdou v červenci a v srpnu? | (47:38) Hry co nemají datum vydání | (1:16:42) Seriály a filmy podle her | (1:26:10) Závěr a pravidlo \#608}{}{}{}
\FC{fight-club-609.mp3}{\FCtitle{\#609}{Detektivové ve hrách}}{2023-01-11}{\Bety, \PHajda, \Luisa, \PMakal}{\FCrule{\#609}{Chcete udělat Patrikovi radost, dejte mu mrtvé děvky a bude spokojený.}{\Bety}}{Jaké herní detektivky máme nejraději? A jak se staví dobře vyšetřování případu do herního prostředí?}{}{}{}
\FC{fight-club-610.mp3}{\FCtitle{\#610}{Pád Ubisoftu}}{2023-01-18}{\OverWatch, \Bety, \Luisa, \PMakal}{\FCrule{\#610}{Jedna špatná hra, dobrý. Dvě špatný hry, horší než EA!}{\Bety}}{Co se stalo s Ubisoftem? Proč poslední dobou kvalita her upadá a jaká bude budoucnost? | (00:00) Úvod | (01:46) Novinky na GAMES.CZ | (04:20) Co hraje Pavel? | (07:38) Co hraje Šárka? | (09:00) Co hraje Adam | (10:00) Co hraje Bětka? | (12:31) Doba, kdy byl Ubisoft ještě dobrý | (22:24) Assassin's Creed | (26:27) Prince of Persia a Splinter Cell remake | (28:47) Proč velký tituly nestíhají? | (29:57) Problémy ve vývoji | (34:25) Nákup Tencentu | (37:10) Just Dance | (38:46) Dotaz od fanouška: Kdy bude ubisoft+ ? | (42:15) : Dotaz od fanouška Hell's Kitchen a kolik vydělávají hry jako Silent Hunter? | (44:19) Kdo bude následníkem Ubisoftu? | (53:31) Dotazy od fanoušků | (1:01:52) Závěr a pravidlo \#610}{}{}{}
\FC{fight-club-611.mp3}{\FCtitle{\#611}{O hraní her s Martinem Veselovským}}{2023-01-25}{\Ulrik, \PMakal, \MVeselovsky}{\FCrule{\#611}{Když nemám v ruce kvér a nestřílím, tak je něco špatně.}{\Ulrik}}{Dlouho jste si psali o hosta, tak jsme si k nám do studia pozvali redaktora a moderátora, Martina Veselovského, který mimo neústupné rozhovory hraje ve svém volném čase hry. Jaké? A na jaké platformě? A co ho bavilo dříve a co teď. To všechno se dozvíte v našem podcastu. | (00:00) Úvod a představení hosta | (02:23) Call of Duty: Mobile | (06:20) Investice do her | (08:51) Co hraje Aleš? | (10:15) Co hraje Pavel? | (12:54) Interakční hry Firewatch | (14:20) Hry od Amanita Design | (16:33) Žánr střílečky | (21:30) Jak sleduje herní svět pan Veselovský? | (24:47) Multitasking | (28:56) Herní filmy a seriály | (31:30) VR | (33:29) Jestli pan Veselovský hraje se svými dětmi? | (35:29) Diablo | (36:56) Jak tráví volný čas pan Veselovský? | (41:19) Deskovky a Sci-fi | (48:24) Politika ve hrách | (55:00) Debata o rozhovorech | (1:00:55) Objektivita u médií, který jsou veřejnoprávní | (1:05:13) Mafie a Kingdom Come: Deliverance | (1:06:18) Jestli Babiš přijde do DVTV? | (1:06:57) Nejhorší hosti? | (1:11:54) Newsroom | (1:12:30) Crowdfunding | (1:18:25) Paywall | (1:27:18) Ústup čtení | (1:32:40) Nejlepší host? | (1:35:25) Jestli pan Veselovský má vzor a jak se připravujete? | (1:13:49) Dívá se pan Veselovský na konkurenci? | (1:37:56) Arcane | (1:39:00) Závěr}{}{}{}
\FC{fight-club-612.mp3}{\FCtitle{\#612}{O herních generálech a prezidentech}}{2023-02-01}{\Ulrik, \VPechacek, \Luisa}{\FCrule{\#612}{Neštvěte Aleše nebo se vám rozřeže.}{\VPechacek}}{O víkendu skončila volba prezidenta, a tak není od věci si připomenout všechny naše oblíbené herní generály, generálky nebo prezidenty a prezidentky. | (00:00) Úvod | (09:20) Co hraje Vašek? | (12:00) Co hraje Šárka? | (13:00) Co hraje Aleš? | (15:29) Das Schwarze Auge | (19:10) : Dotaz od fanouška: Kolik práce má průměrný stálý člen redakce? | (19:50) Hogwarts Legacy | (23:12) Videoherní prezidenti | (34:33) Tropico a Suzerain | (39:30) Saints Row 4 | (41:11) Série Red Alert | (44:55) Command \&{} Conquer: | (51:08) Mass Effect | (54:57) Generál Shepherd | (56:11) Hitman | (58:21) Company of Heroes 3 | (1:00:22) Dotaz od fanouška: Jak jsme spokojení s konzolí Switch? | (1:06:00) Závěr a pravidlo \#612}{}{}{}
\FC{fight-club-613.mp3}{\FCtitle{\#613}{O Hogwarts Legacy a dalších hrách}}{2023-02-08}{\VPechacek, \PHajda, \Luisa, \PMakal}{\FCrule{\#613}{Hej to je můj čaj, Avadaa Kedavra.}{\VPechacek}}{V pátek vychází Hogwarts Legacy, venku už máme recenzi. Jak si ovšem nové dobrodružství ze světa čar a kouzel vede v porovnání s jinými hrami s Harry Potterem? | (00:00) Úvod | (01:22) Hogwarts Legacy | (48:20) Forspoken Dead Space | (51:00) Season: A Letter to the Future | (57:50) Závěr a pravidlo \#613}{}{}{}
\FC{fight-club-614.mp3}{\FCtitle{\#614}{Proč tady není Fight Club \#614?}}{2023-02-15}{\Ulrik, \PHajda, \Luisa}{\FCrule{\#614}{Špatná čeština je taki čeština.}{\Luisa}}{Potkaly nás technické potíže, a tak vysílání Fight Clubu \#614 musíme bohužel odsunout na zítra.}{}{}{}
\FC{fight-club-615.mp3}{\FCtitle{\#615}{O PlayStation VR 2 a nových hrách}}{2023-02-22}{\Ulrik, \Bety, \Luisa, \PMakal}{\FCrule{\#615}{Chcete oplzlé ruské roboty, nechoďte na Twitter, zahrajte si Atomic Heart.}{\Bety}}{Pavel otestovat pořádně nový PlayStationn VR 2, a tak nám přijde povyprávět, jaký nový headset pro virtuální realitu je, ale i jaké jsou hry, které na PSVR2 vyšly. | (00:00) Úvod | (04:07) Co hraje Aleš? | (04:42) Co hraje Šárka? | (05:23) Co hraje Bětka? | (11:22) PlayStation VR2 | (32:20) Ryū ga Gotoku Ishin! | (44:55) Dotaz od fanouška: Jestli v PlayStation VR2 je možnost ovládání pro leváky? | (46:38) Dotaz od fanouška: Kolik vyšlo Yakuz? | (48:43) Atomic Heart | (1:11:32) Dotazy z emailu | (1:15:23) Závěr a pravidlo \#615}{}{}{}
\FC{fight-club-616.mp3}{\FCtitle{\#616}{O nejlepších a nejhorších DLC}}{2023-03-01}{\Bety, \Luisa, \PHajda, \PMakal}{\FCrule{\#616}{Prosím neochutnávejte hry, pěkně je vylízejte.}{\Bety}}{Jaké datadisky nás nejvíce bavily? A které jsou naopak zbytečné? A u jakých her bychom si DLC přáli? | (00:00) Úvod | (01:20) Co hraje Patrik? | (09:57) Co hraje Pavel? | (18:29) Dotaz od fanouška: Jaký PC gamepad doporučujeme? | (21:10) Co hraje Bětka? | (30:00) Hlavní téma | (35:14) Co by mohlo být DLC v Elden Ringu? | (36:39) Oblíbené DLC podle Šárky | (41:41) Oblíbené DLC podle Pavla | (43:55) Oblíbené DLC podle Bětky | (47:44) Dotaz od fanouška: Používá se ještě pojem datadisk? | (50:20) The Last of Us Left Behind | (52:33) Jagged Alliance 2 Unfinished Business | (53:46) "První" DLC: Brnění pro koně v Oblivionu? | (54:55) Patrikův vztah k DLC | (57:35) Red Dead Redemption: Undead Nightmare | (58:42) Far Cry 3: Blood Dragon | (1:00:48) Vadí nám DLC?! | (1:12:55) Dotazy z emailu | (1:17:10) Závěr a pravidlo \#616}{}{}{}
\FC{fight-club-617.mp3}{\FCtitle{\#617}{O hrách s Čestmírem Strakatým}}{2023-03-08}{\Bety, \PMakal, \CStrakaty}{\FCrule{\#617}{Heroin vždycky s mírou, pak to jde.}{\Bety}}{Do Fight Clubu jsme si pozvali dalšího hosta, tentokrát je jím moderátor a novinář Čestmír Strakatý. Tématem budou nejen hry. | (00:00) Úvod | (06:20) Je Čestmír hráč? | (08:54) Zábavný kvíz | (15:31) World of Warcraft a MMORPG | (25:44) Do jaké hry Čestmír propadl? | (33:20) Hraje Čestmír se svojí ženou? | (36:18) Čestmír a seriály | (1:00:26) Dotaz od fanouška: Jestli někdo z Čestmírova kolektivu hraje hry? | (1:01:46) Na čem Čestmír hraje? | (1:11:52) Na co se Čestmír těší? | (1:13:58) Dotaz od fanouška: Plánuje Čestmír vlastní značku čaje? | (1:17:04) Jak Čestmír přistupuje k práci? | (1:25:05) Čte Čestmír recenze? | (1:28:22) Dotazy fanoušků | (1:41:54) Závěr a pravidlo  \#617}{}{}{}
\FC{fight-club-618.mp3}{\FCtitle{\#618}{Jak připravujeme Games.cz}}{2023-03-15}{\Ulrik, \VPechacek, \PHajda, \Luisa, \PMakal}{\FCrule{\#618}{Všechny nás spojují bolavá záda a dětská traumata.}{\PMakal}}{Jak se píší novinky, recenzují hry, editují články, ale i vede tým? O tom všem a mnohem dalším si pokecáme v dnešním Fight Clubu. | (00:00) Úvod | (01:42) Co hraje Pavel? | (05:25) Co hraje Patrik? | (08:00) Co hraje Aleš? | (09:28) Co hraje Vašek? | (11:34) Co hraje Šárka? | (12:27) Dotaz od fanouška: Jaký máme vztah k Ubisoftu? | (15:03) Jak nám začíná práce? | (18:31) Jaký máme role v redakci? | (20:06) Jak vypadá Alešův pracovní den? | (27:35) Novinky | (44:46) Dotaz od fanouška: Kdo tráví v práci nejvíce času? | (46:27) Porušili jsme někdy embargo? | (50:28) Recenzování | (1:20:44) Dotaz od fanouška: Hádáme se o recenze? | (01:25) ::17 Dotaz od fanouška: Vrátí se Hardwareclub? | (1:26:32) Dotaz od fanouška: Na kterou hru se teď v prvním půlroku 2023 nejvíce těšíme? | (1:30:26) Dotaz od fanouška: Neměli jste někdy problém, že jste byli zaměněni s nějakým recenzentem? | (1:32:02) Poděkování | (1:33:27) Dotaz od fanouška: Plánujete přejit místo číselného hodnoceni na šestihran? | (1:35:48) Dotaz od fanouška: Plánujeme v blízké budoucnosti další sraz s diváky? | (1:36:45) Závěr a pravidlo \#618}{}{}{}
\FC{fight-club-619.mp3}{\FCtitle{\#619}{O Redfall a upírech ve hrách obecně}}{2023-03-22}{\Bety, \VPechacek, \Luisa}{\FCrule{\#619}{Orgasmus je kamarádem smradu.}{\VPechacek}}{Šárka si byla v Polsku vyzkoušet Redfall, skvělá příležitost se podívat na téma upírů ve hrách. Pojďme si tak říci, kdo je náš oblíbený hemoglobinový someliér. | (00:00) Úvod | (01:58) Co hraje Šárka? | (02:34) Stmívání | (05:48) Co hraje Vašek? | (07:50) Co hraje Bětka? | (11:50) Upíři jsou LGBT literatura | (14:57) Eragon | (16:00) Hlavní téma: Upíři ve hrách | (18:32) The Witcher 3 Blood and Wine | (25:50) Vampyr | (30:13) Upíři ve světě Harryho Pottera | (32:34) Je Drákula dobrej záporák? | (34:04) Evil West | (35:26) Redfall | (49:50) Vampire: The Masquerade – Bloodhunt | (51:17) V Rising | (52:00) Vampire Survivors | (52:58) Co se Bětka dozvěděla/zajímavosti | (55:53) BloodRayne | (56:06) The Order: 1886 | (57:04) Závěr a pravidlo \#618}{}{}{}
\FC{fight-club-620.mp3}{\FCtitle{\#620}{O Tekkenu a dalších tématech}}{2023-03-29}{\Ulrik, \Bety, \PHajda, \PMakal}{\FCrule{\#620}{Pokud chcete fotku s Pavlem, tak odteď pouze v bojové pozici.}{\Bety}}{Pavel se byl v Lyonu podívat na nového Tekkena, tak vám povypráví o svých dojmech | (00:00) Úvod | (01:21) Co hraje Patrik? | (10:27) Co hraje Pavel? | (22:41) Co hrála Bětka | (28:11) Co hrál Aleš? | (34:00) Diablo IV | (39:40) Tekken 8 | (1:06:47) The Last of us na pc | (1:08:13) Oznámení a pravidlo \#620}{}{}{}
\FC{fight-club-621.mp3}{\FCtitle{\#621}{O smrtkách, které nelze zapomenout}}{2023-04-05}{\Ulrik, \PMakal, \JMalcharek}{\FCrule{\#621}{Kde Šárka není, ani smrt nebere.}{\PMakal}}{Šárka si zahrála Have a Nice Death, navíc se blíží Velikonoce, svátky smrti a vzkříšení. Ideální čas připomenout si naše oblíbené herní smrtky a smrťáky! | (00:00) Úvod | (05:33) Co hraje Kuba? | (08:16) The Last of Us na PC | (11:37) Crime Boss: Rockay City | (21:45) Co hrál Aleš? | (24:18) Sons Of The Forest | (32:43) Super Mario Bros | (41:39) Smrtky ve hrách | (1:08:04) Smrtky v kultuře | (1:17:24) Dotazy od fanoušků | (1:22:56) Závěr a pravidlo \#621}{}{}{}
\FC{fight-club-622.mp3}{\FCtitle{\#622}{O herních světech, kde je náročné přežít}}{2023-04-12}{\Ulrik, \VPechacek, \Luisa, \PMakal}{\FCrule{\#622}{Máte radši obrazy od Prickassa nebo Titiana?}{\VPechacek}}{Fanoušci herních survivalů si momentálně mohou užívat předběžný přístup do Sons of the Forest, ale světů, ve kterých je těžké přežít, známe stovky. A mnohé patří do žánrů, od kterých bychom to ani náhodou nečekali. Právě o přežívání ve virtuálních dimenzích si budeme tento týden povídat | (00:00) Úvod | (05:48) Sons of the Forest | (24:21) Have a Nice Death | (27:35) Sherlock Holmes The Awakened | (32:00) Like a Dragon: Ishin! | (33:12) Dotazy od fanoušků | (34:17) Chtěli bychom žít ve hře Animal Crossing? | (36:12) V jakých světech bychom nechtěli žít? | (1:10:35) Dotazy od fanoušků | (1:19:13) Závěr a pravidlo \#622}{}{}{}
\FC{fight-club-623.mp3}{\FCtitle{\#623}{O hrách, které prošly vývojovým peklem}}{2023-04-20}{\Ulrik, \VPechacek, \PMakal}{\FCrule{\#623}{Buďte jako Duke, když někdo vyvíjí hru patnáct let, hoďte po něm hovnem.}{\VPechacek}}{Při příležitosti vydání Dead Island 2 si budeme povídat o hrách, které měly složitou cestu na svět. U některých se čekání vyplatilo, u jiných zase ne!¨ | (00:00) Úvod | (06:27) Co hraje Aleš? | (10:44) Co hraje Pavel? | (17:29) Hlavní téma | (18:20) Dead Island | (22:55) Duke Nukem Forever | (35:50) L.A. Noire | (40:32) Aliens: Colonial Marines | (46:48) Diablo 3 | (51:08) Star Citizen | (55:45) Six Days in Fallujah | (58:22) Beyond Good and Evil 2 | (1:02:52) The Last Guardian | (1:04:43) StarCraft: Ghost | (1:07:44) Prey | (1:09:39) Dotazy | (1:18:53) Závěr a pravidlo \#623}{}{}{}
\FC{fight-club-624.mp3}{\FCtitle{\#624}{O hrách a Japonsku s Naomi Adachi}}{2023-04-26}{\Ulrik, \Luisa, \NAdachi}{\FCrule{\#624}{V Japonsku i silná noha může být poetická.}{\Ulrik}}{Ve studiu nás tento týden navštíví moderátorka, herečka a modelka Naomi Adachi, která si s Alešem a Šárkou popovídá o hrách a o japonské kultuře | (00:00) Úvod | (02:20) Jak se Naomi dostala ke hraní? | (05:09) První zkušenost s hraním? | (06:10) Jaká je Naomi hráčka? | (14:00) O jménu Adachi | (16:03) Naomi a popkultura | (36:29) Co by si Naomi ráda zahrála? | (43:50) Animefest a Comic-Con | (47:00) Naomi a hry | (1:03:18) Japonsko a hry | (1:09:20) Kdo je Naomi? | (1:10:18) Chtěla by Naomi vytvořit hru? | (1:12:45) Jaký žánr Naomi nemá ráda? | (1:15:10) Jaká hra baví Naomi? | (1:16:54) Open world hry nebo menší hry | (1:18:19) Vr a Naomi | (1:22:13) Závěr a pravidlo \#624}{}{}{}
\FC{fight-club-625.mp3}{\FCtitle{\#625}{O nezvládnutých vydáních her}}{2023-05-04}{\PHajda, \Luisa, \PMakal}{\FCrule{\#625}{To, že jste konzolista, není svobodná volba.}{\PHajda}}{V poslední době se zase roztrhl pytel s nepovedenými vydáními očekávaných her. Jaké jsou ty nejhorší launche herní historie a kolik her se z nich dokázalo vykoupit? | (00:00) Úvod | (02:09) Co hrála Šárka? | (04:06) Co hrál Patrik? | (12:05) GTA IV | (19:05) Sherlock Holmes The Awakened | (20:46) Hlavní téma | (22:49) Star Wars Jedi: Survivor a Redfall | (27:43) Battlefield 2042 | (31:08) No Man's Sky | (37:00) Cyberpunk 2077 | (41:44) Anthem | (45:02) Fallout 76 | (46:59) Batman Arkham Knight | (49:14) Mortal Kombat X | (50:25) Star Wars Battlefront II | (52:59) Balan Wonderworld | (55:54) Street Fighter V | (57:04) We Happy Few | (59:26) GTA: The Trilogy | (1:03:34) Diablo 3 | (1:05:49) Marvel's Avengers | (1:09:17) Artifact | (1:11:44) Assassin's Creed Unity | (1:15:04) BABYLON S FALL | (1:17:14) Kingdom Come: Deliverance | (1:19:25) Wild Hearts | (1:20:35) Dotazy | (1:26:03) Závěr a pravidlo \#625}{}{}{}
\FC{fight-club-626.mp3}{\FCtitle{\#626}{Kam kráčíš, Xboxe?}}{2023-05-10}{\Ulrik, \Luisa, \PMakal}{\FCrule{\#626}{Empatii ve hrách nechceme.}{\PMakal}}{Po mizerném launchi Redfallu poskytl Phil Spencer několik velmi upřímných rozhovorů. Jaká je dle jeho názoru budoucnost Xboxu a co si o tom myslíme my?}{}{}{}
\FC{fight-club-627.mp3}{\FCtitle{\#627}{O Zeldě a herních pískovištích}}{2023-05-18}{\Ulrik, \PHajda, \Luisa, \PMakal}{\FCrule{\#627}{Novou Zeldu, ne děkuji, raději voblejt lajf.}{\PHajda}}{The Legend of Zelda: Tears of the Kingdom momentálně vládne světu, a tak se o ní samozřejmě budeme bavit ve Fight Clubu. Zároveň se ale zamyslíme i nad sandboxovou hratelností v jiných titulech. | (00:00) Úvod | (03:44) Co hraje Patrik? | (06:24) Co hraje Aleš? | (14:40) Co hraje Pavel? | (22:14) The Legend of Zelda: Tears of the Kingdom | (53:31) Dotazy | (1:03:45) Závěr a pravidlo \#627}{}{}{}
\FC{fight-club-628.mp3}{\FCtitle{\#628}{Co přinesla PlayStation Showcase?}}{2023-05-26}{\OverWatch, \Luisa, \PMakal}{\FCrule{\#628}{Nejlepší herní děníček je plný nálepek a razítek kytiček.}{\Luisa}}{Společnost Sony se po dlouhém čekání konečně vytasila s PlayStation Showcase, akcí, na které ukázala plány své konzolové platformy pro příští měsíce a roky. Jak se nám akce líbila? | Timestampy k včerejšímu fc | (00:00) Úvod | (01:31) Co hraje Šárka? | (06:00) Co hraje Pavel? | (07:35) Jak zapůsobila PlayStation Showcase? | (16:18) Na co se těšíme? | (29:31) Assassin's Creed Mirage | (34:00) Spider-Man 2 | (39:21) Alan Wake 2 | (42:49) Microsoft | (49:01) Dotazy | (1:02:31) Handheld od Sony | (1:03:17) Závěr a pravidlo \#628}{}{}{}
\FC{fight-club-629.mp3}{\FCtitle{\#629}{Herní zážitky ze Středozemě}}{2023-05-31}{\Ulrik, \VPechacek, \PHajda, \PMakal}{\FCrule{\#629}{Koně rozsekat, pavouka postříkat.}{\VPechacek}}{Minulý týden vyšel The Lord of the Rings: Gollum a okamžitě se zapsal mezi nejhorší hry poslední doby. Na jaké herní návštěvy Středozemě vzpomínáme rádi a jaké za moc nestály? O tom si popovídáme v aktuálním Fight Clubu | (00:00) Úvod | (04:20) Jaká je naše oblíbená hra? | (24:35) Lego LOTR/ Hobbit | (28:30) Shadow of Mordor/ Shadow of War | (39:30) Gollum | (46:19) Co nás baví na LOTR? | (48:57) LOTR: Return to Moria | (50:10) Nový LOTR MMORPG | (54:30) Magic: The Gathering | (1:01:28) Co hraje Vašek? | (01:03) Co hraje Patrik? | (1:04:26) Co hraje Pavel? | (1:08:46) Dotazy | (1:13:24) Závěr a pravidlo \#629}{}{}{}
\FC{fight-club-630.mp3}{\FCtitle{\#630}{Diablo a hry z pekla}}{2023-06-08}{\Ulrik, \PHajda, \Luisa, \PMakal}{\FCrule{\#630}{Když v pekle zbraň, tak voodoo panenka. Když bradavky, tak s jazýčkem.}{\Ulrik}}{Tento týden si nemůžeme povídat o čemkoli jiném než o Diablu. Ale protože je tohohle akčního RPG teď všude plno, rozšíříme náš záběr i na další hry, ve kterých se podíváte do pekla | (00:00) Úvod | (04:12) Diablo IV | (41:01) Blood | (43:58) Clive Barker's Jericho | (46:00) Constantine | (47:12) Dante's Inferno | (54:28) Hades | (55:53) Metal Hellsinger a Hellblade | (58:49) Painkiller | (59:52) Saints Row Gat out of Hell | (1:01:52) Doom | (1:04:06) The Darkness | (1:06:13) Hellgate: London | (1:07:39) Agony | (1:09:03) Darksider | (1:10:42) Co hraje Patrik? | (1:13:14) Co hraje Šárka? | (1:19:51) Závěr a pravidlo \#630}{}{}{}
\FC{fight-club-631.mp3}{\FCtitle{\#631}{O našem nadšení z odhalených her}}{2023-06-14}{\VPechacek, \PHajda, \Luisa}{\FCrule{\#631}{Čím je váš svět větší, tím víc Gamesy brečí.}{\VPechacek}}{Letošní Ne-E3 byla dostatečně plodná na to, aby si na ní každý z nás našel to svoje. Jaké hry se nám líbily nejvíc? A stýská se nám po opravdové E3? | (00:00) Úvod | (04:20) Fable | (07:00) Star Wars Outlaws | (12:40) Assassin's Creed Mirage | (16:45) Avatar: Frontiers of Pandora | (19:00) Starfield | (27:07) Avowed | (29:47) Cyberpunk 2077: Phantom Liberty | (31:50) Microsoft Flight Simulator 2024 | (35:34) Hellblade II | (36:40) Nová hra od Daedalic | (37:54) Last Train Home | (42:21) Like a Dragon: Infinite Wealth a FINAL FANTASY VII REBIRTH | (44:40) Spider-Man 2 | (51:20) Prince of Persia The Lost Crown | (52:54) Halls Of Torment | (55:10) Diablo IV | (1:01:40) Dotazy | (1:07:20) Závěr a pravidlo \#631}{}{}{}
\FC{fight-club-632.mp3}{\FCtitle{\#632}{O nekonečné sérii Final Fantasy}}{2023-06-22}{\Ulrik, \Luisa, \PMakal}{\FCrule{\#632}{Final Fantasy XVI je dobrá i bez diverzity a inkluze.}{\PMakal}}{I kdyby měla série Final Fantasy jen těch 16 číslovaných dílů, patřila by mezi rekordmany. Jenže je k nim ještě zapotřebí připočíst celou řadu spinoffů a levobočků, které za šestatřicet let nasbírala. O většině z nich si budeme povídat v dnešním Fight Clubu. | (00:00) Úvod | (04:23) Hlavní téma Final Fantasy | (55:41) Co hraje Aleš? | (58:45) Co hraje Pavel? | (1:01:57) Historky z MHD | (1:03:45) Závěr a pravidlo \#632}{}{}{}
\FC{fight-club-633.mp3}{\FCtitle{\#633}{O nadcházející herní náloži}}{2023-06-28}{\Ulrik, \OverWatch, \PHajda, \Luisa}{\FCrule{\#633}{Podzim, heslo zní: řítí se na nás.}{\Ulrik}}{Pokud jsme předchozí roky brblali, že nemáme co hrát, letos podobný nedostatek rozhodně nehrozí. Od srpna do listopadu se na nás řítí taková spousta her, že se v nich doslova budeme topit. Co všechno letos na podzim vyjde a na co se z té nabídky nejvíc těšíme? | (00:00) Úvod | (02:27) Co vychází v srpnu? | (07:54) Co vychází v září? | (21:48) Co vychází v říjnu? | (41:37) Co vychází v listopadu? | (48:30) Co vychází v prosinci? | (50:46) Tři nejočekávanější hry pro Patrika? | (51:30) Tři nejočekávanější hry pro Aleše? | (51:57) Tři nejočekávanější hry pro Adama? | (52:17) Tři nejočekávanější hry pro Šárku? | (53:45) Dotazy | (58:20) Co hraje Patrik? | (1:00:25) Co hraje Aleš? | (1:02:32) Co hraje Adam? | (1:10:28) Závěr a pravidlo \#633}{}{}{}
\FC{fight-club-634.mp3}{\FCtitle{\#634}{O výsledku právních tahanic Microsoftu}}{2023-07-13}{\OverWatch, \VPechacek, \PHajda, \Luisa}{\FCrule{\#634}{Když někoho máte dost a cítíte k němu zlost, honem s ním na Černý Most.}{\VPechacek}}{Akvizice giganta Activision Blizzard se Microsoftu v průběhu času poněkud zkomplikovala, nyní ale slaví dílčí vítězství. V aktuálním Fight Clubu se budeme bavit o průběhu a hlavně důsledcích soudních jednání. | (00:00) Úvod | (03:13) Hlavní téma | (34:36) Co hraje Patrik? | (42:03) Co hraje Adam? | (45:05) Co hraje Šárka? | (46:25) Co hraje Vašek? | (51:55) Dotazy | (1:02:29) Závěr a pravidlo \#634}{}{}{}
\FC{fight-club-635.mp3}{\FCtitle{\#635}{O návratech herních legend}}{2023-07-21}{\Ulrik, \PHajda, \Luisa, \PMakal}{\FCrule{\#635}{Moly ano, moly né (Molyneux).}{\PHajda}}{Herní scéna je plná legendárních jmen, která se podílela na tvorbě slavných a velevážených značek a pokládala základy celých žánrů. Čas od času se tato jména s větším či menším úspěchem vrátí z odpočinku zpět do činné služby, a právě o tom si dnes budeme povídat. | (00:00) Úvod | (04:55) O návratech herních legend | (56:05) Co hraje Aleš? | (56:32) Co hraje Pavel? | (59:14) Dotazy | (1:11:25) Závěr a pravidlo \#635}{}{}{}
\FC{fight-club-636.mp3}{\FCtitle{\#636}{O erotice ve hrách}}{2023-07-26}{\VPechacek, \Luisa, \PMakal, \JMalcharek}{\FCrule{\#636}{Když je vaše holka v tom, čas na jebačku jak hrom.}{\VPechacek}}{Erotika je ve hrách přítomná už od od samých počátků domácí videoherní zábavy. My se v aktuálním Fight Clubu pobavíme o tom, kterým autorům se ji podařilo implementovat dobře a kde nám z ní naopak bylo trapně. | (00:00) Úvod | (04:20) Erotika ve hrách | (55:03) Dotazy | (1:04:47) Co hraje Pavel? | (1:06:08) Co hraje Šárka? | (1:07:29) Co hraje Jakub? | (1:10:19) Závěr a pravidlo \#636}{}{}{}
\FC{fight-club-637.mp3}{\FCtitle{\#637}{Jak zvládáme naše očekávání?}}{2023-08-03}{\Ulrik, \PHajda, \Luisa, \PMakal}{\FCrule{\#637}{A teď všichni: Wobbly Life, Wobbly Life, Wobbly Life, pojďte, Wobbly Life, Wobbly Life, Wobbly Life,...}{\PHajda}}{Léto je v polovině, naditá podzimní sezóna se blíží, a tím pádem se také pomalu rozjíždí pověstný hype train, tedy vlak, který místo uhlí jezdí na marketing a naleštěné trailery. Jak zvládáme krotit svá očekávání a profesionálně upozaďovat svá fanouškovství? | (00:00) Úvod | (03:00) Jak zvládáme naše očekávání? | (52:56) Co hraje Pavel? | (54:41) Co hraje Šárka? | (55:13) Co hraje Aleš? | (57:01) Co hraje Patrik? | (59:40) Dotazy | (1:08:11) Závěr a pravidlo \#637}{}{}{}
\FC{fight-club-638.mp3}{\FCtitle{\#638}{Všichni hrají Baldur's Gate 3}}{2023-08-09}{\VPechacek, \PMakal, \TTousek}{\FCrule{\#638}{Když jsou spolu do pokoje hobgoblin a zlobryně, klepejte.}{\VPechacek}}{V polovině léta, které by mělo být tradiční okurkovou sezónou, sestoupil mezi hráče RPG bůh a s nemilosrdnou razancí se začal ucházet o titul hry roku. Jaké jsou naše prozatímní zážitky z Baldur's Gate 3?}{}{}{}
\FC{fight-club-639.mp3}{\FCtitle{\#639}{Potkáme se u Kolína}}{2023-08-17}{\VPechacek, \PHajda, \Luisa, \PMakal}{\FCrule{\#639}{Na gramofonu v Patrikově Mustangu bude hrát techno.}{\VPechacek}}{Gamescom začíná už příští týden, takže je nejvyšší čas seznámit vás s našimi plány a také si popovídat o tom, co od téhle obří herní akce čekáme}{}{}{}
\FC{fight-club-640.mp3}{\FCtitle{\#640}{Naše zážitky z Gamescomu}}{2023-08-24}{\VPechacek, \PHajda, \Luisa, \PMakal}{\FCrule{\#640}{Když prodáváš otroky, nezapomeň na klec.}{\VPechacek}}{Patrik, Pavel, Šárka a Vašek se vydali do Kolína nad Rýnem, kde probíhá každoroční obří herní výstava Gamescom, a přímo z místa pro vás připravili podcast. Co v prvních dvou dnech zažili, co se jim líbilo nejvíc a na co byste se měli začít těšit?}{}{}{}
\FC{fight-club-641.mp3}{\FCtitle{\#641}{O návratech z cest i cestách samotných}}{2023-08-30}{\Ulrik \VPechacek, \PHajda, \PMakal}{\FCrule{\#641}{Nepláču, protože to končí, usmívám se, protože se to stalo.}{\VPechacek}}{Naše čtyřčlenná výprava se vrátila z Gamescomu, vedle toho má ale ještě jednu důležitou novinku. Vašek Pecháček odchází z Games.cz, a tak se sním pojďme rozloučit.}{}{}{}
\FC{fight-club-642.mp3}{\FCtitle{\#642}{O Starfieldu bez rukaviček}}{2023-09-06}{\PHajda, \Luisa, \PMakal, \JMalcharek}{\FCrule{\#642}{Přes Bethesdu ke hvězdám.}{\JMalcharek}}{Podzimní sezónu otevřelo masivní vermírné RPG Starfield od Bethesdy. Už si můžete přečíst naši recenzi a určitě jste viděli i to, jak se poslední opus Todda Howarda hraje. Teď si o svých zážitcích z vesmíru popovídá i zbytek redakce. | (00:00) Úvod | (01:55) Starfield | (1:12:16) Dotazy | (1:16:15) Závěr a pravidlo \#642}{}{}{}
\FC{fight-club-643.mp3}{\FCtitle{\#643}{My a Asassin's Creed}}{2023-09-13}{\Ulrik, \PHajda, \PMakal, \JMalcharek}{\FCrule{\#643}{Jsou fanoušci Assing Creedu 3 vůbec lidi?}{\PMakal}}{Patrik si mohl zahrát Assassin's Creed Mirage, a tak nám povypráví o svých zážitcích z Bagdádu. Kromě toho si ale zhodnotíme celou sérii a podělíme se o své vzpomínky na starší díly. | (00:00) Začátek | (03:45) Patrik o AC Mirage | (50:00) Dotazy (Horizon, Spiderman, Bojovky) | (58:00) IPhone 15 | (1:07:20) Nové zpoplatnění Unity | (1:27:00) Lies of PI | (1:37:25) Závěr}{}{}{}
\FC{fight-club-644.mp3}{\FCtitle{\#644}{Naše podzimní očekávání a povídání o Cyberpunku a Xboxu}}{2023-09-20}{\OverWatch, \Luisa, \PMakal, \JMalcharek}{\FCrule{\#644}{Nový X-Box? Ne, raději komín z Osvětimi.}{\JMalcharek}}{Letošní podzimní sezóna je stejně nabitá jako zbytek roku, a jak to zatím vypadá, budeme mít až do prosince o zábavu postaráno. Na co se těšíme nejvíc, a co naopak podle nás propadne? | (00:00) Úvod | (02:55) Podzimní očekávání | (51:56) Microsoft leak | (57:45) Cyberpunk 2077: Phantom Liberty | (1:11:40) Dotazy | (1:12:55) Závěr a pravidlo \#644}{}{}{}
\FC{fight-club-645.mp3}{\FCtitle{\#645}{O Cyberpunk 2077: Phantom Liberty a návratu do Night City}}{2023-09-28}{\Luisa, \PMakal, \JMalcharek}{\FCrule{\#645}{Se sexuálními predátory nikdo nechce pracovat.}{\PMakal}}{Skoro tři roky po vydání Cyberpunku 2077 vychází první a jediné rozšíření s názvem Phantom Liberty. Jak se vývojářům z CD Projektu povedlo upravit hru v rámci Update 2.0 a jaký je nový příběh? O tom všem si budeme povídat v aktuálním Fight Clubu. | (00:00) Úvod | (02:25) Cyberpunk 2077: Phantom Liberty | (41:06) Co se stalo za minulý týden? | (43:58) : Co hrajeme? | (46:42) Dotazy | (1:07:43) Závar a pravidlo \#645}{}{}{}
\FC{fight-club-646.mp3}{\FCtitle{\#646}{O hrách s novinářem Jakubem Szántó}}{2023-10-05}{\Ulrik, \PMakal, \JSzanto}{\FCrule{\#646}{Když do boje, tak jedině s Oops!... I did it again.}{\PMakal}}{Novináře, reportéra a zahraničního zpravodaje Jakuba Szántó jistě není třeba představovat. Specialista na Blízký východ, válečné zpravodajství a autor několika knih je ale mimo jiné také vášnivým hráčem, a to už od dob osmibitů. O zajímavá témata v tomto díle Fight Clubu jistě nebude nouze. | (00:00) Úvod | (02:16) Začátky v hraní | (18:00) Fallout | (26:25) Válečný hry | (33:45) Karetní hry a deskové hry | (36:31) Call of duty mobile a další hry | (46:41) Kultura v Izraeli | (1:11:19) Hry ze současných válek | (1:16:39) Závěr a pravidlo \#646}{}{}{}
\FC{fight-club-647.mp3}{\FCtitle{\#647}{O našich nejtěžších recenzích}}{2023-10-11}{\Ulrik, \Luisa, \PMakal, \JMalcharek}{\FCrule{\#647}{Spíš než sex.cz čtěte games.cz, nepřijde si pro vás policie.}{\JMalcharek}}{Jsou hry, které se hrají s radostí a recenze na ně se píší skoro samy. A pak jsou tu ty druhé. Právě o nich si budeme povídat v dnešním Fight Clubu. | (00:00) Úvod | (02:20) Steam next fest | (17:30) Dotaz – Forza Motorsport | (21:58) Nová PS5 | (29:40) O našich nejtěžších recenzích | (57:50) Dotazy (hodnocení, trofeje, platina…) | (1:14:20) Poděkování}{}{}{}
\FC{fight-club-648.mp3}{\FCtitle{\#648}{O komiksech ve hrách}}{2023-10-19}{\PHajda, \Luisa, \PMakal, \JMalcharek}{\FCrule{\#648}{Štít je ňůďo.}{\JMalcharek}}{Spider-Man 2 bere herní svět útokem, a tak zavzpomínáme na naše oblíbené i neoblíbené herní adaptace komiksových příběhů. | (00:00) Úvod | (01:25) Úvod 2. pokus | (03:10) Spiderman 2 | (26:45) Super Mario Wonder | (32:55) Gothic | (39:30) Hlavní téma | (1:15:34) Závěr}{}{}{}
\FC{fight-club-649.mp3}{\FCtitle{\#649}{O dokončení akvizice Activision Blizzard}}{2023-10-26}{\Ulrik, \OverWatch, \Luisa, \JMalcharek}{\FCrule{\#649}{Počítač chlazený lidskou krví? Ne, počkejte raději na cumputer.}{\Luisa}}{Po dlouhých peripetiích se Microsoftu definitivně podařilo uzavřít akvizici vydavatelství Activision Blizzard. Co to znamená do budoucna a jaká jsou naše očekávání se dozvíte v novém Fight Clubu | (00:00) Úvod | (02:12) Hlavní téma | (29:55) Novinky (Kampaň pro Diablo, Alan Wake 2) | (51:42) Dotazy | (1:03:35) Co hrajeme? | (1:09:35) Závěr a pravidlo \#649}{}{}{}
\FC{fight-club-650.mp3}{\FCtitle{\#650}{O našich nejmilovanějších herních systémech}}{2023-11-02}{\Ulrik, \Luisa, \PMakal}{\FCrule{\#650}{My jsme tým hry.}{\PMakal}}{Většina z nás v redakci hraje víc než dvě dekády (slovutný kmet Aleš dokonce víc než tři), a tak nám za tu dobu prošlo pod rukama velké množství různých herních systémů. S jakým jsme si užili nejvíc zábavy a jaký nás nakonec zklamal? | (00:00) Úvod | (04:20) Hlavní téma | (1:07:20) Závěr a pravidlo \#650}{}{}{}
\FC{fight-club-651.mp3}{\FCtitle{\#651}{Osmdesátkoví hrdinové ve hrách}}{2023-11-09}{\Ulrik, \PHajda, \PMakal, \JMalcharek}{\FCrule{\#651}{Patrik je Terminátor.}{\JMalcharek}}{Robocop: Rogue City je až překvapivě povedená střílečka, a tak se společně pobavíme o dalších zářezech osmdesátkových filmových hrdinů v herním světě. | (00:00) Úvod | (04:02) GTA IV | (10:17) Legend of Zelda film | (14:08) Call of Duty | (20:54) Příběh s UltraWide monitorem | (28:58) Hlavní téma | (1:06:58) Co hraje Patrik? | (1:09:52) Dotazy | (1:31:10) Závěr}{}{}{}
\FC{fight-club-652.mp3}{\FCtitle{\#652}{S Pavlem Barákem o Game Developers Session}}{2023-11-16}{\PMakal, \JMalcharek, \PBarak}{\FCrule{\#652}{Stavte se na GDS-ko, možná nad mísou chlebíčků potkáte Briana Farga.}{\JMalcharek}}{Ve studiu nás tento týden navštíví předseda Asociace českých herních vývojářů Pavel Barák a společně si popovídáme hlavně o akci Game Developers Session, jejíž další ročník proběhne již velmi brzy | (00:00) Úvod a představení hosta | (01:29) Game Developers Session | (06:00) Hosti na Game Developers Session | (20:04) Dotaz: Je možnost se na konferenci poohlídnout po práci v herním byznyse? | (27:10) Na co se tam můžeme těšit? | (33:00) Naše cíle do budoucna? | 41: 36 Český herní průmysl | 45: 31 Asociace českých herních vývojářů | (1:04:25) Co hraje Pavel Barák? | (1:09:55) Závěr}{}{}{}
\FC{fight-club-653.mp3}{\FCtitle{\#653}{Half-Life slaví 25. narozeniny}}{2023-11-23}{\Ulrik, \Luisa, \PMakal, \JMalcharek}{\FCrule{\#653}{Nezlobte nebo dostanete Spáčidlem.}{\JMalcharek}}{Jedna z nejdůležitějších stříleček historie slaví 25. narozeniny, a tak společně zavzpomínáme na své zážitky v kůži Gordona Freemana či sličné Alyx a pobavíme se o tom, kdy se tedy konečně dočkáme toho třetího dílu | (00:00) Úvod | (02:39) Remaster Last of Us Part II | (10:27) Suicide Squad | (15:30) Cyberpunk 2077 Definitive Edition | (26:29) Hlavní téma Half-Life slaví 25. narozeniny | (56:02) Co hraje Aleš? | (57:15) Co hraje Šárka? | (59:49) Co hraje Jakub? | (1:03:57) Dotazy | (1:17:11) Závěr a pravidlo \#653}{}{}{}
\FC{fight-club-654.mp3}{\FCtitle{\#654}{O vlacích ve hrách}}{2023-11-30}{\Ulrik, \PHajda, \PMakal, \JMalcharek}{\FCrule{\#654}{Š, š, š, š, hůů.}{\PMakal}}{Tento týden vychází nová česká hra Last Train Home, ve které hraje vlak důležitou roli. Ale v české (a nejen v ní) herní historii najdeme her, nad kterými by zaplesal nejeden šotouš, celou řadu. | (00:00) Úvod | (01:54) Spor mezi Valve vs Wildfire | (06:00) Finanční výsledky CD Projektu | (09:59) Last train home | (29:30) Hlavní téma | (50:22) Jak trávíme čas ve vlacích? | (56:02) Kde bychom chtěli vlaky? | (58:44) Co hraje Patrik? | (1:03:55) Co hraje Aleš? | (1:05:02) Co hraje Jakub? | (1:08:43) Co hraje Pavel? | (1:14:55) Závěr a pravidlo \#654}{}{}{}
\FC{fight-club-655.mp3}{\FCtitle{\#655}{Co čekáme od Game Awards}}{2023-12-07}{\PHajda, \Luisa, \PMakal, \JMalcharek}{\FCrule{\#655}{Kdo se s námi nedívá, není náš kamarád ani aspirant.}{\PMakal}}{Tento týden dojde k vyhlášení nejlepších her roku na již tradiční akci Game Awards, kromě rozdávání sošek se ale můžeme těšit na množství trailerů a oznámení. Jaká máme letos očekávání?}{}{}{}
\FC{fight-club-656.mp3}{\FCtitle{\#656}{Ohlížíme se za letošním rokem}}{2023-12-13}{\Luisa, \PMakal, \JMalcharek}{\FCrule{\#656}{Je lepší zimní deprese, než mít očního herpese.}{\PMakal}}{Předposlední Fight Club tohoto roku je za námi, a tak je na čase rekapitulovat a bilancovat. Co se nám nejvíce líbilo? Co se povedlo, a co naopak ne? To vše probereme v dnešním podcastu.}{}{}{}
\FC{fight-club-657.mp3}{\FCtitle{\#657}{Předvánoční radovánky}}{2023-12-20}{\Ulrik, \Luisa, \PMakal}{\FCrule{\#657}{All I Want for Christmas Is privatizace České pošty.}{\PMakal}}{Poslední Fight Club letošního roku se ponese v uvolněné atmosféře radosti z konce roku. Pojďme si vyprávět o tom, co plánujeme dělat během volných dní a co si přejeme pod stromeček!}{}{}{}
\FC{fight-club-658.mp3}{\FCtitle{\#658}{Jaký bude rok 2024?}}{2024-01-03}{\Ulrik, \PHajda, \Luisa, \PMakal, \JMalcharek}{\FCrule{\#658}{Patrikovy horké úchopy vždy potěší.}{\PMakal}}{Kola herního průmyslu se sice budou teprve pomalu roztáčet, ale my máme své tipy a očekávání, o které se s vámi chceme hned z kraje roku podělit.}{}{}{}
\FC{fight-club-659.mp3}{\FCtitle{\#659}{Perský princ se vrací}}{2024-01-11}{\Ulrik, \PHajda, \Luisa, \JMalcharek}{\FCrule{\#659}{Tak co to dneska bude? Princ nebo Cornetto?}{\PHajda}}{Ubisoft konečně vydává nové pokračování slavné série Prince of Persia a my zavzpomínáme na naše zážitky ze starších dílů}{}{}{}
\FC{fight-club-660.mp3}{\FCtitle{\#660}{Bossové ve hrách}}{2024-01-17}{\OverWatch, \Luisa, \PMakal, \JMalcharek}{\FCrule{\#660}{Nejlepší bossové jsou v Bullerbynu.}{\PMakal}}{Jací jsou naši nejoblíbenější bossové? Kdo nám nejvíc zatopil? A jakým bossem bychom byli my? O tom všem si povídáme v dnešním Fight Clubu}{}{}{}
\FC{fight-club-661.mp3}{\FCtitle{\#661}{O japonských hrách na hrdiny}}{2020}{\Ulrik, \Luisa, \JMalcharek}{\FCrule{\#661}{Nelevelujte Eris.}{\Ulrik}}{Tento týden vychází Like a Dragon: Infinite Wealth, a to je ideální příležitost, aby se někteří členové redakce vyznali ze své lásky k JRPG. }{}{}{}
\FC{fight-club-662.mp3}{\FCtitle{\#662}{O hranicích mezí inspirací a plagiátem}}{2020}{\PHajda, \Luisa, \PMakal}{\FCrule{\#662}{Ty nejlepší věci se většinou dějí před nebo po vysílání.}{\PMakal}}{Jak to vypadá, když to vývojáři s poctami svým vzorům trochu přeženou. }{}{}{}
\FC{fight-club-663.mp3}{\FCtitle{\#663}{Skončí Microsoft jako Sega?}}{2024-02-08}{\Ulrik, \Luisa, \PMakal, \JMalcharek}{\FCrule{\#663}{Co je šeptem, to je Philem Spencerem.}{\JMalcharek}}{Herním světem se šíří zvěsti o tom, že Microsoft hodlá své exkluzivity nakonec portovat i na ostatní platformy. S čím dál větším důrazem na Game Passy to v budoucnu mohlo znamenat konec konzolového hardware od Microsoftu.}{}{}{}
\FC{fight-club-664.mp3}{\FCtitle{\#664}{Hry pod pirátskou vlajkou}}{2024-02-15}{\PHajda, \Luisa, \PMakal}{\FCrule{\#664}{Mít lemura je lepší než mít vlka.}{\PMakal}}{Nebojte, tentokrát si nebudeme povídat o problémech s kradením her, raději probereme naše nejoblíbenější tituly s piráty, korzáry a bukanýry!}{}{}{}
\FC{fight-club-665.mp3}{\FCtitle{\#665}{Tak jde čas s Larou Croft}}{2024-02-21}{\OverWatch, \PMakal, \JMalcharek}{\FCrule{\#665}{Hranatá prsa jsou nesmrtelná.}{\PMakal}}{K příležitosti vydání remasterů prvních tří dílů série Tomb Raider vzpomínáme na naše zkušenosti se slavnou archeoložkou a sdílíme svá přání do budoucna}{}{}{}
\FC{fight-club-666.mp3}{\FCtitle{\#666}{Peklo ve hrách}}{2024-02-28}{\Ulrik, \PHajda, \PMakal, \JMalcharek}{\FCrule{\#666}{Peklo jsou Ti druzí.}{\PMakal}}{Jaké jiné téma by měl mít Fight Club s pořadovým číslem 666?}{}{}{}
\FC{fight-club-667.mp3}{\FCtitle{\#667}{Realismus a hry}}{2024-03-07}{\PHajda, \Luisa, \PMakal, \JMalcharek}{\FCrule{\#667}{Zabíjet lidi ve spánku je v pohodě.}{\PMakal}}{Jak moc realismus hrám prospívá, a jaká míra už zážitku začíná škodit? O tom si budeme povídat v aktuálním Fight Clubu}{}{}{}
\FC{fight-club-668.mp3}{\FCtitle{\#668}{Karetní hry}}{2024-03-15}{\OverWatch, \PHajda, \PMakal}{\FCrule{\#668}{V zájmu produktivity celého lidstva by měli Balatro zakázat.}{\PHajda}}{Karetních her nebo aspoň miniher je pořádná spousta a momentálně navíc frčí Balatro, a tak se Patrik rozhodl, že nás trochu vzdělá. A vy můžete být u toho!}{}{}{}
\FC{fight-club-669.mp3}{\FCtitle{\#669}{Limity ve hrách}}{2024-03-21}{\Ulrik, \Luisa, \PHajda, \JMalcharek}{\FCrule{\#669}{Njevětší omezení jsou daně z příjmu fyzických osob.}{\Luisa}}{Tento týden vychází řada velkých her, ve Fight Clubu si ale budeme povídat o Dragon's Dogma 2 a předevších o různých omezeních, které nám hry staví do cesty. Může jít o limit nosnosti postavy, správu staminy a další podobné herní mechanismy.}{}{}{}
\FC{fight-club-670.mp3}{\FCtitle{\#670}{Zmrtvýchvstání ve hrách}}{2024-03-28}{\Ulrik, \Luisa, \PMakal, \JMalcharek}{\FCrule{\#670}{Sebevražednému oddílu by nepomohl ani Ježíš.}{\JMalcharek}}{Velikonoce se blíží, a tak jsme se rozhodli probrat různé způsoby oživování ve hrách.}{}{}{}
\FC{fight-club-671.mp3}{\FCtitle{\#671}{Proč nesnášíme apríl?}}{2024-04-04}{\Ulrik, \Luisa, \PMakal, \JMalcharek}{\FCrule{\#671}{Příští Fight Club nebude... apríl, ha ha ha.}{\JMalcharek}}{První duben je den, kdy se na sociálních sítích herní firmy předhánějí v co nejvypečenějších vtípcích. Bohužel je někdy opravdu těžké rozeznat, zda jde o pravdivou zprávu nebo takzvaný humor. Právě o tom si budeme povídat v novém Fight Clubu.}{}{}{}
\FC{fight-club-672.mp3}{\FCtitle{\#672}{O radioaktivním spadu}}{2024-04-11}{\Ulrik, \PHajda, \PMakal, \JMalcharek}{\FCrule{\#672}{Válka, válka je stále stejná.}{\PMakal}}{Na Amazon Prime tento týden dorazí seriálová adaptace Falloutu, a než její tvůrci stihnout naštvat všechny fanoušky původních dílů, popovídáme si o tom, co pro nás tohle legendární RPG znamená.}{}{}{}
\FC{fight-club-673.mp3}{\FCtitle{\#673}{O hrách, které tu už měly být}}{2024-04-17}{\PHajda, \Luisa, \PMakal, \JMalcharek}{\FCrule{\#673}{Pozor na laviny!}{\PMakal}}{Některým oznámeným hrám, případně vytouženým pokračováním, to na svět trvá až příliš dlouho. Právě o nich si budeme povídat v dnešním Fight Clubu}{}{}{}
\FC{fight-club-674.mp3}{\FCtitle{\#674}{O Kingdom Come: Deliverance 2 s designérem Ondřejem Bittnerem}}{2024-04-24}{\PHajda, \PMakal, \OBittner}{\FCrule{\#674}{Ti nejlepší slavíci nejsou z Madridu, ale z Ratají.}{\PMakal}}{Kingdom Come: Deliverance 2 je od minulého týdne jasnou skutečností. V aktuálním Fight Clubu si o hře budeme povídat s designérem Ondřejen Bittnerem z Warhorse}{}{}{}
\FC{fight-club-675.mp3}{\FCtitle{\#675}{O PR katastrofách v herním průmyslu}}{2024-05-08}{\Ulrik, \Luisa, \PMakal}{\FCrule{\#675}{Hráče není radno nasírat.}{\PMakal}}{Helldivers 2 se z miláčka milionů stali zdrojem velmi negativních vášní. Jaké známe další příklady, kdy herní firmy nezvládli komunikaci, případně se ke hráčům zachovaly nefér? O tom si budeme povídat tento týden}{}{}{}
\FC{fight-club-676.mp3}{\FCtitle{\#676}{Jsou AAA hry na ústupu?}}{2024-05-16}{\PHajda, \Luisa, \PMakal, \JMalcharek}{\FCrule{\#676}{Dobrá hratelnost překoná i hnusnou grafiku.}{\PMakal}}{Ne, nebojte, určitě si nemyslíme, že největší produkci odzvonilo. Jen nám první měsíce letošního roku připomněly, že menší tituly s daleko nižším rozpočtem dokáží chytnout za srdíčko často mnohem víc a taky se v poslední době vyrojilo několik zajímavých indie her, které strkají AAA konkurenci do kapsy.}{}{}{}
\FC{fight-club-677.mp3}{\FCtitle{\#677}{Co čekáme od Summer Game Festu?}}{2024-05-22}{\Ulrik, \Luisa, \PMakal, \JMalcharek}{\FCrule{\#677}{Věra je nová Zdena.}{\JMalcharek}}{Summer Game Fest se blíží a my jsme ready na zásobu nových her. Co bychom si od show Geoffa Keighleyho přáli nejvíc?}{}{}{}
\FC{fight-club-678.mp3}{\FCtitle{\#678}{O hrách, které zestárly dobře}}{2024-05-31}{\OverWatch, \PHajda, \PMakal}{\FCrule{\#678}{Nové nemusí byt vždycky lepší.}{\PMakal}}{Některé hry vzdorují proudu času lépe a jiné hůře. O nejlepších příkladech si budeme povídat v tomto díle Fight Clubu.}{}{}{}
\FC{fight-club-679.mp3}{\FCtitle{\#679}{Jaké hry budou na letošních konferencích chybět?}}{2024-06-06}{\Ulrik, \PHajda, \JMalcharek}{\FCrule{\#679}{Aleš má rád Necromundu.}{\JMalcharek}}{Tento týden proběhne hned několik herních konferencí a určitě se na nich ukáže spousta zajímavých her. Některé dlouho očekávané ale určitě budou chybět}{}{}{}
\FC{fight-club-680.mp3}{\FCtitle{\#680}{Co nás nejvíc oslovilo na Summer Game Festu?}}{2024-06-13}{\Ulrik, \Luisa, \JMalcharek}{\FCrule{\#680}{Microsoft na pomoc na záchranu.}{\JMalcharek}}{Máme za sebou prakticky vše, co letošní Summer Game Fest alias nE3 2024 nabídl a my se vrátíme k tomu nejlepšímu...a nejhoršímu.}{}{}{}
\FC{fight-club-681.mp3}{\FCtitle{\#681}{Elden Ring a ti druzí ze Souls}}{2024-06-20}{\PHajda, \Luisa, \PMakal}{\FCrule{\#681}{Zážitek nemusí být vždy příjemný, hlavně když je silný.}{\PMakal}}{Tentokrát si sedneme a probereme nejen jak dopadl nejnovější datadisk pro Elden Ring, ale také vlastně jak jsou na tom Souls-like hry v současnosti.}{}{}{}
\FC{fight-club-682.mp3}{\FCtitle{\#682}{U jakých her se nejlépe zchladíte?}}{2024-06-27}{\Ulrik, \Luisa, \JMalcharek}{\FCrule{\#682}{Bruslení s řemdihy je nejlepší.}{\JMalcharek}}{Jaké hry jsou nejlepší, abyste se v létě pořádně ochladili? Na tuto palčivou otázku vám odpoví podcast Fight Club s pořadovým číslem 682. Poskytneme vám hromadu tipů, jak přežít úmorné počasí, když zrovna z jakéhokoli důvodu musíte zůstat doma.}{}{}{}
\FC{fight-club-683.mp3}{\FCtitle{\#683}{Jazyky ve hrách}}{2024-07-04}{\Ulrik, \Luisa, \JMalcharek}{\FCrule{\#683}{Jazyky jsou skvělé hlavně místo bradavek.}{\JMalcharek}}{Blíží se svátek příchodu slovanských věrozvěstů Cyrila a Metoděje, a proto se podíváme na hry, které zajímavě pracují s jazykem.}{}{}{}
\FC{fight-club-684.mp3}{\FCtitle{\#684}{Ke kterým hrám se rádi vracíme?}}{2024-07-11}{\Ulrik, \Luisa, \JMalcharek}{\FCrule{\#684}{Halí belí, casus belli.}{\JMalcharek}}{Když zrovna nejsme zahlceni přívalem nových a krásných her, k jakým titulům se nejraději vracíme a čím vyplňujeme herní sucho? Na to odpoví v novém díle Fight Clubu Aleš, Šárka a Jakub.}{}{}{}
\FC{fight-club-685.mp3}{\FCtitle{\#685}{Freestyle o všem, co hrajeme a sledujeme}}{2024-07-18}{\Ulrik, \Luisa, \JMalcharek}{\FCrule{\#685}{Aleš se rozkládá.}{\JMalcharek}}{Protože vychází hry, filmy a seriály, které nemají jednotné téma, budeme si povídat o všem, co hrajeme a sledujeme v popkulturním freestyle. Na čem byla Šárka v kině? Co recenzoval Aleš? A u jakou betu zkoušel Jakub? O tom se pobavíme v dnešním Fight Clubu.}{}{}{}
\FC{fight-club-686.mp3}{\FCtitle{\#686}{Loučení s Adamem, Patrikem a Pavlem}}{2024-07-24}{\Ulrik, \OverWatch, \PHajda, \PMakal}{\FCrule{\#686}{Byla to jízda. Sbohem a šáteček. Když vlákna, tak herní.}{\PHajda{} + \PMakal{} + \OverWatch}}{Gamesy si koncem června prošly neveselou personální změnou a musely se rozloučit s trojicí redaktorů, kolegů a kamarádů.  Adam, Patrik a Pavel se ale přijdou rozloučit pořádně, zavzpomínají na staré dobré časy a hlavně prozradí, na čem teď pracují a kde je nadále můžete sledovat.}{}{}{}
\FC{fight-club-687.mp3}{\FCtitle{\#687}{S Alžbětou Trojanovou}}{2024-08-01}{\Bety, \Luisa, \JMalcharek}{\FCrule{\#687}{Vyhlížejte Bloody Hearts of Europe.}{\JMalcharek}}{Do Fight Clubu se na kus řeči vrací na návštěvu Bětka Trojanová, moderátorka a bývalá redaktorka Games. Budeme se bavit o tom, co momentálně hraje, jakým projektům se věnuje, a která hra ji v poslední době učarovala. Společnost jí bude dělat Šárka s Jakubem.}{}{}{}
\FC{fight-club-688.mp3}{\FCtitle{\#688}{S Michalem Holánem}}{2024-08-08}{\Ulrik, \JMalcharek, \MHolan, \MTucan, \JVana}{\FCrule{\#688}{Kudy se de do těch Budějovic?}{\JMalcharek}}{Druhý týden po sobě přivítáme v redakci Fight Clubu hosta. Bude jím herec a dabér Michal Holán, který si zahrál například v seriálu Policie Modrava. Nás bude ale zajímat hlavně jeho práce z dabérské kukaně. Kromě namlouvání postav velkých seriálů jako The Walking Dead nebo Breaking Bad, propůjčil svůj hlas i do her, jmenovitě World of Warships. Vyzpovídá ho Aleš s Jakubem.}{}{}{}
\FC{fight-club-689.mp3}{\FCtitle{\#689}{Mají kamenné prodejny her budoucnost?}}{2024-08-15}{\Ulrik, \Luisa, \JMalcharek}{\FCrule{\#689}{Hentai figurky letí.}{\JMalcharek}}{Uhýbají fyzické prodeje her těm digitálním? A pokud ano, jaká v tom případě čeká budoucnost kamenné prodejny s hrami? O tom vyzpovídá v dnešním Fight Clubu Jakub, Šárka a Aleš hosta, kterým je Petr Kratochvíl z řetězce obchodů JRC a Smarty}{}{}{}
\FC{fight-club-690.mp3}{\FCtitle{\#690}{Jsme na Gamescomu}}{2024-08-22}{\Ulrik, \Luisa, \JMalcharek}{\FCrule{\#690}{...}{\JMalcharek}}{Největší evropský videoherní festival si samozřejmě nemůžeme nechat ujít. Z našeho skromného ubytování v Kolíně nad Rýnem se v improvizovaném studiu pokusíme s Vámi spojit ve čtvrtek večer a sdělit nejčerstvější dojmy z nových her, které jsme si na Gamescomu vyzkoušeli, a z atmosféry kolínského výstaviště.}{}{}{}
\FC{fight-club-691.mp3}{\FCtitle{\#691}{O hrách s novinářem Vietem Tranem}}{2024-08-29}{\Luisa, \JMalcharek, \VietTran}{\FCrule{\#691}{Vyhlížejte zlatou éru.}{\JMalcharek}}{K dnešnímu Fight Clubu se k Šárce a Jakubovi přidá také hráč a novinář Viet Tran, který pro CNN Prima News zpracovává nejen zpravodajská témata, ale i zprávy z videoherního průmyslu.}{}{}{}
\FC{fight-club-692.mp3}{\FCtitle{\#692}{Na co se těšíme v podzimním návalu}}{2024-09-05}{\Ulrik, \Luisa, \JMalcharek}{\FCrule{\#692}{Cibule je šalotka je cibule proi lepší lidi. Ne, šalotka je cibule pro lepší lidi. Aňa Šalotka je cibule a má vrstvy. Aňa Šalotka je Yennefer pro lepší lidi. Yennefer je cibule pro Aňu, ne... Yenefer je šalotka Witchera.}{\JMalcharek{} + \Luisa{} + \Ulrik}}{Končí nám letní sucho a s ním se povážlivě otřásá přehrada plná her, která hrozí, že nás zavalí. Aleš, Šárka a Jakub už poctivě recenzují první podzimní vlaštovky a s nadějí vyhlíží všechny ty velké a krásné hry, které se chystají.}{}{}{}
\FC{fight-club-693.mp3}{\FCtitle{\#693}{Nejen o Star Wars Outlaws s Lukášem Grygarem}}{2024-09-11}{\Ulrik, \Luisa, \Janchor}{\FCrule{\#693}{I krize středního věku je příležitost.}{\Janchor}}{Do Fight Clubu po dlouhé době zavítá novinář Lukáš Grygar, veterán herní žurnalistiky a věčný přítel redakce. S Alešem a Šárkou mimo jiné probere, jak se mu líbí nejnovější akční adventura ze světa Hvězdných válek, co dalšího teď hraje a co obecně dělá se svým životem, když už primárně nepokrývá hry.}{}{}{}
\FC{fight-club-694.mp3}{\FCtitle{\#694}{O Crusader Kings, fotbale a podcastech s Vaškem Pecháčkem}}{2024-09-19}{\Ulrik, \Luisa, \VPechacek}{\FCrule{\#694}{Když na deadline kašleš, zbičuje je Ašleš.}{\VPechacek}}{Do Fight Clubu tentokrát zavítá Vašek Pecháček, bývalý člen redakce Games.cz a stávající hvězda podcastů Tam a zase zpátky nebo Kontrapresink. Popovídáme si s ním nejen o Prstenech moci a fotbale, ale i o tom, jestli má ve svém nabitém programu vůbec ještě čas hrát videohry.}{}{}{}
\FC{fight-club-695.mp3}{\FCtitle{\#695}{Je Zelda lepší než Link?}}{2024-09-26}{\Ulrik, \Luisa, \JMalcharek}{\FCrule{\#695}{Aleš asimiluje neštovice.}{\JMalcharek}}{Když vyjde nová Zelda, je to vždycky velká událost a nejinak tomu je i v případě Echoes of Wisdom. Konečně tak dostaneme odpověď na to, jestli je Zelda skutečně lepší než hrdinný Link.}{}{}{}
\FC{fight-club-696.mp3}{\FCtitle{\#696}{Ubisoft má problémy}}{2024-10-03}{\Luisa, \JMalcharek}{\FCrule{\#696}{Intimčo s Jakubem a Šárkou za každého počasí.}{\JMalcharek}}{Francouzský gigant Ubisoft se potýká s problémy. Star Wars Outlaws se špatně prodávají, střílečka XDefiant nedokázala přitáhnout požadovaný počet hráčů a Assassin's Creed Shadows musel být odložený na únor, ve kterém je navíc obrovský přetlak titulů. Jakou čeká firmu budoucnost?}{}{}{}
\FC{fight-club-697.mp3}{\FCtitle{\#697}{Nová nálož hororů}}{2024-10-10}{\Ulrik, \Luisa, \JMalcharek}{\FCrule{\#697}{Ve svých neklidných snech vidím to město... Ostrava.}{\JMalcharek}}{Říjen je pro mnohé měsícem hororu a stejně tak jej vnímají i herní vydavatelé. Dočkali jsme se skvělého Silent Hillu 2, remaku Until Dawn, chystá se hra podle filmů Tiché místo. Jakub, Šárka a Aleš ale zavzpomínají i na další hry, které je kdy děsily.}{}{}{}
\FC{fight-club-698.mp3}{\FCtitle{\#698}{Které hry nám letos udělaly radost?}}{2024-10-17}{\Ulrik, \JMalcharek}{\FCrule{\#698}{Naposled se špatným zvukem.}{\JMalcharek}}{Rok se přehoupl do poslední čtvrtiny, a tak si Aleš s Jakubem připomenou hry, které jim zatím letos udělaly radost, a na které by se při každoročním zúčtování nemělo zapomenout.}{}{}{}
\FC{fight-club-699.mp3}{\FCtitle{\#699}{Kdo potřebuje herní média?}}{2024-10-24}{\Ulrik, \Luisa, \JMalcharek}{\FCrule{\#699}{Jsme pro aby poplatky Games byly dány zákonem.}{\Luisa}}{Herní průmysl zažívá krizi, která se přeneseně dotýká i herních médií. Ostatně i my v Games jsme si letos prošli obrovskou personální změnou a zeštíhlování redakcí jde vidět na každém rohu. Proto se Aleš, Jakub a Šárka pobaví o tom, jakou roli vlastně herní média mají a jak se mění.}{}{}{}
\FC{fight-club-700.mp3}{\FCtitle{\#700}{Odpovídáme na vše, co jste kdy chtěli vědět}}{2024-10-31}{\Ulrik, \Luisa, \JMalcharek}{\FCrule{\#700}{Víte někdo jestli má Šárka sklep?}{\JMalcharek}}{Standardně kulaté podcasty slavíme v poměrně velkém stylu a setkáváme se během nich i s vámi, naším věrným čtenářstvem, diváctvem, posluchačstvem. Letos ale naši redakci postihly nešťastné personální změny, kdy jsme se bohužel museli rozloučit s našimi dlouholetými kolegy a kamarády – Adamem, Patrikem a Pavlem. A protože je tahle rána i pro nás pořád ještě poměrně čerstvá, usoudili jsme, že na bujaré večírky ještě není ta správná doba. | Sedmistovku nicméně nenecháme projít jen tak bez povšimnutí, jen její příchod pojmeme v trochu komornějším, umírněnějším duchu. Budeme odpovídat na vaše dotazy, které jste nám v posledních dvou týdnech zaslali do mailu, ankety na webu, nebo je klidně i položíte živě během vysílání. Můžete se nás zeptat úplně na cokoliv, na libovolné téma, klidně i mimo hry, na co v klasickém Fight Clubu bez jubilea pravděpodobně nepřijde řada, anebo bychom odmávli jako příliš velký off-topic.}{}{}{}
\FC{fight-club-701.mp3}{\FCtitle{\#701}{O novém Dragon Age s Karlem Drdou}}{2024-11-06}{\Ulrik, \Luisa, \KDrda}{\FCrule{\#701}{Všichni jsme adoptivní synové Karla Drdy.    }{\Luisa}}{K debatě o pokračování kultovní RPG série Dragon Age jsme si pozvali odborníka na slovo vzatého – bývalého redaktora Games Karla Drdu. Diskutovat s ním o kvalitách The Veilguard bude hlavně Šárka, která má hru už dohranou, a Aleš, který nové dobrodružství v Thedasu teprve začíná.}{}{}{}
\FC{fight-club-702.mp3}{\FCtitle{\#702}{Na výlet do virtuální Prahy}}{2024-11-14}{\Ulrik, \Luisa}{\FCrule{\#702}{Max Bolest žere jenom lentilky.}{\Ulrik}}{Praga Mater Urbium. Matka měst patří mezi nejkrásnější města na světě a ročně láká tisícovky turistů. Je ale stejně hojně využívána a zobrazena ve hrách? O virtuálním obraze Prahy se pobaví Aleš se Šárkou.}{}{}{}
\FC{fight-club-703.mp3}{\FCtitle{\#703}{O novém Stalkerovi}}{2024-11-21}{\Ulrik, \Luisa, \JMalcharek}{\FCrule{\#703}{Lepší než kafe z nekra je hrát Stalkera.}{\JMalcharek}}{Po deseti letech zkroušeného vývoje vychází dlouho očekávaný Stalker 2 a Jakub, Aleš a Šárka se tím pádem nemohou bavit o ničem jiném. Zavítejte s námi do černobylské Zóny.}{}{}{}
\FC{fight-club-704.mp3}{\FCtitle{\#704}{Hraní na mobilech}}{2024-11-28}{\Ulrik, \Luisa, \JMalcharek}{\FCrule{\#704}{Ani Rammstein, ani U2, nejlepší je Mewtwo.}{\JMalcharek}}{Redakci okouzlila nejnovější sběratelská mánie Pokémon Trading Card Game Pocket, a tak Aleš, Šárka a Jakub zavzpomínají na další mobilní hity. Třeba na hodiny utopené v Marvel Snap!, tisíce v Magic: The Gathering nebo jak už poker kvůli Balatra nejsme schopni hrát normálně. Probereme taky jaký monetizační model nám u kapesních titulů nejvíce vyhovuje.}{}{}{}
\FC{fight-club-705.mp3}{\FCtitle{\#705}{Radíme, co pod stromeček}}{2024-12-05}{\Ulrik, \Luisa, \JMalcharek}{\FCrule{\#705}{Když nemáte jabko, dejte do salátu pomeranč.}{\JMalcharek}}{Přišel ten čas... shonu, stresu a vymýšlení, čím obdařit své blízké na Vánoce. Aleš, Jakub a Šárka proto dají ve Fight Clubu hlavy dohromady a podělí se o své tipy na dárky pod stromeček. Probereme samozřejmě i čerstvé herní novinky a jakým hrám se momentálně věnujeme.}{}{}{}
\FC{fight-club-706.mp3}{\FCtitle{\#706}{Co čekáme od The Game Awards}}{2024-12-12}{\Ulrik, \Luisa, \JMalcharek}{\FCrule{\#706}{Raději než hrát Bloodlines 2 bych se koupal v řece Ganze.}{\JMalcharek}}{Závěrečná událost roku, The Game Awards, s sebou krom vyhlášení těch nejlepších a nejkrásnějších her přináší také celou řadu oznámení. Jaké vývojáře přivítá Geoff Keighley na pódiu? Jakých trailerů se dočkáme? A kdo si odnese cenu nejvyšší? To bude věštit z křišťálové koule Jakub, Šárka a Aleš.}{}{}{}
\FC{fight-club-707.mp3}{\FCtitle{\#707}{Předvánoční zhodnocení}}{2024-12-19}{\Ulrik, \Luisa, \JMalcharek}{\FCrule{\#707}{Lepší než Geralt na střeše je COnnor v hrsti.}{\JMalcharek}}{Týden před svátky by se mohlo zdát, že budeme prázdně tlachat o tom, kdo doma dělá, jaký bramborový salát a kolik vaječňáku se vypije při pečení cukroví, ale není tomu tak! Uplynulé The Game Awards nám nadělily spoustu videoherních témat, kterým se podíváme na zoubek včetně třeba první ukázky ze Zaklínače 4. Rozloučit se před svátky přijde Aleš, Šárka a Jakub.}{}{}{}
\FC{fight-club-708.mp3}{\FCtitle{\#708}{Věštíme budoucnost}}{2025-01-09}{\Ulrik, \Luisa, \JMalcharek}{\FCrule{\#708}{Ať je teplo nebo zima, stále se ohřívá klima.}{\Ulrik}}{Redakce se po svátcích vrací v plné polní a kromě povídání o tom, jak jsme trávili Vánoce a Silvestra, si načrtneme, jak bude vypadat letošní rok. Kdo se zmýlí? Kdo bude pesmista, a kdo naopak vidí herní průmysl pozitivně? O tom se pobavíme v prvním Fight Clubu letošního roku.}{}{}{}
\FC{fight-club-709.mp3}{\FCtitle{\#709}{Do historie s Civilizací}}{2025-01-16}{\Ulrik, \JMalcharek}{\FCrule{\#709}{Kdyby místo Half=Life 3 vyšel raději Alyx 2.}{\JMalcharek}}{Nesmrtelná tahovka Civilization zanedlouho opět ožije sedmým dílem. U nás už hru recenzujeme, takže se pobavíme o tom, co nový díl dělá jinak. Zavzpomínáme ale taky na starší díly a probereme naše oblíbené vojevůdce. V klubovně se navíc Jakub, Šárka a Aleš rozpovídají o tom, co se jim přihodilo v běžném životě.}{}{}{}
\FC{fight-club-710.mp3}{\FCtitle{\#710}{Switch předá korunu nástupci}}{2025-01-23}{\Ulrik, \Luisa, \JMalcharek}{\FCrule{\#710}{Ať je léto nebo zima, Switch 2 bude určo prima.}{\JMalcharek}}{Konečně jsme se dočkali. Po osmi letech se fantasticky hravá konzole Nintendo Switch dočká nástupce. Za tu dobu nám hybridní konzolka přirostla k srdci (a prstům) a proto zavzpomínáme na posledních takřka dekádu, kdy nám cesty zpříjemňovalo hraní na Switchi a taky na to, co bychom od Switche 2 čekali.}{}{}{}
\FC{fight-club-711.mp3}{\FCtitle{\#711}{Renesance Microsoftu znamená zkázu Xboxu?}}{2025-01-30}{\Luisa, \JMalcharek}{\FCrule{\#711}{Ať je léto nebo zima, bez X-Boxu nebude to prima.}{\JMalcharek}}{Všechno je Xbox, jsme i my Xbox? Nová cesta Microsoftu kráčí multiplatformními hrami a cloudovým hraním, ostatně to dosvědčují i hry ohlášení na Xbox Developers Directu. Co to ale znamená pro konzoli jako takovou? O tom se baví Aleš, Šárka a Jakub.}{}{}{}
\FC{fight-club-712.mp3}{\FCtitle{\#712}{Sedm let čekání na Kingdom Come: Deliverance 2 u konce!}}{2025-02-06}{\Ulrik, \Luisa, \JMalcharek}{\FCrule{\#712}{Chcete být drsní hoši z Ostravy, bijte se za mlýny a šňupejte skořici.}{\Ulrik}}{Je to tady! Vychází Kingdom Come: Deliverance 2, a proto nemůžeme Fight Club zasvětit ničemu jinému než středověkému RPG z českých luhů a hájů. Jak hru hodnotí v zahraničí? Co na ni říkáme? Jaký má úspěch? To vše Aleš, Jakub a Šárka proberou v podcastu.}{}{}{}
\FC{fight-club-713.mp3}{\FCtitle{\#713}{S Tobim za oponou Kingdom Come: Deliverance 2}}{2025-02-20}{\Ulrik, \Luisa, \JMalcharek, \Tobi}{\FCrule{\#713}{Audentes fortuna iuvat.}{\JMalcharek}}{Tento týden nás ve studiu navštíví PR manažer studia Warhorse, Tobias Stolz-Zwilling. Spolu s Alešem a Šárkou nahlédnou pod pokličku vývoje Kingdom Come: Deliverance 2. Jaké strasti a starosti s druhým dílem přišly? Kolik práce je na hře ještě čeká? A jak hodnotí vydání hry a její úspěch? O tom se pobavíme ve Figh Clubu s pořadovým číslem 713.}{}{}{}
\FC{fight-club-714.mp3}{\FCtitle{\#714}{Fenomén jménem Monster Hunter}}{2025-02-27}{\Luisa, \JMalcharek}{\FCrule{\#714}{Ať si monstrum nebo lovec, byl bys lepší hlavonožec.}{\Luisa}}{Když vychází nový Monster Hunter, je to ve světě vždycky velká událost. Série akční her o boji s obřími bestiemi se celosvětově prodalo přes sto milionů kopií. S vydáním Monster Hunter Wilds se proto Jakub, Aleš a Šárka pobaví o tom, proč jsou hry tak úspěšné a proč se u nás o nich více nemluví.}{}{}{}
\FC{fight-club-715.mp3}{\FCtitle{\#715}{Kde je konec split screenu?}}{2025-03-06}{\Ulrik, \Luisa, \JMalcharek}{\FCrule{\#715}{Lepší než-li Brutal Assault je Lidl assault.}{\Ulrik}}{S vydáním další kooperativní hry na jedné obrazovce od Josefa Farese si klademe otázku, kam se vlastně split screen poděl. Režim, který byl přítomný v prvních dvou konzolových generacích, je v dnešní době spíše vzácnost. Má split screen své místo mezi hrami? Jak na něj vzpomínáme? O tom se baví Aleš, Šárka a Jakub.}{}{}{}
\FC{fight-club-716.mp3}{\FCtitle{\#716}{Bude Death Stranding 2 nejlepší hra všech dob?}}{2025-03-13}{\Ulrik, \Luisa}{\FCrule{\#716}{V červnu na pláži možná potkáme i vaši matku.}{\Luisa}}{Informacemi nabité preview k očekávanému Death Stranding 2: On the Beach nás totálně odrovnalo. Kodžimovo veledílo je v traileru tak nabité náznaky a pomrknutími, že nemůžeme jinak a emočně se o novém nášupu emocí pobavit ve Fight Clubu. Tentokrát se u mikrofonů střetne Aleš se Šárkou.}{}{}{}
\FC{fight-club-717.mp3}{\FCtitle{\#717}{Potřebuje Assassin's Creed změnu?}}{2025-03-20}{\Ulrik, \Luisa, \JMalcharek, \PHajda}{\FCrule{\#717}{Jestli někdo zvládně Assassin's Creed, je to jedině Hideo Kojima.}{\JMalcharek}}{S vydáním Assassin's Creed Shadows přichází na přetřes spekulace o budoucnosti série. Jak dlouho akčním adventurám vydrží dech bez dalšího zásadního rebootu a změn hratelnosti? O tom debatuje ve Fight Clubu Jakub s Alešem.}{}{}{}
\FC{fight-club-718.mp3}{\FCtitle{\#718}{O hrách, které jsme docenili až časem}}{2025-03-27}{\Ulrik, \Luisa, \JMalcharek}{\FCrule{\#718}{Ať je léto nebo zima, pokud chcete úspěšné GTA, musíte v něm mít padající sních z borovic.}{\Ulrik}}{K některým hrám jsme si hledali cestu složitěji. Mohlo to být zrovna aktuálním rozpoložením, (ne)chutí, anebo nám prostě nesedly. Zkrátka i kriticky oceňované a objektivně skvělé hitovky nemusíme milovat na první dobrou a přesně o takových titulech se ve Fight Clubu bude bavit Jakub, Aleš a Šárka.}{}{}{}
\FC{fight-club-719.mp3}{\FCtitle{\#719}{Hráli jsme nového Dooma, má stále co říct?}}{2025-04-03}{\Ulrik, \Luisa, \JMalcharek}{\FCrule{\#719}{Přístí zastávka... mohl by pán tam vzadu přestat chcát na ten sloup?}{\JMalcharek}}{Legendární akce Doom se letos přihlásí o slovo v řežbě s podtitulem The Dark Ages, kterou si Šárka jela vyzkoušet až do Frankfurtu. Jak se přírůstek dalšího dílu veleslavné střílečky jeví? A co samotný Doomslayer, nepatří už do důchodu? O tom se Šárka pobaví spolu s Alešem a Jakubem.}{}{}{}
\FC{fight-club-720.mp3}{\FCtitle{\#720}{S Petrem Poláčkem nejen o hrách v šéfredaktorském FIght Clubu}}{2025-04-10}{\Ulrik, \Luisa, \Sleepless}{\FCrule{\#720}{Správny Fight Club končí v 17:20.}{\Sleepless}}{Po delší době v redakci Aleš se Šárkou přivítají hosta. Není jim nikdo jiný než bývalý šéfredaktor Games, stálice Levelu a odborník na hry slovo vzatý, Petr Poláček. Shledání po letech se ponese nejen v duchu vzpomínání, ale dojde řeč i na novinky a samozřejmě taky na hry.}{}{}{}
\FC{fight-club-721.mp3}{\FCtitle{\#721}{Vyzkoušeli jsme si Switch 2. Stojí za to?}}{2025-04-17}{\Luisa, \JMalcharek}{\FCrule{\#721}{This is the way of život.}{\JMalcharek}}{Šárka vyrazila do Londýna, aby si na vlastní ruce vyzkoušela Nintendo Switch 2. Je z konzole nadšená? Jak na ní vypadají staré i nové hry? A co joy-cony? Všechny aspekty druhé generace hybridní konzole probere ve Fight Clubu spolu s Jakubem a Alešem. Bude Switch 2 stejný hit jako jednička?}{}{}{}
\FC{fight-club-722.mp3}{\FCtitle{\#722}{Má Oblivion co říct i dnes?}}{2025-04-24}{\Ulrik, \Luisa, \JMalcharek}{\FCrule{\#722}{Lepší remaster v hrsti než Oblivion co bys kamenem dohodil.}{\Luisa}}{Bethesda konečně oznámila oficiální remaster/remake The Elder Scrolls IV: Oblivion. Aleš, Šárka a Jakub na legendární RPG zavzpomínají a taky se pobaví o tom, co by od moderní verze očekávali.}{}{}{}
\FC{fight-club-723.mp3}{\FCtitle{\#723}{Svědčí velký počet remasterů a remaků o vyčerpání nápadů?}}{2025-05-01}{\Ulrik, \JMalcharek}{\FCrule{\#723}{Raději než remaster si dejte pořádný remasturbr.}{\JMalcharek}}{S každým rokem přibývá remasterů a remaků starých her, které na nás vydavatelé valí. Ať už jde i o dekády staré tituly jako v případě The Elder Scrolls IV: Oblivion Remaster nebo novější tituly jako je Days Gone Remastered, může oprašování a znovuprodávání starých her svědčit o kreativním vyčerpání? O tom bude diskutovat Jakub s Alešem.}{}{}{}
\FC{fight-club-724.mp3}{\FCtitle{\#724}{Jak se pozná dobrá Grand strategie?}}{2025-05-08}{\Ulrik, \JMalcharek}{\FCrule{\#724}{Lepší než Kobylisy je šlápnout do krysy.}{\JMalcharek}}{Paradox hodlá oznámit tajuplný Project Caser, což nás přivádí na myšlenku, v čem jsou vlastně Grand strategie tak zábavné, návykové a jak do nich nejlépe proniknout. Historickou přednášku na téma: Kterak se nejlépe dobývá svět povede zkušený stratég Aleš a bedlivě mu bude naslouchat Jakub.}{}{}{}
\FC{fight-club-725.mp3}{\FCtitle{\#725}{Co čekáme od nového dílu Mafie?}}{2025-05-15}{\Ulrik, \JMalcharek}{\FCrule{\#725}{Chcete pořádnou mafii, počkejte do voleb.}{\Ulrik}}{Konečně se nám představila nová Mafia: Old Country, což je samozřejmě skvělá příležitost pobavit se o tom, co od nového dílu kultovní série očekáváme a jak na ni vzpomínáme. Do nostalgie se pohrouží Aleš s Jakubem.}{}{}{}
\FC{fight-club-726.mp3}{\FCtitle{\#726}{O Šárčiných cestách Japonskem}}{2025-05-22}{\Luisa, \JMalcharek}{\FCrule{\#726}{Misuri nagasu.}{\Luisa}}{Šárka se vrátila po skoro třítýdenní dovolené v Japonsku, odkud si přivezla kromě spousty suvenýrů taky neméně zážitků. A to i herních! Jakubovi s Alešem ve Fight Clubu poví, co ji překvapilo, jaké bylo Nintendo muzeum v Kjótu, jak se spalo v kapslích a jestli už je ráda, že je doma.}{}{}{}
\FC{fight-club-727.mp3}{\FCtitle{\#727}{Co hrajeme a co sledujeme}}{2025-05-29}{\Ulrik, \Luisa, \JMalcharek}{\FCrule{\#727}{Star Wars by byly mnohem lepší, kdyby v nich byly světelné jeskyňky.}{\JMalcharek}}{V netradičně tradičním formátu se vrátíme ke hrám, které hrajeme a filmům či seriálům, na které koukáme. Můžete se těšit na soulsovky, bagety, farmaření i diktátory. Tentokrát v plné sestavě Šárka, Aleš a Jakub.}{}{}{}
\FC{fight-club-728.mp3}{\FCtitle{\#728}{Co čekáme od Summer Game Festu}}{2025-06-05}{\Ulrik, \Luisa}{\FCrule{\#728}{Když odsunete GTA, vytvoříte více životního prostoru pro jiné hry.}{\Ulrik}}{Aleš se Šárkou se pohrouží do spekulací o tom, co bychom mohli vidět na letošním Summer Game Festu, který se odvysílá už tento pátek (a my budeme u toho!). Mezi rekordním počtem partnerů najdeme třeba po dlouhé době Nintendo, ale taky Valve. Očekáváme taky spoustu přesných dat vydání po tom, co se odložilo GTA VI.}{}{}{}
\FC{fight-club-729.mp3}{\FCtitle{\#729}{Co nám na Summer Game Festu udělalo radost}}{2025-06-12}{\Ulrik, \Luisa, \JMalcharek}{\FCrule{\#729}{Nevím sice čím by se bojovalo v Call of Duty Prehistoric, vím ale, že by tam byla Helena Vondráčková.}{\JMalcharek}}{Víkend naditý hrami je za námi, můžeme na chvíli vydechnout a vstřebat všechna ta oznámení, která se udála. Pobavíme se o tom, co nám udělalo největší radost, co nás naopak zklamalo a na co se nejvíce těšíme. K Šárce s Alešem se připojí Jakub, který možná přihodí i nějaké ty dojmy z hraní na Nintendu Switch 2.}{}{}{}
\FC{fight-club-730.mp3}{\FCtitle{\#730}{Death Stranding 2 přichází}}{2025-06-26}{\Ulrik, \Luisa}{\FCrule{\#730}{Když přejíždíte koaly, svět vás dislajkuje.}{\Ulrik}}{Hideo Kodžima a jeho tým nám tento týden po šesti letech nadělují pokračování simulátoru poslíčka. Death Stranding 2: On the Beach je přístupnější a svižnější, ale pořád po čertech podivná záležitost, takže hezky probereme všechny časodeště, miminka v lahvích, vyvržené věci a bosé slečny pobíhající dehtem.}{}{}{}
\FC{fight-club-731.mp3}{\FCtitle{\#731}{Do jaké hry na dovolenou?}}{2025-07-03}{\Ulrik, \Luisa, \JMalcharek}{\FCrule{\#731}{Lepší je trávit dovolenou v Death Stranding, aspoň vám Němec nesahá na mořskou okurku.}{\JMalcharek}}{Začínají prázdniny, horka, období dovolených. Tak si dáme téma odlehčeně letní – do jaké hry bychom vyrazili na dovolenou? Je ideální představa volna nošení balíků v pustém světě Death Stranding? Relaxování na pláži s Tifou a Aerith? Nebo naopak návštěva mrazivé Antarktidy v The Thing? O tom debatuje Šárka, Jakub a Aleš.}{}{}{}
\FC{fight-club-732.mp3}{\FCtitle{\#732}{Superhrdinské hry}}{2025-07-10}{\Ulrik, \Luisa, \JMalcharek}{\FCrule{\#732}{Hru se Supermanem sice nemám, ale takové cipoviny jako Suicide Squad by jim šlo.}{\JMalcharek}}{Jak nabírají na popularitě filmy s maskovanými superhrdiny, zákonitě se jejich oblíbenost promítá i do her. Kdysi spíše okrajový žánr konečně získal pozornost minstreamových studií, které dokáží tvořit skvělé pecky, ale i nepochopitelné propadáky. Ty nejlepší, nejhorší a chystané superhrdinské hry si propere Aleš, Jakub a Šárka.}{}{}{}
\FC{fight-club-733.mp3}{\FCtitle{\#733}{Kdo po Mass Effectu převezme pochodeň?}}{2025-07-17}{\Luisa, \JMalcharek}{\FCrule{\#733}{Lepší než Mass Effect od Bioware, je ječné zrno v oku.}{\JMalcharek}}{Mass Effect udělal za svou brilantní sérií space oper rozkydlou tečku v podobě Andromedy. Od té doby, aby člověk pořádně epické vesmírné sci-fi pohledal. Na obzoru se ale rýsuje slibný Exodus nebo bombastický The Expanse. Oba se spoustou styčných bodů s Mass Effectem. Podaří se jim prorazit a třeba pokořit chystaný Mass Effect 4? A co jsou vlastně nutné prvky, které by dobrá space opera měla mít? O tom se baví Šárka, Jakub a Aleš.}{}{}{}
\FC{fight-club-734.mp3}{\FCtitle{\#734}{Naše oblíbené herní soundtracky}}{2025-07-24}{\Luisa, \JMalcharek}{\FCrule{\#734}{Odpočívej v pokoji, Ozzy.}{\JMalcharek}}{Hry jsou audiovizuální médium, a v dnešním Fight Clubu se proto Šárka s Jakubem zaměří hlavně na tu audio stránku, konkrétně na soundtracky. On i letošní rok totiž přihrál hned několik hudebních podkresů, které stojí za to, pobavíme se ale i o našich nejoblíbenějších soundtracích vůbec.}{}{}{}
\FC{fight-club-735.mp3}{\FCtitle{\#735}{Jak vzniká herní recenze?}}{2025-07-31}{\Ulrik, \Luisa, \JMalcharek}{\FCrule{\#735}{Když je hra rozbitá při vydání, je to součást zážitku.}{\Luisa}}{Recenze – náš i váš nejoblíbenější novinářský žánr. Aleš, Šárka a Jakub si ve Fight Clubu popovídají o tom, jaké mají při hraní rituály, osvědčené postupy, aby hru co nejlépe zhodnotili a napsali co nejlepší text. Řeč dojde ale i na další témata jakým jsou seriály podle her, které se množí jak houby po dešti. Pobavíme si taky o tom, co hrajeme.}{}{}{}
\FC{fight-club-736.mp3}{\FCtitle{\#736}{O Kromlechu s Pavlem Makalem}}{2025-08-07}{\Luisa, \JMalcharek, \PMakal}{\FCrule{\#736}{Jak si usteleš, tak si taky Kromlech-neš.}{\Luisa}}{Do Fight Clubu se po roce vrací náš bývalý kolega Pavel Makal, kterého mezitím kariérní proudy donesly na křeslo PR zástupce v českém studiu Perun Creative. A proto se kromě uplynulého roku pobavíme hlavně o keltském akčním RPG Kromlech, na kterém Peruni pilně pracují.}{}{}{}
\FC{fight-club-737.mp3}{\FCtitle{\#737}{O úpadku slavných studií}}{2025-08-14}{\Ulrik, \Luisa}{\FCrule{\#737}{Úpadek je přiřozenou součástí přírody i herního businessu.}{\Ulrik}}{Kdysi legendy, dnes symbol korporátní hamižnosti a her bez špetky invence. Aleš a Šárka si popovídají o úpadku dříve slavných studií, která přišla o své kouzlo a duši zaprodala mamonu.}{}{}{}
\FC{fight-club-738.mp3}{\FCtitle{\#738}{O Gamescomu z Kolína}}{2025-08-21}{\Luisa, \JMalcharek}{\FCrule{\#738}{.}{\JMalcharek}}{Nastal herní svátek roku a my u něj nesmíme chybět. Šárka a Jakub se vydali na kolínské výstaviště na největší evropský veletrh Gamescom. Hrají nové hry, viděli dosud neoznámené novinky, pořádají rozhovory a stýkají se se zajímavými lidmi z herního vývoje. V neposlední řadě taky fasují spoustu herního merche, se kterým se taky tradičně pochlubí.}{}{}{}
\FC{fight-club-739.mp3}{\FCtitle{\#739}{Dozvuky Gamescomu}}{2025-08-28}{\Ulrik, \Luisa, \JMalcharek}{\FCrule{\#739}{Největším nepřítelem herního redaktora je Deutche Bahn.}{\Luisa}}{Největší videoherní veletrh Gamescom je za námi – hráli jsme spoustu her, zažili ještě více zážitků a hodně cestovali. O tom, co všechno jsme v Kolíně nad Rýnem viděli si povídá Jakub se Šárkou za zvídavých dotazů Aleše.}{}{}{}
\FC{fight-club-740.mp3}{\FCtitle{\#740}{Jaká je budoucnost Metal Gearu?}}{2025-09-04}{\Luisa, \JMalcharek}{\FCrule{\#740}{Alfa, beta, gama delta, Bi Boss pláče, Snake je meta.}{\JMalcharek}}{Metal Gear Solid Delta ve velkém stylu vrátil na výsluní slavnou špionážní sérii. Konami pomalu oživuje značku, která dlouho ležela u ledu kvůli odchodu tvůrce Hidea Kodžimy. Jaká by podle nás měla být její budoucnost? O tom se pobaví Jakub se Šárkou v dnešním Fight Clubu.}{}{}{}
\FC{fight-club-741.mp3}{\FCtitle{\#741}{O metroidvaniích: Od Metroidu po Silksong}}{2025-09-12}{\Ulrik, \Luisa, \JMalcharek}{\FCrule{\#741}{Vyšel Silksonk. Veselé Vánoce. Ho, ho, ho, Hollow Knight.}{\JMalcharek}}{Sedm let jsme čekali na vytoužený Hollow Knight: Silksong, který opět pozvedl žánr metroidvanií na výsluní. Žánr, který nese jméno po slavné hře na NES už od roku 1986. O Silksongu a všech důležitých zástupcích se pobaví Aleš, Jakub a Šárka.}{}{}{}
\FC{fight-club-742.mp3}{\FCtitle{\#742}{Jsou dlouhé early accessy ke škodě?}}{2025-09-18}{\Ulrik, \Luisa, \JMalcharek}{\FCrule{\#742}{Tento Fight Club byl v Early Accessu. 1.0 vyjde příští týden.}{\JMalcharek}}{Předběžný přístup: ano, nebo ne? A jak dlouhý? Šárka, Jakub a Aleš debatují nad hrami, které světlo světa spatří ještě před vydáním plné verze. Má čím dál populárnější praktika ve videoherním průmyslu místo? A pokud ano, jakých her by se měla týkat?}{}{}{}
\FC{fight-club-743.mp3}{\FCtitle{\#743}{Hororová renesance}}{2025-09-25}{\Luisa, \JMalcharek}{\FCrule{\#743}{Bafiky, baf, nesežrals ho?}{\JMalcharek}}{Hororová série Silent Hill zažívá recenzi. Po dlouhých letech, kdy značka ležela ladem a rodila jen nepodarky, jsme dostali fantastický remake Silent Hill 2 a neméně povedený Silent Hill f. Okrajový žánr tak slaví renesanci návratem slavných značek, třeba v podobě nového Resident Evilu Requiem, které jsou silnější než kdy dřív. O hororových hrách debatuje Šárka, Jakub a Aleš.}{}{}{}
\FC{fight-club-744.mp3}{\FCtitle{\#744}{Dawn of War a Warhammer ve hrách}}{2025-10-02}{\Ulrik, \JMalcharek}{\FCrule{\#744}{Dnes ve Fight Clubu Warhammer. Přístí týden možná Šárkhammer.}{\JMalcharek}}{Games Workshop rozdává licence na svoji wargamingovou stolní hru Warhammer 40,000 vlevo vpravo. Výsledkem je nejen spousta knih, seriálů a chystaných filmů, ale taky her. Ať už jde o návrat strategie Dawn of War IV, akční řež Space Marine nebo RPG Dark Heresy, Warhammer 40K se stává popkulturním fenoménem, který hrozí pohltit i Hvězdné války. O tématu se baví ti nejpovolanější – Aleš s Jakubem.}{}{}{}
\FC{fight-club-745.mp3}{\FCtitle{\#745}{Herní adaptace dobré i špatné}}{2025-10-09}{\Ulrik, \Luisa, \JMalcharek}{\FCrule{\#745}{Jindrův meč je schovaný na Královehradecku.}{\Luisa}}{Tvůrce Zaklínače Andrzej Sapkowski se nechal slyšet, že žádná adaptace nemůže být tak dobrá, jak původní dílo. Proto se ve Fight Clubu zaměříme na adaptace slavných knih, komiksů a filmů... co dělá adaptaci dobrou? A které se naopak vůbec nepovedly? O tom se pobaví Aleš, Šárka a Jakub.}{}{}{}
\FC{fight-club-746.mp3}{\FCtitle{\#746}{Moderní válka}}{2025-10-16}{\Ulrik, \Luisa, \JMalcharek}{\FCrule{\#746}{Call of Duty nebo Battlefield, pravá bitva se bojuje s Nintendem.}{\JMalcharek}}{S vydáním Battlefieldu 6 se nemůžeme bavit o ničem jiném než o nejnovějších explozivních obrazech moderní války. Řeč ale přijde i na vojenské střílečky obecně, a taky na to, co momentálně hrajeme. U mikrofonů se potká Šárka, Jakub a Aleš.}{}{}{}
\FC{fight-club-747.mp3}{\FCtitle{\#747}{Bloodlines 2 a jiní krvesajové}}{2025-10-23}{\Ulrik, \Luisa, \JMalcharek}{\FCrule{\#747}{Nešťouchejte do Kuby nebo se mu splaší kladiva.}{\Luisa}}{Bloodlines 2 se po letech vývojového pekla dostaly na světlo světa a stejně jako skutečný upír na slunečních paprscích shořel. O jeho pohnutém osudu a dalších upírských hrách si poví Jakub, Šárka a Aleš.}{}{}{}
\FC{fight-club-748.mp3}{\FCtitle{\#748}{Soumrak gaučové kooperace}}{2025-10-30}{\Ulrik, \Luisa}{\FCrule{\#748}{Do každé rodiny zasloužíme gaučový coop a hezké sofa. A hodiny.}{\Ulrik{} + \Luisa}}{Jen tak si s přáteli sednout před televizi a zahrát si společně nějakou hru – kdysi poměrně běžný úkaz. Dnes čím dál vzácnější, a dokonce i tradičně kooperační tituly se hraní na jedné obrazovce zbavují. Mají ještě tyto hry budoucnost? A jaké vám doporučíme, pokud chcete třeba zábavu ve větším kolektivu? Debatuje Šárka s Alešem.}{}{}{}
\FC{fight-club-749.mp3}{\FCtitle{\#749}{Handheldy a seriály}}{2025-11-06}{\Ulrik, \Luisa, \JMalcharek}{\FCrule{\#749}{Víte jak to bolí telekynetickou holí? Jestli ne, tak se dívejte na Zaklínače.}{\JMalcharek}}{V redakci teď všichni hrajeme na nějaké formě handheldu. Ať už to je MSI Claw, Legion S Go nebo Switch 2, pobavíme se o tom, jaké jsou naše zážitky s novou generací hardwaru. Aleš, Jakub a Šárka pak proberou nejnovější seriálovou produkci v čele se 4. sérií Zaklínače a To: Vítejte v Derry.}{}{}{}
\FC{fight-club-750.mp3}{\FCtitle{\#750}{Anno 117 a Řím ve hrách}}{2025-11-13}{\Ulrik, \Luisa, \JMalcharek}{\FCrule{\#750}{Radši než 117-krát Anno, 117-krát ánus.}{\JMalcharek}}{Nový díl budovatelské strategie Anno nás vezme nejdále do historie do období starověkého Říma. Při té příležitosti se Aleš, Jakub a Šráka pobaví nejen o tom, jak se nejnovější dílo hraje, ale taky na další hry, které zajímavé a populární období ztvárňují.}{}{}{}
\FC{fight-club-751.mp3}{\FCtitle{\#751}{Kdo vyhraje The Game Awards 2025?}}{2025-11-20}{\Ulrik, \Luisa}{\FCrule{\#751}{Hebká ouška, krásný ocásek, Umamusume je vaše derby.}{\Ulrik}}{Geoff Keighley v pondělí nominace na letošní The Game Awards, a tak se Šárka s Alešem pobaví o kandidátech na letošní hry roku v čele s rekordmanem Clair Obscur: Expedition 33, ale také o českém zástupci v podobě Kingdom Come: Deliverance II, které soupeří rovnou ve třech kategoriích.}{}{}{}
\FC{fight-club-752.mp3}{\FCtitle{\#752}{O budoucnosti AI ve hrách}}{2025-11-27}{\Luisa, \JMalcharek}{\FCrule{\#752}{Lepší než umělá inteligence je palcátem do palice.}{\JMalcharek}}{Čím dál více herních společností investuje do rozvoje umělé inteligence. Co to bude znamenat pro hry a herní průmysl jako takový? V čem si myslím, že je AI ve vývoji přínostná? A dokáže nahradit lidskou kreativitu? O tom se pobaví Jakub se Šárkou.}{}{}{}
\FC{fight-club-753.mp3}{\FCtitle{\#753}{Ohlédnutí za rokem s Karlem Drdou}}{2025-12-04}{\Luisa, \JMalcharek, \KDrda}{\FCrule{\#753}{Nejlepší houbový krém Karlovi vaří Dominik Cumberbearn.}{\Luisa}}{Na prozatimní ohlédnutí za letošním rokem jsme si přizvali hosta s neotřelými názory a hbitým ostrovtipem. V dnešním Fight Clubnu redakci ve složení Šárka a Jakub navštíví člověk, jehož není nutné příliš představovat – Karel Drda. Spolu se pobaví o nejpamátnějších herních zážitcích roku 2025}{}{}{}
\FC{fight-club-754.mp3}{\FCtitle{\#754}{Kočky ve hrách}}{2025-12-11}{\Ulrik, \Luisa, \JMalcharek}{\FCrule{\#754}{Uchošťoury do koček nepatří.}{\JMalcharek}}{Měli jsme možnost intenzivně týden hrát novou hru Edmunda McMillena – Mewgenics. A už teď můžeme říct, že se v ní rýsuje indie miláček příštího roku. Nebudeme se ale bavit jen o křížení koček v neskutečně hlubokém taktickém RPG, ale o hrách s kočkami obecně.}{}{}{}
\FC{fight-club-755.mp3}{\FCtitle{\#755}{Čas vánoční}}{2025-12-18}{\Ulrik, \Luisa, \JMalcharek}{\FCrule{\#755}{Dárky rozdává Ježíšek, rány rozdáváme my.}{\Ulrik}}{Rok se chýlí ke konci a s ním se na svátky odebírá i naše redakce. Jaké jsou naše plnány na Vánoce? Které herní resty hodláme dohnat? A očekáváme třeba nějaké herní dárky? O tom se baví Jakub, Aleš a Šárka.}{}{}{}
\FC{fight-club-756.mp3}{\FCtitle{\#756}{Jak řešíme herní backlog?}}{2026-01-09}{\Luisa, \JMalcharek}{\FCrule{\#756}{Lepší než žába v hrsti je chobotnice na talíři.}{\JMalcharek}}{Vánoční svátky pravidelně využíváme k odškrtávání her z rostoucího seznamu titulů, na které přes rok nebyl čas. Jak se nám to povedlo letos? A jak obecně k hernímu backlogu přistupujeme? O tom se baví v prvním díle Fight Clubu v roce 2026 Šárka a Jakub}{}{}{}
\FC{fight-club-757.mp3}{\FCtitle{\#757}{Věštíme rok 2026}}{2026-01-15}{\Ulrik, \Luisa, \JMalcharek}{\FCrule{\#757}{Ať už se těšíte v roce 2026 na cokoliv, pamatujte, že zklamání je součást života.}{\Ulrik}}{Začátek roku nahrává k tradičnímu pohledu do křišťálové koule a našim predikcím, co se v nadcházejícím roce stane. Věštíme zásadní změny v herním průmyslu? Nástup nových žánrů? A jak si myslím, že dopadne GTA VI? Do tarotu pohlédnou a interpretují výsledky Aleš, Šárka a Jakub. | (00:00) Úvod | (02:00) Predikce pro rok 2026 | (31:50) Hardwarová krize | (41:39) Co hrajeme | (46:02) Baldur's Gate III}{}{}{}
\FC{fight-club-758.mp3}{\FCtitle{\#758}{Mlžné město Silent Hill}}{2026-01-22}{\Ulrik, \Luisa, \JMalcharek}{\FCrule{\#758}{Nejhorší noční můry nečekají v Silent Hillu, ale v kinosálu.}{\Luisa}}{Šárka s Jakubem viděli Silent Hill: Noční Můry a vzhledem k momentálnímu mlžnému počasí to berou jako ideální záminku pobavit se o slavné hororové značce v širších souvislostech. Co nás čeká, zavzpomínat na zlatou éru PlayStationých klasik i načrtnout směrým, kterým bychom byli rádi, aby se série ubírala. Jakuba se Šárkou pak doplní Aleš. | (00:00) Úvod | (00:45) Silent Hill: Noční můry | (07:10) Silent Hill téma | (26:02) Nejhorší filmy | (38:35) Xbox Showcase a Fable | (44:35) Co hrajeme (Wobbly Life!!!) | (1:01:25) Dotazy}{}{}{}
\FC{fight-club-759.mp3}{\FCtitle{\#759}{Masakr v Ubisoftu}}{2026-01-29}{\Ulrik, \Luisa, \JMalcharek}{\FCrule{\#759}{Microsoft, Ubisoft, FromSoft... Proč nikod není hard?}{\JMalcharek}}{Ubisoft odstartoval nový rok masivním propouštěním, zavíráním studií a rušením projektů, které mimo jiné odskákal i remake Prince of Persia: The Sands of Time. Kromě toho se pobavíme taky o očekávaných hrách, které se ukázaly na Xbox Directu. Ve studiu se sejde Šárka, Aleš a Jakub. | (00:00) Úvod | (02:50) Škrty v Ubisoftu a zavírání studií | (23:00) Budoucnost Ubisoftu a jejich značek | (38:10) Highguard | (48:00) Co hrajeme (KCD2, Mafia, Suzerain) | (59:14) Dotazy}{}{}{}
\FC{fight-club-760.mp3}{\FCtitle{\#760}{O první desítce letošního roku a nové Yakuze}}{2026-02-12}{\Ulrik, \Luisa, \JMalcharek}{\FCrule{\#760}{Lepší než tasemnice je kočičí množenice.}{\JMalcharek}}{Únor teprve začal a máme tady první desítkovou recenzi. Proto si ve Fight Clubu probereme Mewgenics, dojde taky na novou Yakuzu Kiwami 3 anebo bizár od Sudy51 Romeo is a Dead Man. U mikrofonů se sejde Jakub, Šárka a Aleš. | (00:00) Úvod | (01:15) Mewgenics | (13:22) Baldur's Gate III vsuvka | (20:10) Yakuza Kiwami 3 | (25:20) Romeo is a Dead Man | (33:55) Ace Combat 7 | (41:01) Borderlands 4 | (53:04) Co bude na State of Play? | (57:50) Hejt na Kingdom Hearts | (1:01:54) Pokračujeme ve State of Play}{}{}{}
\FC{fight-club-761.mp3}{\FCtitle{\#761}{Co se urodilo na State of Play 2026}}{2026-02-19}{\Ulrik, \Luisa, \JMalcharek}{\FCrule{\#761}{Nejlepším vyvrcholením byl Dead or Alive 6.}{\JMalcharek}}{Sony odvisílalo první letošní State of Play, tak si shrneme, co všechno jsme na něm viděli, které novinky nás učarovaly, a co naopak chybělo. U mikrofonů se sejde Jakub, Šárka a Aleš. | (00:00) Úvod | (00:25) State of Play 2026 | (41:29) Co hrajeme (Baldur's Gate 3, KCD2, Dispatch) | (50:28) Co jsme viděli (Predátor, Rytíř Sedmi království, The Pitt)}{}{}{}
\FC{fight-club-762.mp3}{\FCtitle{\#762}{Rekviem za Resident Evil}}{2026-02-26}{\Ulrik, \Luisa, \JMalcharek}{\FCrule{\#762}{Taťka Leon vám doveze nákup v Poršáku.}{\JMalcharek}}{Na scénu se vrací legendární hororová série Resident Evil s novým dílem Requiem, a to je ideální příležitost, jak si pokecat nejen o posledním dílu, ale taky o Resi obecně a načrtnout směr, kterým kráčí. U mikrofonů se sejde Jakub, Šárka a Aleš. | (00:00) Úvod | (00:45) Kdo tady smrdí? | (02:37) Resident Evil Requiem | (26:15) Konec Phila Spencera v Microsoftu | (32:55) Daniel Vávra a konec vývoje her | (39:55) GTA za 100 dolarů? | (44:06) Co hrajeme? (Borderlands 4, Enshrouded, Sons of Sparta) | (1:12:00) Aleš potřebuje pomoc s Goat Simulatorem 3}{}{}{}
\FC{fight-club-763.mp3}{\FCtitle{\#763}{Herní degustační menu}}{2026-03-05}{\Ulrik, \Luisa, \JMalcharek}{\FCrule{\#763}{V březnu vsaďte na tři věci: pokéši, tenis, toxický penis.}{\Ulrik}}{Oddávali jsme se spoustě her. Ochutnávali dema ze Steam Next Festu, oddávali se server slamu Marathonu, máme rozpracované recenzování Star Trek: Voyager a Mario Tenis Fever a ve volných chvílích hrajeme třeba Pokémony nebo Reigns: Witcher. O tom všem si popovídá Aleš, Jakub a Šárka. | (00:00) Úvod | (01:33) Marathon | (07:14) Konec Hughguard | (09:33) Zero Parades | (15:46) Lego Voyagers | (17:30) Laysara: Summit Kingdom | (19:50) Enshrouded | (23:35) Pokémon | (33:06) Mario Tennis Fever | (39:05) Steam Hardware průzkum | (50:00) Na jaké hry se v březnu těšíme | (58:17) Závěr a dotazy}{}{}{}
\FC{fight-club-764.mp3}{\FCtitle{\#764}{My a multiplayer}}{2026-03-12}{\Luisa, \JMalcharek}{\FCrule{\#764}{Když Dlouhé Bidlo běží marathon, nestačí mu ani Rychlonožka.}{\JMalcharek}}{V dnešním Fight Clubu se pobavíme o téma, o které se otíráme jen vyjímečně – o multiplayeru. U jakých her jsme začínali? Které nás téměř stály studia? Co Alešovy staré reflexy? A jaký je vlastně Marathon? Na to se pokusí od mikrofonů odpovědět Jakub, Šárka a Aleš. | (00:00) Úvod | (01:25) Marathon | (24:45) Multiplayerové lásky (Overwatch, DotA, World of Warcraft) | (44:50) Pokémon Pokopia | (56:58) Dotazy}{}{}{}
\FC{fight-club-765.mp3}{\FCtitle{\#765}{Proč nehrajeme Crimson Desert?}}{2026-03-26}{\Ulrik, \Luisa, \JMalcharek}{\FCrule{\#765}{Chcete-li karmínový dezert, zahrajte si Zeldu.}{\JMalcharek}}{Herním světem cloumá open world akční RPG Crimson Desert. My jsme jej ale hisptersky ještě nehráli. Proč? A jakým hrám jsme místo toho dali přednost? O tom se pobaví Jakub, Aleš a Šárka. | (00:00) Úvod | (01:20) Proč nehrajeme Crimson Desert? | (10:25) Masters of Albion | (15:35) Marathon | (31:00) Star Trek: Voyager - Across the Unknown | (36:40) Fatal Frame II: Crimson Butterfly Remake | (44:35) Spasitel | (49:40) Seriálové okénko}{}{}{}
\FC{fight-club-766.mp3}{\FCtitle{\#766}{Posun času jako herní mechanika}}{2026-04-02}{\Luisa, \JMalcharek}{\FCrule{\#766}{Dočkej klasu jako husa času. Honk.}{\Luisa}}{Máme za sebou tradiční posun na letní čas, tedy ten horší z posunů. To nás vedlo k nápadu pobavit se o hrách, které nějakým zajímavým nebo unikátním způsobem s mechanikou posouvání času nebo různých časových rovin pracují. V časovém zmatení se pousí přijít včas Šárka, Jakub a Aleš. | (00:00) Úvod | (00:40) Čas jako herní mechanika | (24:45) Gray Zone Warfare | (31:03) Heroes of Science and Fiction | (35:05) Reigns: The Witcher | (42:02) Odbočka k Final Fantasy | (48:25) Situace ve Warhorse}{}{}{}
\FC{fight-club-767.mp3}{\FCtitle{\#767}{Je Starfield po letech lepší?}}{2026-04-08}{\JMalcharek, \Luisa, \Ulrik}{\FCrule{\#767}{Lepší než kajuty plné bramor ve Starfieldu jsou brambory s ručičkami}{\Ulrik}}{S vydáním na PlayStation se vesmírný magnum opus od Bethesdy dočkal taky nového obsahu a updatů. Šárka se ve Starfieldu vydala mezi hvězdy v recenzi, a hledala tam odpověď na otázku nejdůležitější: Je Starfiled po skoro třech letech konečně už dobrá hra? Zpovídat ji bude Jakub s Alešem.}{}{}{}
\FC{fight-club-768.mp3}{\FCtitle{\#768}{O sci-fi akci Pragmata}}{2026-04-15}{\JMalcharek, \Luisa, \Ulrik}{\FCrule{\#768}{Ještě, že se to jmenuje Pragmata a ne pregnata.}{\JMalcharek}}{Capcom se nebojí experimentovat a tvořit nové značky, jak je vidno z aktuální sci-fi akce Pragmata. Jak se originální vesmírná arkáda hraje? O tom si popovídá Jakub, Aleš a Šárka.}{}{}{}
\FC{fight-club-769.mp3}{\FCtitle{\#769}{Novinky na Games}}{2026-04-22}{\JMalcharek, \Luisa, \Ulrik}{\FCrule{\#768}{Nezapomeňte jaké je pravidlo klubu rváčů.}{\JMalcharek}}{S měnícími se vodami internetu se měníme i my. Games procházejí evolucí a o novinkách, které chystáme, si popovídáme v dnešním Fight Clubu. Co jsme si připravili? Jak to bude v praxi fungovat? Kdy, kde, kdo, proč a za kolik? Řekne Aleš, Šárka a Jakub.}{}{}{}



\newpage
\section*{Social club}

{\color{white}
Social club je jeden z dvojice nových formátů vzniklých oddělením od původního 
Fight Clubu. Bohužel se nedočkal příliš velkého počtu dílů.
}

\hspace{0pt}  % very 1st minipage more to the left problem hack
\SC{social-club-1.mp3}{\SCtitle{\#1}{Volejte řediteli}}{2014-04-11}{\Sleepless, \Luke, \Blixa}{Odpovědi na vaše dotazy z chatu a emailu. | {\em (2:15) Goat simulator a Psiculus | (7:00) Stárnutí v herním průmyslu | (15:30) Wildstar | (16:40) Distribuce EA her od ABC data | (19:15) Kde hledat spolupráci na vývoji her | (21:40) Nejvíce peněz vydaných za hru | (24:40) První dohraná hra a připomínání soutěže | (26:35) Dostupné scénáře Shadowrunu a další série | (29:10) Musí programátoři umět matematiku nebo ne?}}{0:48:17}{21033268}{}
\SC{social-club-2.mp3}{\SCtitle{\#2}{Bacha na vyhazovače Bacha!}}{2014-04-18}{\Sleepless, \Luke, \Blixa}{Odpovědi na vaše dotazy z mailu a chatu. | {\em (1:00) Shovívavost a nahota | (2:50) Servisní záležitosti | (9:30) Klenoty mezi japonskými RPG | (12:00) Prodejní čísla od Cinemaxu | (12:40) Plány na E3 | (14:20) Dotaz na stav The Keep pro 3DS | (15:05) Soukromí Chrise Sawyera a Twitery dalších vývojářů | (23:00) Za jakých podmínek by si Sleepless nechal oholit hlavu | (24:40) Návrat k původnímu formátu | (26:20) Kdy zas uvidíme Lukáše? Kreativita u vystřihování soutěžních otázek | (27:30) O vybití herního hněvu | (29:00) Bach vyhazovačem | (30:05) S.T.A.L.K.E.R Lost Alpha | (30:40) Indie hra k Levelu | (33:05) Dotaz gamesplay s Janchorem na Hotline Miami 2\,\dots | (35:00)Kam se budou zvát klasičtí hosté | (36:20) Miloš Bohoněk by chtěl zítra udělat jiný GamesPlay | (36:45) Na čem se podílí Diesel\,\dots | (38:00) Hraní více hráčů | (40:35) Draken se nemůže vyjádřit | (41:45) Odhad ceny covermountu k časopisu (Fallout 1) | (42:40) Konec dotazů }}{0:47:41}{21451876}{}
\SC{social-club-3.mp3}{\SCtitle{\#3}{Grafická katovna}}{2014-04-25}{\Sleepless, \Luke, \Janchor}{Odpovědi na vaše dotazy z mailu a chatu. | {\em (0:50) Vyignorované servisní okénko a oznámená uzávěrka LEVELu (MiK spamuje) | ()  | }}{1:03:05}{28504492}{}
\SC{social-club-4.mp3}{\SCtitle{\#4}{Blázen, pesimista, hister a diplomat}}{2014-05-10}{\Sleepless, \Luke, \Janchor, \Blixa}{Odpovědi na vaše dotazy z mailu a chatu.}{0:59:52}{26794800}{}
\SC{social-club-5.mp3}{\SCtitle{\#5}{Slovenský koutek}}{2014-05-17}{\Sleepless, \Luke, \Janchor}{Odpovědi na vaše dotazy z mailu a chatu.}{0:57:47}{23576272}{}
\SC{social-club-6.mp3}{\SCtitle{\#6}{Dva v tom}}{2014-05-23}{\Sleepless, \Luke}{Odpovědi na vaše dotazy z mailu a chatu.}{0:36:38}{14935012}{}
\SC{social-club-7.mp3}{\SCtitle{\#7}{Club Synth}}{2014-05-30}{\Sleepless, \Luke, \Janchor, \Blixa}{Odpovědi na vaše dotazy z mailu a chatu.}{0:55:52}{23193988}{}




\newpage
\section*{Tea club}

{\color{white}
Tea club je jeden z formátů, který byl oddělen od původního Fight Clubu. Jedná
se o relativně krátký podcast a vidcast, jenž je založen na rozhovorech 
s českými a slovenskými herními tvůrci. Jedná se o občasník.
}

\hspace{0pt}  % very 1st minipage more to the left problem hack
\TC{tea-club-01.mp3}{\TCtitle{\#1}{Fiola Soft}}{2014-04-10}{\Luke, \Fiola}{S Filipem Kraucherem s Fiola Softu si povídáme hlavně o čerstvě oznámené plošinovce Blackhole.}{0:29:26}{14964584}{}
\TC{tea-club-02.mp3}{\TCtitle{\#2}{Hyperbolic Magnetism}}{2014-05-08}{\Luke, \JIlavsky, \VHrincar}{Kdyby existovala statistika o počtu vyprodukovaných her za rok, programátorské duo Hyperbolic Magnetism by se s přehledem umístilo na prvních místech. Vladimír „Loki“ Hrinčár a Honza „Split“ Ilavský programují jak o život a pravidelně se hlásí na vývojářské maratony Global Game Jam a Ludum Dare. S jejich tituly se nejčastěji můžete setkat na mobilech, ale i na PC, které ve své hyperproduktivitě nezanedbávají. Příklady z poslední doby mohou být třeba Heterochroma z Global Game Jam nebo Honzova samostatná Temple of Ra z Ludum Dare.}{0:34:15}{19255532}{}
\TC{tea-club-03.mp3}{\TCtitle{\#3}{}}{2014-05-24}{\Luke, \JPavlasek}{O strategické tahovce Cubesis, ve které ovládáte ekosystém planety, jste možná neslyšeli. Nic si z toho nedělejte, my také ne. Proto jsme do dalšího Tea Clubu pozvali jednoho ze spoluautorů, aby nám o hře, jejíž vývoj trval sedm let, něco pověděl.}{0:29:23}{17355276}{}
\TC{tea-club-04.mp3}{\TCtitle{\#4}{Cinemax}}{2014-06-03}{\Luke, \LMacura, \MLinda}{Lukáš Macura a Martin Linda vyprávějí o vzniku konceptu a vymýšlení hádanek a soubojů pro krokovací dungeon The Keep na Nintendo 3DS a PC. Kromě toho jsme povídají o situaci na 3DS trhu a hledání toho správného publika. Probrali jsme také, jak se testují hry nejen v Cinemaxu, ale i v široširém světě.}{0:22:00}{11567484}{}
\TC{tea-club-05.mp3}{\TCtitle{\#5}{Future Factory}}{2014-06-07}{\Luke, \VGersl}{Po zkušenostech s prací v herním vývoji založil Vláďa Geršl studio Fun 2 Robots a připravuje hru, která se snad bude líbit všem příznivcům taktických akcí a literatury. Že to nejde dohromady? Future Factory obsahuje svižnou akci a příběh inspirovaný Čapkovým R.U.R. }{0:24:43}{12749724}{}
\TC{tea-club-06.mp3}{\TCtitle{\#6}{Novus Inceptio}}{2014-06-30}{\Luke, \MMastny}{MMORPG Novus Inceptio s procedurálně generovanou krajinou nenabídne tisíce naskriptovaných úkolů a předem připravená města a vesnice osázená vesničany. Herní obsah budou tvořit samotní hráči svojí aktivitou, která začíná přežívání v divočině a končí třeba stavbou města, volením starosty a pronajímáním pozemku.}{0:26:43}{13034020}{}
\TC{tea-club-07.mp3}{\TCtitle{\#7}{}}{2014-07-03}{\Luke, \MJezek}{Ve studiu 2K prý měli plyšového ježka zavřeného v lahvi. Když se jí Matouš Ježek pokusil otevřít, rychle ho zastavili s tím, že ježek se nikdy nesmí dostat ven. Jak to dopadlo s plyšákem nevíme, ale cestu Matouše ze studia 2K můžeme dobře sledovat. Po odchodu založil v Brně studio Trickster Arts a loni vydal pozitivně přijatou mobilní hru Hero of Many, která vyšla i pro OUYA a Amazon Fire a v současnosti se připravuje i pro Windows.}{0:24:43}{11794000}{}
\TC{tea-club-08.mp3}{\TCtitle{\#8}{Factorio}}{2014-08-03}{\Luke, \MKovarik, \TKozelek}{Michal Kovařík a Tomáš Kozelek nám v půl hodinovém rozhovoru popsali celou cestu jejich projektu Factorio – od prvních krůčků, zoufalství a naivity až k funkčnímu „garážovému“ týmu, který jde cílevědomě za svou vizí. Kluci mluví i o vlídném přijetí komunitou anebo o běhání jako očistném procesu v dobách, kdy jde vše do háje.}{0:30:57}{15993392}{}
\TC{tea-club-09.mp3}{\TCtitle{\#9}{Překlady her do češtiny}}{2014-10-25}{\Luke,\,\dots}{V Tea Clubu jsme se tentokrát zaměřili na slasti a strasti amatérského překládání her, konkrétně Divinity: Original Sin. Michal a Vítek, kteří šéfují týmu překladatelů a vedou komunitu na sociálních sítí, mluví o náročnosti překladu ukecaného RPGčka.}{0:16:48}{8066480}{}
\TC{tea-club-10.mp3}{\TCtitle{\#10}{SIMT MHD}}{2014-12-09}{\Luke, \TFaina}{Každý den hráči najezdí v autobusech a trolejbusech v SIMT téměř dva tisíce kilometrů, nabírají a vykládají cestující a snaží se dodržovat jízdní řád. Podle Tomáše Fainy, který nám o hře povídal v Tea Clubu, jim to docela jde a jsou kultivovaní. No bus trolling in da štrýt! Uznejte sami – v kolika hrách vám cestující nenastoupí, když máte zasviněný autobus? Takže pěkně hadru a Savo a drhněte!}{0:19:20}{9485636}{}
\TC{tea-club-11.mp3}{\TCtitle{\#11}{Wayward Terran Frontier}}{2014-12-29}{\Luke, \JOrszulika}{Wayward Terran Frontier je projektem Američana George Hultgrena a Ostravaka Honzy Orszulika, který se podílí na vzhledu hry a herních pravidlech. Na první pohled jde sice o jednoduchou 2D hru ve vesmíru, ale jakmile se podíváte zblízka, nestačíte se divit: rozlehlý vesmír, design interiéru vlastní lodi a základny, abordace nepřátel, přestřelky v koridorech, destrukce jednotlivých částí nepřátelského vybavení a spousta dalšího.}{0:16:26}{7809584}{}
\TC{tea-club-12.mp3}{\TCtitle{\#12}{Medieval Engineers}}{2015-01-25}{\Luke, \MRosa}{Marek Rosa ze studia Keen Software House (Miner Wars 2081, Space Engineers) vypráví o složitostech při vytváření středověké inženýrské simulace a také zmiňuje spoustu možností, které hru možná čekají.}{0:28:17}{13337132}{}
\TC{tea-club-13.mp3}{\TCtitle{\#13}{The Great Wobo Escape}}{2015-04-11}{\Luke, \RKajan, \MWilczak}{Pavel si tentokrát pod povídal s šéfem studia Gamifi.cc Rudou Kajanem a s grafikem Martinem Wilczakem, jak jejich hru The Great Wobo Escape inspirovaly klasiky jako Abe's Odyssee, Another World nebo film Wall-E. Dozvíte se ale samozřejmě mnohem víc.}{0:18:35}{8060756}{}
\TC{tea-club-14.mp3}{\TCtitle{\#14}{Capricorn 7}}{2015-04-26}{\Luke, \DKos}{Možná by si Future Cop: LAPD zasloužila zařadit mezi retro harampádí. Vyšla před necelými dvaceti lety na PlayStation a dnes vypadá stejně jako všechny rané 3D hry: staře a nekoukatelně. Kdo si na ní pamatuje, ten ale nehodnotí grafiku, nýbrž vynikající městskou akci s mechwarriorem, kterou Electronic Arts dokázala přinést jak v kampani, tak v multiplayeru. David Kos šel ještě dál. Podle vzoru Future Cop vystavěl v Unity vlastní Capricorn 7 a doufá, že uspěje na IndieGoGo.}{0:23:53}{10284908}{}
\TC{tea-club-15.mp3}{\TCtitle{\#15}{Rememoried}}{2015-07-12}{\Luke, \VKoudelka}{Jestli jste před viděli trailer na českou hru Rememoired, možná jste se podivili. Co to má být? To je hra? A o čem? Prý se v ní bude pracovat se sny, tvrdila tisková zpráva. Z toho jsme nikdo nebyli moudřejší. Proto v novém Tea Clubu přijal pohovku pro hosta Vláďa Koudelka, který jako samojediný autor za Rememoired stojí.}{0:22:35}{9876452}{}
\TC{tea-club-16.mp3}{\TCtitle{\#16}{Jindřich Skeldal}}{2015-11-17}{\Luke, \Skeldal}{Brány Skeldalu: 7 mágů sice čerpá jednak z klasických dungeonů jako Wizardry, Legend of Grimrock nebo Eye of the Beholder, ale mytologie hry je ovlivněna další Jindřichovou tvorbou. Kromě her jsou to i knihy, z nichž více jak tisíci stránkový román Brány Skeldalu čeká na vydání, nebo legendární westerny jako Sedm statečných. Z naší půl hodinky se toho zkrátka dozvíte hodně.}{0:34:14}{14997884}{}
\TC{tea-club-17.mp3}{\TCtitle{\#17}{Logistics}}{2016-01-23}{\Luke, \JZeleny}{Hrajete Open Transport Tycoon Deluxe a je vám líto, že nevychází moderní hry, které by na tradici dopravních tycoonů navazovaly? Neplačte! Logistics od českého programátora Honzy Zeleného představuje lákavou alternativu.}{0:34:56}{18110478}{}
\TC{tea-club-18.mp3}{\TCtitle{\#18}{Planet Nomads}}{2016-02-28}{\Luke, \DMaslovsky, \PRuzicka}{Kombinace stavění domů, ale i vozidel, prozkoumávání procedurálně generované planety a přežívání mezi agresivními zvířaty přinesla ostravskému studiu Craneballs téměř sto tisíc liber na Kickstarteru. O obsahu, a také o komplikacích vývoje, přišli do našeho studia popovídat Dan Maslovský a Petr Růžička, pro které je Planet Nomads jednak úspěšným projektem, a pak premiérou na PC.}{0:31:03}{13062142}{}
\TC{tea-club-19.mp3}{\TCtitle{\#19}{The Fall of the Fly}}{2016-04-24}{\Luke, \JValenta}{I když se snažím poctivě a metodicky kopírovat, tak mi z toho vzniká něco bláznivého, říká v jednu chvíli Honza Valenta, autor projektu Fall of the Fly a zároveň host dalšího Tea Clubu. Jeho hra je na první pohled skutečně... jiná. Někdo by možná řekl šílená a nesmyslná, ale i to je záměr. Přeci si nemyslíte, že uprostřed Dutozemě by měl být normální svět obydlený běžnými smrtelníky?}{0:20:00}{9596480}{}
\TC{tea-club-20.mp3}{\TCtitle{\#20}{Vaporum}}{2016-05-22}{\Luke, \TRepta, \MZajacik}{Je tomu už více než rok, co jsme psali o slovenském krokovacím dungeonu Vaporum. Proniknete v něm do vysoké věže s vnitřnostmi z parních strojů a arteriemi z párou natlakovaného potrubí. K čemu slouží? Odpověď znají mechanické příšery postávající na každém rohu, ale ty vám jí neřeknou. Jste pro ně vetřelec, lidský ztroskotanec, který se shodou okolností stává hrdinou. Kromě příběhu klade Vaporum důraz na neobvyklé vylepšování hrdiny a souboje, které se mají vymykat (nejen) moderním krokovacím klasikám jako Legend of Grimrock. Takže: Držte si gombíky, jdeme krokovat!}{0:39:37}{18375116}{}
\TC{tea-club-21.mp3}{\TCtitle{\#21}{The Cycle}}{2016-06-01}{\Luke, \MBudik}{Přibližně čtyři roky pracuje Marek Budík na střílečce The Cycle, která vypadá jako by vypadla z geometrova šíleného snu. Procedurálně generované úrovně jsou tu osazené nepřáteli, kteří čekají na příchod hrdiny s velkým mečem, aby ho zabili – nebo byli zabiti. Rychlé reakce jsou pro přežití výhodou, hrát se však dá i pomaleji a s rozmyslem. Na The Cycle kromě svižné akce potěší i grafická stylizace, která odkazuje k minimalismu a geometrickému vidění světa. Marek v herním návrhu odráží i to, co se naučil během studií průmyslového designu, a nebojí se o tom povídat.}{0:20:24}{9388904}{}
\TC{tea-club-22.mp3}{\TCtitle{\#22}{Clayburger o Hopply}}{2016-06-11}{\Luke, \PKnut, \MPavelka}{Pamětníci budou hned doma: Vždyť jde o Perestrojku, fenomén učeben informatiky v první polovině devadesátých let! Peter a Matěj inspiraci nezastírají. Jde jim především o to udělat perfektní skákací rychlík, který perfektně funguje a má výraznou výtvarnou stránku. Jak se to dělá? I to uvidíte na vlastní oči během našeho nového Tea Clubu o hře Hopply.}{0:21:45}{9913208}{}
\TC{}{\TCtitle{\#23}{Battle for Nethervein}}{2016-08-19}{\Luke, \JVanecek}{Honza Vaněček a jeho studio From the Void rozhodně nejsou začátečníci. Honza je české deskovkové komunitě dobře znám svojí karetní hrou Shuffle Heroes a předchozí McJohny's. V našem rozhovoru se kromě Battle for Nethervein rozpovídal o své minulosti a o deskovkové komunitě v Česku, i o náročnosti při prosazování českých deskovek v zahraničí či o komplikacících distribuce anebo o pořádání deskovkových game jamů.}{0:20:24}{}{}
\TC{tea-club-24.mp3}{\TCtitle{\#24}{Morning Men}}{2016-10-03}{\Luke, \LNizky}{Když se podíváte na obrázky ze slovenské adventury Morning Men, hodí vám do očí písek z pouštní Duny a na solar udeří direktem od M\oe{}bia. A není to náhoda. K jeho práci se výtvarníci ze studia Pixel Federation hrdě hlásí -- a ne pouze k němu\,\dots}{0:19:18}{9888260}{}
\TC{tea-club-24.mp3}{\TCtitle{\#24}{Morning Men}}{2016-10-03}{\Luke, \LNizky}{Když se podíváte na obrázky ze slovenské adventury Morning Men, hodí vám do očí písek z pouštní Duny a na solar udeří direktem od M\oe{}bia. A není to náhoda. K jeho práci se výtvarníci ze studia Pixel Federation hrdě hlásí -- a ne pouze k němu\,\dots}{0:19:18}{9888260}{}
\TC{tea-club-25.mp3}{\TCtitle{\#25}{Dimension Brothers}}{2016-11-20}{\Luke, \MBerlinger}{Jakou hru mám hrát s přítelkyní nebo rodiči? Taková otázka není až tak vzácná. Položil si jí i Michal Berlinger, a protože je programátor, sám si odpověděl hrou Dimension Brothers. V lokální kooperativní skákačce si dva hráči pomáhají překonat neviditelné překážky, které ale musí nejprve sami objevit. Kvůli červeným a modrým brýlím totiž vidí každý z nich herní svět jinak. Michal v Tea Clubu nepovídá jen o prototypování, zastaví se i u problematiky přístupnosti lokálních multiplayerů a zabrousí na půdu gamejamů.}{0:21:51}{10107644}{}
\TC{tea-club-26.mp3}{\TCtitle{\#26}{Vikings – Wolves of Midgard}}{2016-11-26}{\Luke, \PNagy}{Když vám vypálí vesnici a máte štěstí, tak utečete. Když máte smůlu, tak ne. A když mají smůlu oni, tak vypálili vesnici klanu Ulfungů a je úplně jedno, že jsou to obři a mají velké kyje a temnou magii. Takový je námět akční RPG od slovenského studia Games Farm, jehož šéfa Petera Nagyho jsme zastihli s kamerou a mikrofonem během pražské Game Developers Session. Peter nám toho o běsnících vikinzích vyprávěl skoro celou ságu!}{0:18:17}{10601168}{}
\TC{tea-club-27.mp3}{\TCtitle{\#27}{Children of the Galaxy}}{2017-04-14}{\Luke, \FDusek}{Již čtyři roky vyvíjí Filip Dušek 4X strategii Children of the Galaxy, ve které se obrací na velikány žánru jako je Stallaris, Galactic Civilization nebo Endless Space. Zároveň přidává řadu systémů, které objevování, bitvy a kolonizování osvěžují. Jeden příklad za všechny: Tahové souboje s jednotkami, kterým se v průběhu tažení vylepšují statistiky. Pro Filipa jde sice o první vlastní hru, ale zkušenosti nasbíral již ve společnosti 2K a studiem žánrových spřízněnců. V našem Tea Clubu se zmiňuje i o tom, proč je výhodné dělat programátorské video streamy.}{0:30:31}{15628278}{}
\TC{tea-club-28.mp3}{\TCtitle{\#28}{LostHero}}{2018-07-09}{\Luke, }{Povídání o vývoji Lost Hero}{0:24:55}{10768712}{}
% \TC{tea-club-29.mp3}{\TCtitle{\#29}{}}{2018--}{\Luke, }{}{}{}{}



\newpage
\section*{Hardware club}

{\color{white}
Hardware club je novým formátem nastartovaným v průběhu roku 2018.
}

\hspace{0pt}  % very 1st minipage more to the left problem hack
\FC{hardware-club-001.mp3}{\HCtitle{\#1}{Tak nám ten rok pěkně začal}}{2018-02-01}{\Sleepless, \Janchor, \JSvak, \JZubec}{}{(02:58) Spectre/Meltdown | (21:16) Nedostatek grafických karet na trhu | (35:43) Notebook s GTX 1080 Max-Q (Asus Zephyrus) | (54:42) CES 2018 a GeForce NOW}{1:03:52}{41250060}{}
\FC{hardware-club-002.mp3}{\HCtitle{\#2}{Procesory}}{2018-02-12}{\Sleepless, \Janchor, \JSvak, \JZubec}{}{(02:04) Jak jsou na tom dnes procesory? (Ryzen/Coffee Lake) | (16:37) Přetaktování CPU | (20:00) Upgrade, nebo počkat? | (30:24) Operační paměti}{0:55:03}{36100368}{}
\FC{hardware-club-003.mp3}{\HCtitle{\#3}{AMD Raven Ridge}}{2018-02-22}{\Sleepless, \Janchor, \JSvak, \JZubec}{}{(03:02) Přestupy Intel/AMD | (10:58) Trh GPU a AMD | (20:11( APU AMD (Raven Ridge) | (45:09) Cenové srovnání: APU vs. CPU+GPU}{1:00:16}{40688460}{}
\FC{hardware-club-004.mp3}{\HCtitle{\#4}{Přetaktování}}{2018-03-09}{\Sleepless, \Janchor, \JSvak, \JZubec}{}{(03:25) Přetaktování | (09:56) Odbočka – frekvence vs. jádra | (15:22) Na čem jde přetaktovat a jak to funguje?  | (24:44) Náklady na OC I. | (32:50) Odbočka: 4K + VR | (46:43) Náklady na OC II. + podvoltování | (56:48) Dotazy}{1:11:56}{29674152}{}
% \FC{hardware-club-005.mp3}{\HCtitle{\#5}{Výrobní procesy, taktování AMD, delid}}{2018-03-22}{\Sleepless, \Janchor, \JSvak, \JZubec}{}{(01:17) Konec křemíku: 7nm výrobní proces a dále | (14:08) VSUVKA: Kvantové počítače | (20:37) Výrobní procesy – pokračování, pozice Intelu | (33:46) Taktování AMD, výhled s novou generací | (46:15) Delid procesoru | (56:52) Dotazy}{}{}{}
\FC{hardware-club-006.mp3}{\HCtitle{\#6}{Ray-tracing, notebook HP Omen X}}{2018-04-05}{\Sleepless, \Janchor, \JSvak, \JZubec}{}{Ray-tracing – v herní branži periodicky se objevující termín dostal novou krev na letošním GDC. Nvidia a další technologičtí giganti oznámili, že se na tento „svatý grál“ vykreslování hodlají v blízké budoucnosti zaměřit. Společně si toto téma probereme. V druhé půli se koukneme na momentálně recenzovaný notebook HP Omen X s GTX 1080. Jak vypadá v porovnání s dříve testovaným Asus ROG Zephyrus?}{1:11:03}{46254580}{}
\FC{hardware-club-007.mp3}{\HCtitle{\#7}{Ryzen 2700X \& 2600X, osmijádro od Intelu a pokažené Atomy}}{2018-04-19}{\Janchor, \JSvak, \JZubec}{}{(01:11) Druhá generace procesorů Ryzen | (14:31) Nově objevené chyby platformy Intel | (27:16) Klepy o osmijádru Intel | (31:25) Alternativa z diskuze: starší serverový Xeon? | (38:19) Ryzen 7 2700X | (48:13) Dotazy}{1:10:40}{47498688}{}
\FC{hardware-club-008.mp3}{\HCtitle{\#8}{Grafické karty}}{2018-05-03}{\Sleepless, \Janchor, \JSvak, \JZubec}{}{Trh s GPU regeneruje - Kdy jít do koupě? - Kvantum dotazů}{1:13:23}{48592152}{}
\FC{hardware-club-009.mp3}{\HCtitle{\#9}{Kdy vyjdou nové herní konzole a další novinky}}{2018-05-18}{\Sleepless, \Janchor, \JSvak, \JZubec}{}{(01:40) Novinky | (03:14) Nvidia Turing již v červenci? | (14:50) AMD Vega 20 a 7nm proces | (20:05) Co by měl vlastně přinést 7nm proces? | (26:26) Herní konzole – kdy by mohly reálně vyjít? | (50:18) Dotazy}{1:05:34}{28564504}{}
\FC{hardware-club-010.mp3}{\HCtitle{\#10}{Monitory}}{2018-05-31}{\Sleepless, \Janchor, \JSvak, \JZubec}{}{(00:57) Téma: monitory | (02:58) CRT? | (03:49) Hlas mimo kameru | (04:33) LCD TV vs. LED TV | (05:49) Typy LCD panelů | (06:43) Odezva | (12:28) Jindrův favorit mezi monitory | (19:08) IPS vs. TN | (22:54) VA panel, pozorovací úhly | (25:37) Srovnání panelů - pokračování | (34:09) Obnovovací frekvence | (43:29) GSync/FreeSync | (50:10) Dotaz: levné monitory | (54:52) Dotaz: OLED/MicroLED}{1:04:22}{24719436}{}
\FC{hardware-club-011.mp3}{\HCtitle{\#11}{Monitory - HDR, ultrawide a naše prognózy}}{2018-06-21}{\Janchor, \JSvak, \JZubec}{}{(01:33) Monitory podruhé | (03:25) Mini LED | (10:41) Co je to HDR? | (11:30) HDR: foto vs. video | (14:03) Jak HDR funguje? | (17:17) Monitor 4:3 HDR? | (21:04) Dotaz: HDR settings ve hrách | (22:46) Podmínky HDR – jas, gamut, barevná hloubka, HDR formáty | (26:28) Režie HDR obrazu | (29:35) Podpora HDR | (38:50) Ultrawide | (47:51) 4K HDR 144 Hz G-Sync displeje a jejich problémy | (56:25) Dotazy | (1:09:52) Jak budou vypadat monitory za 10 let? | (1:19:08) Jindrova sestava | (1:20:40) [soupis sestavy]}{1:24:59}{32265144}{}
\FC{hardware-club-012.mp3}{\HCtitle{\#12}{Navrhujeme herní PC do 15\,000, 30\,000 až 50\,000+ Kč}}{2018-07-03}{\Sleepless, \Janchor, \JSvak, \JZubec}{}{(03:19) Herní PC za 15\,000\,Kč | (23:52) Sestava za 30\,000\,Kč | (43:52) High-end kolem 50\,000\,Kč | (47:54) Herní dělo pro přetaktování | (53:36) Workstation za 100\,000\,Kč}{1:12:47}{31475068}{}
\FC{hardware-club-013.mp3}{\HCtitle{\#13}{Virtuální realita s Markem Ďurišem}}{2018-07-19}{\Janchor, \JSvak, \MDuris}{}{(05:02) Představení VR headsetů a naše zkušenosti | (11:28) Hardwarové požadavky na VR (PC) | (17:12) Sledování pozice | (19:59) Rozdíly v headsetech? | (22:18) Prodeje VR | (24:18) Kam s kabely? | (27:57) Bezdrátové VR | (29:58) Nevýhody a kompromisy VR | (43:25) VR a obsah pro dospělé | (45:18) VR pro mobily | (47:53) Killer hra pro VR? | (51:25) Dotazy | (1:01:15) Budoucnost VR}{1:08:47}{29154306}{}
\FC{hardware-club-014.mp3}{\HCtitle{\#14}{Chlazení}}{2018-08-02}{\Sleepless, \JSvak, \JZubec}{}{(02:09) Airflow | (06:58) Konstrukce PC skříně | (11:46) Vzduchové chlazení | (18:18) Chladič CPU + jak fungují heatpipes? | (26:14) Chlazení GPU | (35:09) Vodní chlazení | (53:32) Chlazení olejem nebo pasivně? | (58:15) Tipy na chladiče | (1:20:36) Dotazy}{1:28:58}{39573786}{}
\FC{hardware-club-015.mp3}{\HCtitle{\#15}{Herní notebooky}}{2018-08-16}{\Sleepless, \Janchor, \JSvak, \JZubec}{}{(04:24) Proč si lidé kupují herní notebooky? | (10:58) Nevýhody herních notebooků | (14:25) Externí GPU? | (16:08) Co už je herní notebook? A ceny? | (25:55) Nepřehlednost nabídky | (33:57) Poměr cena/výkon desktop vs. notebook | (38:20) Undervolting | (46:12) Modulární konfigurace | (51:15) Dotazy}{1:11:21}{30678093}{}
\FC{hardware-club-015-spec.mp3}{\HCtitle{SPECIÁL}{Nvidia Geforce RTX (Gamescom 2018)}}{2018-08-23}{\Janchor, \JSvak, \JZubec}{}{(03:04) RTX 2080Ti, 2080, 2070 | (04:55) Ray tracing a Turing | (16:18) Ceny | (20:52) Úvaha nad strategií Nvidia | (26:59) Prosadí se ray tracing? | (52:05) Parametry čipů Turing}{1:06:49}{29474397}{}
\FC{hardware-club-016.mp3}{\HCtitle{\#16}{Zdražování}}{2018-09-06}{\Janchor, \JSvak, \JZubec}{}{ (02:57) Drahé novinky ze stáje Nvidia | (15:54) Intel a nebezpečí nedostatku výrobních kapacit | (24:55) i9-9900K / i7-9700K | (28:54) Global Foundries a zastavený vývoj 7nm výrobního procesu | (32:37) BFGD | (38:17) Příležitost pro AMD? | (41:24) Stav na trhu s RAM | (45:09) Nová konzole v Číně | (47:20) AMD GPU na 7nm | (48:50) Prolomení bobříka mlčení | (50:29) Dotazy}{1:03:01}{27692316}{}
\FC{hardware-club-017.mp3}{\HCtitle{\#17}{GeForce RTX – shrnutí recenzí a naše dojmy}}{2018-09-20}{\Janchor, \JSvak, \JZubec}{}{Kde je ray tracing? | Konečně grafika pro 4K / 60 FPS? | GeForce RTX 2080 od Gigabyte}{0:54:31}{23630013}{}
\FC{hardware-club-018.mp3}{\HCtitle{\#18}{SSD disky}}{2018-10-04}{\Janchor, \JSvak, \JZubec}{}{(02:50) Co je to SSD? | (09:34) Životnost | (12:49) Rozhraní | (23:25) Paměťové čipy | (46:48) Spolehlivost | (55:00) Kapacita vs. rychlost}{1:14:31}{32863401}{}
\FC{hardware-club-019.mp3}{\HCtitle{\#19}{Nvidia RTX 2070, RAID}}{2018-10-18}{\Janchor, \JSvak, \JZubec}{}{(02:32) Doháníme dotazy | (19:06) RTX 2070: výkon, cena | (34:49) RAID | (48:19) RAID jako záloha? | (50:27) RAID 5, 6, 10 | (55:02) Jiné řešení redundance}{1:06:22}{29451888}{}
\FC{hardware-club-020.mp3}{\HCtitle{\#20}{Intel Core i9-9900K a i7-9700K}}{2018-11-01}{\Sleepless, \JSvak, \JZubec}{}{(01:52) i9-9900K a i7-9700K - parametry a výkon | (19:17) Provozní vlastnosti nových osmijader | (29:33) TDP a vliv základních desek | (39:18) Vyplatí se nové procesory? | (55:50) Opravy Meltdown/Spectre | (1:03:03) Dotazy}{1:23:26}{37838925}{}
\FC{hardware-club-021.mp3}{\HCtitle{\#21}{Externí GPU}}{2018-11-15}{\Janchor, \JSvak, \JZubec}{}{(02:36) Co je to externí grafika? | (06:48) Jak vyřešit připojení externí GPU? | (14:10) Existuje prostor pro eGPU? Výhody a nevýhody | (38:46) Novinka: RTX v Battlefield V | (59:40) Nadějný výhled AMD | (1:00:40) Dotazy | (1:07:21) Upoutávka na příště}{1:08:39}{31751073}{}
\FC{hardware-club-022.mp3}{\HCtitle{\#22}{Navrhujeme herní počítače pod stromeček}}{2018-11-29}{\Janchor, \JSvak, \JZubec}{}{(04:14) Sestava \uv{z nouze ctnost} (8\,450\,Kč) | (15:38) Sestava \uv{i Skrblík by se zaradoval} (16\,600--17\,500\,Kč) | (28:10) Sestava MID aka střední třída (23\,700--30\,900\,Kč) | (42:36) Sestava HIGH aka nejvyšší třída včetně OC (56\,600--66\,600\,Kč)}{1:05:02}{30793392}{}
\FC{hardware-club-023.mp3}{\HCtitle{\#23}{AMD a Intel v roce 2019, konkurence ARMu a další novinky}}{2018-12-13}{\Janchor, \JSvak, \JZubec}{}{(16:49) Drby o procesorech AMD Ryzen 3000 | (38:49) Nové Radeony příští rok? | (47:47) Intel a nová 10nm jádra Sunny Cove | (57:17) Ohrožení Intelu procesory ARM | (1:07:06) Dotazy / očekávání roku 2019}{1:13:12}{34443324}{}
\FC{hardware-club-024.mp3}{\HCtitle{\#24}{To nejdůležitější z CES 2019}}{2019-01-14}{\Sleepless, \Janchor, \JSvak, \JZubec}{}{(02:30) Nvidia GeForce RTX 2060: výkon a cena | (10:39) FreeSync na Nvidia kartách? | (26:17) AMD na CES: prezentace Lisy Su | (31:47) AMD a nová konfigurace CPU | (34:10) Jak budou vypadat nové Ryzeny? | (39:49) Stane se z AMD druhý Intel? | (49:23) Nová grafika AMD Radeon VII: výkon, cena, dostupnost | (54:01) Dotazy I | (1:04:09) Intel na závěr | (1:09:14) Dotazy II}{1:13:11}{36866961}{}
\FC{hardware-club-025.mp3}{\HCtitle{\#25}{Mini PC}}{2019-01-24}{\Sleepless, \Janchor, \JSvak, \JZubec}{}{(05:33) Co je Mini PC | (21:48) Chlazení | (27:25) Grafika vs iGPU | (37:25) Pasivní HTPC | (47:45) Intel NUC | (50:17) AMD APU}{1:07:56}{32853753}{}
\FC{hardware-club-026.mp3}{\HCtitle{\#26}{AI neboli umělá inteligence}}{2019-02-14}{\Janchor, \JSvak, \JZubec}{}{(01:53) Průmyslová revoluce v4.0? | (04:46) Vývoj AI | (09:55) Hardware AI | (24:45) Princip AI | (1:03:30) Současné aplikace AI a potenciální problémy | (1:54:27) AI jako další vývojový stupeň člověka?}{2:08:21}{62102196}{}
\FC{hardware-club-027.mp3}{\HCtitle{\#27}{AI neboli umělá inteligence – pokračování}}{2019-02-28}{\Janchor, \JSvak, \JZubec}{}{(01:02) Budoucnost je dnes! | (05:36) Osud lidstva v éře AI? | (16:00) Superinteligence | (23:54) Co náš čeká v nejbližších letech? | (33:13) Pozitivní/negativní vize světa s AI | (49:25) AI ve sci-fi dílech}{1:28:53}{42222816}{}
\FC{hardware-club-028.mp3}{\HCtitle{\#28}{Vyplatí se dnes hotové sestavy z obchodu?}}{2019-03-14}{\Janchor, \JSvak, \JZubec}{}{(05:13) Co se změnilo na trhu s hotovými PC sestavami? | (09:09) Srovnání nejprodávanějších PC sestav v ČR | (26:26) Výhody a nevýhody hotových sestav | (41:47) Kvalita sestavení | (54:31) Konfigurovatelná PC | (57:56) Rady na závěr}{1:08:14}{32157936}{}
\FC{hardware-club-029.mp3}{\HCtitle{\#29}{Ray tracing pro všechny, výkon GTX 1660 / 1660 Ti a další novinky}}{2019-03-28}{\Janchor, \JSvak, \JZubec}{}{(01:23) Ray tracing bez karet Nvidia RTX | (13:51) GTX 1660 / 1660 Ti: výkon ve Full HD | (18:45) Poměr cena/výkon | (25:06) Výkon ve QHD | (32:04) Spotřeba | (33:57) Ryzen 3000: 16/32 jader!}{0:52:01}{24322909}{}
\FC{hardware-club-030.mp3}{\HCtitle{\#30}{Kolik paměti je potřeba v herním PC?}}{2019-04-15}{\Janchor, \JSvak, \JZubec}{}{(02:20) RAM vs VRAM | (13:14) Vliv na výkon | (15:24) Paměť a nové hry | (31:34) Kolik RAM/VRAM mít | (40:04) Aktuální stav trhu | (44:15) Časování a latence}{1:22:07}{39592728}{}
\FC{hardware-club-031.mp3}{\HCtitle{\#31}{Stavíme herní PC (jarní aktualizace)}}{2019-04-25}{\Janchor, \JSvak, \JZubec}{}{(01:46) PC do 10\,000\,Kč | (14:00) Herní PC do 15\,000\,Kč | (31:28) Střední třída do 30\,000\,Kč | (46:15) Dělo za 40\,000\,Kč s RTX 2080 | (55:35) Nejvýkonnější herní PC za 68\,000\,Kč}{1:19:47}{39644100}{}
\FC{hardware-club-032.mp3}{\HCtitle{\#32}{Jubilejní s Michalem Rybkou}}{2019-05-16}{\Janchor, \JSvak, \JZubec, \Wolf}{}{(00:27) Představení hosta | (02:10) Téma: Google Stadia a streamované hraní | (55:40) Várka dotazů | (1:22:06) Soutěž | (1:24:19) Nečekaný vstup}{1:29:28}{42652512}{}
\FC{hardware-club-033.mp3}{\HCtitle{\#33}{Computex 2019}}{2019-05-30}{\Sleepless, \Janchor, \JSvak, \JZubec}{}{(01:05) Nové procesory AMD Ryzen 3000 | (25:41) Ryzeny 3000 a nové konzole | (34:58) Nové grafiky AMD Navi | (53:37) Dotazy (57:35) Procesory Intel 10. generace}{1:03:35}{31435596}{}
% \FC{hardware-club-034.mp3}{\HCtitle{\#34}{Xbox Scarlett, 16jádro AMD a nové grafické karty (#E3 2019)}}{2019-06-14}{\JSvak, \JZubec}{}{(02:19) AMD Ryzen 9 3950X | (22:04) Nové Radeony | (26:44) RX 5700 XT | (37:06) RX 5700 (bez XT) | (45:34) Xbox Scarlett | (1:00:09) Nvidia SUPER}{1:20:18}{43288388}{}
% \FC{hardware-club-035.mp3}{\HCtitle{\#35}{Hraní na Linuxu v roce 2019}}{2019-06-27}{\Janchor, \JSvak, \JZubec}{}{(01:38) Hraní a Linux: oxymóron? | (04:41) Stav Linuxu jako herní platformy | (15:36) Jak a co si lze na Linuxu zahrát | (29:46) SteamPlay/Proton | (43:09) ProtonDB | (46:41) Problémy Linuxu}{1:18:07}{38465424}{}
% \FC{hardware-club-036.mp3}{\HCtitle{\#36}{Výkon procesorů Ryzen 3000 a GPU Navi}}{2019-07-11}{\Janchor, \JSvak, \JZubec}{}{(03:22) Přehled modelů | (15:54) Výkon Ryzen 3000 | (25:22) Přetaktování | (28:51) Single thread | (30:10) Herní testy, vliv RAM | (50:47) Spotřeba | (59:49) Nové grafiky AMD RX 5700 / 5700 XT}{1:24:39}{40874796}{}
% \FC{hardware-club-037.mp3}{\HCtitle{\#37}{Paměti hrozí zdražováním, testy RTX 2080 Super a další novinky}}{2019-07-25}{\Sleepless, \Janchor, \JSvak, \JZubec}{}{(02:22) Situace na trhu s RAM, spor Japonsko vs. Korea | (14:41) Koupit operační paměť teď, nebo čekat? | (20:14) SpaceX | (44:06) Herní AI Google | (46:38) GeForce RTX 2080 Super: stojí za to? | (51:16) Nové Raspberry Pi}{0:59:58}{28244412}{}
% \FC{hardware-club-038.mp3}{\HCtitle{\#38}{Jak si vedou starší procesory v roce 2019?}}{2019-08-15}{\Janchor, \JSvak, \JZubec}{}{(01:22) Výkon procesorů od 486 | (20:24) Vývoj od Sandy Bridge dále | (30:05) Srovnání v Cinebench | (41:05) Herní výkon napříč generacemi | (45:28) GTA V}{1:03:30}{31561344}{}
% \FC{hardware-club-039.mp3}{\HCtitle{\#39}{Nové herní sestavy s Ryzeny a AMD Navi}}{2019-08-29}{\JSvak, \JZubec}{}{1:50 Změny na trhu | 14:53 Základní herní PC za 12000 Kč | 25:39 Počítač za 17000 Kč | 43:31 Herní PC za cca 30000 Kč | 57:26 Sestava High aka co nejvyšší výkon kolem 40000 Kč | 1:06:29 Sestava TOP OC}{1:23:07}{42933024}{}
% \FC{hardware-club-040.mp3}{\HCtitle{\#40}{O nejnáročnějších PC hrách}}{2019-09-12}{\JSvak, \JZubec}{}{5:13 90. léta: DOOM, Descent, Wing Commander, Quake I \&{} II, Total Annihilation, Unreal, Half Life, Outcast | 47:44 Rok 2000+: Giants: Citizen Kabuto, Morrowind, Far Cry, Half Life2, DOOM 3, Fear, TES V: Oblivion | 57:52 Rok 2007 a legendární Crysis | 1:04:07 2008+: GTA IV, S.T.A.L.K.E.R.: Call of Pripyat, Metro Last Light, Crysis 3}{1:25:49}{45336744}{}
% \FC{hardware-club-041.mp3}{\HCtitle{\#41}{O bazarovém hardwaru}}{2019-09-26}{\JSvak, \JZubec}{}{(02:53) Výhody a nevýhody nákupu z bazaru | (10:28) Jaké komponenty z bazaru koupit a jaké raději ne? | (38:52) Kde nakupovat/prodávat + jak se vyhnout podvodníkům | (57:21) Co můžeme dnes najít v bazaru a za kolik? Co se vyplatí? | (1:28:29) Naše zkušenosti, nejlepší/nejhorší koupě z bazaru}{1:34:44}{48602052}{}
% \FC{hardware-club-042.mp3}{\HCtitle{\#42}{O herních klávesnicích}}{2019-10-17}{\JSvak, \JZubec}{}{(02:48) Co je to herní klávesnice a jaké má mít parametry? | (13:38) Podsvícení kláves | (28:51) Mechanické klávesnice| (48:47) Proč jsou mechanické klávesnice hlučné? | (55:18) Exotické a „sci-fi“ klávesnice}{1:27:22}{44062164}{}
% \FC{hardware-club-043.mp3}{\HCtitle{\#43}{O  budoucnosti výroby hardwaru}}{2019-11-01}{\JSvak, \JZubec}{}{(03:21) GeForce GTX 1660 Super | (10:58) Intel Core i9-9900KS | (14:42) Zvýšené TDP a háček v 5GHz boostu | (27:28) Výroba mikročipu – vývoj a dnešní stav | (44:40) EUVL a technologie ASML | (59:33) Budoucnost litografie}{1:25:18}{45489168}{}
% \FC{hardware-club-044.mp3}{\HCtitle{\#44}{Jak měřit FPS ve hrách}}{2019-11-14}{\JSvak, \JZubec}{}{(02:11) Co je to FPS a jak to měřit? | (08:42) Pokročilé možnosti | (22:03) FPS versus frametimes | (38:25) Novinky – faul marketingu Intel + další zranitelnosti CPU | (54:32) AMD Radeon 5500 | (59:54) Drby o Nvidia Ampere}{1:11:58}{40026996}{}
% \FC{hardware-club-045.mp3}{\HCtitle{\#45}{Neznakomp Vašek staví počítač}}{2019-12-05}{\JSvak, \JZubec, \VPechacek}{}{Překoná Václav Pecháček nástrahy konektorů a drátoví?}{1:33:54}{}{}
% \FC{hardware-club-046.mp3}{\HCtitle{\#46}{AMD a Intel na CES 2020 a další best of veletrhu}}{2020-01-09}{\JSvak, \JZubec}{}{Nejzajímavější představení z veletrhu CES 2020 v Las Vegas. Threadripper, APU Renoir a další novinky od AMD a Intel.}{1:18:08}{49921776}{}
% \FC{hardware-club-047.mp3}{\HCtitle{\#47}{Kolik jader procesoru je už moc?}}{2020-01-23}{\JSvak, \JZubec}{}{S vydáním Threadripperu 3990X s 64 jádry / 128 vlákny se na jazyk znova krade otázka, zda je tak rapidní vývoj v počtu jader u spotřebitelských procesorů vůbec k něčemu.}{1:09:28}{40501404}{}
% \FC{hardware-club-048.mp3}{\HCtitle{\#48}{Zpackané vydání Radeonu RX 5600 XT}}{2020-02-06}{\JSvak, \JZubec}{}{Nová grafická karta AMD Radeon RX 5600 XT dostala pod tlakem konkurence na poslední chvíli injekci s vyšším výkonem. Co taková změna způsobila a proč AMD raději kartu nezlevnilo?}{1:00:15}{31151844}{}
% \FC{hardware-club-049.mp3}{\HCtitle{\#49}{Cena PlayStation 5 a dopad koronaviru na HW}}{2020-02-20}{\JSvak, \JZubec}{}{Oproti minulé generaci se na startu prodeje PlayStationu 5 dočkáme zdražení, vyplývá ze zprávy Bloomberg, která odhalila pravděpodobnou výrobní cenu nové konzole. Nemůže za to třeba koronavirus?}{0:56:37}{27834444}{}
% \FC{hardware-club-050.mp3}{\HCtitle{\#50}{O modování her}}{2020-03-05}{\JSvak, \JZubec}{}{Jedni modování odsuzují, pro druhé je to prakticky další náplň hry. Jisté ale je, že mody jsou recept na dlouhověkost her: Skyrim, Half-Life, DOOM, nebyly by bez modů už v propadlišti dějin?}{1:22:55}{43613532}{}
% \FC{hardware-club-051.mp3}{\HCtitle{\#51}{Nové procesory Intel Comet Lake-S: 10/20 jader, ale pořád 14 nm}}{2020-04-30}{\JSvak, \JZubec}{}{Na vydání nových desktopových procesorů Intel se čeká už dlouho. Poslední refresh vyšel před dvěma lety a mezitím AMD zamíchalo kartami. Co novinky přinesou a vrátí Intel zase na trůn PC builderů?}{1:14:05}{33560208}{}
% \FC{hardware-club-052.mp3}{\HCtitle{\#52}{Nové GPU Ampere od Nvidia}}{2020-05-14}{\JSvak, \JZubec}{}{Jensen Huang představil novou architekturu GPU Ampere, ač zatím pouze jako výpočetní čip pro strojové učení. Epic ukázal demo Unreal Engine 5. Oboje si v 52. HWC probereme.}{0:54:40}{26402544}{}
% \FC{hardware-club-053.mp3}{\HCtitle{\#53}{}}{2020--}{\JSvak, \JZubec}{}{}{}{}{}
% \FC{hardware-club-054.mp3}{\HCtitle{\#54}{}}{2020--}{\JSvak, \JZubec}{}{}{}{}{}
% \FC{hardware-club-055.mp3}{\HCtitle{\#55}{}}{2020--}{\JSvak, \JZubec}{}{}{}{}{}
% \FC{hardware-club-056.mp3}{\HCtitle{\#56}{}}{2020--}{\JSvak, \JZubec}{}{}{}{}{}
% \FC{hardware-club-057.mp3}{\HCtitle{\#57}{}}{2020--}{\JSvak, \JZubec}{}{}{}{}{}
% \FC{hardware-club-058.mp3}{\HCtitle{\#58}{}}{2020--}{\JSvak, \JZubec}{}{}{}{}{}
% \FC{hardware-club-059.mp3}{\HCtitle{\#59}{}}{2020--}{\JSvak, \JZubec}{}{}{}{}{}
% \FC{hardware-club-060.mp3}{\HCtitle{\#60}{}}{2020--}{\JSvak, \JZubec}{}{}{}{}{}
% \FC{hardware-club-061.mp3}{\HCtitle{\#61}{}}{2020--}{\JSvak, \JZubec}{}{}{}{}{}
% \FC{hardware-club-062.mp3}{\HCtitle{\#62}{}}{2020--}{\JSvak, \JZubec}{}{}{}{}{}
% \FC{hardware-club-063.mp3}{\HCtitle{\#63}{}}{2020--}{\JSvak, \JZubec}{}{}{}{}{}
% \FC{hardware-club-064.mp3}{\HCtitle{\#64}{}}{2020--}{\JSvak, \JZubec}{}{}{}{}{}
% \FC{hardware-club-065.mp3}{\HCtitle{\#65}{}}{2020--}{\JSvak, \JZubec}{}{}{}{}{}
% \FC{hardware-club-066.mp3}{\HCtitle{\#66}{}}{2020--}{\JSvak, \JZubec}{}{}{}{}{}
% \FC{hardware-club-067.mp3}{\HCtitle{\#67}{}}{2020--}{\JSvak, \JZubec}{}{}{}{}{}
% \FC{hardware-club-068.mp3}{\HCtitle{\#68}{}}{2020--}{\JSvak, \JZubec}{}{}{}{}{}
% \FC{hardware-club-069.mp3}{\HCtitle{\#69}{}}{2020--}{\JSvak, \JZubec}{}{}{}{}{}



\end{document}
